\documentclass[leqno]{article}

\usepackage[english]{babel}

\usepackage[utf8]{inputenc}
\usepackage{amsmath}
\usepackage{amssymb}
\usepackage{amsfonts}
\usepackage{enumerate}
\usepackage{vmargin}
\usepackage[all]{xy}
\usepackage{mathrsfs}
\usepackage{mathtools}
\usepackage{lmodern}
\usepackage{slashed}
\usepackage{xurl}
\usepackage[colorlinks=true,linkcolor=blue,pagebackref=true,hyperfootnotes=false]{hyperref}

\usepackage{array}
\setmarginsrb{3cm}{3cm}{3.5cm}{3cm}{0cm}{0cm}{1.5cm}{3cm}










\usepackage{comment}


\IfFileExists{accents.sty}{\usepackage{accents}}{}





\newlength{\dhatheight}
\newcommand{\doublehat}[1]{
    \settoheight{\dhatheight}{\ensuremath{\hat{#1}}}
    \addtolength{\dhatheight}{-0.2ex}
    \hat{\vphantom{\rule{1pt}{\dhatheight}}
    \smash{\hat{#1}}}}



\footskip1.3cm

\newcommand{\A}{\ensuremath{\mathcal{A}}}
\newcommand{\Cbb}{\ensuremath{\mathbb{C}}}
\newcommand{\C}{\mathrm{C}}

\newcommand{\D}{\mathrm{D}}
\newcommand{\E}{\ensuremath{\mathbb{E}}}
\newcommand{\F}{\mathcal{F}}
\newcommand{\N}{\ensuremath{\mathbb{N}}}
\newcommand{\B}{\mathrm{B}} 
\let\H\relax 
\newcommand{\H}{\mathrm{H}}
\newcommand{\I}{\mathrm{I}} 
\newcommand{\II}{\mathrm{II}} 
\newcommand{\K}{\ensuremath{\mathbb{K}}}
\newcommand{\KP}{\ensuremath{\mathbb{KP}}}
\let\L\relax 
\newcommand{\L}{\mathrm{L}}
\let\O\relax 
\newcommand{\J}{\mathcal{J}} 
\newcommand{\O}{\ensuremath{\mathbb{O}}}
\let\P\relax 

\newcommand{\P}{\mathsf{P}} 
\newcommand{\Q}{\mathrm{Q}}
\newcommand{\IV}{\mathrm{IV}} 
\newcommand{\V}{\mathrm{V}} 
\newcommand{\VI}{\mathrm{VI}} 
\newcommand{\SG}{\mathrm{SG}} 
\newcommand{\BMO}{\mathrm{BMO}} 
\newcommand{\Mult}{\mathrm{Mult}}
\newcommand{\scr}{\mathscr}
\newcommand{\SU}{\mathrm{SU}}
\newcommand{\Tron}{\mathcal{T}}
\newcommand{\SO}{\mathrm{SO}}
\newcommand{\Aut}{\mathrm{Aut}} 
\newcommand{\QG}{\ensuremath{\mathbb{G}}}
\newcommand{\M}{\mathrm{M}}
\newcommand{\CAR}{\mathrm{CAR}} 
\newcommand{\CB}{\mathrm{CB}} 
\newcommand{\diag}{\mathrm{diag}} 
\newcommand{\RCD}{\mathrm{RCD}}
\newcommand{\dist}{\mathrm{d}}

\newcommand{\meas}{\mu}
\newcommand{\CC}{\mathrm{CC}}
\newcommand{\HS}{\mathrm{HS}}
\newcommand{\absco}{\mathrm{absco}} 
\newcommand{\NP}{\mathrm{NP}}
\newcommand{\PPT}{\mathrm{PPT}}
\newcommand{\inv}{\diamond}
\newcommand{\Mat}{\ensuremath{\mathbb{M}}}

\let\div\relax
\newcommand{\div}{\mathrm{div}}
\let\cal\relax
\newcommand{\cal}{\mathcal}
\newcommand{\conv}{\mathrm{conv}}
\newcommand{\Z}{\ensuremath{\mathbb{Z}}}
\newcommand{\Zc}{\mathrm{Z}}
\newcommand{\R}{\ensuremath{\mathbb{R}}}

\newcommand{\T}{\ensuremath{\mathbb{T}}}
\newcommand{\U}{\mathcal{U}}
\newcommand{\W}{\mathrm{W}}
\newcommand{\sgn}{\mathrm{sgn}}
\let\j\relax 
\newcommand{\j}{\frak{j}} 
\let\S\relax 
\newcommand{\S}{\mathrm{S}} 
\newcommand{\vect}{\ensuremath{\mathop{\rm Span }\nolimits}}

\newcommand{\clos}{\ensuremath{\mathop{\rm closure }\nolimits}}
\newcommand{\Id}{\mathrm{Id}}
\newcommand{\PL}{\mathrm{PL}}
\newcommand{\VN}{\mathrm{VN}} 
\newcommand{\Asym}{\mathrm{Asym}} 
\newcommand{\Sym}{\mathrm{Sym}} 
\newcommand{\Reg}{\mathrm{Reg}} 
\newcommand{\la}{\langle}
\newcommand{\ra}{\rangle}
\newcommand{\EA}{\mathrm{EA}} 
\newcommand{\CPAP}{\mathrm{CPAP}}
\newcommand{\CD}{\mathrm{CD}}
\newcommand{\flip}{\mathrm{flip}}
\newcommand{\CBAP}{\mathrm{CBAP}}
\newcommand{\MK}{\mathrm{MK}} 
\newcommand{\CCAP}{\mathrm{CCAP}}
\newcommand{\ccb}{\mathrm{ccb}} 
\newcommand{\cocb}{\mathrm{cocb}}
\renewcommand{\leq}{\ensuremath{\leqslant}}
\renewcommand{\geq}{\ensuremath{\geqslant}}
\newcommand{\qed}{\hfill \vrule height6pt  width6pt depth0pt}
\newcommand{\bnorm}[1]{ \big\| #1  \big\|}
\newcommand{\Bnorm}[1]{ \Big\| #1  \Big\|}
\newcommand{\bgnorm}[1]{ \bigg\| #1  \bigg\|}
\newcommand{\Bgnorm}[1]{ \Bigg\| #1  \Bigg\|}
\newcommand{\norm}[1]{\left\Vert#1\right\Vert}
\newcommand{\Wa}{\mathcal{W}} 
\newcommand{\mk}{\medskip}
\newcommand{\xra}{\xrightarrow}
\newcommand{\co}{\colon}
\newcommand{\otp}{\widehat{\ot}}
\newcommand{\otpb}{\hat{\ot}}
\newcommand{\ot}{\otimes}
\newcommand{\ovl}{\overline}
\newcommand{\otvn}{\ovl\ot}

\newcommand{\ul}{\mathcal{U}}
\newcommand{\dec}{\mathrm{dec}}
\newcommand{\pdec}{\mathrm{pdec}}
\newcommand{\reg}{\mathrm{reg}}
\newcommand{\Dec}{\mathrm{Dec}}
\newcommand{\cb}{\mathrm{cb}}
\newcommand{\fin}{\mathrm{fin}} 
\newcommand{\TRO}{\mathrm{TRO}}
\newcommand{\Ch}{\mathrm{Ch}}
\newcommand{\PDec}{\mathrm{PDec}}
\newcommand{\QWEP}{\mathrm{QWEP}}
\newcommand{\dsp}{\displaystyle}
\newcommand{\op}{\mathrm{op}} 
\let\i\relax 
\newcommand{\i}{\mathrm{i}}
\let\j\relax 
\newcommand{\j}{\mathrm{j}} 
\let\k\relax 
\newcommand{\k}{\mathrm{k}} 
\newcommand{\loc}{\mathrm{loc}}
\newcommand{\ov}{\overset}
\newcommand{\sa}{\mathrm{sa}}
\newcommand{\JW}{\mathrm{JW}}
\newcommand{\JC}{\mathrm{JC}}
\newcommand{\JB}{\mathrm{JB}}
\newcommand{\kR}{\mathfrak{R}}
\newcommand{\JBW}{\mathrm{JBW}}
\newcommand{\Vol}{\mathrm{Vol}}
\newcommand{\fix}{\mathrm{fix}}
\newcommand{\ORNP}{\mathrm{ORNP}}
\newcommand{\MD}{\mathrm{MD}} 
\newcommand{\LMD}{\mathrm{LMD}} 
\newcommand{\RMD}{\mathrm{RMD}} 
\newcommand{\w}{\mathrm{w}} 
\newcommand{\wo}{\mathrm{wo}} 
\newcommand{\leftsupp}{{\mathrm{leftsupp}}}
\newcommand{\epsi}{\varepsilon}
\renewcommand{\d}{\mathop{}\mathopen{}\mathrm{d}} 
\newcommand{\e}{\mathrm{e}} 
\renewcommand{\d}{\mathop{}\mathopen{}\mathrm{d}} 
\newcommand{\bra}[1]{\langle{#1}|}
\newcommand{\ket}[1]{|{#1}\rangle}
\newcommand{\Pol}{\mathrm{Pol}} 

\def\kb{ \ket \bra }
\newcommand{\proj}[1]{ | #1 \kb  #1|}
\DeclareMathOperator{\Span}{span} 
\DeclareMathOperator{\Lip}{\mathrm{Lip}} 
\DeclareMathOperator{\lip}{\mathrm{lip}} 
\DeclareMathOperator{\aut}{Aut} 
\DeclareMathOperator{\supp}{supp} 
\DeclareMathOperator{\Fix}{Fix} 
\DeclareMathOperator{\tr}{Tr} 
\DeclareMathOperator{\Tr}{Tr} 
\let\ker\relax 
\DeclareMathOperator{\ker}{Ker} 
\DeclareMathOperator{\Ran}{Ran} 
\DeclareMathOperator{\dom}{dom} 
\DeclareMathOperator{\rank}{rank} 
\let\Re\relax 
\DeclareMathOperator{\Re}{Re} 
\let\Im\relax 
\DeclareMathOperator{\Im}{Im} 
\DeclareMathOperator{\card}{card} 
\DeclareMathOperator{\Irr}{Irr} 
\DeclareMathOperator{\Rep}{Rep} 
\DeclareMathOperator{\Ind}{Ind} 
\DeclareMathOperator*{\esssup}{esssup} 
\DeclareMathOperator*{\Lie}{Lie} 
\DeclareMathOperator{\diam}{diam}
\DeclareMathOperator{\Sp}{Sp}


\selectlanguage{english}
\newtheorem{thm}{Theorem}[section]
\newtheorem{defi}[thm]{Definition}
\newtheorem{prop}[thm]{Proposition}
\newtheorem{conj}[thm]{Conjecture}
\newtheorem{Property}[thm]{Property}
\newtheorem{cor}[thm]{Corollary}
\newtheorem{lemma}[thm]{Lemma}
\newtheorem{fact}[thm]{Fact}
\newtheorem{remark}[thm]{Remark}
\newtheorem{example}[thm]{Example}
\newtheorem{prob}[thm]{Problem}


\newenvironment{proof}[1][]{\noindent {\it Proof #1} : }{\hbox{~}\qed
\smallskip
}








\usepackage{tocloft}
\setlength{\cftbeforesecskip}{0pt}
		
\numberwithin{equation}{section}
\usepackage[nottoc,notlot,notlof]{tocbibind}






\let\OLDthebibliography\thebibliography
\renewcommand\thebibliography[1]{
  \OLDthebibliography{#1}
  \setlength{\parskip}{0pt}
  \setlength{\itemsep}{0pt plus 0.3ex}
}




\newcommand\reallywidehat[1]{\arraycolsep=0pt\relax
\begin{array}{c}
\stretchto{
  \scaleto{
    \scalerel*[\widthof{\ensuremath{#1}}]{\kern-.5pt\bigwedge\kern-.5pt}
    {\rule[-\textheight/2]{1ex}{\textheight}} 
  }{\textheight} 
}{0.5ex}\\           
#1\\                 
\rule{-1ex}{0ex}
\end{array}
}


\emergencystretch=3em

\begin{document}
\selectlanguage{english}
\title{\bfseries{The Haagerup center at central charge eight via conformal nets}}
\date{}
\author{\bfseries{C\'edric Arhancet}}
\maketitle




\begin{abstract}
Let $S\subset V_{E_6\oplus A_2}$ be the unitary vertex operator algebra of central charge eight with currents at levels $(3,13)$ constructed in our previous work. An exact presentation of a full corner of its Zhu algebra gives strong rationality, the twelve simple modules and the lattice branching rules. Using these results, we identify the representation category of its conformal net with the Haagerup center $Z(\cal H_3)$. We construct a holomorphic diagonal extension of $\cal A_S\otimes \cal A_S$ and apply the vacuum trace class theorem of Henriques and Tener to prove trace class heat operators in every sector of $\cal A_S$. Reconstruction gives a unitary holomorphic vertex operator algebra containing $S\ot S$ as mutual commutants. The coset theorems of Gui then give the complete unitarity of $S$ and a braided equivalence $\Rep(S)\cong Z(\cal H_3)$, so that $S$ is a Haagerup algebra. Using results of Dong, Jiao and Xu, we also identify the Huang--Riesen algebra with the twin of $S$ and describe their common fixed point algebra using the intersection of the two binary lattices of discriminant $-39$.
\end{abstract}




\makeatletter
 \renewcommand{\@makefntext}[1]{#1}
 \makeatother
 \footnotetext{
 2020 {\it Mathematics subject classification:}
 17B69, 46L37, 81T40, 18M20, 11H56.
\\
{\it Key words}: vertex operator algebra, Haagerup subfactor, conformal net, Drinfeld center, hyperbolic genus, simple current extension, complete unitarity.}

{
  \hypersetup{linkcolor=blue}
 \tableofcontents
}


%%%%%%%%%%%%%%%%%%%%%%%%%%%%%%%%%%%%%%%%%%%%%%%%%%%%%%%%%%%%%%%%%%%%%%%%%%%%%
\section{Introduction}
\label{sec:intro}

\paragraph{Background.}
The Haagerup subfactor \cite{Haa94,AsH99} is the smallest finite depth subfactor of index greater than four, with index $\frac{5+\sqrt{13}}2$, and its even parts $\cal H_1$, $\cal H_2$ and the category $\cal H_3$ of \cite[Definition 3.19 p.~23]{GrS12} are the three Morita equivalent Haagerup fusion categories of \cite{GrS12}. Their common Drinfeld center $Z(\cal H_3)$ is a unitary modular tensor category of rank twelve whose modular data was computed by Izumi \cite{Izu01} and by Evans and Gannon \cite{EvG11}, who conjectured that it is realized by a rational conformal field theory of central charge eight with six currents. Explicit candidates were announced in October 2026 by Riesen \cite{Rie26} and by Huang \cite{Hua26}, with currents $\mathfrak{sl}_2 \oplus \mathfrak{sl}_2$ at levels $(1,39)$, and by Jiao \cite{Jia26}, with currents at levels $(3,13)$, all as vertex subalgebras of lattice vertex operator algebras.

In \cite{Arh26} we constructed the algebra $S \subset V' = V_{E_6 \oplus A_2}$ with currents at levels $(3,13)$ from the weight-one and weight-two data forced by the modular data, we proved that it is simple, unitary, $C_2$-cofinite, strongly generated in weights at most three, and isomorphic to the algebra of Jiao, and we developed the structure theory of \emph{Haagerup algebras}, the positive strongly rational vertex operator algebras of central charge eight with representation category $Z(\cal H_3)$: their weight-one Lie algebras, their lowest weight spaces, their Heisenberg commutants and their lattices, and a twin construction exchanging the two binary lattices of discriminant $-39$. In \cite{ArhII} we studied the conformal net $\cal A_S$ and used the hypothesis
\[
\mathrm{(H_J)} \co \text{$S$ is a Haagerup algebra},
\]
for the results specific to $S$. Appendix \ref{app:zhu} computes a full corner of Zhu's algebra $A(S)$, proves strong rationality and positivity, determines the twelve simple modules and gives the complete restrictions of the $V'$-modules. Theorem \ref{thm:algebraic-F4} deduces from it an algebraic description of $\Rep(S)$ as the Drinfeld center of a pseudo-unitary fusion category, which is not used in the proof of Theorem \ref{thm:HJ-main}.

\paragraph{The analytic proof.}
The algebraic inputs are stated in Section \ref{sec:inputs}. The unitary realization of $S$ and its weight-one Lie algebra (I1) are given by \cite[Theorems 10.10 and 10.21]{Arh26}. The calculation of Zhu's algebra in Appendix \ref{app:zhu} supplies strong rationality and positivity (F1), the twelve simple modules and their lowest weight data (F2), and the lattice restrictions (F3).

Theorem \ref{thm:net} (Theorem B) identifies the completely rational net category $\Rep(\cal A_S)$ with $Z(\cal H_3)$ using (I1), (F1) and (F3). For each simple $S$-summand $M$ of the two lattice extensions of $V'$, \cite[Proposition 5.2]{ArhII} gives an invariant scalar product on $\ovl M \boxtimes M$ for which this fusion product is a transport module in the sense of \cite[Definition 3.3 p.~26]{Ten24}, and this is the input of the finite-index criterion of Tener. The recognition theorem of \cite[Theorem 6.1]{ArhII} then applies to the unitary fusion categories of the net extensions.

Section \ref{sec:diagonal} constructs the holomorphic diagonal net $\cal B$ with the same stress-energy tensor as the represented product net. The Hamiltonian identity of Proposition \ref{prop:net-B} lets us apply the vacuum trace theorem of Henriques and Tener \cite{HeT26} to $\cal B$ and deduce trace class for every sector of $\cal A_S$ by compression. Reconstruction \cite{HeT25} gives a unitary holomorphic vertex operator algebra containing $S \ot S$, with the prescribed scalar products on both copies of $S$. Theorem \ref{thm:HJ-main} applies the coset comparison of Section \ref{sec:unitarity} to obtain $\mathrm{(H_J)}$, complete unitarity and the VOA--net equivalence from (I1) and (F1)--(F3). Since (I1) and (F1)--(F3) are established for the explicit algebra $S$ in Theorems \ref{thm:inputs-ArhI} and \ref{thm:zhu-inputs}, Theorem \ref{thm:HJ-main} proves that $S$ is a Haagerup algebra, that $S$ is completely unitary and that $\Rep^u(S) \cong \Rep(\cal A_S) \cong Z(\cal H_3)$ as unitary modular tensor categories. This is the main result of the article. Corollary \ref{cor:S-consequences} identifies the lattice labels across the equivalence and the third holomorphic net extension.

\paragraph{The twin identification.}
Let $\mathfrak h \subset S_1$ be the Cartan subalgebra spanned by the Cartan currents $p$, $b$ of levels $3$ and $13$, let $\Omega_S = \operatorname{Com}_S(M(\mathfrak h))$ be the Heisenberg commutant, and let $L_{\mathrm{II}}$ and $L_{\mathrm I}$ be the two even lattices of discriminant $-39$ of \cite[Lemma 12.1]{Arh26}, with Gram matrices $\left[\begin{smallmatrix}6&3\\3&8\end{smallmatrix}\right]$ and $\left[\begin{smallmatrix}2&1\\1&20\end{smallmatrix}\right]$. In Section \ref{sec:genus} we show that the construction of \cite[Section 12]{Arh26} applies to $S$ under (I1) and (F1)--(F4): $S = \bigoplus_{\alpha} \Omega_S^{\tau(\alpha)} \ot V_{\alpha+L_{\mathrm{II}}}$ over $\alpha \in L_{\mathrm{II}}'/L_{\mathrm{II}} \cong \Z_{39}$, and exchanging $L_{\mathrm{II}}$ for $L_{\mathrm I}$ gives the twin $S^\natural$. The key observation is that the involutions $\sigma$ of $S$ and $\tau$ of $U$ used in \cite{DJX26} are the characters $\e^{2\pi\i(s,\cdot)}$ and $\e^{2\pi\i(p,\cdot)}$ of $L_{\mathrm{II}}$ and $L_{\mathrm I}$, where $s = \frac{p+b}4$ is the Cartan current of level one, and that both have kernel $N = L_{\mathrm I} \cap L_{\mathrm{II}}$, a lattice of discriminant $-156$ whose discriminant form is the orthogonal sum of the metric group of $Z(\mathrm{Vec}_{\Z_2})$ and of a nondegenerate form on $\Z_{39}$. Hence the common fixed point algebra $C$ of \cite[Theorem 7.3]{DJX26} is the extension of $\Omega_S \ot V_N$ by the thirty-nine simple currents indexed by the $39$-part, and $S$, $U$ are its two extensions by the two bosons of $N'/N$.

\begin{thm}[Theorem A]
\label{thm:intro-twin}
Assume \textup{(I1)}, \textup{(I3)}, \textup{(F1)}--\textup{(F4)}, \textup{(D1)} and \textup{(D2)}. Then the algebra $U$ of Huang and Riesen is isomorphic to the twin $S^\natural$ of $S$. Consequently:
\begin{enumerate}
\item $U$ is a simple, self-contragredient, positive, strongly rational vertex operator algebra of CFT type and central charge eight, and $\Rep(U)$ is ribbon equivalent to $\Rep(S)$, a pseudo-unitary modular tensor category of rank twelve.
\item The Heisenberg commutant of $U$ is isomorphic to $\Omega_S$, its associated lattice is $L_{\mathrm I}$, and $U_1 = \mathfrak{sl}_2 \oplus \mathfrak{sl}_2$ with currents at levels $(1,39)$. In particular $U \not\cong S$.
\item $S$ and $U$ belong to the same hyperbolic genus, every strongly rational vertex operator algebra in this genus is isomorphic to $S$ or to $U$, and $S$, $U$ are the only positive strongly rational algebras of central charge eight with $\Rep$ ribbon equivalent to $\Rep(S)$ and Heisenberg commutant isomorphic to $\Omega_S$.
\item The common fixed point algebra $C = S^\sigma \cong U^\tau$ is strongly rational, and $\Rep(C) \cong Z(\mathrm{Vec}_{\Z_2}) \boxtimes \Rep(S)$ has rank $48$.
\end{enumerate}
\end{thm}

The lattice exchange of \cite[Theorem 12.3]{Arh26} constructs the twin. Theorem A identifies it with the explicit case (I) algebra using (D1) and (D2). It gives a second proof of the strong rationality of $U$, through the extension theory of \cite{HKL15,CKM22} rather than the rationality criterion of \cite[Theorem 3.1]{DJX26}, and a proof of $U_1 \cong \mathfrak{sl}_2 \oplus \mathfrak{sl}_2$ from the root system of the twin, without the weight-one computations of \cite{Hua26}. It also identifies the algebras themselves, whereas \cite[Proposition 9.2]{DJX26} identifies their module categories. The proof uses two uniqueness statements for simple current extensions with prescribed summands \cite[Proposition 5.3]{DoM04}, the quantum Galois theory of Dong and Mason \cite{DoM97}, the regularity of fixed point subalgebras \cite{CaM16}, and the finite identities of \cite{DJX26}.

\begin{thm}[Theorem D]
\label{thm:intro-consequences}
Assume \textup{(I1)}, \textup{(I3)}, \textup{(F1)}--\textup{(F4)}, \textup{(D1)} and \textup{(D2)}. Then $S$ and $U$ are Haagerup algebras, every Haagerup algebra is isomorphic to $S$ or to $U$, the Haagerup algebras form a single hyperbolic genus, every simple $S$-module is unitarizable, $S$ is completely unitary with $\Rep^u(S) \cong \Rep(\cal A_S) \cong Z(\cal H_3)$ as unitary modular tensor categories, the third holomorphic extension of $\cal A_S$ is isomorphic to $\cal A_{E_8}$, and the twin $U$ is completely unitary with $\Rep^u(U) \cong \Rep(\cal A_U) \cong Z(\cal H_3)$.
\end{thm}

Theorem \ref{thm:consequences} gives the detailed consequences in Theorem D. It collects the statements of \cite[Theorem 12.4]{Arh26} and \cite[Theorems 7.2, 7.17, 8.3 and 8.6, Corollaries 7.19 and 8.5, Proposition 7.10]{ArhII} which become available once $\mathrm{(H_J)}$ is known, with Theorem A in place of the hypothesis $\mathrm{(H_R)}$. The assertions of Theorem D on $S$, on its modules and on the net $\cal A_S$ only use (I1) and (F1)--(F3), through Theorem \ref{thm:HJ-main} and Corollary \ref{cor:S-consequences}. The inputs (I3), (F4), (D1) and (D2) enter only through Theorem A, for the assertions on $U$ and on the classification. All these inputs are established for the explicit algebras, in \cite{Arh26}, in Appendix \ref{app:zhu} and in \cite{DJX26}, see Section \ref{subsec:DJX-verification}, so that Theorems A and D hold for $S$ and $U$ without further hypothesis.

\paragraph{The algebraic comparison and the supporting calculations.}
Section \ref{sec:algebraic-route} collects the pseudo-unitary center description (F4), the algebraic construction from a braid-reversing equivalence, the compatible unitary structure found by exact lattice calculation and the equivalent formulations of Theorem \ref{thm:unitarity-equivalences} (Theorem C). Section \ref{sec:DJX} compares the results with those of Dong, Jiao and Xu and verifies the finite identities used for the twin theorem. Section \ref{sec:certificates} indexes the certificates. Appendix \ref{app:zhu} contains the proofs of the algebraic inputs.

\paragraph{Conventions.}
Vertex operator algebras are over $\Cbb$, of CFT type unless stated otherwise, and strongly rational means simple, rational, $C_2$-cofinite, self-contragredient and of CFT type, so that the module category is a modular tensor category \cite{Hua08}. Positive means that every simple module other than the algebra has positive lowest conformal weight. We write $\Rep(V)$ for the category of ordinary modules, $M \boxtimes N$ for the fusion product, $\cal C^{\mathrm{rev}}$ for a braided category with the inverse braiding, $\ovl{\cal C}$ for the complex conjugate category, and $Z(\cal F)$ for the Drinfeld center. The labels $0, b, a, c_0, c_1, c_2, d_1, \ldots, d_6$ of the simple objects of $Z(\cal H_3)$ and of the simple $S$-modules are those of \cite{EvG11} and \cite[Table 2]{Arh26}, with $d = \frac{3+\sqrt{13}}2$. The two simple $S$-modules labelled $a$ and $c_0$ in (F2) have the same lowest weight and the same lowest weight space, and they are not distinguished by the computation of Appendix \ref{app:zhu}. Before Section \ref{sec:center}, only the unordered pair $\{M_a,M_{c_0}\}$ is used, and the intrinsic labels are fixed in Section \ref{sec:center}, after Theorem \ref{thm:HJ-main}. Page references to arXiv preprints refer to the versions indicated in the bibliography.


%%%%%%%%%%%%%%%%%%%%%%%%%%%%%%%%%%%%%%%%%%%%%%%%%%%%%%%%%%%%%%%%%%%%%%%%%%%%%
\section{Inputs and their status}
\label{sec:inputs}

We state the lattice and representation-theoretic results used in the analytic argument. The proofs for the Zhu algebra and the lattice branching are in Appendix \ref{app:zhu}. The vertex operator algebra $S$ is defined in \cite[Theorems 10.9 and 10.21]{Arh26} as the vertex subalgebra of $V' = V_{E_6 \oplus A_2}$ generated by a Lie algebra $U_1 \cong \mathfrak{sl}_2 \oplus \mathfrak{sl}_2$ of currents at levels $(3,13)$ and by an explicit $120$-dimensional space $U_2 \subset V'_2$ defined over $\K = \Q(\sqrt{13},\sqrt{-3})$. We write $\mathfrak g' \oplus \mathfrak g''$ for $S_1$, $A = L_{\mathfrak{sl}_2}(3,0) \ot L_{\mathfrak{sl}_2}(13,0)$ for the affine subalgebra generated by $S_1$, which is simple since it is a unitary subalgebra of $(S,H_\chi)$ by \cite[Theorem 10.10]{Arh26} and \cite[Propositions 5.3 and 5.23 pp.~34 and 43]{CKLW18}, and $W_g$, $g \in \Z_3 \times \Z_3$, for the nine simple $V'$-modules, with $W_{(0,0)} = V'$.


%%%%%%%%%%%%%%%%%%%%%%%%%%%%%%%%%%%%%%%%%%%%%%%%%%%%%%%%%%%%%%%%%%%%%%%%%%%%%
\subsection{The unitary lattice realization of \texorpdfstring{$S$}{S}}

\begin{thm}[{\cite{Arh26}}]
\label{thm:inputs-ArhI}
\begin{enumerate}
\item[\textup{(I1)}] $S$ is a simple unitary vertex operator algebra of CFT type and central charge eight, self-contragredient, with the conformal vector of $V'$, and $S$ is a unitary subalgebra of $V'$ for the scalar product $H_\chi$ and its associated PCT operator $\theta$ \cite[Theorem 10.10]{Arh26}. Moreover $S_1 = U_1$, so that $\dim S_1 = 6$ \cite[Theorem 10.21]{Arh26}.
\item[\textup{(I2)}] $S$ is strongly generated in weights at most three, $\dim S_2 = 120$, $\dim S_3 = 956$, $\dim S_4 = 6010$, and $S$ is $C_2$-cofinite \cite[Proposition 10.17 and Theorem 10.21]{Arh26}. The weight-four certificate condition called (F4) in \cite[Definition 10.16]{Arh26} is unrelated to the input (F4) of Theorem \ref{thm:algebraic-F4}.
\end{enumerate}
\end{thm}

These statements are computer-assisted, with exact certificates over $\K$ \cite[Appendix A]{Arh26}. They make no hypothesis on the representation category of $S$.

The strong locality theorem and the submodule lemma of \cite[Theorem 4.5 and Lemma 4.7]{ArhII} apply to the unitary conformal embedding in (I1). We use the general transport and recognition results of \cite[Proposition 5.2 and Theorem 6.1]{ArhII} in Section \ref{sec:net}.


%%%%%%%%%%%%%%%%%%%%%%%%%%%%%%%%%%%%%%%%%%%%%%%%%%%%%%%%%%%%%%%%%%%%%%%%%%%%%
\subsection{Rationality, simple modules and lattice branching}

\begin{thm}
\label{thm:zhu-inputs}
\begin{enumerate}
\item[\textup{(F1)}] The vertex operator algebra $S$ is strongly rational and every simple $S$-module other than $S$ has positive lowest conformal weight. Its Zhu algebra is semisimple.
\item[\textup{(F2)}] $S$ has exactly twelve simple modules $M_z$ up to isomorphism, labelled by $z \in \{0,b,a,c_0,c_1,c_2,\allowbreak d_1,\ldots,d_6\}$ with $M_0 = S$. Its Zhu algebra has dimension $36625$. The lowest weight spaces of the simple modules $R_z$, as modules over $S_1$, are the twelve spaces of \cite[Table 2]{Arh26}, and their lowest conformal weights are
\[
h_0 = 0,
\quad
h_b = h_a = h_{c_0} = 1,
\quad
h_{c_1} = \tfrac13,
\quad
h_{c_2} = \tfrac23,
\]
\[
(h_{d_1},\ldots,h_{d_6}) = \big(\tfrac{6}{13},\tfrac{11}{13},\tfrac{15}{13},\tfrac{5}{13},\tfrac{7}{13},\tfrac{8}{13}\big).
\]
\item[\textup{(F3)}] $V'|_S \cong S \oplus M_b$, $W_{(0,\pm1)}|_S \cong M_{c_1}$, $W_{(\pm1,0)}|_S \cong M_{c_2}$, and the four modules $W_g$ with both coordinates of $g$ nonzero restrict to simple $S$-modules $X_+ = W_{(1,1)}|_S \cong W_{(-1,-1)}|_S$ and $X_- = W_{(1,-1)}|_S \cong W_{(-1,1)}|_S$ of lowest weight one, with $\{X_+,X_-\} = \{M_a,M_{c_0}\}$. The two holomorphic lattice extensions $H_\pm \cong V_{E_8}$ of $V'$ restrict to $S$ as $H_\pm|_S \cong S \oplus M_b \oplus 2X_\pm$, and we put $L_\pm=H_\pm|_S$.
\end{enumerate}
\end{thm}

\begin{proof}
The semisimplicity and rationality in (F1) are Theorem \ref{zhu:thm:Z} and Corollary \ref{zhu:cor:rational}, and positivity follows from Proposition \ref{zhu:prop:exact-h}. Theorem \ref{zhu:thm:twelve}, Table \ref{zhu:table:blocks} and Corollary \ref{zhu:cor:dimension} give (F2). Propositions \ref{zhu:prop:restrictions} and \ref{zhu:prop:a-c0} give the restrictions in (F3), and the two lattice decompositions in Section \ref{zhu:subsec:lagrangian} give the restrictions of $H_\pm$.
\end{proof}

Recall that $\mathrm{(H_J)}$ is the hypothesis of \cite[Section 1]{Arh26} that $S$ is a Haagerup algebra, that is, $S$ is positive, strongly rational and $\Rep(S)$ is braided equivalent to $Z(\cal H_3)$. By (F1), it is equivalent here to the braided equivalence $\Rep(S) \cong Z(\cal H_3)$. The vertex operator algebra $S$ satisfies (I1) and (F1)--(F3) by Theorems \ref{thm:inputs-ArhI} and \ref{thm:zhu-inputs}. In the statements below, these labels only indicate which of these properties are used in the proof.


%%%%%%%%%%%%%%%%%%%%%%%%%%%%%%%%%%%%%%%%%%%%%%%%%%%%%%%%%%%%%%%%%%%%%%%%%%%%%
\subsection{Scope of the algebraic inputs}
\label{subsec:audit-inputs}

Appendix \ref{app:zhu} derives exact relations in a full corner of Zhu's algebra $A(S)$ from products of vertex operators in $V'$ and deduces that this corner, of dimension $762$, is semisimple with twelve simple blocks. Hence $A(S)$ is semisimple of dimension $36625$. The computations are listed in Section \ref{zhu:app:certificates}. The proofs of Theorem \ref{thm:net} and of Proposition \ref{prop:sector-trace} use only (I1), (F1) and (F3). The count of simple modules and their lowest weights in (F2) are used from Lemma \ref{lem:lattice-sectors} on.

%%%%%%%%%%%%%%%%%%%%%%%%%%%%%%%%%%%%%%%%%%%%%%%%%%%%%%%%%%%%%%%%%%%%%%%%%%%%%
\subsection{Additional inputs for the twin identification}
\label{subsec:twin-inputs}

The twin theorem uses the algebraic center description (F4) of Theorem \ref{thm:algebraic-F4} and the additional inputs (I3), (D1) and (D2). The finite identities for (D1) and (D2) are verified independently in Section \ref{subsec:DJX-verification}.

\paragraph{(I3).} $S$ is isomorphic to the vertex operator algebra constructed by Jiao \cite{Jia26}, through an explicit isomorphism which identifies the generators \cite[Theorem 10.24]{Arh26}.

Let $U^{\mathrm R}$ and $U^{\mathrm H}$ be the vertex subalgebras of $V'$ generated by the explicit vectors of Riesen \cite{Rie26} and of Huang \cite{Hua26}, with currents of levels $(1,39)$, and let $V$ be the algebra of Jiao, so that $V \cong S$ by (I3).

\begin{thm}[{\cite{DJX26}}]
\label{thm:inputs-DJX}
\begin{enumerate}
\item[\textup{(D1)}] $U^{\mathrm R}$ and $U^{\mathrm H}$ are isomorphic \cite[Section 4.5 and Proposition 4.1]{DJX26}. We write $U$ for either.
\item[\textup{(D2)}] Let $\tau = \exp(\frac{\pi\i}{2}(3H_L+H_R)_{(0)})$ and $\sigma = \exp(\frac{\pi\i}{2}(P_H+B_H)_{(0)})$, where $H_L$, $H_R$ and $P_H$, $B_H$ are the Cartan currents of $U$ and of $V$. These are involutions of $U$ and of $V$, and there is an isomorphism of vertex operator algebras $\Phi \co U^\tau \to V^\sigma$ induced by an isometry of lattices, which maps Cartan currents to Cartan currents \cite[Theorem 7.3]{DJX26}.
\end{enumerate}
\end{thm}

The other results of \cite{DJX26}, on the strong rationality of $U$, $V$, $U^\tau$ and on their module categories and conformal nets, are discussed in Section \ref{sec:DJX}. We do not use any statement of \cite{Jia26,Rie26,Hua26} beyond the definitions of the explicit vertex operator algebras.


%%%%%%%%%%%%%%%%%%%%%%%%%%%%%%%%%%%%%%%%%%%%%%%%%%%%%%%%%%%%%%%%%%%%%%%%%%%%%
\section{The conformal net of \texorpdfstring{$S$}{S}}
\label{sec:net}

For the net of $S$, the inputs are the unitary conformal embedding and weight-one data (I1), strong rationality (F1) and the branching of the lattice modules (F3). The proof uses \cite[Proposition 5.2]{ArhII} on the lattice summands to apply Tener's finite-index criterion, and \cite[Theorem 6.1]{ArhII} to recognize the fusion categories of the net extensions.


%%%%%%%%%%%%%%%%%%%%%%%%%%%%%%%%%%%%%%%%%%%%%%%%%%%%%%%%%%%%%%%%%%%%%%%%%%%%%
\subsection{Setting}
\label{subsec:net-setting}

We use the explicit model of $V' = V_{E_6 \oplus A_2}$ of \cite[Section 10.1]{Arh26} and the scalar product $H_\chi$ of \cite[Theorem 10.10]{Arh26}, for which $S$ is a unitary subalgebra of $V'$ in the sense of \cite[Definition 5.22]{CKLW18}, containing the conformal vector $\omega$ of $V'$. Let $\Lambda_\pm = \la(1,\pm1)\ra \subset L^\circ/L \cong \Z_3 \times \Z_3$ be the two isotropic lines of the discriminant form of $L = E_6 \oplus A_2$, and let $H_\pm = H_{\Lambda_\pm} \cong V_{E_8}$ be the two lattice extensions of $V'$ \cite[Lemma 4.1]{ArhII}, equipped with the scalar products $\{\cdot|\cdot\}_{\Lambda_\pm}$ of \cite[Lemma 4.2]{ArhII}, which extend $H_\chi$. Let $\cal A_S \subset \cal A_{V'} \subset \cal A_{H_\pm} \cong \cal A_{E_8}$ be the conformal nets of \cite{CKLW18}. By \cite[Theorem 4.5 and Proposition 4.3]{ArhII}, whose proofs make no hypothesis on $S$ beyond (I1), these are covariant subnets with trivial cosets and irreducible inclusions, and $[\cal A_{H_\pm} : \cal A_{V'}] = 3$. For an $S$-submodule $M$ of $H_\pm$, we write $\pi_M$ for the representation of $\cal A_S$ on the closure of $M$ given by \cite[Lemma 4.7]{ArhII}, and we put $\pi_0 = \pi_S$, $\pi_b = \pi_{N_b}$ where $N_b$ is the orthogonal complement of $S$ in $V'$, and $\pi_\pm = \pi_{W_{(1,\pm1)}}$, where $W_{(1,\pm1)} \subset H_\pm$. Let $\Theta_\pm$ be the dual canonical endomorphism of $\cal A_S \subset \cal A_{H_\pm}$, let $\theta_{V'}$ be that of $\cal A_S \subset \cal A_{V'}$, and let $\cal F^{\mathrm{net}}_\pm = \Rep(\cal A_S)_{\Theta_\pm}$ be the category of modules over $\Theta_\pm$ in $\Rep(\cal A_S)$. We label $Z(\cal H_3)$ as in \cite{EvG11}, see \cite[Section 2.2]{ArhII}, and we write $L_2 = 0 \oplus b \oplus 2a$ and $L_3 = 0 \oplus b \oplus 2c_0$ for the Lagrangian algebras of \cite[Proposition 3.3 and Corollary 3.8]{ArhII}, with $Z(\cal H_3)_{L_i} \cong \cal H_i$.


%%%%%%%%%%%%%%%%%%%%%%%%%%%%%%%%%%%%%%%%%%%%%%%%%%%%%%%%%%%%%%%%%%%%%%%%%%%%%
\subsection{Statement}

\begin{thm}[Theorem B]
\label{thm:net}
Assume \textup{(I1)}, \textup{(F1)} and \textup{(F3)}.
\begin{enumerate}
\item The conformal net $\cal A_S$ is completely rational. The inclusions $\cal A_S \subset \cal A_{V'}$ and $\cal A_S \subset \cal A_{H_\pm}$ have finite index. The vacuum representation of $\cal A_{V'}$ restricts to $\cal A_S$ as $\pi_0 \oplus \pi_b$, and that of $\cal A_{H_\pm}$ as $\pi_0 \oplus \pi_b \oplus 2\pi_\pm$, where $\pi_0$, $\pi_b$, $\pi_+$, $\pi_-$ are pairwise inequivalent irreducible representations with statistical dimensions $d(\pi_0) = 1$, $d(\pi_b) = d^2 = 3d+1$ and $d(\pi_\pm) = 3d+2$. Moreover
\[
[\cal A_{V'} : \cal A_S] = 3d+2,
\quad
[\cal A_{H_\pm} : \cal A_S] = 3(3d+2),
\quad
\mu_{\cal A_S} = 9(3d+2)^2.
\]
\item For each sign, $\Theta_\pm$ is a Lagrangian algebra of $\Rep(\cal A_S)$, the category $\cal F^{\mathrm{net}}_\pm$ is a unitary fusion category of rank six with Frobenius--Perron dimension $3(3d+2)$, $\Rep(\cal A_S)$ is braided equivalent to $Z(\cal F^{\mathrm{net}}_\pm)$, and $\cal F^{\mathrm{net}}_\pm$ is equivalent to $\cal H_2$ or to $\cal H_3$. Its pointed part, which corresponds to the intermediate net $\cal A_{V'}$, is $\mathrm{Vec}_{\Z_3}$ with trivial associator.
\item $\Rep(\cal A_S)$ is equivalent to $Z(\cal H_3)$ as a unitary modular tensor category. A braided equivalence $\Psi \co \Rep(\cal A_S) \to Z(\cal H_3)$ is unique up to natural isomorphism, and it satisfies $\Psi(\pi_b) \cong b$ and $\{\Psi(\pi_+),\Psi(\pi_-)\} = \{a,c_0\}$. Writing $\pi^{\mathrm{net}}_a$ and $\pi^{\mathrm{net}}_{c_0}$ for the representations among $\pi_\pm$ sent to $a$ and $c_0$, and $\Theta^{\mathrm{net}}_a$, $\Theta^{\mathrm{net}}_{c_0}$ for the corresponding algebras $\Theta_\pm$, we have $\Psi(\Theta^{\mathrm{net}}_a) \cong L_2$, $\Psi(\Theta^{\mathrm{net}}_{c_0}) \cong L_3$, $\cal F^{\mathrm{net}}_a \cong \cal H_2$ and $\cal F^{\mathrm{net}}_{c_0} \cong \cal H_3$.
\item Up to isomorphism of extensions, $\cal A_S$ has exactly three irreducible local extensions $\cal B$ with $\mu_{\cal B} = 1$. Their dual canonical endomorphisms are $\Theta_1 \cong \pi_0 \oplus \pi_b \oplus \pi_+ \oplus \pi_-$, $\Theta^{\mathrm{net}}_a$ and $\Theta^{\mathrm{net}}_{c_0}$, with $\Rep(\cal A_S)_{\Theta_1} \cong \cal H_1$, and the extensions attached to $\Theta_\pm$ are the nets $\cal A_{H_\pm} \cong \cal A_{E_8}$. The only intermediate nets between $\cal A_S$ and $\cal A_{H_\pm}$ are $\cal A_S$, $\cal A_{V'}$ and $\cal A_{H_\pm}$, and there is no intermediate net strictly between $\cal A_S$ and the extension attached to $\Theta_1$. The Haagerup subfactor arises from $\cal A_S$ in the sense of \cite[Definition 4.1]{Bis17a}. There is a proper action of the hypergroup $K = \{k_0,k_1\}$, $k_1k_1 = d^{-2}k_0+(1-d^{-2})k_1$, on $\cal A_{V'} = \cal A_{E_6,1} \ot \cal A_{SU(3),1}$ with $\cal A_S = \cal A_{V'}^K$, of weight $3d+2$ and parameter $w = d^2$, which solves \cite[Problem 4.27]{Bis17b}. Finally $\Rep(\cal A_S)_{\theta_{V'}}$ is a near-group fusion category of type $\Z_3^2+9$.
\end{enumerate}
\end{thm}

The net argument alone does not identify the third holomorphic extension with $\cal A_{E_8}$. That identification follows from Theorem \ref{thm:HJ-main}, see Corollary \ref{cor:S-consequences}(2). The labels $a$ and $c_0$ of part 3 are defined on the net side and are not compared with the labels of the simple $S$-modules $X_\pm$ of (F3), see Remark \ref{rem:net-labels}.


%%%%%%%%%%%%%%%%%%%%%%%%%%%%%%%%%%%%%%%%%%%%%%%%%%%%%%%%%%%%%%%%%%%%%%%%%%%%%
\subsection{Proof of Theorem \ref{thm:net}}

The proof follows \cite[Sections 5--7]{ArhII}. We indicate, for each step, which statement of \cite{ArhII} is used, where its proof uses the Haagerup hypothesis, and by what this hypothesis is replaced. Two conventions of Tener are used: a VOA with bounded localized vertex operators, in the sense of \cite[Definition 4.12 p.~42]{Ten19}, and, for a unitary module $M$ of such a VOA, the representation $\pi^M$ of \cite[Definition 5.1 p.~49]{Ten19} and \cite[Definition 2.18 p.~20]{Ten24}, when it exists.

\emph{Step 0: inputs from \cite[Section 4]{ArhII}.} Let $H$ be one of $H_\pm$. By \cite[Proposition 4.3(1)]{ArhII}, $H$ is strongly local, has bounded localized vertex operators, its nets in the sense of \cite{Ten19} and of \cite{CKLW18} are isomorphic, and $\cal A_H \cong \cal A_{E_8}$ is completely rational with $\mu_{\cal A_H} = 1$ \cite[Theorem 7.4 p.~71, Proposition 7.5 and Corollary 7.7 p.~72]{Ten19}, \cite[Corollary 3.19]{DoX06}. By (I1) and the proof of \cite[Theorem 4.5]{ArhII}, $S$ is a unitary subalgebra of $H$ in the sense of \cite[Definition 2.3 p.~12]{Ten19}, since the PCT operator of $H$ restricts to that of $V'$, which preserves $S$, and since $L(1)S \subset S$, and $S$ is a conformal subalgebra with trivial coset \cite[Proposition 2.7]{Ten19}, \cite[Remark 5.30]{CKLW18}. By \cite[Proposition 4.16 p.~43]{Ten19}, $S$ has bounded localized vertex operators. By (F1) and \cite[Theorem 4.5 p.~9]{ABD04}, $S$ is regular. By regularity and the finite-dimensional homogeneous spaces of $H$ for the common conformal vector, $H$ is a finite direct sum of simple $S$-modules, which we choose orthogonal for its invariant scalar product.

\emph{Step 1: the representations $\pi^M$.} Let $M$ be an $S$-submodule of $H$, with the restricted scalar product. By \cite[Theorem 5.6 p.~52]{Ten19}, applied to $V = H$, to the unitary subalgebra $S$ and to the unitary $H$-module $H$ itself, whose representation $\pi^H$ is the vacuum representation, the representation $\pi^M$ exists. Since the PCT operator $\theta_H$ preserves $S$, the subspace $\theta_H(M)$ is an $S$-submodule of $H$, unitarily isomorphic to the conjugate module $\ovl M$, so that $\pi^{\ovl M}$ exists.

\emph{Step 2: the transport module for $(\ovl M,M)$.} Let $M$ be a simple $S$-summand of $H$. By Step~0 and Step~1, $S$ is a simple regular unitary vertex operator algebra with bounded localized vertex operators, and the unitary $S$-modules $\ovl M$ and $M$ satisfy the hypotheses of \cite[Proposition 5.2]{ArhII}, with $\ovl M$ and $M$ in the roles of $M$ and $N$ there, since $\pi^M$ exists. Hence $\ovl M \boxtimes M$ admits an invariant scalar product for which $(\ovl M \boxtimes M,\cal Y_\boxtimes)$ is a transport module for $\ovl M$ and $M$ in the sense of \cite[Definition 3.3 p.~26]{Ten24}. Since the invariant scalar product of $M$ identifies $\ovl M$ with the contragredient module $M'$, rigidity gives $\operatorname{Hom}_S(S,\ovl M \boxtimes M) \cong \operatorname{End}_S(M) \neq 0$, so that $\ovl M \boxtimes M$ contains the vacuum module $S$. This construction uses neither the modules $M_{d_l}$ nor $\mathrm{(H_J)}$.

\emph{Step 3: finite index and complete rationality.} By skew symmetry in $H$, the projection $p_MY^H(m,z)s$ onto $M$, for $m\in M$ and $s\in S$, is $\cal Y_M^+(m,z)s$. The common conformal vector makes the power shift in \cite[Theorem 5.15]{Ten19} zero, so $\cal Y_M^+\in\cal I_{\mathrm{loc}}$. Steps~1 and~2 provide the representations $\pi^M$ and $\pi^{\ovl M}$ and the transport module $\ovl M \boxtimes M$, which contains $S$. Thus \cite[Theorem 4.19 p.~55]{Ten24} applies directly and gives finite index for each $\pi^M$.

Write $\pi_H$ for the restriction of the vacuum of $\cal A_H^{\mathrm T}$ to $\cal A_S^{\mathrm T}$, with the superscript denoting the nets in the sense of \cite{Ten19}. By the finite decomposition in Step~0, $\pi_H$ is a finite sum of the sectors just obtained and has finite index. For a complementary interval $I'$, locality gives
\[
\pi_H(\cal A_S^{\mathrm T}(I))
 \subset \cal A_H^{\mathrm T}(I)
 \subset \pi_H(\cal A_S^{\mathrm T}(I'))'.
\]
The exterior inclusion of factors has finite index by the index of $\pi_H$. A finite-index conditional expectation from $\pi_H(\cal A_S^{\mathrm T}(I'))'$ onto $\pi_H(\cal A_S^{\mathrm T}(I))$ restricts to a conditional expectation from $\cal A_H^{\mathrm T}(I)$ onto the same lower factor. The restriction has finite index by \cite[Remark 3.13(i) p.~227]{BDH88}. This gives finite index for the subnet inclusion, as in the last step of the proof of \cite[Corollary 4.24]{Ten24}. Since the net of $H$ is completely rational, \cite[Theorem 24 p.~20]{Lon03} gives complete rationality of $\cal A_S^{\mathrm T}$, and \cite[Lemma 7.6 p.~72]{Ten19} identifies it with $\cal A_S$. By \cite[Theorem 4.5]{ArhII}, $\cal A_{V'}$ and $\cal A_H$ are irreducible local extensions of $\cal A_S$, so $[\cal A_{V'}:\cal A_S]$ and $[\cal A_H:\cal A_S]$ are finite by \cite[Proposition 2.5(2)]{ArhII}. The specific branching in (F3) enters in Step~4.

\emph{Step 4: decompositions.} By (F3) and the orthogonality argument of \cite[proof of Theorem 5.6(2)]{ArhII}, the four indicated $S$-summands of $H$ form an orthogonal direct sum. By \cite[Lemma 4.7(1)]{ArhII}, the vacuum representation of $\cal A_{H_\pm}$ restricts to $\cal A_S$ as $\pi_S \oplus \pi_{N_b} \oplus \pi_{W_{(1,\pm1)}} \oplus \pi_{W_{(-1,\mp1)}}$, and that of $\cal A_{V'}$ as $\pi_S \oplus \pi_{N_b}$. By (F3) and \cite[Lemma 4.7(2),(3)]{ArhII}, these representations are irreducible, $\pi_{W_{(-1,\mp1)}} \cong \pi_{W_{(1,\pm1)}} = \pi_\pm$, and $\pi_0$, $\pi_b$, $\pi_\pm$ are pairwise inequivalent. Applying \cite[Lemma 4.7(3)]{ArhII} to $H_+$ and $H_-$, which contain $S$ as a unitary subalgebra with the same conformal vector and with scalar products both restricting to $H_\chi$, and using $X_+ \not\cong X_-$ from (F3), we obtain $\pi_+ \not\cong \pi_-$. This is \cite[Theorem 5.6(2)]{ArhII}, with (F3) in place of \cite[Lemma 5.4]{ArhII}, which was the only use of $\mathrm{(H_J)}$ there.

\emph{Step 5: indices.} Put $\delta = d(\pi_b)$. By \cite[Proposition 2.5(1),(3)]{ArhII} and Step 4, $\theta_{V'} = \pi_0 \oplus \pi_b$, so that $[\cal A_{V'} : \cal A_S] = 1+\delta$ and $\mu_{\cal A_S} = (1+\delta)^2\mu_{\cal A_{V'}} = 9(1+\delta)^2 = [\cal A_{H_\pm} : \cal A_S]^2$ \cite[Proposition 4.3(2)]{ArhII}, \cite[Proposition 24]{KLM01}. Hence $[\cal A_{H_\pm} : \cal A_S] = d(\Theta_\pm) = 1+\delta+2d(\pi_\pm) = 3(1+\delta)$ and $d(\pi_\pm) = 1+\delta$. This is \cite[Theorem 5.6(3)]{ArhII}, whose proof is unchanged.

\emph{Step 6: the Lagrangian algebras $\Theta_\pm$.} By \cite[Proposition 2.5(3)]{ArhII} and the irreducibility of $\cal A_S \subset \cal A_{H_\pm}$, $\Theta_\pm$ is a commutative $Q$-system with $\dim\operatorname{Hom}(\mathrm{id},\Theta_\pm) = 1$, hence a connected étale algebra in the unitary modular tensor category $\Rep(\cal A_S)$ \cite[Theorem 33 and Corollary 39]{KLM01}. Since $\operatorname{FPdim}\Rep(\cal A_S) = \mu_{\cal A_S} = d(\Theta_\pm)^2$ \cite[Theorem 33]{KLM01}, $\Theta_\pm$ is Lagrangian \cite[Definition 4.6]{DMNO13}, and $\Rep(\cal A_S) \cong Z(\cal F^{\mathrm{net}}_\pm)$ as braided fusion categories \cite[Corollary 4.1]{DMNO13}, with $\operatorname{FPdim}\cal F^{\mathrm{net}}_\pm = 3(1+\delta)$ \cite[Lemma 3.11]{DMNO13}. The category $\cal F^{\mathrm{net}}_\pm$ is the category of modules over the $Q$-system $\Theta_\pm$ in the $C^*$-tensor category $\Rep(\cal A_S)$, equivalently the category of $\alpha$-induced sectors of $\cal A_S(I) \subset \cal A_{H_\pm}(I)$, a unitary fusion category \cite[Section 4.6]{BKLR15}, \cite[Proposition 4.5]{Bis17a}. Its rank is $\dim\operatorname{End}(\Theta_\pm) = 1+1+4 = 6$ by \cite[Lemma 3.1]{ArhII} and Step 4. The intermediate net $\cal A_{V'}$ gives a connected étale subalgebra $\theta_{V'} \subset \Theta_\pm$ \cite[Proposition 4.11]{BKLR15}, which corresponds by \cite[Lemma 3.1]{ArhII} and \cite[Theorem 4.10]{DMNO13} to a fusion subcategory $\cal P_\pm$ of $\cal F^{\mathrm{net}}_\pm$ with $\operatorname{FPdim}\cal P_\pm = 3$, which is pointed \cite[Corollary 8.30]{ENO05}. This is \cite[Theorem 5.6(4)]{ArhII}, whose proof is unchanged.

\emph{Step 7: $\delta > 1$ and the three invertible objects.} Suppose that $\delta = 1$. By \cite[Proposition 5.7]{ArhII}, which uses only Steps 3--5 and \cite[Theorem 4.22 and Section 4.6]{Bis17b}, we have $\cal A_S = \cal A_{V'}^K$ for a hypergroup $K = \{k_0,k_1\}$ with $k_1k_1 = k_0$, and the argument of \cite[Proposition 5.8]{ArhII}, based on the Kadison--Schwarz inequality, \cite[Proposition 4.21]{Bis17b} and \cite[Theorem 6.9, Theorem 7.5 and Proposition 7.6]{CKLW18}, produces a unitary automorphism $g$ of $V'$ of order two with $S = (V')^g$. Then $S_1 = (V'_1)^g$ is the fixed point subalgebra of a nontrivial involution of $\mathfrak e_6 \oplus \mathfrak{sl}_3$ preserving both ideals, so that $\dim S_1 \geq 39$ by the classification of involutions of simple Lie algebras \cite[Chapter X]{Hel78}, which contradicts $\dim S_1 = 6$. Hence $\delta > 1$. Since $\cal F^{\mathrm{net}}_\pm$ has rank six and contains the pointed subcategory $\cal P_\pm$ of order three, its group of invertible objects has order $3$ or $6$, and order $6$ would make $\cal F^{\mathrm{net}}_\pm$ pointed with $\operatorname{FPdim} = 6 = 3(1+\delta)$, that is, $\delta = 1$. Hence $\cal F^{\mathrm{net}}_\pm$ has exactly three invertible simple objects, which generate $\cal P_\pm$. This is \cite[Proposition 5.8]{ArhII}, with $\dim S_1 = 6$ in place of the weight-one classification of \cite{Arh26} for Haagerup algebras.

\emph{Step 8: trivial associator of $\cal P_\pm$.} The proof of \cite[Lemma 6.10]{ArhII} applies verbatim: it uses Step 6, \cite[Corollary 4.1 and Theorem 4.10]{DMNO13}, \cite[Proposition 2.2(5)]{ArhII}, \cite[Proposition 6.4]{BKL15}, the two Lagrangian algebras of $\Rep(\cal A_{V'})$ given by the lattice nets $\cal A_{H_\pm}$ \cite[Proposition 4.3(2) and Theorem 4.5]{ArhII}, which are not isomorphic because their restrictions to $\cal A_S$ are $\pi_0 \oplus \pi_b \oplus 2\pi_+$ and $\pi_0 \oplus \pi_b \oplus 2\pi_-$ with $\pi_+ \not\cong \pi_-$ by Step 4, \cite[Proposition 4.8]{DMNO13} and \cite[Lemma 6.7(2)]{ArhII}. Hence $\cal P_\pm \cong \mathrm{Vec}_{\Z_3}$ with trivial associator.

\emph{Step 9: recognition.} By Steps 6 and 7, $\cal F^{\mathrm{net}}_\pm$ is a unitary fusion category with six simple objects, three of which are invertible. By Step 4, $\operatorname{End}(\Theta_\pm) \cong \Cbb \oplus \Cbb \oplus \M_2(\Cbb)$, and since $\Theta_\pm$ corresponds to the induction $I(1)$ \cite[Lemma 3.1]{ArhII} and $K(\cal F^{\mathrm{net}}_\pm) \ot \Cbb \cong \operatorname{End}(I(1))$ \cite[Corollary 2.16]{Ost15}, the Grothendieck ring of $\cal F^{\mathrm{net}}_\pm$ is noncommutative. By Step 8 and \cite[Theorem 6.1]{ArhII}, a categorical statement whose part 1 uses unitarity only through the spherical structure with positive dimensions and whose part 2 uses the $C^*$-structure through \cite[Theorems 5.9 and 9.1]{Izu18}, the category $\cal F^{\mathrm{net}}_\pm$ is equivalent to $\cal H_2$ or to $\cal H_3$. Comparing $\operatorname{FPdim}\cal F^{\mathrm{net}}_\pm = 3(1+\delta)$ with $\dim\cal H_i = 3(3d+2)$ gives $\delta = 3d+1 = d^2$, which completes the proof of parts 1 and 2. By Step 6 and \cite[Theorem 3.1]{ENO11}, $\Rep(\cal A_S) \cong Z(\cal F^{\mathrm{net}}_\pm) \cong Z(\cal H_3)$ as braided fusion categories, and such an equivalence is naturally isomorphic to a unitary braided equivalence \cite[Theorem 5]{Reu23}. This recovers the conclusion of \cite[Theorem 6.11]{ArhII} for $S$ from (I1), (F1) and (F3).

\emph{Step 10: labels.} Let $\Psi \co \Rep(\cal A_S) \to Z(\cal H_3)$ be a braided equivalence. It is unique up to natural isomorphism, since the Brauer--Picard group of $\cal H_3$ is trivial \cite[Theorem 1.5]{GrS12}, \cite[Theorem 1.1]{ENO10}. The first part of the proof of \cite[Lemma 7.1]{ArhII} applies: $d(\pi_b) = 3d+1$ forces $\Psi(\pi_b) \cong b$ by the dimensions of $Z(\cal H_3)$, $\Psi(\Theta_\pm) \cong 0 \oplus b \oplus 2\Psi(\pi_\pm)$ is a Lagrangian algebra of $Z(\cal H_3)$, hence $\Psi(\pi_\pm) \in \{a,c_0\}$ by \cite[Lemma 3.2]{ArhII}, and $\Psi(\pi_+) \not\cong \Psi(\pi_-)$ by Step 4. By \cite[Proposition 3.3]{ArhII} and the first two paragraphs of the proof of \cite[Proposition 3.6]{ArhII}, which are categorical, $L_2$ and $L_3$ are the only Lagrangian algebras of $Z(\cal H_3)$ with these objects, and $Z(\cal H_3)_{L_i} \cong \cal H_i$. This proves part 3.

\emph{Step 11: consequences.} The first assertions of part 4 are \cite[Theorem 7.2(2),(3)]{ArhII}, with Step 10 in place of \cite[Lemma 7.1]{ArhII} and without \cite[Proposition 3.5]{ArhII}, which is not needed: the irreducible local extensions of $\cal A_S$ with $\mu = 1$ correspond to the Lagrangian algebras of $\Rep(\cal A_S)$ \cite[Proposition 2.5(1),(2)]{ArhII}, \cite[Theorem 4.9]{LoR95}, \cite[Theorems 3.2 and 3.9]{CGGH23}, of which there are exactly three by Step 10 and \cite[Proposition 3.3]{ArhII}, with the stated objects. The algebras $\Theta_\pm$ are two of them by Step 6, their extensions are $\cal A_{H_\pm} \cong \cal A_{E_8}$ \cite[Theorem 4.5]{ArhII}, and the intermediate nets correspond to the fusion subcategories of $\cal H_i$ through \cite[Proposition 4.11]{BKLR15}, \cite[Lemma 3.1]{ArhII} and \cite[Theorem 4.10]{DMNO13}. The assertion on the Haagerup subfactor is \cite[Corollary 7.4]{ArhII}: $\Rep(\cal A_S) \cong Z(\cal H_3)$ is the quantum double of the Haagerup subfactor and \cite[Proposition 4.7(1)]{Bis17a} applies. The hypergroup assertion is \cite[Proposition 7.5 and Corollary 7.6]{ArhII} with $\delta = d^2$ from Step 9 and \cite[Proposition 4.3(2)]{ArhII} for $\cal A_{V'} \cong \cal A_{E_6,1} \ot \cal A_{SU(3),1}$. The last assertion is \cite[Remark 7.8]{ArhII}, a computation in $Z(\cal H_3)$ with the étale algebra $0 \oplus b$, which uses Step 9 and the fusion rules of $Z(\cal H_3)$.
\qed

\begin{remark}\normalfont
\label{rem:third-extension}
The extension $\cal B_1$ attached to $\Theta_1$ is holomorphic of central charge eight. Theorem \ref{thm:net} identifies its module category over the inclusion, but does not itself classify the vacuum net. Theorem \ref{thm:HJ-main} subsequently supplies complete unitarity and the VOA--net comparison. Then \cite[Proposition 7.10]{ArhII} identifies $\cal B_1$ with $\cal A_{E_8}$, as in Corollary \ref{cor:S-consequences}(2). This uses the classification of the corresponding holomorphic VOA, not an assumed classification of all holomorphic nets at charge eight.
\end{remark}

\begin{remark}\normalfont
\label{rem:net-labels}
The labels $a$ and $c_0$ of part 3 are intrinsic to the net through $\Psi$. The signs in $X_\pm$ refer to the lattice coordinates of (F3). Corollary \ref{cor:S-consequences}(1) compares the intrinsic labels of the modules and sectors.
\end{remark}


%%%%%%%%%%%%%%%%%%%%%%%%%%%%%%%%%%%%%%%%%%%%%%%%%%%%%%%%%%%%%%%%%%%%%%%%%%%%%
\section{The diagonal holomorphic extension from the conformal net}
\label{sec:diagonal}

Theorem \ref{thm:net} identifies the representation category of $\cal A_S$ with $Z(\cal H_3)$ without assuming $\mathrm{(H_J)}$. We construct the diagonal holomorphic extension on this net side, prove the trace class property of every sector using \cite{HeT26}, and reconstruct a unitary holomorphic vertex operator algebra containing $S\ot S$. The compatible unitary structure comes from the inclusion of nets. This supplies an analytic alternative to the lattice computation of \cite[Section 7.5]{ArhII}.

We use the notation of Section \ref{sec:net}. The twelve irreducible representations of $\cal A_S$ are labelled $\pi_z$, $z \in \{0,b,a,c_0,c_1,c_2,d_1,\ldots,d_6\}$, through $\Psi(\pi_z) \cong z$, so that $\pi_a = \pi^{\mathrm{net}}_a$ and $\pi_{c_0} = \pi^{\mathrm{net}}_{c_0}$. We write $\cal A_2 \ov{\mathrm{def}}{=} \cal A_S \ot \cal A_S$ for the tensor product net on $\cal H_S \ot \cal H_S$, which is the net of the strongly local vertex operator algebra $(S \ot S, H_\chi \ot H_\chi)$ \cite[Corollary 8.2 p.~62]{CKLW18}. For a representation $(K,\pi)$ of a conformal net, the local unitaries $\pi_I(U(\gamma))$, $\gamma \in \mathrm{Diff}(I)$, define a projective representation of the universal cover of $\mathrm{Diff}(S^1)$ on $K$ \cite[Section 2 p.~7--8]{HeT25}, \cite[Theorem 11]{Hen19}, whose generator of rotations $L_0$ is positive \cite[Proposition 3.8]{Wei06}, and $K^{\mathrm{f.e.}}$ denotes the algebraic direct sum of the eigenspaces of $L_0$.


%%%%%%%%%%%%%%%%%%%%%%%%%%%%%%%%%%%%%%%%%%%%%%%%%%%%%%%%%%%%%%%%%%%%%%%%%%%%%
\subsection{The braid-reversing equivalence of \texorpdfstring{$Z(\cal H_3)$}{Z(H3)} and the Lagrangian algebra \texorpdfstring{$\Theta_\Delta$}{Theta}}
\label{subsec:lagrangian-diagonal}

\begin{lemma}
\label{lem:sigma}
Assume \textup{(I1)}, \textup{(F1)}, \textup{(F3)}. There is a braided equivalence $\sigma$ from $\Rep(\cal A_S)^{\mathrm{rev}}$ to $\Rep(\cal A_S)$, unique up to natural isomorphism, with $\sigma \circ \sigma^{\mathrm{rev}} \cong \mathrm{id}$. It fixes $\pi_0$, $\pi_b$, $\pi_a$, $\pi_{c_0}$, exchanges $\pi_{c_1}$ and $\pi_{c_2}$, and maps $\pi_{d_l}$ to $\pi_{d_{l'}}$ with $l' \equiv \pm 5l \bmod 13$, that is, $d_1 \leftrightarrow d_5$, $d_2 \leftrightarrow d_3$, $d_4 \leftrightarrow d_6$.
\end{lemma}

\begin{proof}
The first half of the proof of \cite[Proposition 7.12(1)]{ArhII} concerns $Z(\cal H_3)$ alone: for a unitary fusion category $\cal C$, inverting half-braidings gives $Z(\cal C^{\ot\mathrm{op}}) \cong Z(\cal C)^{\mathrm{rev}}$ and $f \mapsto f^*$ composed with duality gives $\cal C^{\ot\mathrm{op}} \cong \ovl{\cal C}$, while the even parts of the Haagerup subfactor satisfy $\ovl{\cal H_i} \cong \cal H_i$ \cite[Theorem 4.1 p.~21]{JMS14}, since the conjugate planar algebra has the same principal graphs. With $Z(\cal H_1) \cong Z(\cal H_3)$ \cite[Theorem 3.1 p.~9]{ENO11}, this gives a braided equivalence $R \co Z(\cal H_3)^{\mathrm{rev}} \to Z(\cal H_3)$, and $\sigma = \Psi^{-1} \circ R \circ \Psi^{\mathrm{rev}}$ with $\Psi$ from Theorem \ref{thm:net}. Uniqueness and $\sigma \circ \sigma^{\mathrm{rev}} \cong \mathrm{id}$ follow from the triviality of the Brauer--Picard group of $\cal H_3$ \cite[Theorem 1.5 p.~5]{GrS12}, \cite[Theorem 1.1 p.~3]{ENO10}. The action on the labels is computed in the second half of the proof of \cite[Proposition 7.12(1)]{ArhII}, which uses only the dimensions, twists and self-duality of the simple objects of $Z(\cal H_3)$ and its Lagrangian algebras \cite[Proposition 3.3]{ArhII}.
\end{proof}



\begin{lemma}
\label{lem:tensor-nets}
Let $\cal A$ be a completely rational conformal net. Then $\cal A \ot \cal A$ is completely rational, $\mu_{\cal A \ot \cal A} = \mu_{\cal A}^2$, and $\rho \boxtimes \rho' \mapsto \rho \ot \rho'$, $t \boxtimes t' \mapsto t \ot t'$, is a braided equivalence of unitary modular tensor categories $\Rep(\cal A) \boxtimes \Rep(\cal A) \to \Rep(\cal A \ot \cal A)$.
\end{lemma}

\begin{proof}
Complete rationality is \cite[Lemma 25 p.~21]{Lon03}. For irreducible localized endomorphisms $\rho_i$, $\rho'_i$, the endomorphism $\rho_1 \ot \rho_2$ is irreducible, since $(\rho_1 \ot \rho_2)(\cal A \ot \cal A)'' = \B(\cal H_{\rho_1}) \otvn \B(\cal H_{\rho_2})$, and $\rho_1 \ot \rho_2 \cong \rho'_1 \ot \rho'_2$ forces $\rho_i \cong \rho'_i$, by restriction to both $\cal A\ot 1$ and $1\ot\cal A$: each restriction is an amplification of one irreducible factor, and inequivalent irreducible factor representations are disjoint. Every factor representation of $\cal A$ is of type I \cite[Corollary 14 p.~14]{KLM01}, so every irreducible localized endomorphism of $\cal A \ot \cal A$ is of the form $\rho_1 \ot \rho_2$ \cite[Lemma 27 p.~22]{KLM01}, and every finite-index representation is a finite sum of irreducibles \cite[Corollary 39 p.~27]{KLM01}. The functor is thus an equivalence of semisimple categories, strictly monoidal, and braided, since the DHR braiding of $\rho_1 \ot \rho_2$ and $\rho'_1 \ot \rho'_2$ is computed with charge transporters $u \ot u'$ and equals $\epsi(\rho_1,\rho'_1) \ot \epsi(\rho_2,\rho'_2)$. Finally $\mu_{\cal A \ot \cal A} = \sum d(\rho_1)^2d(\rho_2)^2 = \mu_{\cal A}^2$ \cite[Theorem 33 p.~24]{KLM01}.
\end{proof}

\begin{lemma}
\label{lem:theta-delta}
Let $\cal C$ be a non-degenerate braided fusion category and $\sigma \co \cal C^{\mathrm{rev}} \to \cal C$ a braided equivalence. Then $\Theta_\Delta \ov{\mathrm{def}}{=} \bigoplus_{z \in \Irr(\cal C)} z \boxtimes \sigma(z^*)$ carries the structure of a connected étale Lagrangian algebra of $\cal C \boxtimes \cal C$. For $\cal C = \Rep(\cal A_S)$ and $\sigma$ as in Lemma \ref{lem:sigma}, the object $\Theta_\Delta \cong \bigoplus_z \pi_z \boxtimes \sigma(\pi_z)$ of $\Rep(\cal A_S) \boxtimes \Rep(\cal A_S) \cong \Rep(\cal A_2)$ is a connected étale Lagrangian algebra with twelve pairwise non-isomorphic simple summands and $\operatorname{FPdim}\Theta_\Delta = \mu_{\cal A_S} = 9(3d+2)^2$.
\end{lemma}

\begin{proof}
Let $G \co \cal C \boxtimes \cal C^{\mathrm{rev}} \to Z(\cal C)$ be the braided equivalence of \cite[(11) p.~5]{DMNO13} and $F$ its composition with the forgetful functor, a central functor with $F(X \boxtimes Y) \cong X \ot Y$. Its right adjoint $I$ makes $I(1)$ a connected étale algebra \cite[Lemma 3.5 p.~8]{DMNO13}, with $\operatorname{Hom}(x \boxtimes y,I(1)) \cong \operatorname{Hom}_{\cal C}(x,y^*)$ for simple $x$, $y$, so $I(1) \cong \bigoplus_z z \boxtimes z^*$ and $\operatorname{FPdim} I(1) = \operatorname{FPdim}\cal C = (\operatorname{FPdim}(\cal C \boxtimes \cal C^{\mathrm{rev}}))^{1/2}$ \cite[(8) p.~5]{DMNO13}: $I(1)$ is Lagrangian \cite[Definition 4.6 p.~17]{DMNO13}. The braided equivalence $\mathrm{id} \boxtimes \sigma$ preserves connectedness, commutativity, separability \cite[Remark 2.6(ii) p.~5]{DMNO13} and dimensions, and maps $I(1)$ to an algebra with underlying object $\Theta_\Delta$. For $\Rep(\cal A_S)$, the simple objects of $Z(\cal H_3)$ are self-dual \cite[Section 2.2]{ArhII}, $z \mapsto (z,\sigma(z))$ is injective, Lemma \ref{lem:tensor-nets} applies, and $\operatorname{FPdim}\Theta_\Delta = \sum_z d(\pi_z)^2 = \mu_{\cal A_S}$ \cite[Theorem 33 p.~24]{KLM01}.
\end{proof}


%%%%%%%%%%%%%%%%%%%%%%%%%%%%%%%%%%%%%%%%%%%%%%%%%%%%%%%%%%%%%%%%%%%%%%%%%%%%%
\subsection{The holomorphic net \texorpdfstring{$\cal B$}{B}}
\label{subsec:holomorphic-net}


\begin{prop}
\label{prop:net-B}
Assume \textup{(I1)}, \textup{(F1)} and \textup{(F3)}. The diagonal algebra of Lemma \ref{lem:theta-delta}, with its specified multiplication, determines an irreducible local extension $\cal B$ of $\cal A_2$, unique up to isomorphism for this algebra. Its vacuum representation restricts as
\begin{equation}
\label{eq:HB-decomposition}
\cal H_{\cal B} = \bigoplus_z \cal H_{\pi_z} \ot \cal H_{\sigma(\pi_z)}.
\end{equation}
The extension is completely rational, $[\cal B:\cal A_2] = \mu_{\cal A_S}$, $\cal A_2(I)' \cap \cal B(I) = \Cbb$ and $\mu_{\cal B} = 1$. It is diffeomorphism covariant of central charge $16$, with local implementing unitaries in the represented subnet $\cal A_2(I)$. Its conformal Hamiltonian satisfies
\begin{equation}
\label{eq:diagonal-Hamiltonian}
L_0^{\cal B}|_{\cal H_{\pi_z}\ot\cal H_{\sigma(\pi_z)}} = \overline{L_0^{\pi_z}\ot 1+1\ot L_0^{\sigma(\pi_z)}}.
\end{equation}
The closure in this formula is the nonnegative self-adjoint tensor sum. No discreteness or finite multiplicity assumption on the sector Hamiltonians is used.
\end{prop}

\begin{proof}
Transport the canonical diagonal algebra $I(1)$ by $\mathrm{id}\boxtimes\sigma$ as in Lemma \ref{lem:theta-delta}. Its algebra structure, not merely its underlying object, is part of this construction. By \cite[Theorems 3.2 and 3.9]{CGGH23}, it is isomorphic to a normalized commutative haploid $Q$-system in $\Rep(\cal A_2)$. Its twist is the identity: the canonical positive spherical structures are preserved by $\sigma$, and the twists of $\pi_z$ and $\sigma(\pi_z)$ are inverse. Equivalently one can use \cite[Proposition 2.22]{CGGH23}.

The extension theorem of \cite[Theorems 2.6, p.~11, and 5.7, p.~26]{GuiF25}, in the Connes fusion formulation of \cite[Theorem 4.9]{LoR95}, now applies. Here the algebra of Lemma \ref{lem:theta-delta}, defined in the DHR category $\Rep(\cal A_2)$, is transported to the category of $\cal A_2$-modules with Connes fusion through the equivalence of \cite[Theorem 6.6 p.~77]{Gui21a}, which is braided or braid-reversing (Remark \ref{rem:braiding-conventions}), and commutativity and haploidy are preserved in both cases. The smaller net $\cal A_2$ is an irreducible conformal net, the $Q$-system is commutative and haploid and its twist is trivial. The result gives an irreducible local conformal extension on the representation space of the $Q$-system, with vacuum vector given by its unit. Its restriction is \eqref{eq:HB-decomposition}. The vacuum is cyclic, and the extension has a genuine positive-energy M\"obius representation and a projective diffeomorphism representation. The index is the dimension of the $Q$-system, namely $\mu_{\cal A_S}$. By \cite[Theorem 24]{Lon03} and \cite[Proposition 24]{KLM01}, the extension is completely rational and $\mu_{\cal B} = \mu_{\cal A_2}/[\cal B:\cal A_2]^2 = 1$.

In \cite[Theorem 2.6]{GuiF25}, covariance of the extension is implemented by the canonical covariance on the $\cal A_2$-module of the $Q$-system. In particular its local diffeomorphism implementers are $\pi_I(U_{\cal A_2}(\gamma))$ up to phase for $\gamma\in\mathrm{Diff}(I)$. These implementers give the same projective representation of the covering diffeomorphism group, since the local diffeomorphisms generate it projectively. On the universal covering M\"obius group the genuine lifts coincide. Indeed two such lifts of the same projective action differ by a continuous character, and the connected covering group has perfect Lie algebra. Thus the rotation generator has the same normalization as in $\pi$. The local stress tensor is the represented stress tensor of $S\ot S$, so the Virasoro action has central charge $16$ and is the same in the extension and in $\pi$. For a diffeomorphism fixing an interval pointwise, the implementer on the vacuum subnet belongs projectively to the algebra of the complementary interval by locality and Haag duality, in the closed-interval convention of \cite[Section 2 p.~7]{HeT25}. Its image under $\pi$ commutes with the extension algebra of the original interval by locality. This gives the local triviality in \cite[Definition 2.3]{HeT26}. The local conformal extension in \cite[Theorem 5.7]{GuiF25} has a cyclic vacuum and a unique M\"obius invariant vector.

Each summand $\cal H_{\pi_z} \ot \cal H_{\sigma(\pi_z)}$ is invariant under the M\"obius representation of $\pi$, which commutes with $\pi(\cal A_2)'$. On this summand, both this representation and the tensor product of the M\"obius representations of $\pi_z$ and $\sigma(\pi_z)$ are continuous unitary representations of the universal covering group of $\mathrm{PSL}(2,\R)$ implementing the M\"obius covariance of the irreducible representation $\pi_z \ot \sigma(\pi_z)$ of $\cal A_2$. Hence they differ by a continuous character of this group, which is trivial since its Lie algebra is perfect, as in the previous paragraph. Thus the two representations coincide. Stone's theorem and the joint spectral calculus for the two commuting one-parameter groups of rotations then give \eqref{eq:diagonal-Hamiltonian} and show that the summand projections reduce $L_0^{\cal B}$. Moreover $\e^{2\pi\i L_0^{\cal B}} = 1$. Indeed, by \eqref{eq:diagonal-Hamiltonian}, on each summand this operator is the product of the scalars $\e^{2\pi\i L_0^{\pi_z}}$ and $\e^{2\pi\i L_0^{\sigma(\pi_z)}}$, which are, for a fixed orientation of the DHR braiding, the twists of $\pi_z$ and of $\sigma(\pi_z)$ by the conformal spin and statistics theorem \cite{GuL96}, and these twists are inverse to each other. Hence the projective representation of the universal cover of $\mathrm{Diff}(S^1)$ on $\cal H_{\cal B}$ is a projective representation of $\mathrm{Diff}(S^1)$.
\end{proof}


\begin{remark}\normalfont
\label{rem:diff-covariance}
Diffeomorphism covariance of $\cal B$ can also be obtained from \cite[Proposition 3.7 p.~19]{Car04}, which applies to every local extension of a Virasoro net in the sense of \cite[Definition 3.1 p.~14]{Car04}, that is, to every irreducible local M\"obius covariant net containing a Virasoro subnet such that the von Neumann algebra generated by its local algebras contains the M\"obius unitaries, and requires neither finite index nor irreducibility of the inclusion. A third proof uses the structure $\cal B(I) = \cal A_2(I)v$ given by a $Q$-system localized in a subinterval \cite[Lemma 15 p.~14]{Lon03} and the localization of charged intertwiners. The same proposition is used in \cite[proof of Proposition 4.8 p.~28]{RTT22} for the nets of AQFT-local vertex operator algebras.
\end{remark}


%%%%%%%%%%%%%%%%%%%%%%%%%%%%%%%%%%%%%%%%%%%%%%%%%%%%%%%%%%%%%%%%%%%%%%%%%%%%%
\subsection{The trace class theorem and the sector Hamiltonians}
\label{subsec:hypothesis-D}

The definition of a conformal net in \cite[Definition 2.3, pp.~7--8]{HeT26} requires isotony, locality, a strongly continuous projective diffeomorphism representation extending a genuine M\"obius representation, local triviality of diffeomorphisms, nonnegative energy, a cyclic vacuum and uniqueness of the M\"obius invariant vector. It does not assume finite-dimensional energy eigenspaces. Proposition \ref{prop:net-B} verifies these requirements for $\cal B$. The proof of \cite[Theorem 4.1]{HeT26} uses the split property, which holds for $\cal B$ since $\cal B$ is completely rational.

\begin{thm}[Henriques--Tener, {\cite[Theorem 4.1, pp.~22--23]{HeT26}}]
\label{thm:HT-trace}
For a conformal net satisfying these axioms, its vacuum Hamiltonian satisfies $\operatorname{Tr}(\e^{-\beta L_0}) < \infty$ for every $\beta > 0$.
\end{thm}

This result concerns vacuum representations. Its use for the nonvacuum sectors of $\cal A_S$ requires the extension $\cal B$ and the following elementary observation.

\begin{lemma}
\label{lem:tensor-trace}
Let $A$ and $B$ be nonzero bounded positive operators on nonzero Hilbert spaces. If $A\ot B$ is trace class, then $A$ and $B$ are trace class. Consequently, if $H_1,H_2$ are nonnegative self-adjoint operators and $\e^{-\beta\overline{H_1\ot1+1\ot H_2}}$ is trace class, then both $\e^{-\beta H_i}$ are trace class.
\end{lemma}

\begin{proof}
Choose a unit vector $\eta$ with $\langle B\eta,\eta\rangle > 0$ and let $J_\eta\xi = \xi\ot\eta$. Then $J_\eta^*(A\ot B)J_\eta = \langle B\eta,\eta\rangle A$. The trace class operators form a two-sided ideal, so $A$ is trace class. Interchanging the factors proves the assertion for $B$. For the last assertion, the joint spectral calculus gives $\e^{-\beta\overline{H_1\ot1+1\ot H_2}} = \e^{-\beta H_1}\ot\e^{-\beta H_2}$, and each factor is nonzero.
\end{proof}

\begin{prop}
\label{prop:sector-trace}
Assume \textup{(I1)}, \textup{(F1)} and \textup{(F3)}. For every irreducible representation $\pi$ of $\cal A_S$ and every $\beta > 0$, $\operatorname{Tr}(\e^{-\beta L_0^\pi}) < \infty$. In particular $L_0^\pi$ has compact resolvent, pure point spectrum and finite-dimensional eigenspaces.
\end{prop}

\begin{proof}
Theorem \ref{thm:net} and Proposition \ref{prop:net-B} construct $\cal B$ without a spectral finiteness assumption. Apply Theorem \ref{thm:HT-trace} to its vacuum representation. Each summand projection in \eqref{eq:HB-decomposition} reduces $L_0^{\cal B}$, and the compression of its heat operator to that summand is trace class. By \eqref{eq:diagonal-Hamiltonian} and functional calculus, this compression is $\e^{-\beta L_0^{\pi_z}}\ot\e^{-\beta L_0^{\sigma(\pi_z)}}$. Lemma \ref{lem:tensor-trace} gives the assertion for every $z$. Finally a compact heat operator bounds the rank of every bounded spectral interval of its generator. Since the heat operator is injective with dense range, its spectral decomposition gives a complete eigenbasis for the generator, with finite multiplicities and eigenvalues tending to infinity. This also gives compact resolvent.
\end{proof}

\begin{remark}\normalfont
\label{rem:status-D}
Theorem \ref{thm:HT-C} below only requires that $L_0^\pi$ has discrete spectrum and finite-dimensional eigenspaces, which is weaker than the trace class property of Proposition \ref{prop:sector-trace}. Theorem \ref{thm:HT-trace} concerns vacuum representations, and we do not apply it to the representations $\pi_z$ themselves. The proof of Proposition \ref{prop:sector-trace} uses the diffeomorphism covariance of $\cal B$ and the identity \eqref{eq:diagonal-Hamiltonian}, but neither \textup{(F2)} nor \textup{(F4)}.
\end{remark}

We now use the reconstruction results of \cite{HeT25}. Their finite-energy hypotheses hold for $\cal A_S$, $\cal A_2$, $\cal B$ and the representations considered below by Proposition \ref{prop:sector-trace} and Theorem \ref{thm:HT-trace}.

\begin{thm}[{\cite{HeT25}}]
\label{thm:HT-AB}
Let $\cal A$ be a conformal net. The space $V_{\cal A}$ carries a natural structure of unitary vertex operator algebra \cite[Theorems A and B p.~3, Theorems 7.11 p.~59 and 7.13 p.~61]{HeT25}, with vacuum $\Omega$, conformal vector $\nu = L_{-2}\Omega$ whose field is the stress-energy tensor of $\cal A$, the scalar product of $\cal H_0$, and PCT operator $\Theta$ the modular conjugation of $(\cal A(I_-),\Omega)$ for the lower half circle $I_-$ \cite[(34) p.~23]{HeT25}. For $v \in V_{\cal A}$ and $f \in C^\infty(S^1)$ supported in $I$, the smeared field $Y(v,f)$, the closure of $\int_{S^1} f(z)Y(z^{L_0}v,z)\frac{\d z}{2\pi\i z}$ on $V_{\cal A}$ \cite[Definition 8.2 p.~62]{HeT25}, is affiliated with $\cal A(I)$ \cite[Proposition 8.8 p.~65]{HeT25}, and these operators generate $\cal A(I)$ \cite[Theorem 8.9(i) p.~66]{HeT25}. Conversely, let $V$ be a unitary vertex operator algebra which is AQFT-local, that is, $Y(v_1,f_1)$ and $Y(v_2,f_2)$ commute strongly whenever $f_1$, $f_2$ have disjoint supports \cite[Definition 8.3 p.~62]{HeT25}. Then $\cal A_V(I) = W^*\{Y(v,f) : \operatorname{supp} f \subset I\}$ is a conformal net on the completion of $V$, and $V_{\cal A_V} = V$ as unitary vertex operator algebras \cite[Theorem 8.9 p.~66]{HeT25}.
\end{thm}

The last assertion is \cite[Theorem 8.9(ii)]{HeT25}, which rests on \cite[Proposition 4.6 p.~27]{RTT22}. It gives the same space, scalar product and vertex operator algebra structure, and the PCT operator is then determined by the scalar product \cite[Proposition 5.1 p.~33]{CKLW18}.

\begin{thm}[{\cite{HeT25}}]
\label{thm:HT-C}
Let $\cal A$ be a conformal net and $(K,\pi)$ a representation of $\cal A$ in which $L_0$ has discrete spectrum and finite-dimensional eigenspaces. Then $M = K^{\mathrm{f.e.}}$ carries a unique structure of unitary $V_{\cal A}$-module $Y^M$ \cite[Theorem C p.~5, Proposition 9.3 p.~68, Theorem 9.5 p.~70, Corollaries 9.6 and 9.7 p.~72]{HeT25} such that
\begin{equation}
\label{eq:HT-119}
\pi_I(Y(v,f)) = Y^M(v,f)
\end{equation}
for all $v \in V_{\cal A}$ and $f$ supported in an interval $I$, where $Y^M(v,f)$ is the closure on $M$ of the smeared module field, defined from $Y^M$ as in Theorem \ref{thm:HT-AB}, and $\pi_I(Y(v,f))$ is the image of the affiliated operator $Y(v,f)$ under $\pi_I$. Both sides have $M$ as a core. The $L_0$-grading of $M$ is the grading by the zero mode of $Y^M(\nu,z)$, the scalar product of $K$ is invariant, and for $0 < r < |z| < R$ and $a \in M$,
\begin{equation}
\label{eq:HT-115}
Y^M(v,z)a = R^{L_0}\mathbb A^K_{R,r}[v(z)]r^{-L_0}a,
\end{equation}
where $\mathbb A^K_{R,r}[v(z)]$ is the bounded operator on $K$ attached by \cite[(73) p.~43]{HeT25} to the annulus $r \leq |w| \leq R$ with the insertion of $v$ at $z$. For $K = \cal H_0$, $M = V_{\cal A}$ and $Y^M = Y$.
\end{thm}

By Proposition \ref{prop:sector-trace}, Theorem \ref{thm:HT-C} applies to every irreducible representation of $\cal A_S$, and to every finite direct sum of tensor products of such representations as a representation of $\cal A_2$.

\begin{lemma}
\label{lem:HT-irreducible}
Assume \textup{(I1)}, \textup{(F1)}. Let $(K,\pi)$ be a representation of $\cal A_S$ in which $L_0$ has discrete spectrum and finite-dimensional eigenspaces, and let $M = M(\pi) \ov{\mathrm{def}}{=} K^{\mathrm{f.e.}}$, with the unitary module structure of Theorem \ref{thm:HT-C} over $V_{\cal A_S}$, which is $(S,H_\chi)$ by Lemma \ref{lem:AQFT-local} below.
\begin{enumerate}
\item If $N \subset M$ is an $S$-submodule, the orthogonal projection $P$ onto its closure lies in $\pi(\cal A_S)'$. If $\pi$ is irreducible, then $M$ is a simple $S$-module.
\item If $M$ is a simple $S$-module, then $\pi$ is irreducible.
\item If $(K',\pi')$ is another such representation and the modules $M(\pi)$ and $M(\pi')$ are simple and isomorphic, then $\pi \cong \pi'$.
\item If $(K',\pi')$ is another such representation and $\pi \cong \pi'$, then the $S$-modules $M(\pi)$ and $M(\pi')$ are isomorphic.
\end{enumerate}
\end{lemma}

\begin{proof}
1. The submodule $N$ is stable under $L_0 = \nu^M_{(1)}$, so $N = \bigoplus_h N(h)$ with $N(h) = N \cap M(h)$ finite-dimensional, $P$ maps $M(h)$ onto $N(h)$, and $M = N \oplus (N^\perp \cap M)$. The invariance of the scalar product implies that the algebraic adjoint of each mode is a finite linear combination of modes of $L_1^j\theta v$. Since $N$ is a submodule, its graded orthogonal complement is stable under every mode as well. Consequently $P$ commutes with every mode on $M$. For a test function $f$ supported in $I$, the convergent mode expansion gives $PXa = XPa$ on $M$, where $X = \pi_I(Y(v,f))$ and $M$ is a core by Theorem \ref{thm:HT-C}. Passing to graph limits gives $P\operatorname{Dom}(X)\subset\operatorname{Dom}(X)$ and $XP = PX$ there. The same holds for $1-P$. Thus $X$ is affiliated with $\{P\}'$. Since $\pi_I$ is implemented by a unitary \cite[p.~70]{HeT25} and the $Y(v,f)$ generate $\cal A_S(I)$ (Theorem \ref{thm:HT-AB}), the operators $X$ generate $\pi_I(\cal A_S(I))$, and $P \in \pi(\cal A_S)'$. If $\pi$ is irreducible then $P \in \{0,1\}$, and $P = 1$ gives $N(h) = M(h)$ for all $h$ by finite-dimensionality.

2. A projection $T \in \pi(\cal A_S)'$ commutes with the local unitaries $\pi_I(U(\gamma))$, hence with the representation of the universal cover of $\mathrm{Diff}(S^1)$ on $K$ and with $L_0$, so $T$ preserves $M$ and each $M(h)$. From $TXa = XTa$ for $a \in M$ and all $f$ supported in $I$, comparing homogeneous components gives $Tv^M_ma = v^M_mTa$ for all modes. Thus $T|_M$ is an endomorphism of the simple module $M$, whose lowest weight space is finite-dimensional, so $T|_M$ is a scalar and $T \in \{0,1\}$.

3. Let $\phi \co M(\pi) \to M(\pi')$ be an isomorphism. The form $\la\phi a,\phi b\ra$ is an invariant scalar product on the simple unitary module $M(\pi)$ for the same $\theta$, so it equals $\la Ta,b\ra$ with $T$ positive, graded and commuting with all modes, hence a positive scalar. Thus $\phi$ is a positive multiple of a unitary $u \co M(\pi) \to M(\pi')$, which extends to a unitary $K \to K'$ intertwining the smeared module fields on the cores, hence $u\pi_I(Y(v,f))u^* = \pi'_I(Y(v,f))$. The normal representations $\operatorname{Ad}u \circ \pi_I$ and $\pi'_I$ of $\cal A_S(I)$ are implemented by unitaries $uU_I$ and $U'_I$ from $\cal H_S$, and $U'^*_IuU_I$ fixes every $Y(v,f)$, $\operatorname{supp} f \subset I$, so it lies in $\cal A_S(I)'$ and $\operatorname{Ad}u \circ \pi_I = \pi'_I$. Hence $\pi \cong \pi'$.

4. Let $w \co K \to K'$ be a unitary intertwiner. Then $w\pi_I(U(\gamma))w^* = \pi'_I(U(\gamma))$ for $\gamma \in \mathrm{Diff}(I)$, so that $w$ intertwines the representations of the universal cover of $\mathrm{Diff}(S^1)$, hence the operators $L_0$, and $wM(\pi) = M(\pi')$. Moreover $w\pi_I(Y(v,f))w^* = \pi'_I(Y(v,f))$, so that $w$ maps $Y^{M(\pi)}(v,f)$ onto $Y^{M(\pi')}(v,f)$ by \eqref{eq:HT-119}. Comparing homogeneous components as in part 2, $w|_{M(\pi)}$ intertwines all the modes, and it is an isomorphism of $S$-modules.
\end{proof}

\begin{lemma}
\label{lem:lattice-sectors}
Assume \textup{(I1)} and \textup{(F1)}--\textup{(F3)}. For $M \in \{S, N_b, W_{(1,\pm1)}\}$, the representation $\pi_M$ of \cite[Lemma 4.7]{ArhII} satisfies $M(\pi_M) = M$, so that $M(\pi_0) = S$, $M(\pi_b) = M_b$, $M(\pi_\pm) = X_\pm$. For $g \in \{(0,\pm1),(\pm1,0)\}$, let $\pi_{W_g}$ be the representation of $\cal A_{V'}$ on the closure of the unitary $V'$-module $W_g$ with $\pi_{W_g,I}(Y(v,f)) = Y_{W_g}(v,f)$ \cite[Definition 4.22 p.~43 and Example 6.19(b) p.~71]{Gui25}, \cite[Corollary 4.11 p.~21]{CGGH23}, and let $\pi_{c_1}$, respectively $\pi_{c_2}$, be its restriction to $\cal A_S$ for $g \in \{(0,\pm1)\}$, respectively $g \in \{(\pm1,0)\}$. Then $\pi_{c_1}$ and $\pi_{c_2}$ are irreducible, inequivalent to each other and to $\pi_0$, $\pi_b$, $\pi_\pm$, their $L_0$ have discrete spectrum with finite multiplicities, $M(\pi_{c_j}) = M_{c_j}$, and $\Psi(\pi_{c_j}) \cong c_j$, so that the notation agrees with the labelling of the introduction of this section.
\end{lemma}

\begin{proof}
Since $S$ contains the conformal vector of $H_\pm$, the local implementers of diffeomorphisms of $\cal A_{H_\pm}$ lie in $\cal A_S$, so that the conformal Hamiltonian of $\pi_M$ is the restriction of that of $H_\pm$ and the finite-energy space of $\pi_M$ is $M$. By \cite[Lemma 4.7(1)]{ArhII} and its proof, $\pi_M$ maps $Y(w,f)|_{\cal H_S}$ to the closure of the smeared module field of $M$ with core $M$, which is $Y^M(w,f)$, so $M(\pi_M) = M$ by the uniqueness in Theorem \ref{thm:HT-C}. Since $S$ is a unitary subalgebra of $V'$, the restriction of $\pi_{W_g}$ to $\cal A_S(I) \subset \cal A_{V'}(I)$ maps $Y(s,f)$, $s \in S$, to $Y_{W_g}(s,f)$, the smeared field of the $S$-module $W_g|_S$, so $M(\pi_{c_j}) = W_g|_S \cong M_{c_j}$ by (F3), a simple module, and the $L_0$ of $\pi_{W_g}$ is that of the lattice module $W_g$. Irreducibility and the inequivalences follow from Lemma \ref{lem:HT-irreducible}(2),(4), since $S$, $M_b$, $X_\pm$, $M_{c_1}$, $M_{c_2}$ are pairwise non-isomorphic by (F2), (F3). For the orientation of the DHR braiding of Section \ref{sec:net} for which the conformal spin and statistics theorem \cite{GuL96} gives the twist $\e^{2\pi\i L_0}$ of a sector, the twist of $\pi_{c_j}$ is $\e^{2\pi\i h_{c_j}}$ (we use this orientation for $\Psi$ and for the labels $\pi_z$ of this section, which is allowed since Theorem \ref{thm:net} holds for both orientations by Lemma \ref{lem:ZH3-self-reverse}), braided equivalences between unitary modular tensor categories preserve the canonical spherical structures \cite[Proposition 8.23 p.~34]{ENO05}, hence the twists, and $c_j$ is the only simple object of $Z(\cal H_3)$ with twist $\e^{2\pi\i h_{c_j}}$ \cite[Section 2.2]{ArhII}.
\end{proof}

\begin{cor}
\label{cor:tau-bijection}
Assume \textup{(I1)}, \textup{(F1)}--\textup{(F3)}. The map $\tau \co \pi \mapsto [M(\pi)]$ is a bijection from the equivalence classes of irreducible representations of $\cal A_S$ onto the isomorphism classes of simple $S$-modules, with $\tau(\pi_0) = S$, $\tau(\pi_b) = M_b$, $\tau(\pi_\pm) = X_\pm$, $\tau(\pi_{c_j}) = M_{c_j}$, and it maps $\{\pi_{d_1},\ldots,\pi_{d_6}\}$ onto $\{M_{d_1},\ldots,M_{d_6}\}$. Every simple $S$-module admits an invariant scalar product for $\theta$.
\end{cor}

\begin{proof}
By Proposition \ref{prop:sector-trace} and Lemma \ref{lem:HT-irreducible}(1),(3),(4), $\tau$ is a well-defined injection from the twelve classes of Theorem \ref{thm:net} into the twelve classes of (F2), hence a bijection. The values on the six listed representations are Lemma \ref{lem:lattice-sectors}, and the unitarity is that of the modules $M(\pi)$ in Theorem \ref{thm:HT-C}.
\end{proof}


%%%%%%%%%%%%%%%%%%%%%%%%%%%%%%%%%%%%%%%%%%%%%%%%%%%%%%%%%%%%%%%%%%%%%%%%%%%%%
\subsection{The vertex operator algebra of \texorpdfstring{$\cal B$}{B} and the identification of \texorpdfstring{$S \ot S$}{S x S}}
\label{subsec:identification-SS}

\begin{lemma}
\label{lem:AQFT-local}
Let $V$ be a strongly local simple unitary vertex operator algebra \cite[Definition 6.7 p.~53]{CKLW18}. Then $V$ is AQFT-local, the nets of \cite{CKLW18} and \cite{HeT25} coincide, and $V_{\cal A_V} = V$ as unitary vertex operator algebras. In particular $V_{\cal A_S} = (S,H_\chi)$ and $V_{\cal A_2} = (S \ot S, H_\chi \ot H_\chi)$, with PCT operators $\theta$ and $\theta \ot \theta$.
\end{lemma}

\begin{proof}
The smeared fields of \cite[(106) p.~49]{CKLW18} and \cite[Definition 8.2]{HeT25} agree on $V$. For a vector $a$ of weight $d_a$, the smearing of $Y(z^{L_0}a,z)$ is $\sum_n \hat f(n)a_{(n+d_a-1)}$. Since $V$ is strongly local, it is energy-bounded, so that there are $k \in \N$ and $C > 0$ with $\norm{Y(a,f)\xi} \leq C\norm{(1+L_0)^k\xi}$ for any $\xi \in \cal H^\infty$. Let $P_n$ be the spectral projection of $L_0$ for the interval $[0,n]$. For any $\xi \in \cal H^\infty$, we have $P_n\xi \in V$, $P_n\xi \to \xi$ and $(1+L_0)^kP_n\xi \to (1+L_0)^k\xi$, hence $Y(a,f)P_n\xi \to Y(a,f)\xi$. Thus $V$ is a core for the closure of $Y(a,f)$ on $\cal H^\infty$, and the two closures coincide.

Let $f_1,f_2$ have disjoint supports. Write $f_1=\sum_k f_1^{(k)}$ and $f_2=\sum_l f_2^{(l)}$ with supports compactly contained in intervals $J_k$ and $J'_l$, respectively, with $J_k\cap J'_l=\emptyset$ for every $k,l$. Put $N_1=\bigvee_k\cal A_V(J_k)$ and $N_2=\bigvee_l\cal A_V(J'_l)$. For $x\in N_1'$, \cite[Lemma 6.5]{CKLW18} gives the weak commutation identity $(x^*b\mid Y(a,f_1^{(k)})c)=(Y(a,f_1^{(k)})^*b\mid xc)$ for $b,c\in V$. Linearity of smearing and the adjoint formula on $V$ \cite[(95) and (113)]{CKLW18} let us add these identities for $f_1$. For small rotations $|t|<\delta$, each rotated support remains in $J_k$, and the same summed identity holds for the rotated test function. These are the identities needed for the smoothing in the proof of \cite[Lemma 6.5]{CKLW18}. Convolving $\operatorname{Ad}(\e^{\i tL_0})(x)$ against a smooth approximate identity supported in $(-\delta,\delta)$ gives commutation on the finite-energy core for the smoothed operators. Passing to the limit using closedness gives $xY(a,f_1)\subset Y(a,f_1)x$. The argument for $x^*$ gives affiliation to $N_1$. The corresponding argument gives affiliation to $N_2$ for the second field. Strong locality makes $N_1$ and $N_2$ commute.

Thus $V$ is AQFT-local. For a closed interval in the convention of \cite{HeT25}, fields with support in the interior generate the \cite{CKLW18} algebra. Fields supported in the closed interval commute strongly with fields supported in the interior of its complement, so Haag duality places them in affiliation with the same algebra. Hence the two nets agree. By \cite[Theorem 8.9(ii) p.~66]{HeT25}, reconstruction recovers $V$. Apply this to $S$ using \cite[Theorem 4.5]{ArhII} and to $S\ot S$ using \cite[Corollary 8.2 p.~62]{CKLW18}. The tensor product PCT is $\theta \ot \theta$ by \cite[Proposition 5.1 p.~33]{CKLW18}.
\end{proof}

From now on we assume (I1) and (F1)--(F3), and $\cal B$ is the net of Proposition \ref{prop:net-B}. By Proposition \ref{prop:net-B} and Theorem \ref{thm:HT-trace}, $\cal B$ is a conformal net in the sense of \cite[Definition 2.1 p.~6]{HeT25}, and, since $L_0^\pi = L_0^{\cal B}$ by Proposition \ref{prop:net-B}, the representation $\pi$ of $\cal A_2$ in \eqref{eq:HB-decomposition} satisfies the hypothesis of Theorem \ref{thm:HT-C}. Put $W \ov{\mathrm{def}}{=} V_{\cal B} = \cal H_{\cal B}^{\mathrm{f.e.}}$, a unitary vertex operator algebra by Theorem \ref{thm:HT-AB}, with vacuum $\Omega_{\cal B}$, conformal vector $\nu_W = L_{-2}\Omega_{\cal B}$ and PCT operator $\Theta_W$, and let $Y^M$ be the unitary $S \ot S$-module structure on $M = W$ given by Theorem \ref{thm:HT-C} for $\pi$. Let $\cal H'_0 = \cal H_{\pi_0} \ot \cal H_{\pi_0}$, the closure of $\cal A_2(I)\Omega_{\cal B}$, and let $u \co \cal H_S \ot \cal H_S \to \cal H'_0$ be the unitary intertwiner of representations of $\cal A_2$ with $u(\Omega \ot \Omega) = \Omega_{\cal B}$, unique since the vacuum representation is irreducible and occurs once in \eqref{eq:HB-decomposition}. It intertwines the Möbius representations and, by Proposition \ref{prop:net-B}, the stress-energy tensors. Put $W_0 \ov{\mathrm{def}}{=} (\cal H'_0)^{\mathrm{f.e.}} = u(S \ot S) \subset W$.

\begin{lemma}
\label{lem:naturality}
For $v \in S \ot S$ and $a \in W$, $Y^W(uv,z)a = Y^M(v,z)a$.
\end{lemma}

\begin{proof}
In \cite[Sections 5 and 6]{HeT25}, a lifted annulus $\mathbb A$ with disjoint intervals $I_j$ in its interior and $x_j \in \cal A(I_j)$ defines a worm operator $\mathbb A[x_1 \ldots x_n]$ on any representation $K$ of $\cal A$, the composition $\mathbb A_{n+1}x_n \cdots \mathbb A_2x_1\mathbb A_1$ of \cite[(45) p.~26]{HeT25}, where the lifted annuli $\mathbb A_j$ act on $K$ through the Virasoro representation of $K$ and the $x_j$ through the representation of $\cal A$ and changes of parametrization. For a finite-energy vector $v$ inserted at $z$, one chooses a disc $D$ around $z$ in the annulus, intervals $I_1$, $I_2$ covering $\partial D$ and $x_i \in \cal A(I_i)$ with $|v(z)\ra_D = x_1\Omega_D+x_2\Omega_D$ \cite[Lemma 6.1 p.~39 and Corollary 6.6 p.~42]{HeT25}, and puts $\mathbb A[v(z)] = \mathbb A[x_1]+\mathbb A[x_2]$ \cite[(73) p.~43]{HeT25}, independently of all choices \cite[Lemma 6.9 p.~43]{HeT25}.

Apply this to $\cal A_2$ and $\cal B$ on the common space $\cal H_{\cal B}$, writing $\pi(x) \in \cal A_2(I) \subset \cal B(I)$ for the image of $x \in \cal A_2(I)$. Use the same lift of the annulus, the same cuts, and the same boundary parametrizations in both constructions. By Proposition \ref{prop:net-B}, the Virasoro representation on $\cal H_{\cal B}$ is the one for $\pi$, so the lifted annuli act by identical operators by \cite[(23)]{HeT25}. Changes of local parametrization act on local operators by conjugation, in which the phase of the implementer cancels. Hence, for $x_j \in \cal A_2(I_j)$,
\begin{equation}
\label{eq:worm-naturality}
\mathbb A^{\cal B}[\pi(x_1) \ldots \pi(x_n)] = \mathbb A^{\pi}[x_1 \ldots x_n],
\end{equation}
the left side being the worm operator of $\cal B$ on its vacuum representation and the right side that of $\cal A_2$ on $\pi$. Let $v \in S \ot S = V_{\cal A_2}$ and $|v(z)\ra_D = x_1\Omega_D+x_2\Omega_D$ for the net $\cal A_2$. The unitary $u$ extends to the vacuum sectors $\cal H_0(D) = \mathrm{Conf}(\mathbb D,D) \times_{\mathrm{M\ddot ob}} \cal H_0$ of \cite[(4) p.~8]{HeT25}, compatibly with the local algebras and the annuli, so $|(uv)(z)\ra_D = \pi(x_1)\Omega^{\cal B}_D+\pi(x_2)\Omega^{\cal B}_D$, an admissible decomposition for $\cal B$ since $\pi(x_i) \in \cal B(I_i)$. By \cite[Lemma 6.9]{HeT25} for $\cal B$ and by \eqref{eq:worm-naturality}, $\mathbb A^{\cal B}[(uv)(z)] = \sum_i \mathbb A^{\cal B}[\pi(x_i)] = \sum_i \mathbb A^\pi[x_i] = \mathbb A^\pi[v(z)]$. For $0<r<|z|<1$, apply \eqref{eq:HT-115} with outer radius $1$ to the vacuum representation of $\cal B$ and to $\pi$. In both cases the equal bounded annulus operators are applied to the same vector $r^{-L_0}a$, defined for $a\in W=\cal H_{\cal B}^{\mathrm{f.e.}}$. The two fields therefore agree on this domain, and equality of their Laurent coefficients gives the stated identity of vertex operators.
\end{proof}

\begin{lemma}
\label{lem:vacuum-module}
$W_0$ is an $S \ot S$-submodule of $(M,Y^M)$ and $Y^M(v,z)uv' = uY^{S \ot S}(v,z)v'$ for $v, v' \in S \ot S$. For every summand $\cal H_{\pi_z} \ot \cal H_{\sigma(\pi_z)}$ of \eqref{eq:HB-decomposition}, its finite-energy space is a submodule of $(M,Y^M)$, equal to $M(\pi_z) \ot M(\sigma(\pi_z))$ with the tensor product module structure of \cite[Section 4.7]{FHL93}.
\end{lemma}

\begin{proof}
The projection $E$ onto a subrepresentation $K' \subset \cal H_{\cal B}$ lies in $\pi(\cal A_2)'$ and commutes with the modes $v^M_m = \sum_j \pi_{I_j}(Y(v,e_{m-d_v+1}\chi_j))|_M$, where $v$ is homogeneous of weight $d_v$, $v^M_m$ is the mode $v^M_{(m)}$, $e_k(z) = z^k$ and $(\chi_j)$ is a partition of unity subordinate to a cover by intervals $I_j$ \cite[p.~72]{HeT25}, so $(K')^{\mathrm{f.e.}}$ is a submodule whose structure satisfies \eqref{eq:HT-119} for $K'$, hence is the module structure of Theorem \ref{thm:HT-C} for $K'$. For $K' = \cal H'_0$, this representation is the vacuum representation transported by $u$, so the structure on $W_0$ is $uY^{S \ot S}u^{-1}$ by uniqueness and the last assertion of Theorem \ref{thm:HT-C}. For $K' = \cal H_{\pi_z} \ot \cal H_{\sigma(\pi_z)}$, the finite-energy space is the algebraic tensor product $M(\pi_z) \ot M(\sigma(\pi_z))$, since Proposition \ref{prop:sector-trace} gives only finitely many eigenvalues, counted with multiplicity, in each bounded energy interval, and the tensor product module satisfies \eqref{eq:HT-119} for the vectors $v \ot 1$ and $1 \ot v'$, whose smeared fields are $Y(v,f) \ot 1$ and $1 \ot Y(v',f)$. These vectors generate $S \ot S$ and $Y(v \ot v',z)$ is the normally ordered product of $Y(v \ot 1,z)$ and $Y(1 \ot v',z)$, so two module structures agreeing on them coincide.
\end{proof}

\begin{prop}
\label{prop:identification}
Assume \textup{(I1)}, \textup{(F1)}--\textup{(F3)}.
\begin{enumerate}
\item $W_0$ is a vertex subalgebra of $W$ and $u \co S \ot S \to W_0$ is an isomorphism of vertex algebras.
\item $\nu_W = u(\omega \ot 1+1 \ot \omega)$, and $W$ has central charge $16$.
\item $L_1W_0 \subset W_0$, $\Theta_WW_0 = W_0$ and $\Theta_W|_{W_0} = u(\theta \ot \theta)u^{-1}$. Hence $W_0$, $u(S \ot 1)$ and $u(1 \ot S)$ are unitary subalgebras of $W$ \cite[Definition 5.22 p.~43]{CKLW18}, $u$ is a unitary isomorphism from $(S \ot S, H_\chi \ot H_\chi)$ onto $W_0$ with the scalar product of $\cal H_{\cal B}$, and the unitary structure induced by $W$ on $u(S \ot 1) \cong S$ is $H_\chi$.
\end{enumerate}
\end{prop}

\begin{proof}
1. For $v, v' \in S \ot S$, Lemmas \ref{lem:naturality} and \ref{lem:vacuum-module} give $Y^W(uv,z)uv' = Y^M(v,z)uv' = uY^{S \ot S}(v,z)v'$.

2. By Theorem \ref{thm:HT-AB}, $Y^W(\nu_W,z)$ is the stress-energy tensor of $\cal B$, which is that of $\pi$ by Proposition \ref{prop:net-B}, and the latter is $Y^M(\omega \ot 1+1 \ot \omega,z)$ by \cite[Corollary 9.6 p.~72]{HeT25}. Hence $\nu_W = L_{-2}\Omega_{\cal B} = (\omega \ot 1+1 \ot \omega)^M_{(-1)}\Omega_{\cal B} = u(\omega \ot 1+1 \ot \omega)$ by Lemma \ref{lem:vacuum-module}.

3. The operator $L_1$ of $\cal B$ preserves $\cal H'_0$ by Möbius covariance of $\cal A_2 \subset \cal B$. The PCT operator $\Theta_W$ is the modular conjugation $J$ of $(\cal B(I_-),\Omega_{\cal B})$. By the Bisognano--Wichmann property the vacuum modular group of $\cal B(I_-)$ is implemented by the dilations preserving $I_-$. By Proposition \ref{prop:net-B} these are the common dilations of the extension and the represented subnet. Hence they preserve $\cal A_2(I_-)$ on $\cal H_{\cal B}$. Takesaki's theorem gives the vacuum-preserving conditional expectation onto this subalgebra, as in \cite[Lemma 13 p.~14]{Lon03}. The projection onto $\cal H'_0=\overline{\cal A_2(I_-)\Omega_{\cal B}}$ reduces $J$, and $J|_{\cal H'_0}$ is the modular conjugation for the compressed subnet algebra. The unitary $u$ identifies it with $J_{\cal A_2(I_-)}$ in the vacuum representation of the product net. By Theorem \ref{thm:HT-AB} and Lemma \ref{lem:AQFT-local}, $J_{\cal A_2(I_-)}$ restricts to $\theta \ot \theta$ on $S \ot S$. Hence $W_0$ is a unitary subalgebra by \cite[Proposition 5.23 p.~43]{CKLW18}, and so are $u(S \ot 1)$, $u(1 \ot S)$, since $S \ot 1$ and $1 \ot S$ are stable under $\theta \ot \theta$ and $L_1 \ot 1+1 \ot L_1$. The map $u$ preserves scalar products, vertex operators, conformal vectors and PCT operators.
\end{proof}

\begin{remark}\normalfont
\label{rem:naive-uniqueness}
Agreement of two closed affiliated operators on a vacuum-generated dense domain does not by itself prove equality: that domain must be a core for both operators. Lemma \ref{lem:naturality} instead compares the bounded annulus operators in \eqref{eq:worm-naturality}, using their construction from compatible local algebras, vacuum vectors and Virasoro representations. The reconstruction formula then identifies the vertex operators on their specified finite-energy domains.
\end{remark}


%%%%%%%%%%%%%%%%%%%%%%%%%%%%%%%%%%%%%%%%%%%%%%%%%%%%%%%%%%%%%%%%%%%%%%%%%%%%%
\subsection{Decomposition, braid reversal and the lattice}
\label{subsec:gluing}

We keep the hypotheses (I1), (F1)--(F3), and identify $S \ot S$ with $W_0 \subset W$ through $u$.

\begin{prop}
\label{prop:decomposition}
As an $S \ot S$-module, with $\tau$ as in Corollary \ref{cor:tau-bijection} and $M_{\tau(\pi)} \ov{\mathrm{def}}{=} M(\pi)$, a representative of the class $\tau(\pi)$,
\begin{equation}
\label{eq:W-decomposition}
W = \bigoplus_z M_{\tau(\pi_z)} \ot M_{\tau(\sigma(\pi_z))},
\end{equation}
a sum of twelve pairwise non-isomorphic simple modules. The vertex operator algebra $W$ is simple, of CFT type and $\Z$-graded, $h_{\tau(\pi_z)}+h_{\tau(\sigma(\pi_z))} \in \Z$ for all $z$, and $W_1$ is the direct sum of $(S \ot S)_1$ and of $R_{z'} \ot R_{z''}$ for $(z',z'') \in \{(c_1,c_2),(c_2,c_1),(d_1,d_5),(d_5,d_1),(d_4,d_6),(d_6,d_4)\}$, where $R_z$ is the lowest weight space of $M_z$.
\end{prop}

\begin{proof}
The decomposition is Lemma \ref{lem:vacuum-module}, the tensor product of simple modules is simple \cite[Section 4.7]{FHL93}, and the pairs $(\tau(\pi_z),\tau(\sigma(\pi_z)))$ are distinct since $\tau$ is injective. The weight-zero space of $W$ is the space of Möbius invariant vectors, $\Cbb\Omega_{\cal B}$, so $W$ is of CFT type and simple \cite[Proposition 5.3 p.~34]{CKLW18}. The grading is by the $L_0$ of $\cal B$, with integral spectrum since the M\"obius representation of $\cal B$ on $\cal H_{\cal B}$ is a representation of $\mathrm{PSL}(2,\R)$. By Lemma \ref{lem:sigma} and Corollary \ref{cor:tau-bijection}, $\tau(\pi_x) = M_x$ for $x \in \{0,b,c_1,c_2\}$, $\tau$ maps $\{\pi_a,\pi_{c_0}\}$ onto $\{X_+,X_-\} = \{M_a,M_{c_0}\}$ and $\{\pi_{d_1},\ldots,\pi_{d_6}\}$ onto $\{M_{d_1},\ldots,M_{d_6}\}$, while $\sigma$ has the orbits $\{d_1,d_5\}$, $\{d_2,d_3\}$, $\{d_4,d_6\}$. By (F2), the pairs $\{d_l,d_{l'}\}$ with $h_{d_l}+h_{d_{l'}} \in \Z$ are $\{d_1,d_5\}$, $\{d_4,d_6\}$ (sum $1$) and $\{d_2,d_3\}$ (sum $2$), so $\tau$ maps the three $\sigma$-orbits onto these pairs in some order. The weight-one vectors of $M_{z'} \ot M_{z''}$ with $h_{z'}+h_{z''} = 1$ form $R_{z'} \ot R_{z''}$, and $h_z \geq 1$ for $z \in \{b,a,c_0\}$.
\end{proof}

\begin{prop}
\label{prop:braid-reversal}
There is a braid-reversed equivalence $\tau_{\mathrm{CKM}} \co \Rep(S) \to \Rep(S)$, that is, a braided equivalence $\Rep(S) \to \Rep(S)^{\mathrm{rev}}$, with $\tau_{\mathrm{CKM}}(M_{\tau(\pi_z)}) \cong (M_{\tau(\sigma(\pi_z))})'$ for all $z$, where the prime denotes the contragredient module. The vertex operator algebra $W$ is strongly rational and holomorphic.
\end{prop}

\begin{proof}
We apply \cite[Theorem 5.10(2) p.~51]{CKM22}, stated as Main Theorem 3(2) p.~4 there, with $\cal U = \cal V = \Rep(S)$ and $U = V = S$. By (F1), $S$ is simple, self-contragredient and strongly rational, so $\Rep(S)$ is a semisimple ribbon category with the vertex tensor category structure of \cite{Hua08}, closed under contragredients, submodules and quotients. By Propositions \ref{prop:identification} and \ref{prop:decomposition}, $W = \bigoplus_z U_z \ot V_z$ with $U_z = M_{\tau(\pi_z)}$ distinct simple modules, $U_0 = S$, $V_z = M_{\tau(\sigma(\pi_z))}$ and $\dim\operatorname{Hom}_S(S,V_z) = \delta_{z,0}$, since $\sigma(\pi_z) = \pi_0$ only for $z = 0$, and $W$ is a simple $\Z$-graded vertex operator algebra extending $S \ot S$ with the same vacuum and conformal vector. All summands are integrally graded, so we take the entire index set as the even part in \cite[Theorem 5.10(2)]{CKM22}. We view $W$ in $\Rep(S) \boxtimes \Rep(S) \cong \Rep(S \ot S)$ \cite[Theorems 5.5 p.~48 and 5.7 p.~49]{CKM22}. The conclusion is a braid-reversed equivalence from the subcategory generated by the $U_z$ onto that generated by the $V_z$, with $U_z \mapsto V_z'$, and both are $\Rep(S)$ since $\tau$ and $\tau\sigma$ are bijections onto the simple modules. By \cite[Corollary 5.13 p.~53]{CKM22}, $W$ is strongly rational.

The extension algebra $A_W=\bigoplus_z U_z\boxtimes V_z$ of \cite[Theorem 3.2 p.~10]{HKL15} is commutative and connected, since the vacuum summand occurs once. The proof of \cite[Theorem 5.12 p.~52]{CKM22} gives semisimplicity of the category of $A_W$-modules, hence separability by \cite[Proposition 2.7]{DMNO13}. A tensor equivalence preserves Frobenius--Perron dimensions and $\operatorname{FPdim}(V_z')=\operatorname{FPdim}(V_z)$. Thus $\operatorname{FPdim}(A_W)=\sum_z\operatorname{FPdim}(U_z)^2=\operatorname{FPdim}\Rep(S)$, the square root of $\operatorname{FPdim}(\Rep(S)\boxtimes\Rep(S))$. Hence $A_W$ is Lagrangian, its local modules are trivial \cite[Remark 4.3 p.~16]{DMNO13}, and $\Rep(W) \cong \mathrm{Vec}$ by \cite[Theorem 3.4 p.~16]{HKL15}.
\end{proof}

\begin{lemma}
\label{lem:commutants}
$\operatorname{Com}_W(1 \ot S) = S \ot 1$ and $\operatorname{Com}_W(S \ot 1) = 1 \ot S$. With $V = 1 \ot S$, we have $V^c = S \ot 1$, $V^{cc} = V$, both are regular unitary subalgebras of $W$ isomorphic to $S$ with conformal vectors $1 \ot \omega$ and $\omega \ot 1$ in the sense of \cite[Proposition 5.29 p.~45]{CKLW18}, and $\nu_W = \omega \ot 1+1 \ot \omega$.
\end{lemma}

\begin{proof}
A vector of $W$ annihilated by the non-negative modes of $1 \ot S$ is annihilated by $1 \ot L_0 = (1 \ot \omega)_{(1)}$, whose spectrum on $M_{\tau(\pi_z)} \ot M_{\tau(\sigma(\pi_z))}$, $z \neq 0$, lies in $h_{\tau(\sigma(\pi_z))}+\N$ with $h_{\tau(\sigma(\pi_z))} > 0$ by (F1), since $\sigma(\pi_z) \neq \pi_0$. Hence it lies in $S \ot S_0 = S \ot 1$, which commutes with $1 \ot S$. The other commutant is symmetric. By Proposition \ref{prop:identification}(3) and \cite[Proposition 5.29]{CKLW18}, the conformal vector of $1 \ot S$ is the orthogonal projection of $\nu_W$ onto $1 \ot S$, which is $1 \ot \omega$ since $\omega \ot 1 \in S_2 \ot S_0$ is orthogonal to $1 \ot S$. Regularity is (F1) and \cite[Theorem 4.5 p.~9]{ABD04}.
\end{proof}

\begin{prop}
\label{prop:lattice-identification}
There are an even unimodular lattice $L \in \{E_8 \oplus E_8, D_{16}^+\}$ and a unitary isomorphism of vertex operator algebras $\psi \co W \to V_L$, for the unitary structure of Theorem \ref{thm:HT-AB} on $W$ and that of \cite[Theorem 4.12]{DoL14} on $V_L$. Hence $W$ satisfies Condition B of \cite[Section 3.4 p.~44]{Gui26} and Condition I+ of \cite[Definition 6.17 p.~70]{Gui25}, and $\cal B \cong \cal A_{V_L}$ as conformal nets.
\end{prop}

\begin{proof}
By Propositions \ref{prop:decomposition} and \ref{prop:braid-reversal}, $W$ is a $C_2$-cofinite holomorphic vertex operator algebra of CFT type and central charge $16$, so $W \cong V_L$ for one of the two even unimodular lattices of rank $16$ \cite[Theorem 2 p.~2]{DoM04b}. Transporting the unitary structure of $W$ through an isomorphism $\psi_0$ gives a normalized invariant scalar product on $V_L$, and by \cite[Proposition 5.19 p.~40]{CKLW18} there is $h \in \Aut(V_L)$ such that $\psi = h \circ \psi_0$ is unitary for the structure of \cite[Theorem 4.12]{DoL14}. Conditions B and I+ are formulated in terms of the unitary vertex operator algebra, its unitary modules, intertwining operators and net, and are preserved by unitary isomorphisms, as in the proof of \cite[Theorem 7.17]{ArhII}. They hold for $V_L$ by \cite[Theorem 3.46 p.~57]{Gui26} and \cite[Example 6.19(b) p.~71]{Gui25}. Finally $\cal B = \cal A_W$ by \cite[Theorem 8.9(i) p.~66]{HeT25} and $\cal A_W \cong \cal A_{V_L}$ by \cite[Theorem 6.8 p.~53]{CKLW18}.
\end{proof}

\begin{remark}\normalfont
\label{rem:which-lattice}
Under $\mathrm{(H_J)}$, the lattice $D_{16}^+$ is excluded in \cite[proof of Proposition 7.12(2)]{ArhII} by a count of $(\mathfrak g' \ot 1 \oplus 1 \ot \mathfrak g')$-invariant vectors in $W_1$, based on the decompositions of the spaces $R_z$ under $S_1$ of \cite[Table 2]{Arh26}. By Proposition \ref{prop:decomposition}, $W_1$ is the same $S \ot S$-module as there, and these decompositions are part of (F2). Hence the count, recalled in the proof of Proposition \ref{prop:W-pi}(2), applies to $W$ and gives $W \cong V_{E_8 \oplus E_8}$, see Corollary \ref{cor:HJ-conditional}. This refinement is not used below, since both lattices satisfy Proposition \ref{prop:lattice-identification}.
\end{remark}


%%%%%%%%%%%%%%%%%%%%%%%%%%%%%%%%%%%%%%%%%%%%%%%%%%%%%%%%%%%%%%%%%%%%%%%%%%%%%
\subsection{The diagonal extension theorem}
\label{subsec:theorem-N}

\begin{thm}
\label{thm:net-diagonal}
Assume \textup{(I1)}, \textup{(F1)}--\textup{(F3)}.
\begin{enumerate}
\item There is a strongly rational holomorphic unitary vertex operator algebra $W$ of central charge $16$, isomorphic to $V_L$ for an even unimodular lattice $L \in \{E_8 \oplus E_8, D_{16}^+\}$, containing $S \ot S$ as a unitary subalgebra with the same conformal vector, such that the induced scalar product on each of $S \ot 1$ and $1 \ot S$ is $H_\chi$, and
\[
\operatorname{Com}_W(S \ot 1) = 1 \ot S,
\quad
\operatorname{Com}_W(1 \ot S) = S \ot 1.
\]
As an $S \ot S$-module, $W = \bigoplus_z M_{\tau(\pi_z)} \ot M_{\tau(\sigma(\pi_z))}$, with twelve simple summands of multiplicity one, where $\sigma$ is the involution of Lemma \ref{lem:sigma}. The algebra $W$ satisfies Condition B of \cite{Gui26} and Condition I+ of \cite{Gui25}, and its conformal net is the net $\cal B$ of Proposition \ref{prop:net-B}.
\item The map $\tau \co \pi \mapsto [K_\pi^{\mathrm{f.e.}}]$, which sends an irreducible representation of $\cal A_S$ to the simple $S$-module structure of Theorem \ref{thm:HT-C} on its finite-energy space, is a bijection from the equivalence classes of irreducible representations of $\cal A_S$ onto the isomorphism classes of simple $S$-modules, with $\tau(\pi_0) = S$, $\tau(\pi_b) = M_b$, $\tau(\pi_\pm) = X_\pm$, $\tau(\pi_{c_j}) = M_{c_j}$, and every simple $S$-module is unitarizable for the PCT operator of $(S,H_\chi)$.
\item $\Rep(S)$ admits a braid-reversing autoequivalence, so that $\Rep(S)$ and $\Rep(S)^{\mathrm{rev}}$ are braided equivalent.
\end{enumerate}
\end{thm}

\begin{proof}
Part 2 is Corollary \ref{cor:tau-bijection}. For part 1, let $W = V_{\cal B}$ as in Section \ref{subsec:identification-SS}. By Proposition \ref{prop:identification}, $S \ot S$ is a unitary subalgebra of $W$ with the same conformal vector and the stated unitary structure. The decomposition is Proposition \ref{prop:decomposition}, the commutants are Lemma \ref{lem:commutants}, strong rationality and holomorphy are Proposition \ref{prop:braid-reversal}, and the lattice, Conditions B and I+ and the net are Proposition \ref{prop:lattice-identification}. Part 3 is Proposition \ref{prop:braid-reversal}.
\end{proof}

\begin{remark}\normalfont
\label{rem:status-N}
The theorem uses (F2) for the count of twelve simple modules, (F3) for the lattice-sector identifications and (F1) for positivity, regularity and the vertex tensor category. The spectral hypothesis of Theorem \ref{thm:HT-C} is supplied by Proposition \ref{prop:sector-trace}. Neither $\mathrm{(H_J)}$ nor (F4) is used. The inclusion of nets supplies the unitary structure on the reconstructed $W$. Section \ref{sec:algebraic-route} treats the separate algebraic construction.
\end{remark}


%%%%%%%%%%%%%%%%%%%%%%%%%%%%%%%%%%%%%%%%%%%%%%%%%%%%%%%%%%%%%%%%%%%%%%%%%%%%%
\subsection{The reverse braiding of the Haagerup center}

\begin{lemma}
\label{lem:ZH3-self-reverse}
There is a braided equivalence $R \co Z(\cal H_3)^{\mathrm{rev}} \to Z(\cal H_3)$. It is unique up to natural isomorphism, it is a ribbon equivalence for the twist $\theta^{-1}$ on the source, and the induced permutation of the simple objects fixes $0$, $b$, $a$ and $c_0$, exchanges $c_1$ and $c_2$, and maps $d_l$ to $d_{l'}$ with $l' \equiv \pm5l \bmod 13$. Moreover $R(L_i) \cong L_i$ as algebras for the three Lagrangian algebras $L_1$, $L_2$, $L_3$ of $Z(\cal H_3)$.
\end{lemma}

\begin{proof}
The self-reversal and uniqueness are the categorical statement used in Lemma \ref{lem:sigma}, proved in the first half of the proof of \cite[Proposition 7.12(1)]{ArhII}, which concerns $Z(\cal H_3)$ alone. Positive categorical dimensions determine the spherical structure, so the equivalence is ribbon. The three Lagrangian algebras are distinguished by their module categories $\cal H_1,\cal H_2,\cal H_3$. Reversal replaces each by its tensor opposite, equivalent to itself by self-conjugacy of the Haagerup even parts. Thus $R$ fixes each algebra. Comparing $L_2 = 0\oplus b\oplus2a$ gives $R(a) = a$ and then $R(c_0) = c_0$. The remaining labels are determined by their dimensions and inverse twists, as in Lemma \ref{lem:sigma}.
\end{proof}


%%%%%%%%%%%%%%%%%%%%%%%%%%%%%%%%%%%%%%%%%%%%%%%%%%%%%%%%%%%%%%%%%%%%%%%%%%%%%
\section{Complete unitarity from the reconstructed extension}
\label{sec:unitarity}

%%%%%%%%%%%%%%%%%%%%%%%%%%%%%%%%%%%%%%%%%%%%%%%%%%%%%%%%%%%%%%%%%%%%%%%%%%%%%
\subsection{Complete unitarity and the comparison functor}
\label{subsec:complete-unitarity}

We use the notions of Gui: condition B of \cite[Section 3.4 p.~44]{Gui26}, Condition I of \cite[Definition 4.15 p.~41]{Gui25} and Condition I$^+$ of \cite[Definition 6.17 p.~70]{Gui25}, which adds to Condition I the complete rationality of the net and the essential surjectivity of the functor below. For a simple unitary strongly local vertex operator algebra $V$ \cite[Section 6]{CKLW18}, $\Rep^u(V)$ denotes the category of unitary $V$-modules, $\Rep^f(\cal A_V)$ the category of representations of finite statistical dimension of its conformal net, which is $\Rep(\cal A_V)$ in the convention of Section \ref{sec:net} when $\cal A_V$ is completely rational, and
\[
\Pi \co \Rep^u(V) \to \Rep^f(\cal A_V),
\quad
(M,Y_M) \mapsto (\cal H_M,\pi_M),
\]
the functor of Carpi, Weiner and Xu, defined on strongly integrable modules by the condition $\pi_{M,I}(Y(v,f)) = Y_M(v,f)$ for $v \in V$ homogeneous and $f$ supported in $I$ \cite[(3.8) p.~41]{Gui26}, \cite[Definition 4.22 and Theorem 4.24 p.~43]{Gui25}. The category $\Rep^f(\cal A_V)$ carries the braided $C^*$-tensor structure of Connes fusion \cite{Gui21a,Gui26}. The net $\cal A_S$ is that of Section \ref{sec:net}, for the unitary structure $H_\chi$ of (I1).

\begin{thm}
\label{thm:complete-unitarity}
Assume \textup{(I1)} and \textup{(F1)}--\textup{(F3)}. Let $W$ be the reconstructed extension of Theorem \ref{thm:net-diagonal}, or, more generally, a simple holomorphic unitary vertex operator algebra containing $S \ot S$ with the same conformal vector, such that $\operatorname{Com}_W(S \ot \mathbf 1) = \mathbf 1 \ot S$, $\operatorname{Com}_W(\mathbf 1 \ot S) = S \ot \mathbf 1$, $\mathbf 1 \ot S$ is a unitary subalgebra of $W$, and $W$ is unitarily isomorphic to a lattice vertex operator algebra $V_\Gamma$ with the unitary structure of \cite[Theorem 4.12]{DoL14}. Then $S$, with the unitary structure of \textup{(I1)}, satisfies condition B of \cite{Gui26} and Condition I$^+$ of \cite{Gui25}. In particular:
\begin{enumerate}
\item $S$ is completely unitary in the sense of \cite[Definition 1.8 p.~13]{Gui22}: every $S$-module is unitarizable and the invariant sesquilinear forms of \cite{Gui19} on the spaces of intertwining operators are positive definite. The category $\Rep^u(S)$ is a unitary modular tensor category, and the forgetful functor $\Rep^u(S) \to \Rep(S)$ is a braided equivalence.
\item Every unitary $S$-module is strongly integrable, and $\Pi \co \Rep^u(S) \to \Rep(\cal A_S)$ is an equivalence of braided $C^*$-tensor categories, for the braiding of Connes fusion on the target.
\item For every $S$-submodule $M$ of one of the lattice extensions $H_\pm$ of \textup{(F3)}, with the restricted scalar product, $\Pi(M)$ is the representation $\pi_M$ of \cite[Lemma 4.7]{ArhII} on the closure of $M$. In particular $\Pi(S) = \pi_0$, $\Pi(N_b) = \pi_b$ and $\Pi(X_\pm) = \pi_\pm$ in the notation of Section \ref{subsec:net-setting}. The functor $\Pi$ preserves twists, so that $\Pi(M_{c_1})$, $\Pi(M_{c_2})$ and $\Pi(M_{d_l})$ are the irreducible sectors of $\cal A_S$ with twists $\e^{2\pi\i/3}$, $\e^{4\pi\i/3}$ and $\e^{2\pi\i h_{d_l}}$.
\end{enumerate}
\end{thm}

\begin{proof}
We follow the proof of \cite[Theorem 7.17]{ArhII}, verifying the hypotheses of each result used. Put $V = \mathbf 1 \ot S \subset W$. By Lemma \ref{lem:commutants}, or by the mutual-commutant assumption in the more general case, $V^c \ov{\mathrm{def}}{=} \operatorname{Com}_W(V) = S \ot \mathbf 1$ and $V^{cc} = V$. The coset $V^c$ of the unitary subalgebra $V$ is the coset subalgebra of \cite[Example 5.27 p.~44]{CKLW18}, so it is a unitary subalgebra as well.

\emph{Conformal vectors.} By \cite[Propositions 5.29 and 5.31 p.~45]{CKLW18}, whose hypotheses are that $W$ is a simple unitary vertex operator algebra and $V$ a unitary subalgebra, the vector $\nu = e_V\omega_W$ is a conformal vector of $V$, with which $(V,(\cdot|\cdot))$ is a simple unitary vertex operator algebra, and $\omega_W = \nu+\nu'$ with $\nu'$ the corresponding vector of $V^c$. Since $\omega_W = \omega_{S \ot 1}+\omega_{1 \ot S}$, the vector $\nu-\omega_{1 \ot S} = \omega_{S \ot 1}-\nu'$ lies in $(\mathbf 1 \ot S) \cap (S \ot \mathbf 1) \cap W_2 = 0$. Hence $V$ and $V^c$, with the conformal vector of $S$ and the scalar product inherited from $W$, are unitary vertex operator algebras isomorphic to $S$, and $V$ is a unitary sub-VOA of $W$ with commutant $V^c$ in the sense of \cite[Section 3.6.1 p.~49]{Gui26}.

\emph{Regularity.} $S$ is of CFT type, $C_2$-cofinite and rational by (F1), hence regular by \cite[Theorem 4.5]{ABD04}. Thus $V$ and $V^c$ are regular. In \cite{ArhII} this step uses $\mathrm{(H_J)}$, and (F1) is its replacement.

\emph{Condition B for $W$.} The lattice vertex operator algebra $V_\Gamma$, with the unitary structure of \cite{DoL14}, satisfies condition B \cite[Theorem 3.46 p.~57]{Gui26}, whose only hypothesis is that $\Gamma$ is an even positive definite lattice. Condition B is preserved by unitary isomorphisms of unitary vertex operator algebras, by transport of the vertex operators, of the energy bounds and of the smeared fields, and two structures of unitary vertex operator algebra on the simple algebra $W \cong V_\Gamma$ are unitarily isomorphic \cite[Proposition 5.19 pp.~40--41]{CKLW18}. Hence $W$ satisfies condition B.

\emph{Condition B for $S$.} The hypotheses of \cite[Theorem 3.34 p.~51]{Gui26} are: $W$ satisfies condition B, $V$ is a unitary sub-VOA of $W$, $V$ and $V^c$ are regular, and $V^{cc} = V$. All four hold by the preceding paragraphs. Hence $V^c$ satisfies condition B. No energy bound on $S$ or on its modules is assumed: the proof of \cite[Theorem 3.34]{Gui26} derives the energy bounds of the modules and intertwining operators of $V^c$ by compression from those of $W$. By \cite[Proposition 5.19]{CKLW18} again, the isomorphism of vertex operator algebras $S \to V^c = S \ot \mathbf 1$, composed with an automorphism of $S$, is a unitary isomorphism from $S$ with the structure of (I1) onto $V^c$, so $S$ satisfies condition B.

\emph{Condition I$^+$ for $S$.} The algebra $V_\Gamma$ satisfies Condition I$^+$ \cite[Example 6.19(b) p.~71]{Gui25}, hence so does $W$ by the same transport, Condition I$^+$ being a property of the vertex operators, their energy bounds, the local algebras and the functor $\Pi$. The hypotheses of \cite[Theorem 6.18(3) p.~70]{Gui25} are: $W$ is a unitary vertex operator algebra satisfying Condition I$^+$, $V$ is a unitary subalgebra of $W$, and $V$ and $V^c$ are $C_2$-cofinite and rational. They hold by (F1). Hence $V^c$, and by transport $S$, satisfies Condition I$^+$. Condition I$^+$ includes the complete rationality of $\cal A_S$, which was obtained independently in Theorem \ref{thm:net}.

1. Condition B implies condition A of \cite{Gui26}, and \cite[Theorem 3.15 p.~45]{Gui26} gives that $S$ is completely unitary and that $\Rep^u(S)$ is a unitary modular tensor category, the latter being \cite[Theorem 7.9 p.~53]{Gui19}. Since every $S$-module is unitarizable, the forgetful functor $\Rep^u(S) \to \Rep(S)$ is an equivalence of modular tensor categories \cite[Remark 4.3 p.~18]{CGGH23}.

2. By \cite[Theorems 3.16 and 3.17 p.~45]{Gui26}, under condition A every unitary $S$-module is strongly integrable and $\Pi$ is a fully faithful $*$-functor, and by \cite[Corollary 3.23 p.~48]{Gui26} it is an isomorphism onto a braided $C^*$-tensor subcategory of $\Rep(\cal A_S)$ for Connes fusion. The functor of \cite[Theorem 4.24 p.~43]{Gui25} is the same functor, both being characterized by \cite[(3.8) p.~41]{Gui26}, which is \cite[(4.6) p.~43]{Gui25}. By \cite[Theorem 6.14 pp.~68--69 and Definition 6.17 p.~70]{Gui25}, Condition I$^+$ means that $\Pi$ is essentially surjective onto $\Rep^f(\cal A_S)$, which is $\Rep(\cal A_S)$ since $\cal A_S$ is completely rational. Hence $\Pi$ is an equivalence of braided $C^*$-tensor categories.

3. By (I1) and \cite[Theorem 4.5]{ArhII}, $S$ is a unitary subalgebra of $(H_\pm,\{\cdot|\cdot\}_{\Lambda_\pm})$ containing its conformal vector, so the restriction of the scalar product of $H_\pm$ makes $M$ a unitary $S$-module, and $M$ is strongly integrable by part 2. By \cite[Lemma 4.7(1)]{ArhII}, $\pi_M$ is a representation of $\cal A_S$ on $\cal H_M$ which maps the closed operator $Y(v,f)|_{\cal H_S}$ to the closed operator $Y_M(v,f)$, of which $M$ is a core. The vertex operator $Y_M(v,\cdot)$ is the restriction to $M$ of a field of $H_\pm \cong V_{E_8}$, which satisfies energy bounds \cite[Theorem 3.46]{Gui26}, and energy bounds pass to compressions \cite[Proposition 3.31 p.~50]{Gui26}, so $Y_M(v,f)$ is also the closure of the smeared field of $M$ on the smooth vectors. This is the defining property of $\Pi(M)$, which is unique since $\cal A_S(I)$ is generated by the operators $Y(v,f)$ \cite[(3.8) p.~41]{Gui26}. Hence $\Pi(M) = \pi_M$, and in particular $\Pi(S) = \pi_0$, $\Pi(N_b) = \pi_b$, $\Pi(X_\pm) = \pi_{W_{(1,\pm1)}} = \pi_\pm$ by (F3). Finally, $\Pi$ is a braided equivalence between the pseudo-unitary modular tensor categories $\Rep^u(S)$ and $\Rep(\cal A_S)$, so it preserves the spherical structures with positive dimensions \cite[Proposition 8.23 p.~34]{ENO05} and the twists, and by (F2) the twists $\e^{2\pi\i h_z}$ of $M_{c_1}$, $M_{c_2}$ and of the $M_{d_l}$ are pairwise distinct and distinct from $1$. Independently of the orientation of the braiding, $\Pi(M_z)$ acts on the closure of $M_z$ and its conformal Hamiltonian is the closure of $L_0$ on $M_z$ \cite[Remark 4.23 p.~43]{Gui25}, so that $\e^{2\pi\i L_0} = \e^{2\pi\i h_z}$ on $\Pi(M_z)$.
\end{proof}

\begin{remark}\normalfont
\label{rem:braiding-conventions}
Part 2 of Theorem \ref{thm:complete-unitarity} identifies $\Rep(S)$ with $\Rep(\cal A_S)$ equipped with the braiding of Connes fusion. By \cite[Theorem 6.6 p.~77]{Gui21a}, this braided $C^*$-tensor category is equivalent to the category of DHR endomorphisms of $\cal A_S$ with the tensor product $\rho \odot \sigma = \sigma \circ \rho$ and the statistics operator $\widetilde c$ of \cite[(6.1) p.~73]{Gui21a}, and the identity functor with the tensor constraint given by the braiding is a braided equivalence onto the same category with the tensor product $\rho \circ \sigma$ and the braiding $b_{\rho,\sigma} = \widetilde c_{\sigma,\rho}$, by the hexagon identities and the Yang--Baxter equation \cite[proof of Theorem 7.17(3)]{ArhII}. The two orientations in the definition of the DHR braiding give $b$ and $b^{\mathrm{rev}}$. Hence, for the DHR braiding used in Section \ref{sec:net}, Theorem \ref{thm:complete-unitarity} gives a braided equivalence of $\Rep(S)$ with $\Rep(\cal A_S)$ or with $\Rep(\cal A_S)^{\mathrm{rev}}$, which sends $S$, $M_b$, $X_+$, $X_-$ to $\pi_0$, $\pi_b$, $\pi_+$, $\pi_-$ in both cases. The equivalence of Lemma \ref{lem:ZH3-self-reverse} handles the reversed orientation on the net side. Corollary \ref{cor:S-consequences} uses it to match the intrinsic labels.
\end{remark}


%%%%%%%%%%%%%%%%%%%%%%%%%%%%%%%%%%%%%%%%%%%%%%%%%%%%%%%%%%%%%%%%%%%%%%%%%%%%%
\section{The Haagerup property and its consequences for \texorpdfstring{$S$}{S}}
\label{sec:center}

\begin{thm}
\label{thm:HJ-main}
The vertex operator algebra $S$ has the following properties.
\begin{enumerate}
\item $S$ is a Haagerup algebra, that is, $\mathrm{(H_J)}$ holds. In particular $\Rep(S)$ is braided equivalent to $Z(\cal H_3)$.
\item With the scalar product $H_\chi$, $S$ is completely unitary and satisfies Condition I$^+$. The functor $\Pi\co\Rep^u(S)\to\Rep^f(\cal A_S)$ is a braided $C^*$-tensor equivalence for the braiding of Connes fusion and agrees with $M\mapsto\pi_M$ on the lattice summands of Theorem \ref{thm:complete-unitarity}(3).
\item The heat operator $\e^{-\beta L_0^\pi}$ is trace class for every irreducible sector $\pi$ of $\cal A_S$ and every $\beta>0$. There is a holomorphic unitary vertex operator algebra of central charge $16$ containing $S\ot S$ as a unitary conformal subalgebra, with the two copies of $S$ as mutual commutants.
\end{enumerate}
\end{thm}

\begin{proof}
By Theorems \ref{thm:inputs-ArhI} and \ref{thm:zhu-inputs}, $S$ satisfies (I1) and (F1)--(F3). Theorem \ref{thm:net} identifies $\Rep(\cal A_S)$ with $Z(\cal H_3)$, using \cite[Proposition 5.2 and Theorem 6.1]{ArhII}. Proposition \ref{prop:sector-trace} and Theorem \ref{thm:net-diagonal} give part 3. Theorem \ref{thm:complete-unitarity}, applied to the reconstructed extension $W$ of Theorem \ref{thm:net-diagonal}, gives part 2 and the braided equivalence $\Rep^u(S)\cong\Rep(S)$. For the DHR braiding the comparison is braided or braid-reversed by Remark \ref{rem:braiding-conventions}. The center is equivalent to its reverse by Lemma \ref{lem:ZH3-self-reverse}, so in both orientations $\Rep(S)\cong Z(\cal H_3)$ as braided categories. Together with positivity and strong rationality in (F1), this is part 1.
\end{proof}

By (I3), the vertex operator algebra of Jiao \cite{Jia26} is also a Haagerup algebra.

%%%%%%%%%%%%%%%%%%%%%%%%%%%%%%%%%%%%%%%%%%%%%%%%%%%%%%%%%%%%%%%%%%%%%%%%%%%%%
\subsection{The labels and the third holomorphic extension}

After Theorem \ref{thm:HJ-main}, we write $\Phi_0\co\Rep(S)\to Z(\cal H_3)$ for a braided equivalence and use the intrinsic labels $M_z=\Phi_0^{-1}(z)$. The trivial Brauer--Picard group of $\cal H_3$ makes these labels independent of the choice of $\Phi_0$. Since $\Phi_0$ preserves the canonical spherical structures \cite[Proposition 8.23 p.~34]{ENO05}, it preserves the twists $\e^{2\pi\i h_z}$, and the conformal weights in (F2), with $h_{d_l} \equiv \frac{6l^2}{13} \bmod \Z$, show that these intrinsic labels agree with those of (F2) for the modules with nontrivial twist. The vacuum label is $0$, and the summand $N_b \cong M_b$ of $V'|_S$ in (F3) has the intrinsic label $b$ by \cite[Lemma 5.4]{ArhII}, which applies now that $\mathrm{(H_J)}$ holds. The modules $X_\pm$ have trivial twist by (F3) and dimension $d(\pi_\pm) = 3d+2$ by Theorems \ref{thm:net}(1) and \ref{thm:complete-unitarity}(3), so that their intrinsic labels lie in $\{a,c_0\}$ \cite[Section 2.2]{ArhII}, and $X_+ \not\cong X_-$ gives $\{M_a,M_{c_0}\}=\{X_+,X_-\}$ for the intrinsic labels. For the net we use the intrinsic labels of Theorem \ref{thm:net}(3). Write $\cal F_\pm=\Rep(S)_{L_\pm}$ for the module categories of the two lattice extension algebras.

\begin{cor}
\label{cor:S-consequences}
The following assertions hold.
\begin{enumerate}
\item \emph{Labels.} The assertion \textup{(N)} of \cite[Section 7]{ArhII} holds: there is a braided equivalence $\Phi \co \Rep(\cal A_S) \to \Rep(S)$ with $\Phi(\pi_0) \cong S$ and $\Phi(\pi_x) \cong M_x$ for $x \in \{b,a,c_0\}$, for either orientation of the DHR braiding. The label of $X_\pm$ is the label of $\pi_\pm$: $X_\pm \cong M_a$ if and only if $\pi_\pm = \pi^{\mathrm{net}}_a$. The lattice extension $H_\pm$ of $S$ carries $\cal H_2$, that is, $\cal F_\pm \cong \cal H_2$, if and only if the net $\cal A_{H_\pm}$ does, that is, $\cal F^{\mathrm{net}}_\pm \cong \cal H_2$, if and only if $X_\pm \cong M_a$.
\item \emph{The third extension.} $S$ has exactly three holomorphic extensions up to isomorphism of extensions, $H_+$, $H_-$ and an algebra $B_1$ attached to the Lagrangian algebra $L_1 \cong 0 \oplus b \oplus a \oplus c_0$, all isomorphic to $V_{E_8}$. The algebra $B_1$ carries a unique unitary structure extending that of $S$, it is strongly local, its conformal net is isomorphic to $\cal A_{E_8}$, and it is the third extension of $\cal A_S$ of Theorem \ref{thm:net}(4).
\end{enumerate}
\end{cor}

\begin{proof}
1. Let $\Psi\co\Rep(\cal A_S)\to Z(\cal H_3)$ be the equivalence of Theorem \ref{thm:net}(3) for the DHR braiding. Regard the comparison functor $\Pi$ as an equivalence from $\Rep(S)$ to the net category by forgetting the chosen invariant forms. If $\Pi$ is braided for this orientation, put $\widetilde\Pi=\Pi$. If it lands in the category with reversed braiding, put $\widetilde\Pi=\Psi^{-1}R\Psi^{\mathrm{rev}}\Pi$, where $R$ is the equivalence of Lemma \ref{lem:ZH3-self-reverse}. In both cases $\Psi\widetilde\Pi$ is a braided equivalence, hence naturally isomorphic to $\Phi_0$. The correction for the reversed orientation fixes the sectors labelled $0,b,a,c_0$. By Theorem \ref{thm:complete-unitarity}(3), $\Pi(X_\pm)=\pi_\pm$. Consequently $\widetilde\Pi^{-1}$ gives (N), and the intrinsic label of $X_\pm$ agrees with that of $\pi_\pm$.

By \cite[Theorem 3.9]{ArhII}, applied now that $\mathrm{(H_J)}$ holds, $L_\pm$ are Lagrangian algebras of $\Rep(S)$. Under $\widetilde\Pi$ their objects are those of $\Theta_\pm$ by (F3) and Theorem \ref{thm:net}(1). By the categorical classification of \cite[Proposition 3.3]{ArhII}, the algebra structures are determined up to isomorphism by these objects. Thus the module categories agree, and the identification of the $\cal H_2$ extension follows from \cite[Proposition 3.6]{ArhII} and Theorem \ref{thm:net}(3).

2. By \cite[Theorem 3.9]{ArhII} the three holomorphic extensions are $H_+,H_-,B_1$, all isomorphic to $V_{E_8}$. Proposition 7.10 of \cite{ArhII} applies by Theorem \ref{thm:HJ-main}, which provides both $\mathrm{(H_J)}$ and complete unitarity. It gives the unique extending unitary structure on $B_1$, strong locality and $\cal A_{B_1}\cong\cal A_{E_8}$. Part 3 of that proposition identifies the extension of $\cal A_S$ through the comparison in part 1, with underlying object $\Theta_1$ of Theorem \ref{thm:net}(4).
\end{proof}

%%%%%%%%%%%%%%%%%%%%%%%%%%%%%%%%%%%%%%%%%%%%%%%%%%%%%%%%%%%%%%%%%%%%%%%%%%%%%
\section{The twin algebra and the common fixed point of Dong, Jiao and Xu}
\label{sec:genus}

In this section we prove Theorem \ref{thm:intro-twin}. We first recall the lattice data of \cite[Section 12]{Arh26} and we relate them to the involutions of \cite[Section 5]{DJX26}, then we extend the construction of the twin algebra of \cite[Theorem 12.3]{Arh26} to $S$ using only (I1) and (F1)--(F4), without Theorem \ref{thm:HJ-main}, and we identify the twin of $S$ with the algebra $U$ of Huang and Riesen. The identification rests on a single observation: the common fixed point algebra $C = U^\tau \cong S^\sigma$ of \cite[Theorem 7.3]{DJX26} is the extension of $\Omega_S \ot V_N$ by the thirty-nine simple currents indexed by the $39$-part of $N'/N$, where $N = L_{\mathrm I} \cap L_{\mathrm{II}}$, and $S$, $U$ are its two extensions by the two bosons of $N'/N$.


%%%%%%%%%%%%%%%%%%%%%%%%%%%%%%%%%%%%%%%%%%%%%%%%%%%%%%%%%%%%%%%%%%%%%%%%%%%%%
\subsection{Two binary lattices and their common sublattice}
\label{subsec:lattices}

Let $\mathfrak h = \Cbb p \oplus \Cbb b$ be a two-dimensional space with the symmetric bilinear form $(p,p) = 6$, $(b,b) = 26$, $(p,b) = 0$, so that $p$ and $b$ are the Cartan currents of $\mathfrak{sl}_2$ at levels $3$ and $13$, with $(h,h) = 2k$ for the Cartan element $h$ of $L_{\mathfrak{sl}_2}(k,0)$. Following \cite[Section 5]{DJX26}, put
\begin{equation}
\label{eq:lattice-vectors}
\lambda = \frac{p+b}{2},
\quad
s = \frac{p+b}{4},
\quad
t = \frac{13p-3b}{4},
\end{equation}
so that $(\lambda,\lambda) = 8$, $(s,s) = 2$, $(t,t) = 78$, $(s,t) = 0$ and $(s,p) = \frac32$. Thus $s$ and $t$ are the Cartan currents of $\mathfrak{sl}_2$ at levels $1$ and $39$. Consider the three even lattices
\begin{equation}
\label{eq:three-lattices}
L_{\mathrm{II}} = \Z p \oplus \Z\lambda,
\quad
L_{\mathrm I} = \Z s \oplus \Z\tfrac{s+t}{2},
\quad
N = \Z\lambda \oplus \Z 2p,
\end{equation}
where $\frac{s+t}{2} = \frac{7p-b}{4}$. We also use the root lattice $L = E_6 \oplus A_2$ of $V'$ in the ordered root basis $\alpha_1,\ldots,\alpha_8$ of \cite[(4.1)]{DJX26}, in which $p = \rho_1+\rho_2+\rho_3 = (1,2,3,2,1,0,0,0)$ and $b = (3,4,3,4,3,2,2,2)$ are the Cartan currents of $S_1$ \cite[Section 4.3]{DJX26}, and the lattices
\begin{equation}
\label{eq:L0-LH}
L_0 = \{x \in L : (\lambda,x) \in 2\Z\},
\quad
L_H = L_0+\Z s,
\end{equation}
of \cite[Section 5]{DJX26}. Here $\lambda = (2,3,3,3,2,1,1,1) \in L$, and $L_0$ has index two in $L$.

\begin{lemma}
\label{lem:common-sublattice}
\begin{enumerate}
\item The Gram matrices of $L_{\mathrm{II}}$, $L_{\mathrm I}$ and $N$ in the bases \eqref{eq:three-lattices} are $\left[\begin{smallmatrix}6&3\\3&8\end{smallmatrix}\right]$, $\left[\begin{smallmatrix}2&1\\1&20\end{smallmatrix}\right]$ and $\left[\begin{smallmatrix}8&6\\6&24\end{smallmatrix}\right]$, of determinants $39$, $39$ and $156$. Hence $L_{\mathrm{II}}$ and $L_{\mathrm I}$ are the lattices of case \textup{(II)} and of case \textup{(I)} of \cite[Lemma 12.1]{Arh26}. Neither has a proper even overlattice in $\mathfrak h$.
\item $N = L_{\mathrm I} \cap L_{\mathrm{II}}$ has index two in $L_{\mathrm I}$ and in $L_{\mathrm{II}}$, with $L_{\mathrm{II}} = N \sqcup (p+N)$ and $L_{\mathrm I} = N \sqcup (s+N)$. The pairings $(s,\cdot)$ and $(p,\cdot)$ take their values in $\frac12\Z$ on $L_{\mathrm{II}}'$ and on $L_{\mathrm I}'$ respectively, and the characters $x \mapsto \e^{2\pi\i(s,x)}$ of $L_{\mathrm{II}}$ and $x \mapsto \e^{2\pi\i(p,x)}$ of $L_{\mathrm I}$ both have kernel $N$.
\item The discriminant group $N'/N$ has order $156$. Its $2$-part is $\{0, s, p, s+p\}+N$, with $q(s+N) = q(p+N) = 0$ and $q(s+p+N) = \frac12$, where $q(x) = \frac{(x,x)}2 \bmod \Z$. Its $39$-part $P$ is cyclic and contained in $(L_{\mathrm I}' \cap L_{\mathrm{II}}')/N$, and the quotient maps $L_{\mathrm{II}}'/N \to L_{\mathrm{II}}'/L_{\mathrm{II}}$ and $L_{\mathrm I}'/N \to L_{\mathrm I}'/L_{\mathrm I}$ restrict to isomorphisms of metric groups $\pi_{\mathrm{II}} \co P \to L_{\mathrm{II}}'/L_{\mathrm{II}}$ and $\pi_{\mathrm I} \co P \to L_{\mathrm I}'/L_{\mathrm I}$. Moreover
\begin{equation}
\label{eq:B-sets}
\{x \in L_{\mathrm{II}}' : (s,x) \in \Z\} = \{x \in L_{\mathrm I}' : (p,x) \in \Z\} = \{x \in N' : x+N \in P\} \ov{\mathrm{def}}{=} \mathsf B,
\end{equation}
and $L_{\mathrm{II}}' = \mathsf B \sqcup (\mathsf B+p)$, $L_{\mathrm I}' = \mathsf B \sqcup (\mathsf B+s)$.
\item The isometry $\psi_N = \pi_{\mathrm I}\pi_{\mathrm{II}}^{-1} \co L_{\mathrm{II}}'/L_{\mathrm{II}} \to L_{\mathrm I}'/L_{\mathrm I}$ maps the generator $(1,1)$ of \cite[Lemma 12.1]{Arh26} to $22$ times the generator $(1,1)$, and the four isometries $u \in \{4,17,22,35\}$ of \cite[Lemma 12.1(3)]{Arh26} are the composites of $\psi_N$ with the elements $1$, $14$, $25$, $38$ of $O(L_{\mathrm{II}}'/L_{\mathrm{II}})$, which are induced by the Weyl group of $\mathfrak g' \oplus \mathfrak g''$ \cite[Lemma 12.1(4)]{Arh26}.
\item In the root coordinates, $s$ and $t$ have half-integral coordinates, $L_0 = \{x \in L : \sum_i x_i \in 2\Z\}$, and $\lambda$, $2p$, $\frac{t-s}2$ belong to $L_0$, while $p \notin L_0$ and $s \notin L$. The lattice $L_H$ is even, $(s,L_0) \subset \Z$ and $(p,L_0) \subset \Z$. For the lattice intersections, use the rational plane $\mathfrak h_{\Q} = \Q p\oplus\Q b$ inside $\mathfrak h$, and intersection with either plane gives the same lattices. The intersections are
\begin{equation}
\label{eq:intersections}
L \cap \mathfrak h = L_{\mathrm{II}},
\quad
L_H \cap \mathfrak h = L_{\mathrm I},
\quad
L_0 \cap \mathfrak h = N,
\end{equation}
and $E \ov{\mathrm{def}}{=} L \cap \mathfrak h^\perp = L_H \cap \mathfrak h^\perp = L_0 \cap \mathfrak h^\perp$ is a lattice of rank six.
\end{enumerate}
\end{lemma}

\begin{proof}
All assertions are finite computations with rational numbers in the coordinates $(p,b)$ or in the root coordinates. They are carried out by the certified scripts of Section \ref{sec:certificates} (certificate A1), and we only indicate the identifications and the arguments which replace an enumeration. In the coordinates $(m',m'')$ of \cite[Section 12.1]{Arh26} for case (II), the vector $(m',m'')$ is $\frac{m'}{6}p+\frac{m''}{26}b$, so that $L_{\mathrm{II}} = \{(3a,13c) : a \equiv c \bmod 2\}$ is $\Z p \oplus \Z\lambda$, and for case (I) the vector $(m',m'')$ is $\frac{m'}{2}s+\frac{m''}{78}t$, so that $L_{\mathrm I} = \{(a,39c) : a \equiv c \bmod 2\}$ is $\Z s \oplus \Z\frac{s+t}2$. The generators $(1,1)$ are $\frac{p}{6}+\frac{b}{26}$ and $\frac{s}{2}+\frac{t}{78}$, with $q$-values $\frac{4}{39}$ and $\frac{10}{39}$. An even overlattice $M$ of $L_{\mathrm{II}}$ in $\mathfrak h$ satisfies $L_{\mathrm{II}} \subset M \subset L_{\mathrm{II}}'$ and $M/L_{\mathrm{II}}$ is isotropic for $q$. Since $\frac{4k^2}{39} \in \Z$ forces $39 \mid k$, and likewise $\frac{10k^2}{39}$, the only isotropic subgroup of $L_{\mathrm{II}}'/L_{\mathrm{II}} \cong \Z_{39}$ and of $L_{\mathrm I}'/L_{\mathrm I}$ is trivial, which proves the last assertion of part 1. For part 2, the bases $L_{\mathrm I} = \Z s+\Z 2p$ and $L_{\mathrm{II}} = \Z p+\Z 2s$ of Remark \ref{rem:rank-two-analytic} give $L_{\mathrm I} \cap L_{\mathrm{II}} = \Z 2s+\Z 2p = N$, the two indices and the two coset decompositions, and the two characters are computed on these bases from $(s,p) = \frac32$, $(s,2s) = 4$ and $(p,2p) = 12$. The values in $\frac12\Z$ follow from $2s \in L_{\mathrm{II}}$, $2p \in L_{\mathrm I}$, $(s,\frac p6+\frac b{26}) = \frac12$ and $(p,\frac s2+\frac t{78}) = 1$. For part 3, the order-two classes $s+N$ and $p+N$ have $q$-values zero and pairing $\frac12$ modulo $\Z$, since $(s,p) = \frac32$. They generate a nondegenerate subgroup $T$ of order four in $N'/N$. Its orthogonal complement has order $39$ and is the odd-primary subgroup $P$. The characters $x+N\mapsto\e^{2\pi\i(s,x)}$ and $x+N\mapsto\e^{2\pi\i(p,x)}$ are nontrivial of order two, and their common kernel is $P$. This proves the asserted containment in $(L_{\mathrm I}'\cap L_{\mathrm{II}}')/N$. Since $L_{\mathrm I}'/N$ has order $78$ and contains the $39$-part, the restriction of $L_{\mathrm I}'/N \to L_{\mathrm I}'/L_{\mathrm I}$ to the $39$-part is a bijection, and it preserves $q$ since both groups carry the form induced from $N'$. The subgroup $\{x \in L_{\mathrm{II}}' : (s,x) \in \Z\}$ has index two in $L_{\mathrm{II}}'$ and contains $N$, so its image in $N'/N$ has order $39$ and is the unique subgroup $P$ of order $39$ of $N'/N$, and it equals the preimage $\mathsf B$ of $P$ in $N'$ because $P \subset L_{\mathrm{II}}'/N$. The same argument applies to $\{x \in L_{\mathrm I}' : (p,x) \in \Z\}$, which proves \eqref{eq:B-sets}, and $p \notin \mathsf B$, $s \notin \mathsf B$ since $(s,p) = \frac32$. The value $u = 22$ and the identification of the four isometries follow from $22 \cdot \{1,14,25,38\} \equiv \{22,35,4,17\} \bmod 39$. For part 5, $(\lambda,\alpha_i) \in \{\pm1\}$ for the eight simple roots, which gives the description of $L_0$. The lattice $L \cap \mathfrak h$ is an even lattice of $\mathfrak h$ containing $L_{\mathrm{II}}$, since $p, \lambda \in L$, and it equals $L_{\mathrm{II}}$ by part 1. Likewise $L_H$ is even, it contains $s$ and $\frac{s+t}2 = s+\frac{t-s}2$, so $L_H \cap \mathfrak h = L_{\mathrm I}$. Finally $L_0 = L \cap L_H$, since $s \notin L$, so $L_0 \cap \mathfrak h = L_{\mathrm I} \cap L_{\mathrm{II}} = N$, and a vector of $L_H$ orthogonal to $\mathfrak h$ lies in $L_0$, since $(p,x+s) = (p,x)+\frac32 \notin \Z$ for $x \in L_0$.
\end{proof}

\begin{remark}\normalfont
\label{rem:rank-two-analytic}
The rank-two part admits a short exact calculation without a program. Put $v = (s+t)/2 = (7p-b)/4 = 2p-s$. Then $L_{\mathrm I} = \Z s+\Z 2p$, $L_{\mathrm{II}} = \Z p+\Z 2s$, and $p,s$ are rationally independent. Comparing their two coefficients gives $L_{\mathrm I}\cap L_{\mathrm{II}} = \Z 2p+\Z 2s = N$ and both indices are two. The Gram matrix of $N$ is $\left[\begin{smallmatrix}8&6\\6&24\end{smallmatrix}\right]$. Its Smith invariants are $2,78$, so its odd-primary part is cyclic of order $39$. The preceding nondegenerate order-four subgroup gives its orthogonal splitting. For the canonical discriminant isometry, the case-II generator $g = p/6+b/26$ satisfies $(s,g) = 1/2$, so $g+p$ represents its unique lift in $P$. The case-I generator is $g' = s/2+t/78$. Direct subtraction gives $(g+p)-22g' = -10s\in L_{\mathrm I}$, proving the multiplier $22$. The four isometries solve $10u^2\equiv4\pmod{39}$, or $u^2\equiv16\pmod{39}$, and the Chinese remainder theorem gives $u \in \{4,17,22,35\}$. Block A1 of Section \ref{sec:certificates} also checks these identities and the rank-eight coordinates.
\end{remark}

\begin{remark}\normalfont
\label{rem:DJX-lattices}
In \cite[Section 5]{DJX26} the vectors $p$, $b$, $s$, $t$ appear in the root lattice $L$ of $V'$, the lattice $L_0$ is the kernel of the character $x \mapsto \e^{2\pi\i(s,x)} = (-1)^{(\lambda,x)}$ of $L$, and $L_H$ is the image of the copy $L^{\mathrm H}$ of $E_6 \oplus A_2$ in Huang's coordinates under an isometry $R_* \co L^{\mathrm H} \to L_H$ with $R_*\ell = s$, $R_*h = t$ for the Cartan vectors $\ell$, $h \in L^{\mathrm H}$ of Huang's currents of levels $1$ and $39$ \cite[Proposition 5.1 and (6.10)]{DJX26}. The lattices of parts 1 and 2 are the intersections \eqref{eq:intersections} of these rank-eight lattices with the Cartan subalgebra $\mathfrak h$, which is spanned by the Cartan currents of $S$ and, after transport by $R_*$, by those of $U$. The involutions $\sigma$ and $\tau$ of \textup{(D2)} act on the lattice vertex operator algebras $V_L$ and $V_{L_H}$ through the characters $x \mapsto \e^{2\pi\i(s,x)}$ and $x \mapsto \e^{2\pi\i(p,x)}$ of part 2, both trivial on $L_0$ by part 5, and the group $\{0,s,p,s+p\}+N$ of part 3 is the fusion group $T \cong Z(\mathrm{Vec}_{\Z_2})$ of \cite[Lemma 9.1]{DJX26}, with $a = s+N$, $b = p+N$ and the fermion $ab = s+p+N$. Since $t \in s+L_0$ has half-integral root coordinates, $t$ lies in $L_H$ and not in $L$, see item 2 of Section \ref{subsec:DJX-verification}.
\end{remark}


%%%%%%%%%%%%%%%%%%%%%%%%%%%%%%%%%%%%%%%%%%%%%%%%%%%%%%%%%%%%%%%%%%%%%%%%%%%%%
\subsection{The Heisenberg commutant and the twin of \texorpdfstring{$S$}{S}}
\label{subsec:twin}

Section 12 of \cite{Arh26} is written for a Haagerup algebra $U$, that is, a positive strongly rational vertex operator algebra of central charge eight with $\Rep(U) \cong Z(\cal H_3)$. By Theorem \ref{thm:HJ-main}, $S$ is a Haagerup algebra, so that \cite[Proposition 12.2 and Theorem 12.3]{Arh26} apply to $S$. We go through these statements for $U = S$ with $\mathrm{(H_J)}$ replaced by (I1) and (F1)--(F4), so that Theorem \ref{thm:twin-S}, Lemma \ref{lem:fixed-points} and Theorem \ref{thm:identification} do not depend on Sections \ref{sec:diagonal}--\ref{sec:center}. It is convenient to name the class of algebras which replaces the Haagerup algebras.

\begin{defi}
\label{def:type-S}
A vertex operator algebra $W$ is \emph{of type $S$} if it is simple, of CFT type, strongly rational, positive, of central charge eight, and $\Rep(W)$ is ribbon equivalent to $\Rep(S)$.
\end{defi}

By (F1) and (F4), $S$ is of type $S$, and by (F4) and Lemma \ref{lem:no-invertible} the category $\Rep(W)$ of an algebra of type $S$ is pseudo-unitary and has no invertible object besides the unit. By Theorem \ref{thm:HJ-main}, the algebras of type $S$ are exactly the Haagerup algebras. This is not used in Sections \ref{subsec:twin}--\ref{subsec:identification}. Let $\mathfrak h = \Cbb p \oplus \Cbb b \subset S_1 = \mathfrak g' \oplus \mathfrak g''$ be the Cartan subalgebra spanned by the Cartan currents of levels $3$ and $13$, with the normalization of Section \ref{subsec:lattices}. The Heisenberg commutant of $S$ is $\Omega_S = \operatorname{Com}_S(M(\mathfrak h))$, where $M(\mathfrak h)$ is the Heisenberg vertex operator algebra generated by $\mathfrak h$ \cite[Definition 2.12]{MoR24a}. For a strongly rational vertex operator algebra $\Omega$, we write $P(\Omega)$ for the group of isomorphism classes of invertible simple $\Omega$-modules, with the quadratic form $\theta$ given by the twists, and $\Omega^\nu$ for the module of class $\nu$.

\begin{prop}
\label{prop:associated-lattice-S}
Assume \textup{(I1)} and \textup{(F1)}--\textup{(F4)}.
\begin{enumerate}
\item $\Omega_S$ is a strongly rational vertex operator algebra of central charge six with $(\Omega_S)_1 = 0$. The commutant of $\Omega_S$ in $S$ is the lattice vertex operator algebra $V_{L_{\mathrm{II}}}$, with $L_{\mathrm{II}} = L \cap \mathfrak h$ as in \eqref{eq:intersections}, and there is an anti-isometry $\tau \co L_{\mathrm{II}}'/L_{\mathrm{II}} \to P(\Omega_S)$ onto the group of invertible simple $\Omega_S$-modules such that
\begin{equation}
\label{eq:S-decomposition}
S = \bigoplus_{\alpha \in L_{\mathrm{II}}'/L_{\mathrm{II}}} \Omega_S^{\tau(\alpha)} \ot V_{\alpha+L_{\mathrm{II}}}
\end{equation}
as an $\Omega_S \ot V_{L_{\mathrm{II}}}$-module, and $\Rep(\Omega_S) \cong \Rep(S) \boxtimes \ovl{\cal C(L_{\mathrm{II}}'/L_{\mathrm{II}})}$ as ribbon categories.
\item Every simple $\Omega_S$-module has a nonnegative lowest conformal weight, and for $\alpha \neq 0$ the module $\Omega_S^{\tau(\alpha)}$ contains no vacuum-like vector. Consequently a vector of $S$ of $\mathfrak h$-weight $\gamma$ has conformal weight at least $\frac{(\gamma,\gamma)}2$.
\item The set of $\mathfrak h$-weights of $S$ is $L_{\mathrm{II}}'$, and $e^{\lambda} \in V_{L_{\mathrm{II}}} \subset S$ generates the simple current $J = L_{\mathfrak{sl}_2}(3,3) \ot L_{\mathfrak{sl}_2}(13,13)$ of the affine subalgebra $A$.
\end{enumerate}
\end{prop}

\begin{proof}
This is \cite[Proposition 12.2]{Arh26}, whose proof we follow, indicating each place where the hypothesis $\mathrm{(H_J)}$ was used and what replaces it. (a) The dual pair $(\Omega_S, V_{L_S})$, the strong rationality of $\Omega_S$ and the decomposition of $S$ as a simple current extension of $\Omega_S \ot V_{L_S}$ over a subgroup $A_S \subset L_S'/L_S$ use \cite[Definition 2.12 and Proposition 2.14]{MoR24a}, \cite[Theorem 1]{Mas14} and \cite{CKLR19}, which require $S$ to be simple, strongly rational and of CFT type, with $S_1$ reductive and $\mathfrak h$ a Cartan subalgebra of $S_1$. These are (F1) and (I1). (b) The proof of \cite[Corollary 4.10]{MoR24a}, applied with $V = W = S$, gives $A_S = L_S'/L_S$, the anti-isometry $\tau$ and the equivalence $\Rep(\Omega_S) \cong \Rep(S) \boxtimes \ovl{\cal C(L_S'/L_S)}$. It uses three properties of $\Rep(S)$: it has no invertible object besides the unit, which is Lemma \ref{lem:no-invertible} under (F4), it is pseudo-unitary, which is (F4) or, alternatively, follows from the positivity in (F1) by \cite[Lemma 4.2]{DJX13} and \cite{Hua08}, and it is unpointed, which is the first property again. (c) The central charge of $\Omega_S$ is $8-\operatorname{rank}L_S = 6$, and $(\Omega_S)_1 = 0$ because a vector of $(\Omega_S)_1$ has $\mathfrak h$-weight zero in $S_1 = \mathfrak{sl}_2 \oplus \mathfrak{sl}_2$ (I1), hence lies in $\mathfrak h$ and is orthogonal to $\mathfrak h$. (d) Part 2 is proved in \cite{Arh26} from Zhu's theorem for $S$ and for $\Omega_S$, from the nonnegativity of the conformal weights of the simple $S$-modules, which is (F1), and from the extension theory of \cite{HKL15,CKM22}. It does not use the category $Z(\cal H_3)$. (e) Part 3 uses $V'_1 = S_1 \oplus R_b$, which is (F3) and (F2) with $h_b = 1$, the highest weights $(a,b)$, $a \equiv b \bmod 2$, of the $S_1$-summands of $R_b$, which are read in the column (II) of \cite[Table 2]{Arh26} and are part of (F2), the fact that $V'$ is generated by $V'_1$, the roots $(2,0)$, $(0,2)$ of $S_1$, and the occurrence of the $S_1$-module $V(1,1)$ in $S_2$, which is explicit in \cite[Section 10]{Arh26}. Hence the set of $\mathfrak h$-weights of $S$ is the group $\Gamma_0$ of \cite[Proposition 12.2]{Arh26}, which is $L_{\mathrm{II}}'$, and $L_S = L_{\mathrm{II}}$. Since $V_{L_S} \subset S \subset V'$ is the lattice vertex operator algebra of the Cartan subalgebra $\mathfrak h$, it is contained in $V_{L \cap \mathfrak h}$, and $L \cap \mathfrak h = L_{\mathrm{II}}$ by \eqref{eq:intersections}, so that $V_{L_S} = V_{L_{\mathrm{II}}} = V_{L \cap \mathfrak h}$. (f) The statement on $J$ is \cite[Proposition 11.4]{Arh26}, proved there by an explicit lattice word, or \cite[Proposition 12.2(4)]{Arh26}, whose proof uses part 2 and the regularity of $A$ only.
\end{proof}

Let $\psi \co L_{\mathrm{II}}'/L_{\mathrm{II}} \to L_{\mathrm I}'/L_{\mathrm I}$ be an isometry, induced by a ribbon equivalence $\Phi_\psi \co \cal C(L_{\mathrm{II}}'/L_{\mathrm{II}}) \to \cal C(L_{\mathrm I}'/L_{\mathrm I})$ \cite[Section 8.4]{EGNO15}. By \cite{HKL15}, $S$ corresponds to the connected étale algebra $\cal A_S = \bigoplus_\alpha \tau(\alpha) \boxtimes \alpha$ of $\Rep(\Omega_S) \boxtimes \cal C(L_{\mathrm{II}}'/L_{\mathrm{II}}) \cong \Rep(\Omega_S \ot V_{L_{\mathrm{II}}})$, and we define $S^\natural_\psi$ as the vertex operator algebra extension of $\Omega_S \ot V_{L_{\mathrm I}}$ corresponding to the connected étale algebra $(\mathrm{id} \boxtimes \Phi_\psi)(\cal A_S)$ \cite{HKL15,CKM22}. As an $\Omega_S \ot V_{L_{\mathrm I}}$-module,
\begin{equation}
\label{eq:twin-decomposition}
S^\natural_\psi = \bigoplus_{\alpha \in L_{\mathrm{II}}'/L_{\mathrm{II}}} \Omega_S^{\tau(\alpha)} \ot V_{\psi(\alpha)+L_{\mathrm I}}.
\end{equation}

\begin{thm}
\label{thm:twin-S}
Assume \textup{(I1)} and \textup{(F1)}--\textup{(F4)}, and let $\psi$ be as above.
\begin{enumerate}
\item $S^\natural_\psi$ is an algebra of type $S$. It is self-contragredient, its Heisenberg commutant for the Cartan subalgebra $\mathfrak h^\natural = \mathbf 1 \ot (L_{\mathrm I} \ot \Cbb)_{(-1)}\mathbf 1$ of $(S^\natural_\psi)_1$ is $\Omega_S$, and its associated lattice is $L_{\mathrm I}$.
\item The image of $\operatorname{Aut}(\Omega_S)$ in $O(P(\Omega_S))$ is the whole group. Consequently, for $\Omega \cong \Omega_S$ and $M \in \{L_{\mathrm I},L_{\mathrm{II}}\}$, all the algebras of type $S$ with Heisenberg commutant isomorphic to $\Omega$ and associated lattice isometric to $M$ are isomorphic.
\item Up to isomorphism, $S^\natural_\psi$ does not depend on $\psi$. We write $S^\natural$ for it. We have $(S^\natural)^\natural \cong S$ and $S^\natural \not\cong S$, and the algebras of type $S$ with Heisenberg commutant isomorphic to $\Omega_S$ are, up to isomorphism, exactly $S$ and $S^\natural$.
\item $S$ and $S^\natural$ belong to the same hyperbolic genus in the sense of \cite[Definition 4.3]{MoR24a}, and every strongly rational vertex operator algebra in this hyperbolic genus is isomorphic to $S$ or to $S^\natural$. In particular it is of type $S$.
\end{enumerate}
\end{thm}

\begin{proof}
We follow the proof of \cite[Theorem 12.3]{Arh26} with Proposition \ref{prop:associated-lattice-S} in place of \cite[Proposition 12.2]{Arh26}, and we indicate the two places where the Haagerup hypothesis entered through \cite[Theorem 8.5]{Arh26}.

1. The paragraphs \emph{extension and representation category}, \emph{CFT type}, \emph{self-contragredience}, \emph{$C_2$-cofiniteness}, \emph{rationality} and \emph{positivity} of the proof of \cite[Theorem 12.3]{Arh26} use \cite[Theorem 3.2 and Remark 3.3]{HKL15}, \cite{CKM22}, \cite{DMNO13}, \cite{Li94}, \cite{Buh02}, \cite{ABD04}, the positivity of $S$, which is (F1), and parts 1 and 2 of Proposition \ref{prop:associated-lattice-S}. They show that $W = S^\natural_\psi$ is a simple, self-contragredient, strongly rational vertex operator algebra of CFT type and central charge $6+2 = 8$, that $\Rep(W)$ is ribbon equivalent, through $\mathrm{id} \boxtimes \Phi_\psi$, to the category of local $\cal A_S$-modules, that is, to $\Rep(S)$, and that $W$ is positive. Hence $W$ is of type $S$. At this point \cite{Arh26} concludes that $W$ is a Haagerup algebra and invokes \cite[Theorem 8.5]{Arh26} to place it in case (I) or (II). We do not use this theorem here. The paragraph \emph{Heisenberg commutant} of \cite{Arh26} shows, using only $(\Omega_S)_1 = 0$, the reductivity of $W_1$ \cite[Theorem 1]{DoM04} and part 2 of Proposition \ref{prop:associated-lattice-S}, that $\mathfrak h^\natural$ is a Cartan subalgebra of $W_1$, that $\operatorname{Com}_W(M(\mathfrak h^\natural)) = \Omega_S \ot \mathbf 1$ and that $\operatorname{Com}_W(\Omega_S) = V_{L_{\mathrm I}}$. This is part 1. The last sentence of this paragraph in \cite{Arh26}, which identifies the case of $W$ and the Lie algebra $W_1$, is not used here. We determine $W_1$ in Corollary \ref{cor:twin-consequences} below by a different argument.

2. The proof of \cite[Theorem 12.3(2)]{Arh26} uses the anti-isometry $\tau$, the Weyl reflections of $\mathfrak g'$ and $\mathfrak g''$, which are inner automorphisms of $S$ normalizing $\mathfrak h$, preserving $\Omega_S$ and $V_{L_{\mathrm{II}}}$ and inducing $14$ and $25$ on $L_{\mathrm{II}}'/L_{\mathrm{II}}$ by \cite[Lemma 12.1(4)]{Arh26}, and the uniqueness statement \cite[Proposition 5.3]{DoM04} for two simple vertex operator algebras which are simple current extensions of $\Omega \ot V_M$ with the same summands. The decomposition of an algebra $Y$ of type $S$ with Heisenberg commutant $\Omega$ and associated lattice $M$ as $\bigoplus_\alpha \Omega^{\sigma(\alpha)} \ot V_{\alpha+M}$ with $\sigma$ an anti-isometry onto $P(\Omega)$ is given by the proof of \cite[Corollary 4.10]{MoR24a}, applied to $Y$ as in step (b) of the proof of Proposition \ref{prop:associated-lattice-S}: it uses only that $\Rep(Y) \cong \Rep(S)$ is unpointed and pseudo-unitary. No other property of $Z(\cal H_3)$ is used.

3. The independence of $\psi$ and the involutivity follow from parts 1 and 2 as in \cite{Arh26}. The lattices $L_{\mathrm I}$ and $L_{\mathrm{II}}$ are not isometric, and an isomorphism of vertex operator algebras maps a Cartan subalgebra of the weight-one Lie algebra to a Cartan subalgebra, all of which are conjugate under inner automorphisms, so it identifies the Heisenberg commutants and the associated lattices. Hence $S^\natural \not\cong S$. For the last assertion, \cite{Arh26} uses \cite[Proposition 12.2(3)]{Arh26} for an arbitrary Haagerup algebra $W$ with $\Omega_W \cong \Omega_S$, which rests on \cite[Theorem 8.5]{Arh26}. We argue instead as follows. Let $W$ be of type $S$ with $\Omega_W \cong \Omega_S$. By \cite[Corollary 4.10]{MoR24a}, whose hypotheses are that $\Rep(W) \cong \Rep(S)$ is unpointed and pseudo-unitary, that $W$ and $S$ have the same central charge and isomorphic Heisenberg commutants, $W$ and $S$ belong to the same hyperbolic genus. By \cite[Definition 4.6 and Theorem 4.7]{MoR24a} the associated lattice of $W$ is in the genus of $L_{\mathrm{II}}$, hence isometric to $L_{\mathrm I}$ or $L_{\mathrm{II}}$ by \cite[Lemma 12.1(2)]{Arh26}, and part 2 gives $W \cong S$ or $W \cong S^\natural$.

4. This is the proof of \cite[Theorem 12.3(4)]{Arh26}, which uses \cite[Corollary 4.10 and Theorems 4.7, 4.8]{MoR24a} and part 2, with the hypotheses of \cite[Corollary 4.10]{MoR24a} verified as in part 3. The conclusion ``in particular it is a Haagerup algebra'' becomes ``it is of type $S$'', since $\Rep(V) \cong \Rep(S)$ by \cite[Theorem 4.8]{MoR24a} for a strongly rational $V$ in the genus, and the positivity of $V$ follows from its decomposition over $\Omega_S \ot V_M$ and part 2 of Proposition \ref{prop:associated-lattice-S} as in the paragraph \emph{positivity} of part 1.
\end{proof}

\begin{remark}\normalfont
\label{rem:twin-not-covered}
Theorem \ref{thm:twin-S} gives the lattice exchange and the classification in the hyperbolic genus of $S$ under its stated inputs. For classification among all Haagerup algebras we use \cite[Theorem 12.4(1),(4)]{Arh26} with the Haagerup property of $S$ from Theorem \ref{thm:HJ-main}. Theorem \ref{thm:identification} below identifies the explicit case (I) algebra using (D1) and (D2).
\end{remark}


%%%%%%%%%%%%%%%%%%%%%%%%%%%%%%%%%%%%%%%%%%%%%%%%%%%%%%%%%%%%%%%%%%%%%%%%%%%%%
\subsection{The two involutions and the common fixed point algebra}
\label{subsec:involutions}

We now use the lattice $N$ of Lemma \ref{lem:common-sublattice}. Let $B = \Omega_S \ot V_N$. Since $N \subset L_{\mathrm{II}}$ and $N \subset L_{\mathrm I}$ have index two, $B$ is a vertex operator subalgebra of $\Omega_S \ot V_{L_{\mathrm{II}}} \subset S$ and of $\Omega_S \ot V_{L_{\mathrm I}} \subset S^\natural$, with the same conformal vector. It is strongly rational, with $\Rep(B) \cong \Rep(\Omega_S) \boxtimes \cal C(N'/N)$ \cite{FHL93,CKM22}. For $\gamma \in \mathsf B$, in the notation \eqref{eq:B-sets}, we write $\bar\gamma = \gamma+L_{\mathrm{II}} \in L_{\mathrm{II}}'/L_{\mathrm{II}}$, so that $\psi_N(\bar\gamma) = \gamma+L_{\mathrm I}$ by definition of $\psi_N$, and we put
\begin{equation}
\label{eq:M-summands}
M_{\gamma+\epsi s+\delta p+N} \ov{\mathrm{def}}{=} \Omega_S^{\tau(\bar\gamma)} \ot V_{\gamma+\epsi s+\delta p+N},
\quad
\gamma \in \mathsf B,\ \epsi,\delta \in \{0,1\}.
\end{equation}
These are pairwise non-isomorphic invertible simple $B$-modules, indexed by $N'/N = P \oplus \{0,s,p,s+p\}$, and their fusion is $M_x \boxtimes M_y \cong M_{x+y}$, since $\tau$ and $\pi_{\mathrm{II}}$ are homomorphisms and the fusion of lattice modules is given by the addition of cosets \cite{DL93}. The twists are $\theta_{M_x} = \e^{2\pi\i h_x}$, where, for $x = \gamma+\epsi s+\delta p+N$,
\begin{equation}
\label{eq:M-weights}
h_x \equiv h_{\tau(\bar\gamma)}+\frac{(\gamma+\epsi s+\delta p)^2}2 \equiv \epsi\frac{(s,s)}2+\delta\frac{(p,p)}2+\epsi\delta(s,p) \equiv \frac{3\epsi\delta}2 \pmod{\Z},
\end{equation}
because $h_{\tau(\bar\gamma)}+\frac{(\gamma,\gamma)}2 \in \Z$, the weights of $S$ being integral, and $(\gamma,s), (\gamma,p) \in \Z$ for $\gamma \in \mathsf B$ by \eqref{eq:B-sets}.

The automorphisms $\sigma = \exp(2\pi\i s_{(0)})$ of $S$ and $\tau' = \exp(2\pi\i p_{(0)})$ of $S^\natural$ are well defined: $s_{(0)}$ acts on the summand $\Omega_S^{\tau(\alpha)} \ot V_{\alpha+L_{\mathrm{II}}}$ of \eqref{eq:S-decomposition} by the scalar $(s,x)$ on the vectors of $\mathfrak h$-weight $x \in \alpha+L_{\mathrm{II}}$, which lies in $\frac12\Z$ by Lemma \ref{lem:common-sublattice}(2), and likewise for $p_{(0)}$ on \eqref{eq:twin-decomposition}. Since $s = \frac{p+b}4$, the automorphism $\sigma$ is the involution $\exp(\frac{\pi\i}2(P_H+B_H)_{(0)})$ of \textup{(D2)}, and $\tau'$ is the transport of the involution $\tau = \exp(\frac{\pi\i}2(3H_L+H_R)_{(0)})$ of $U$, since $R_*^{-1}p = \frac{3\ell+h}4$ \cite[Section 5]{DJX26}.

\begin{lemma}
\label{lem:fixed-points}
Assume \textup{(I1)} and \textup{(F1)}--\textup{(F4)}, and take $\psi = \psi_N$ in \eqref{eq:twin-decomposition}.
\begin{enumerate}
\item $\sigma$ and $\tau'$ are automorphisms of order two of $S$ and of $S^\natural$. As $B$-modules,
\[
S^\sigma = \bigoplus_{\gamma \in P} M_{\gamma},
\quad
S^- \ov{\mathrm{def}}{=} \{v \in S : \sigma v = -v\} = \bigoplus_{\gamma \in P} M_{\gamma+p},
\]
\[
(S^\natural)^{\tau'} = \bigoplus_{\gamma \in P} M_{\gamma},
\quad
(S^\natural)^- \ov{\mathrm{def}}{=} \{v \in S^\natural : \tau' v = -v\} = \bigoplus_{\gamma \in P} M_{\gamma+s},
\]
where $\gamma$ runs over representatives in $\mathsf B$ of the classes of $P$. In particular $S^\sigma$ and $(S^\natural)^{\tau'}$ are simple current extensions of $B$ with the same summands.
\item Put $C = S^\sigma$. Then $C$ is a simple, self-contragredient, strongly rational vertex operator algebra of CFT type and central charge eight, and $S^-$ is a simple $C$-module, not isomorphic to $C$. The same holds for $(S^\natural)^{\tau'}$ and $(S^\natural)^-$, and $C \cong (S^\natural)^{\tau'}$.
\item Let $X$ be a simple $C$-module whose restriction to $B$ contains a submodule isomorphic to $M_{x}$ for some $x \in \{s, p, s+p\}+N$. Then $X \cong \Ind(M_x) \ov{\mathrm{def}}{=} \bigoplus_{\gamma \in P} M_{\gamma+x}$ as $C$-modules, where $\Ind(M_x)$ carries the induced module structure of \cite{HKL15,CKM22}. In particular $S^- \cong \Ind(M_p)$.
\end{enumerate}
\end{lemma}

\begin{proof}
1. The automorphism $\sigma$ acts on the vectors of $\mathfrak h$-weight $x$ by $\e^{2\pi\i(s,x)} = \pm1$, and $\sigma \neq 1$ since $(s,p) = \frac32$ and $e^p \in V_{L_{\mathrm{II}}} \subset S$. On the summand $\Omega_S^{\tau(\alpha)} \ot V_{\alpha+L_{\mathrm{II}}}$ of \eqref{eq:S-decomposition}, the fixed points of $\sigma$ are $\Omega_S^{\tau(\alpha)} \ot V_{(\alpha+L_{\mathrm{II}}) \cap \mathsf B}$, and by \eqref{eq:B-sets} and $L_{\mathrm{II}} = N \sqcup (p+N)$ the set $(\alpha+L_{\mathrm{II}}) \cap \mathsf B$ is a single coset $\gamma+N$ with $\gamma \in \mathsf B$, $\bar\gamma = \alpha$, while the other coset $\gamma+p+N$ of $\alpha+L_{\mathrm{II}}$ is the $-1$-eigenspace. This gives the decompositions of $S^\sigma$ and $S^-$. For $S^\natural_{\psi_N}$, the summand indexed by $\alpha \in L_{\mathrm{II}}'/L_{\mathrm{II}}$ in \eqref{eq:twin-decomposition} is $\Omega_S^{\tau(\alpha)} \ot V_{\psi_N(\alpha)+L_{\mathrm I}}$, and $\psi_N(\alpha) = \gamma+L_{\mathrm I}$ for the representative $\gamma \in \mathsf B$ of $\alpha$, by definition of $\psi_N = \pi_{\mathrm I}\pi_{\mathrm{II}}^{-1}$. The involution $\tau'$ acts by $\e^{2\pi\i(p,x)}$ and its fixed points in $V_{\gamma+L_{\mathrm I}}$ are $V_{\gamma+N}$, since $L_{\mathrm I} = N \sqcup (s+N)$ and $(p,\gamma) \in \Z$, $(p,s) = \frac32$. This gives the decompositions of $(S^\natural)^{\tau'}$ and $(S^\natural)^-$, with the same $\Omega_S$-labels $\tau(\bar\gamma)$ as for $S$.

2. The algebra $S$ is simple, regular \cite{ABD04}, non-negatively graded, with the nonsingular invariant bilinear form given by its unitary structure (I1). By \cite[Theorem 5.24 p.~33]{CaM16}, the fixed point algebra $C = S^\sigma$ is regular, and by \cite[Theorem 4.4]{DoM97} it is simple, and $S^-$, the $-1$-eigenspace of $\sigma$, is a simple $C$-module not isomorphic to $C$ \cite[Theorem 3]{DoM97}. It is of CFT type, since $C_0 = S_0$, and self-contragredient, since $L(1)C_1 \subset L(1)S_1 = 0$ \cite{Li94}. The algebra $S^\natural$ is simple, strongly rational, of CFT type and self-contragredient by Theorem \ref{thm:twin-S}(1), and the invariant bilinear form of a simple self-contragredient algebra of CFT type is nonsingular \cite{Li94}, so the same references apply to $\tau'$. Finally $C$ and $(S^\natural)^{\tau'}$ are simple vertex operator algebras which are simple current extensions of $B$ with the same pairwise non-isomorphic invertible summands $M_\gamma$, $\gamma \in P$, with fusion $M_\gamma \boxtimes M_{\gamma'} \cong M_{\gamma+\gamma'}$ and multiplicities one. Since $B$ is strongly rational, the space of intertwining operators of type $\binom{M_z}{M_x\,M_y}$ has dimension $\dim\operatorname{Hom}_B(M_x \boxtimes M_y,M_z) = \delta_{z,x+y}$ \cite{HL95,Hua08}. Hence the vertex operators of both algebras map $M_x \ot M_y$ into $M_{x+y}((z))$, both algebras are $P$-graded in the sense of \cite[Definition 5.2]{DoM04}, and the fusion rules satisfy the hypothesis $N_{M_x,M_y}^{M_z} = \delta_{z,x+y}$ of \cite[Proposition 5.3]{DoM04}. Transporting the vertex operators of $(S^\natural)^{\tau'}$ to $C$ through an isomorphism of $B$-modules which is the identity on $M_0 = B$, we obtain from \cite[Proposition 5.3]{DoM04} an isomorphism $C \cong (S^\natural)^{\tau'}$, which acts by a nonzero scalar on each summand $M_\gamma$, as in the proof of \cite[Theorem 12.3(2)]{Arh26}.

3. The monodromy of $M_x$ with the summand $M_\gamma$ of $C$ is the scalar $\theta_{M_{\gamma+x}}\theta_{M_\gamma}^{-1}\theta_{M_x}^{-1} = \e^{2\pi\i(\gamma,x)} = 1$ by \eqref{eq:M-weights} and \eqref{eq:B-sets}. Hence $\Ind(M_x) = C \boxtimes_B M_x = \bigoplus_{\gamma \in P} M_{\gamma+x}$ is a $C$-module \cite{HKL15,CKM22}, a finite direct sum of simple $C$-modules since $C$ is regular. By the adjunction between induction and restriction \cite[Theorem 1.6 p.~7]{KiO02},
\[
\operatorname{End}_C(\Ind(M_x)) \cong \operatorname{Hom}_B(M_x,\Ind(M_x)) = \bigoplus_{\gamma \in P}\operatorname{Hom}_B(M_x,M_{\gamma+x}) = \Cbb,
\]
since $M_{\gamma+x} \cong M_x$ only for $\gamma \in N$. Thus $\Ind(M_x)$ is a simple $C$-module. The same adjunction gives a nonzero $C$-module map $\Ind(M_x) \to X$ from the inclusion $M_x \to X$, and it is an isomorphism since both modules are simple. For $S^-$, take $x = p$ and the submodule $M_p \ni e^p$.
\end{proof}


%%%%%%%%%%%%%%%%%%%%%%%%%%%%%%%%%%%%%%%%%%%%%%%%%%%%%%%%%%%%%%%%%%%%%%%%%%%%%
\subsection{The identification \texorpdfstring{$U \cong S^\natural$}{U = twin of S}}
\label{subsec:identification}

In this subsection (I3) is used to transport the fixed-point statement (D2) from Jiao\textquotesingle s algebra to $S$, including its identified Cartan currents. Let $U = U^{\mathrm H} \subset V_{L^{\mathrm H}}$ be the vertex subalgebra generated by the explicit vectors of Huang and by the conformal vector $\omega$ of $V_{L^{\mathrm H}}$ \cite[(4.6)--(4.9)]{DJX26}, where $L^{\mathrm H}$ is the copy of $E_6 \oplus A_2$ in Huang's coordinates of Remark \ref{rem:DJX-lattices}, and let $\tau$ be the involution of \textup{(D2)}. Following \cite[Sections 5--7]{DJX26}, let $\hat\Phi \co V_{L^{\mathrm H}} \to V_{L_H}$ be a lift of the isometry $R_* \co L^{\mathrm H} \to L_H$ of Remark \ref{rem:DJX-lattices}, extending the lift $\Phi$ of \cite[(6.13)]{DJX26} on $V_{R_*^{-1}L_0}$ \cite[p.~21]{DJX26}, and let $\tilde U = \hat\Phi(U) \subset V_{L_H}$ be the transported algebra. The two copies of $V_{L_0}$ in $V_{L_H}$ and in $V_L = V'$ are identified by the cocycle-compatible isomorphism of \cite[(5.7)]{DJX26}, which maps $P e^x$ to $\eta(x)Pe^x$ for a function $\eta \co L_0 \to \Cbb^\times$. Through this identification, which we suppress from the notation, \cite[Theorem 7.3]{DJX26} states that
\begin{equation}
\label{eq:DJX-fixed-point}
\tilde U^\tau = S^\sigma = C,
\quad
\tilde U = C \oplus \tilde U^-,
\quad
S = C \oplus S^-,
\end{equation}
where $\tilde U^-$ is the $-1$-eigenspace of $\tau$ in $\tilde U$, and $\tau$ acts on $V_{L_H}$ by $\e^{2\pi\i(p,x)}$ on the vectors of $\mathfrak h$-weight $x$. The subalgebra $B = \Omega_S \ot V_N$ of $C$ lies in $V_{L_0}$: indeed $V_N \subset V_{L_0}$ since $N = L_0 \cap \mathfrak h$, and $\Omega_S \subset \operatorname{Com}_{V'}(M(\mathfrak h)) = V_E \subset V_{L_0}$ with $E = L \cap \mathfrak h^\perp = L_0 \cap \mathfrak h^\perp$ by Lemma \ref{lem:common-sublattice}(5).

\begin{thm}
\label{thm:identification}
Assume \textup{(I1)}, \textup{(I3)}, \textup{(F1)}--\textup{(F4)} and \textup{(D2)}. Then $U \cong \tilde U \cong S^\natural$ as vertex operator algebras. More precisely, as $B$-modules,
\[
\tilde U \cong \bigoplus_{\gamma \in P} M_\gamma \oplus \bigoplus_{\gamma \in P} M_{\gamma+s} \cong S^\natural_{\psi_N}.
\]
\end{thm}

\begin{proof}
\emph{Step 1: $U$ is simple.} The lattice vertex operator algebra $V_{L^{\mathrm H}}$ carries the positive definite invariant Hermitian form $(\cdot,\cdot)_{\kappa_H}$ and the antilinear involution $\theta_H$ of \cite[Section 4.4]{DJX26}, in the sense of \cite[Theorem 4.12]{DoL14}, and $\theta_H(U) = U$ by \cite[(6.15)]{DJX26}, an identity on the generators verified in Section \ref{subsec:DJX-verification}. The restriction to $U$ of the bilinear form $(\theta_H a,b)_{\kappa_H}$ is a nondegenerate invariant bilinear form on $U$ for the conformal vector $\omega \in U$ of $V_{L^{\mathrm H}}$ \cite[(4.9)]{DJX26}. A proper ideal $I$ of $U$ contains no nonzero vector of weight zero, and for $a \in I$, $b \in U$ the invariance expresses $(a,b)$ as a combination of pairings $(a',\mathbf 1)$ with $a' \in I$ \cite{FHL93,Li94}, so that $(I,U) = 0$ and $I = 0$. Hence $U$, and $\tilde U \cong U$, are simple vertex operator algebras of CFT type.

\emph{Step 2: the two eigenspaces.} By Step 1 and \cite[Theorems 3 and 4.4]{DoM97}, $\tilde U^\tau = C$ is simple and $\tilde U^-$ is a simple $C$-module, not isomorphic to $C$. The current $E_L = e^\ell$ of Huang is mapped by $\hat\Phi$ to a nonzero multiple of $e^{s}$, since $R_*\ell = s$, and $\tau(E_L) = -E_L$ since $(p,s) = \frac32$. Hence $e^s \in \tilde U^-$. In $V_{L_H}$, the $V_E \ot V_N$-submodule generated by $e^s$ is $V_E \ot V_{s+N}$, because $s \in \mathfrak h$ has zero component along $E$, and the $B$-submodule generated by $e^s$ is therefore $\Omega_S \ot V_{s+N} = M_s$, the vector $e^s = \mathbf 1 \ot e^s$ being a vacuum vector for $\Omega_S \subset V_E$. By Lemma \ref{lem:fixed-points}(3), $\tilde U^- \cong \Ind(M_s) = \bigoplus_{\gamma \in P} M_{\gamma+s}$ as $C$-modules, hence as $B$-modules.

\emph{Step 3: the two extensions of $B$.} By \eqref{eq:DJX-fixed-point}, Lemma \ref{lem:fixed-points}(1) and Step 2, $\tilde U$ and $S^\natural_{\psi_N}$ are simple vertex operator algebras containing $B$ with the same conformal vector, and both decompose as $\bigoplus_{x \in L_{\mathrm I}'/N} M_x$ as $B$-modules, where $L_{\mathrm I}'/N = P \oplus \{0,s\}$ by Lemma \ref{lem:common-sublattice}(3). The summands are pairwise non-isomorphic invertible $B$-modules with fusion $M_x \boxtimes M_y \cong M_{x+y}$ and multiplicities one. As in the proof of Lemma \ref{lem:fixed-points}(2), the fusion rules of $B$ give the $(L_{\mathrm I}'/N)$-gradings and the hypotheses of \cite[Proposition 5.3]{DoM04}, which gives an isomorphism $\tilde U \cong S^\natural_{\psi_N}$ acting by a nonzero scalar on each summand $M_x$ and by the identity on $B$. Finally $S^\natural_{\psi_N} \cong S^\natural$ by Theorem \ref{thm:twin-S}(3).
\end{proof}

\begin{cor}
\label{cor:twin-consequences}
Assume \textup{(I1)}, \textup{(I3)}, \textup{(F1)}--\textup{(F4)}, \textup{(D1)} and \textup{(D2)}. Let $U$ be the algebra of Huang, isomorphic to the algebra $U^{\mathrm R}$ of Riesen by \textup{(D1)}.
\begin{enumerate}
\item $U$ is a simple, self-contragredient, positive, strongly rational vertex operator algebra of CFT type and central charge eight, and $\Rep(U)$ is ribbon equivalent to $\Rep(S)$. In particular $\Rep(U)$ is a pseudo-unitary modular tensor category of rank twelve, with the twists of (F2), the dimensions of Lemma \ref{lem:dimensions} and no invertible object other than the unit by Lemma \ref{lem:no-invertible}.
\item The Heisenberg commutant of $U$ for the Cartan subalgebra $\Cbb H_L \oplus \Cbb H_R$ is isomorphic to $\Omega_S$, its associated lattice is $L_{\mathrm I}$, and $U_1 = \mathfrak{sl}_2 \oplus \mathfrak{sl}_2$ with currents of levels $(1,39)$, so that $\dim U_1 = 6$ and $U \not\cong S$.
\item $S$ and $U$ belong to the same hyperbolic genus, every strongly rational vertex operator algebra in this genus is isomorphic to $S$ or to $U$, and $S$, $U$ are the only algebras of type $S$ with Heisenberg commutant isomorphic to $\Omega_S$.
\item $S^\natural$ embeds in $V_{L_H} \cong V'$ as the algebra $\tilde U$, and $\tilde U \cap V_{L_0} = C = S \cap V_{L_0}$ under the identification of the two copies of $V_{L_0}$ used in Theorem \ref{thm:identification}.
\end{enumerate}
\end{cor}

\begin{proof}
Parts 1, 3 and the first two assertions of part 2 follow from Theorems \ref{thm:identification} and \ref{thm:twin-S}, since $U \cong S^\natural$ is of type $S$, and from Lemmas \ref{lem:dimensions} and \ref{lem:no-invertible} applied to $\Rep(U) \cong \Rep(S)$. For $U_1$, the Lie algebra $U_1$ is reductive \cite[Theorem 1]{DoM04}, and $\Cbb H_L \oplus \Cbb H_R$ is a Cartan subalgebra of $U_1$, so that $U_1$ has rank two. Indeed, by Theorem \ref{thm:twin-S}(1), $\mathfrak h^\natural$ is a Cartan subalgebra of $(S^\natural_{\psi_N})_1$, and it is the Heisenberg space of $V_N \subset B$. The isomorphism $\tilde U \cong S^\natural_{\psi_N}$ of the proof of Theorem \ref{thm:identification} is the identity on $B$, hence maps the Heisenberg space of $V_N \subset \tilde U$, which is $\hat\Phi(\Cbb H_L \oplus \Cbb H_R)$ since $R_*\ell = s$ and $R_*h = t$, onto $\mathfrak h^\natural$. It contains the Lie algebra $\mathfrak{sl}_2 \oplus \mathfrak{sl}_2$ spanned by the six currents of Huang, of rank two, so $U_1$ is semisimple of rank two with Cartan subalgebra $\mathfrak h^\natural$, and its root system $R$ contains the roots $\pm s$ and $\pm\frac{t}{39}$ of the two $\mathfrak{sl}_2$, which are orthogonal with $(s,s) = 2$ and $(\frac t{39},\frac t{39}) = \frac2{39}$ for the invariant form of $U$ on $\mathfrak h^\natural$. On each simple ideal of $U_1$ this form is a nonzero multiple of the Killing form, so the ratio of the squared lengths of two roots of the same simple ideal is $1$, $2$, $\frac12$, $3$ or $\frac13$. The irreducible root systems of rank two, $A_2$, $B_2$ and $G_2$, contain no pair of orthogonal roots with squared length ratio $39$: in $A_2$ there are no orthogonal roots, in $B_2$ orthogonal roots have the same length, and in $G_2$ the ratio is $3$. Hence $R = \{\pm s\} \sqcup \{\pm\frac t{39}\}$ is reducible, $U_1 = \mathfrak{sl}_2 \oplus \mathfrak{sl}_2$, and the levels are $\frac{(s,s)}2 = 1$ and $\frac{(t,t)}2 = 39$. Since $S_1$ has levels $(3,13)$ and an isomorphism of vertex operator algebras preserves the weight-one Lie algebras with their invariant forms, $U \not\cong S$, which also follows from the lattices. Part 4 restates Theorem \ref{thm:identification} and \eqref{eq:DJX-fixed-point}.
\end{proof}

\begin{remark}\normalfont
\label{rem:twin-comparison}
Corollary \ref{cor:twin-consequences} identifies $U$ with $S^\natural$ using the inputs of Theorem \ref{thm:identification}. The strong rationality of $U$ follows from the extension theory of \cite{HKL15,CKM22} applied to $\Omega_S \ot V_{L_{\mathrm I}}$. The equality $U_1 = \mathfrak{sl}_2 \oplus \mathfrak{sl}_2$ is obtained from the root system of the twin, without the weight-one computations of \cite{Hua26} used in \cite[Proposition 9.4]{DJX26}. This describes the structure behind \cite[Theorem 7.3]{DJX26}: $C$ is the extension of $B = \Omega_S \ot V_N$ by the thirty-nine simple currents $M_\gamma$, $\gamma \in P$, and its simple currents $\Ind(M_x)$, $x \in \{s,p,s+p\}+N$, are $\tilde U^-$, $S^-$ and $\Ind(M_{s+p})$. The extensions $C \oplus \tilde U^-$ and $C \oplus S^-$ are $\tilde U \cong U$ and $S$, whereas $\Ind(M_{s+p})$ has half-integral weights by \eqref{eq:M-weights} and gives no vertex operator algebra extension of $C$. Thus the two nontrivial bosons $s+N$ and $p+N$ of $N'/N$ correspond to the two algebras of type $S$ with Heisenberg commutant $\Omega_S$.
\end{remark}

We finally recover the structure of $\Rep(C)$ of \cite[Section 9]{DJX26} under our hypotheses, with a proof of the rank $48$ which does not use \cite[Theorem 3.1]{DJX26}.

\begin{prop}
\label{prop:Rep-C}
Assume \textup{(I1)}, \textup{(I3)}, \textup{(F1)}--\textup{(F4)} and \textup{(D2)}. Put $a = \tilde U^- \cong \Ind(M_s)$ and $b = S^- \cong \Ind(M_p)$ in $\Rep(C)$, as in \cite[(9.1)]{DJX26}. In this proposition, $a$ and $b$ do not denote the labels of $Z(\cal H_3)$.
\begin{enumerate}
\item $a$ and $b$ are distinct simple currents of order two of $\Rep(C)$ with $\theta_a = \theta_b = 1$ and $\theta_{ab} = -1$, and the fusion subcategory $\cal T$ generated by $a$ and $b$ is braided equivalent to $Z(\mathrm{Vec}_{\Z_2})$. It is nondegenerate.
\item $\Rep(C) \cong \cal T \boxtimes \cal D$ as ribbon categories, where $\cal D$ is the Müger centralizer of $\cal T$, and $\Rep(U) \cong \cal D \cong \Rep(S)$ as ribbon categories. In particular $\Rep(C) \cong Z(\mathrm{Vec}_{\Z_2}) \boxtimes \Rep(S)$ has rank $48$ and global dimension $4 \cdot 117d^2$.
\end{enumerate}
\end{prop}

\begin{proof}
1. By Lemma \ref{lem:fixed-points}(3) and Theorem \ref{thm:identification}, $a = \Ind(M_s)$ and $b = \Ind(M_p)$ are simple $C$-modules, and induction is a tensor functor \cite[Theorem 1.6 p.~7]{KiO02}, so that $a \boxtimes_C a \cong \Ind(M_{2s}) = \Ind(M_0) = C$, $b \boxtimes_C b \cong C$ and $a \boxtimes_C b \cong \Ind(M_{s+p})$, since $2s = \lambda$ and $2p$ lie in $N$. The twists are $\e^{2\pi\i h}$ for the lowest weights, and all the weights of $\Ind(M_x)$ are congruent to $h_x$ modulo $\Z$ by \eqref{eq:M-weights}, which gives $\theta_a = \theta_b = 1$ and $\theta_{ab} = -1$. Hence $a \neq b$ and the four objects $C$, $a$, $b$, $ab$ are distinct, with fusion group $\Z_2 \times \Z_2$ and quadratic form $(0,0,0,\frac12)$. This metric group is nondegenerate, and a pointed braided fusion category is determined up to braided equivalence by its pre-metric group \cite[Theorem 8.4.12 p.~206]{EGNO15}, so $\cal T \cong Z(\mathrm{Vec}_{\Z_2})$, whose metric group is the hyperbolic plane over $\Z_2$. This is \cite[Lemma 9.1]{DJX26}, whose proof obtains $a^2 \cong b^2 \cong 1$ from orbifold Schur--Weyl duality and $\theta_{ab}$ from a lattice intertwining operator. Here both follow from induction and \eqref{eq:M-weights}.

2. Since $\cal T$ is a nondegenerate fusion subcategory of the nondegenerate braided fusion category $\Rep(C)$, we have $\Rep(C) \cong \cal T \boxtimes \cal D$ with $\cal D$ the centralizer of $\cal T$ \cite[Theorem 8.21.4 p.~245]{EGNO15}. The vertex operator algebras $S = C \oplus b$ and $\tilde U = C \oplus a$ are the extensions of $C$ by the connected étale algebras $A_b = \mathbf 1 \oplus b$ and $A_a = \mathbf 1 \oplus a$ of $\cal T$ \cite[Theorem 3.2]{HKL15}, and $\Rep(S) \cong \Rep(C)^{\loc}_{A_b}$, $\Rep(U) \cong \Rep(C)^{\loc}_{A_a}$ as ribbon categories \cite[Theorem 3.4]{HKL15} \cite[Theorem 3.65]{CKM22}. For $Y \in \cal D$, the induced module $A_b \ot Y$ is local since $Y$ centralizes $b$, and $Y \mapsto A_b \ot Y$ is a braided tensor functor $\cal D \to \Rep(C)^{\loc}_{A_b}$ \cite[Theorem 1.6 p.~7]{KiO02}. It is fully faithful: $\operatorname{Hom}_{A_b}(A_b \ot Y, A_b \ot Y') \cong \operatorname{Hom}(Y, Y' \oplus b \boxtimes Y')$, and $\operatorname{Hom}(Y, b \boxtimes Y') = 0$ since $Y \in \mathbf 1 \boxtimes \cal D$ and $b \boxtimes Y' \in b \boxtimes \cal D$ have no common simple summand in $\cal T \boxtimes \cal D$. The global dimension of $\Rep(C)^{\loc}_{A_b}$ is $\dim\Rep(C)/\dim(A_b)^2 = \dim\cal T\dim\cal D/4 = \dim\cal D$ \cite[Corollary 3.32]{DMNO13}, so the fully faithful tensor functor is an equivalence \cite[Proposition 6.3.3 p.~125]{EGNO15}. The same argument applies to $A_a$. The twists agree because both sides are the categories of modules of the regular vertex operator algebras $S$, $U$ with the ribbon structure given by $\e^{2\pi\i L(0)}$, and all inclusions preserve the conformal vector. The rank of $\Rep(S)$ is twelve by (F2), and its global dimension is $(\operatorname{FPdim}\cal F)^2 = (9d+6)^2 = 117d^2$, since $\Rep(S) \cong Z(\cal F)$ with $\cal F$ pseudo-unitary by (F4) \cite[Theorem 2.15]{ENO05}, see Section \ref{sec:category}.
\end{proof}


%%%%%%%%%%%%%%%%%%%%%%%%%%%%%%%%%%%%%%%%%%%%%%%%%%%%%%%%%%%%%%%%%%%%%%%%%%%%%
\subsection{Consequences of the Haagerup property}
\label{subsec:consequences}

Theorem \ref{thm:HJ-main} supplies $\mathrm{(H_J)}$ and complete unitarity for $S$. The identification $U\cong S^\natural$ from Theorem \ref{thm:identification} then yields the Haagerup property for $U$.

\begin{thm}
\label{thm:consequences}
Assume \textup{(I1)}, \textup{(I3)}, \textup{(F1)}--\textup{(F4)}, \textup{(D1)} and \textup{(D2)}. Theorem \ref{thm:HJ-main} supplies $\mathrm{(H_J)}$.
\begin{enumerate}
\item \emph{Classification.} Up to isomorphism, there are exactly two Haagerup algebras, $S$, with currents of levels $(3,13)$, and $U \cong S^\natural \cong U^{\mathrm R} \cong U^{\mathrm H}$, with currents of levels $(1,39)$. They have isomorphic Heisenberg commutants and form a single hyperbolic genus, and every strongly rational vertex operator algebra in this genus is isomorphic to $S$ or to $U$. Every Haagerup algebra in case \textup{(II)} of \cite[Theorem 8.5]{Arh26} is isomorphic to $S$, and every Haagerup algebra in case \textup{(I)} is isomorphic to $U$.
\item \emph{Complete unitarity.} $S$, $U$ and $\Omega_S$ are completely unitary, and they satisfy Condition I$^+$ of \cite[Definition 6.17]{Gui25} for every structure of unitary vertex operator algebra. Every simple $S$-module and every simple $U$-module is unitarizable. The functors $\Pi$ give equivalences of unitary modular tensor categories $\Rep^u(S) \cong \Rep^f(\cal A_S) \cong Z(\cal H_3)$ and $\Rep^u(U) \cong \Rep^f(\cal A_U) \cong Z(\cal H_3)$, and the correspondence of \cite[Theorem 4.10 p.~21]{CGGH23} between simple CFT-type extensions of $S$, respectively $U$, and irreducible local extensions of $\cal A_S$, respectively $\cal A_U$, applies.
\item \emph{Labels.} The assertions of Corollary \ref{cor:S-consequences}(1) hold: the assertion \textup{(N)} of \cite[Section 7]{ArhII} holds, the label of $X_\pm$ is the label of $\pi_\pm$, and $\cal F_\pm \cong \cal H_2$ if and only if $\cal F^{\mathrm{net}}_\pm \cong \cal H_2$, if and only if $X_\pm \cong M_a$.
\item \emph{The third extension.} The assertions of Corollary \ref{cor:S-consequences}(2) hold: the third holomorphic extension $B_1$ of $S$ is isomorphic to $V_{E_8}$, it carries a unique unitary structure extending that of $S$, and its conformal net is the third extension of $\cal A_S$ of Theorem \ref{thm:net}(4).
\end{enumerate}
\end{thm}

\begin{proof}
1. Apply the classification in \cite[Theorem 12.4(1),(4)]{Arh26} with $\mathrm{(H_J)}$ supplied by Theorem \ref{thm:HJ-main}. It gives the two classes represented by $S$ and $S^\natural$ and their common Heisenberg commutant. Theorem \ref{thm:identification} and \textup{(D1)} identify the twin with $U^{\mathrm H}$ and $U^{\mathrm R}$, and Corollary \ref{cor:twin-consequences} gives its levels. Theorem \ref{thm:twin-S}(4) gives the assertion about the strongly rational algebras in their hyperbolic genus.

2. By Theorem \ref{thm:HJ-main}, $S$ is completely unitary and satisfies Condition I$^+$, and $\Pi \co \Rep^u(S) \to \Rep^f(\cal A_S)$ is an equivalence of braided $C^*$-tensor categories, hence of unitary modular tensor categories. The passage to $U$ follows the proofs of \cite[Theorem 8.3(1) and Corollary 8.5(1)]{ArhII}, with $S$ in the role of the Haagerup algebra there and its twin $S^\natural \cong U$, and with Theorem \ref{thm:HJ-main} in place of \cite[Theorem 7.17]{ArhII}: the Heisenberg commutant $\Omega_S$ is completely unitary by \cite[Lemma 8.2]{ArhII}, whose hypotheses are the complete unitarity of $S$ and the structure of $S$ over $\Omega_S \ot V_{L_{\mathrm{II}}}$ of Proposition \ref{prop:associated-lattice-S}, the lattice algebra $V_{L_{\mathrm I}}$ is completely unitary \cite[Corollary 3.32 p.~49]{Gui22}, so $\Omega_S \ot V_{L_{\mathrm I}}$ is completely unitary \cite[Proposition 3.31 p.~48]{Gui22}, and its simple CFT-type extension $S^\natural$ with the same conformal vector, Theorem \ref{thm:twin-S}(1), carries a unique unitary structure extending the tensor product structure and is completely unitary \cite[Theorem 4.7 p.~19]{CGGH23}. Condition I$^+$ passes from $S$ to $\Omega_S$ by \cite[Theorem 6.18(3) p.~70]{Gui25}, since, for a $\theta$-stable Cartan subalgebra \cite[Lemma 8.1]{ArhII}, $V_{L_{\mathrm{II}}}$ is a regular unitary subalgebra of $S$ with commutant $\Omega_S$ \cite[proof of Lemma 8.2]{ArhII}, to $\Omega_S \ot V_{L_{\mathrm I}}$ by \cite[Theorem 6.18(1) and Example 6.19(b) pp.~70--71]{Gui25}, and to its unitary extension $S^\natural$ by \cite[Theorem 6.18(2) p.~70]{Gui25}, as in the proof of \cite[Corollary 8.5(1)]{ArhII}. Condition I$^+$ includes the complete rationality of $\cal A_U$ and the essential surjectivity of $\Pi$ onto $\Rep^f(\cal A_U)$ \cite[Theorem 6.14 pp.~68--69]{Gui25}. The braided equivalence $\Rep^u(U) \cong \Rep(U) \cong \Rep(S) \cong Z(\cal H_3)$ is naturally isomorphic to a unitary braided equivalence \cite[Theorem 5]{Reu23}, which gives $\Rep^u(U) \cong Z(\cal H_3)$ as unitary modular tensor categories. Since two unitary structures on a simple vertex operator algebra are unitarily isomorphic \cite[Proposition 5.19 pp.~40--41]{CKLW18}, these properties hold for every unitary structure. The unitarizability of the simple modules is part of complete unitarity. The hypotheses of \cite[Theorem 4.10 p.~21]{CGGH23}, namely that $S$ is regular, completely unitary, of CFT type, strongly local, with $\cal A_S$ completely rational and $\Rep^u(S) \cong \Rep^f(\cal A_S)$ as unitary modular tensor categories, are satisfied, and likewise for $U$, whose strong locality follows from Condition I$^+$ \cite[Remark 4.16 p.~41]{Gui25}. This is \cite[Corollary 7.19(3)]{ArhII} for $S$.

3--4. These are the statements of Corollary \ref{cor:S-consequences}.
\end{proof}

\begin{remark}\normalfont
\label{rem:other-consequences}
By Theorem \ref{thm:HJ-main}, the statements of \cite[Section 12]{Arh26} not covered by Theorem \ref{thm:twin-S}, namely \cite[Theorem 12.4 and Remark 12.6]{Arh26}, hold as stated, see Remark \ref{rem:twin-not-covered}, and \cite[Theorem 8.6(2)]{ArhII}, the analogue of (N) for a Haagerup algebra embedded in $V'$ as a unitary subalgebra containing its conformal vector, applies to $U = U^{\mathrm H} \subset V_{L^{\mathrm H}}$, which contains the conformal vector by \cite[(4.9)]{DJX26} and satisfies these hypotheses for the unitary structure of \cite[Section 4.4]{DJX26} by the identity $\theta_H(U) = U$ of \cite[(6.15)]{DJX26}, verified in Section \ref{subsec:DJX-verification}. The assertions of Theorem \ref{thm:net}(4) on the Haagerup subfactor, the hypergroup and the near-group category, which are \cite[Corollaries 7.4 and 7.6 and Remark 7.8]{ArhII}, do not depend on $\mathrm{(H_J)}$.
\end{remark}


%%%%%%%%%%%%%%%%%%%%%%%%%%%%%%%%%%%%%%%%%%%%%%%%%%%%%%%%%%%%%%%%%%%%%%%%%%%%%
\section{Comparison with the work of Dong, Jiao and Xu}
\label{sec:DJX}

The preprint \cite{DJX26} appeared while the present work was in progress. We describe what it proves, which of its statements we use and how we verified them, and what is new here.


%%%%%%%%%%%%%%%%%%%%%%%%%%%%%%%%%%%%%%%%%%%%%%%%%%%%%%%%%%%%%%%%%%%%%%%%%%%%%
\subsection{\texorpdfstring{The results of \cite{DJX26}}{The results of Dong, Jiao and Xu}}

Let $U$ be the vertex subalgebra of $V' = V_{E_6 \oplus A_2}$ generated by the explicit vectors of \cite[Section 2]{Hua26}, with currents of levels $(1,39)$, and let $V$ be the vertex subalgebra generated by the nine vectors of \cite[Section 3]{Jia26}, with currents of levels $(3,13)$, both written in the root basis and with the cocycle of \cite[Section 4.1]{DJX26}. The main results of \cite{DJX26} are the following.
\begin{enumerate}
\item The algebra of Riesen \cite{Rie26} is isomorphic to $U$ \cite[Proposition 4.1]{DJX26}, through an explicit identification of generators and of the two commuting Virasoro vectors $e$, $f$ of central charges $\frac{70}{31}$ and $\frac{2400}{1271}$.
\item The involutions $\tau$ and $\sigma$ of \textup{(D2)} satisfy $\Phi(U^\tau) = V^\sigma$, where $\Phi$ is the lift of an isometry between two rank-eight lattices $L_H$ and $L_J$ with the common index-two sublattice $L_0$ \cite[Sections 5--6 and Theorem 7.3]{DJX26}. The proof is a finite set of mode identities together with a lattice-shift lemma \cite[Lemma 7.1]{DJX26}.
\item A positive conformal-vector identity $\omega = \omega_L+\frac{3403}{4563}\omega_R+\frac{25}{6}\Pi_R\eta_3+\frac{25}{26}\Pi_R\eta_{13}$ holds in $U$ \cite[Proposition 8.1]{DJX26}, where $\eta_3$, $\eta_{13}$ are the conformal vectors of the parafermion algebras $K(\mathfrak{sl}_2,3)$, $K(\mathfrak{sl}_2,13)$ transported from $V^\sigma$ and $\Pi_R$ is an average over the compact group of the level-$39$ currents. With a rationality criterion \cite[Theorem 3.1]{DJX26}, this gives the strong rationality of $U$, then of $U^\tau$ by \cite{CaM16}, then of $V$ by \cite[Example 4.2]{McR22} \cite[Corollary 8.2]{DJX26}.
\item $\Rep(U^\tau) \cong T \boxtimes D$ with $T \cong Z(\mathrm{Vec}_{\Z_2})$ and $\Rep(U) \cong D \cong \Rep(V)$ as ribbon categories \cite[Lemma 9.1 and Proposition 9.2]{DJX26}, so that $\Rep(U^\tau) \cong Z(\mathrm{Vec}_{\Z_2}) \boxtimes Z(\cal H_3)$, of rank $48$, if $\Rep(U)$ or $\Rep(V)$ is ribbon equivalent to $Z(\cal H_3)$ \cite[Corollary 9.3]{DJX26}, and $U \not\cong V$ \cite[Proposition 9.4]{DJX26}.
\item The conformal nets of $U$, $V$ and $U^\tau$ are completely rational, and $\cal A_U^\tau \cong \cal A_{U^\tau} \cong \cal A_V^\sigma$ \cite[Theorem 10.1 and Corollary 10.2]{DJX26}.
\end{enumerate}
\begin{remark}\normalfont
\label{rem:net-DJX}
The complete rationality of $\cal A_S$ is also proved in \cite[Theorem 10.1]{DJX26}, by the same chain of results of Tener and Longo, with the strong rationality of $S$ taken from \cite[Corollary 8.2]{DJX26} instead of (F1). The proof there applies \cite[Corollary 4.24]{Ten24} as stated, whereas the proof of \cite[Corollary 4.24]{Ten24} uses a unitary structure on $M \boxtimes \ovl M$, which is supplied by \cite[Proposition 5.2]{ArhII}, see Step 2 of the proof of Theorem \ref{thm:net}. Parts 2--4 of Theorem \ref{thm:net} go beyond \cite{DJX26}, whose Remark 10.3 points out that complete rationality alone does not identify the representation category of the net.
\end{remark}


%%%%%%%%%%%%%%%%%%%%%%%%%%%%%%%%%%%%%%%%%%%%%%%%%%%%%%%%%%%%%%%%%%%%%%%%%%%%%
\subsection{Independent verification of the finite identities}
\label{subsec:DJX-verification}

We use \cite[Theorem 7.3]{DJX26} in Section \ref{sec:genus}. Its proof rests on the explicit vectors, on finite mode identities and on \cite[Lemma 7.1]{DJX26}. We verified all finite identities with an exact implementation of the lattice vertex operator algebra $V_L$, $L = E_6 \oplus A_2$, in the basis and with the cocycle of \cite[Section 4.1]{DJX26}, over the field $\K = \Q(\sqrt{13},\sqrt{-3})$, written independently of the implementation used in \cite{Arh26} (certificate A4 of Section \ref{sec:certificates}). The checks are the following, with the number of verified identities in parentheses.
\begin{enumerate}
\item The lattice data of \cite[Sections 4.1--4.4]{DJX26} and the $\mathfrak{sl}_2$-triples of levels $1$, $39$, $3$, $13$, the highest weight vectors $u$, $u_{2,6}$, $u_{1,7}$, the unitary structures and the adjoints ($79$).
\item The lattices $L_0$, $L_H$, the vectors $s$, $t$, the isometry $T$ and the characters of $\sigma$ and $\tau$ of \cite[Section 5]{DJX26} ($40$). We note that $t$ has half-integral coordinates in the root basis, so that $t \in s+L_0 \subset L_H$ and $t \notin L$, and the sentence ``$t$ is a lattice vector'' of \cite[Table 1]{DJX26} has to be read with $L_H$ in place of $L$.
\item The isometries $R$, $R_*$, the lift $\Phi$ of \cite[(6.13)]{DJX26}, the identities \cite[(6.14) and (6.15)]{DJX26}, the first of which is \cite[Proposition 6.2]{DJX26}, and $\Phi\theta_H = \theta_J\Phi$ ($56$).
\item The identities \cite[(7.6), (7.7), (7.9), (7.10), (7.11)]{DJX26} and the membership conditions of \cite[Proposition 7.2]{DJX26} ($54$). The identity (7.9) is checked as an equality of vectors. The vector $Y$ of \cite[(7.8)]{DJX26} is $(Z_0)_{(1)}X+\frac43X$, as in \cite[Appendix A.2]{DJX26}, and not $(Z_0)_{(1)}X+\frac34X$ as printed: only the former gives $(Y,Y) = 104$ and (7.9).
\item The parafermion conformal vectors, the Casimir projections, the vectors $e$, $f$ of Riesen and the identities \cite[(4.19)--(4.22), (8.1), (8.2), (8.4)]{DJX26} of Proposition 8.1, as equalities of vectors ($56$), and the data of \cite[Proposition 4.1 and Appendix A.3]{DJX26} for the identification of the algebra of Riesen with $U$ ($35$).
\end{enumerate}
These checks are grouped into six scripts. We also read the proof of \cite[Lemma 7.1]{DJX26} step by step and recomputed the scalar factor of its product formula. The hypothesis ``every generator of $A$ lies in $W+S_uW$'' is to be read with $W+S_uW+S_{-u}W$, since some generators lie in $S_{-u}W$ and the equality $S_{-u}W = S_uW$ is only obtained in the course of the proof. The lemma is applied in $V_{L_H}$ with a bimultiplicative cocycle whose restriction to $L_0$ is cohomologous, but not equal, to the restriction of the cocycle of $V_L$, which is accounted for by the cocycle-compatible identification of $V_{L_0}$ in \cite[(5.7)]{DJX26}. The proof of \cite[Theorem 7.3]{DJX26} uses no property of $U$, $V$ or of the algebra of Riesen beyond their definitions: neither rationality, nor characters, nor the generation statements of \cite{Hua26,Jia26,Rie26}. We therefore regard \cite[Proposition 4.1, Theorem 7.3 and Proposition 8.1]{DJX26} as established.


%%%%%%%%%%%%%%%%%%%%%%%%%%%%%%%%%%%%%%%%%%%%%%%%%%%%%%%%%%%%%%%%%%%%%%%%%%%%%
\subsection{What is new here}

\begin{enumerate}
\item The strong rationality of $V \cong S$ is proved in Corollary \ref{zhu:cor:rational} through the semisimplicity of Zhu's algebra and in \cite[Corollary 8.2]{DJX26} through the positive conformal-vector identity and \cite[Theorem 3.1]{DJX26}. The two proofs are independent. The identity \cite[(8.2)]{DJX26} is verified above, and we have not checked the proof of \cite[Theorem 3.1]{DJX26}, which is a general statement. The classification of the simple modules, (F2), is not obtained in \cite{DJX26}.
\item Theorem \ref{thm:intro-twin} identifies the algebras $U$ and $S^\natural$, not only their module categories as in \cite[Proposition 9.2]{DJX26}, and it explains the common fixed point of \cite[Theorem 7.3]{DJX26} as the lattice realization of the twin construction of \cite[Section 12]{Arh26}, through the common index-two sublattice $N$ of Lemma \ref{lem:common-sublattice}. The rank of $\Rep(U^\tau)$ is $48$ by (F2) and Proposition \ref{prop:Rep-C}.
\item Theorem \ref{thm:net} identifies $\Rep(\cal A_S)$ with $Z(\cal H_3)$ as a unitary modular tensor category and realizes the Haagerup subfactor from $\cal A_S$, whereas \cite[Theorem 10.1]{DJX26} proves the complete rationality of the nets. The proof of \cite[Theorem 10.1]{DJX26} applies \cite[Corollary 4.24]{Ten24} as stated, whose proof uses a unitary structure on the fusion products $M \boxtimes \ovl M$ which is not among its hypotheses, see Remark \ref{rem:net-DJX}. This structure is provided by \cite[Proposition 5.2]{ArhII}.
\item Theorem \ref{thm:HJ-main} identifies $\Rep(S)$ with $Z(\cal H_3)$ relative to the explicit inputs (I1), (F1)--(F3), and the twin identification then gives the same conclusion for $U$. The additional ingredient is the trace class theorem applied to the diagonal net, followed by reconstruction and the coset comparison.
\end{enumerate}


%%%%%%%%%%%%%%%%%%%%%%%%%%%%%%%%%%%%%%%%%%%%%%%%%%%%%%%%%%%%%%%%%%%%%%%%%%%%%
\section{The algebraic route and the equivalent criteria}
\label{sec:algebraic-route}


The Lagrangian algebras of $\Rep(S)$ arising from the lattice extensions give a separate algebraic description of the category. This description is used for the twin theorem and for the comparison with the algebraic construction of the diagonal extension.

\begin{thm}
\label{thm:algebraic-F4}
\textup{(F4)}. Each $L_\pm$ carries the structure of a Lagrangian algebra in $\Rep(S)$. The category $\Rep(S)$ is a pseudo-unitary modular tensor category, braided equivalent to the Drinfeld center $Z(\cal F_\pm)$ of the fusion category $\cal F_\pm = \Rep(S)_{L_\pm}$. The category $\cal F_\pm$ is pseudo-unitary, has six simple objects, three of which are invertible, with the fusion rules \eqref{eq:haagerup-rules} of the Haagerup categories $\cal H_2$ and $\cal H_3$, and its pointed part is $\mathrm{Vec}_{\Z_3}$ with trivial associator.
\end{thm}

\begin{proof}
Proposition \ref{zhu:prop:pseudo} and Theorem \ref{zhu:thm:categorical}(1)--(3) give the assertions.
\end{proof}


%%%%%%%%%%%%%%%%%%%%%%%%%%%%%%%%%%%%%%%%%%%%%%%%%%%%%%%%%%%%%%%%%%%%%%%%%%%%%
\subsection{Consequences of (F4) for the representation category}
\label{sec:category}

The following properties of $\Rep(S)$ follow from (F1)--(F3) of Theorem \ref{thm:zhu-inputs} and (F4) of Theorem \ref{thm:algebraic-F4}. They are used in the twin theorem and the algebraic construction of the diagonal extension. In this subsection, $\cal F = \cal F_+$ is the pseudo-unitary fusion category of Theorem \ref{thm:algebraic-F4}, with simple objects $1$, $\alpha$, $\alpha^2$, $\rho$, $\alpha\rho$, $\alpha^2\rho$ and fusion rules
\begin{equation}
\label{eq:haagerup-rules}
\rho\alpha \cong \alpha^2\rho,
\quad
\rho^2 \cong 1 \oplus \rho \oplus \alpha\rho \oplus \alpha^2\rho,
\end{equation}
and $\Rep(S)$ is braided equivalent to the Drinfeld center $Z(\cal F)$. We write $d = \frac{3+\sqrt{13}}{2}$, so that $d^2 = 3d+1$ and $\operatorname{FPdim}\cal F = 3+3d^2 = 9d+6$. We do not assume that $\cal F$ is unitary.


%%%%%%%%%%%%%%%%%%%%%%%%%%%%%%%%%%%%%%%%%%%%%%%%%%%%%%%%%%%%%%%%%%%%%%%%%%%%%
\subsubsection{Invertible objects}

\begin{lemma}
\label{lem:no-invertible}
Let $\cal F$ be a fusion category with the fusion rules \eqref{eq:haagerup-rules}. Then the universal grading group of $\cal F$ is trivial, and the unit is the only invertible object of $Z(\cal F)$ up to isomorphism.
\end{lemma}

\begin{proof}
Let $U(\cal F)$ be the universal grading group of $\cal F$ \cite[Definition 4.14.2]{EGNO15}, and write $g_X$ for the degree of a simple object $X$, using multiplicative notation since the grading group need not be abelian. The summands $1$ and $\rho$ of $\rho\ot\rho$ give $g_\rho^2 = 1$ and $g_\rho^2 = g_\rho$, hence $g_\rho = 1$. The summand $\alpha\rho$ then gives $g_\alpha = 1$. The simple objects $\alpha$ and $\rho$ generate $\cal F$, so $U(\cal F)$ is trivial. By \cite[Proposition 4.14.3]{EGNO15}, the group of tensor automorphisms of the identity functor of $\cal F$ is the character group of $U(\cal F)$, hence trivial.

Let $(X,\gamma)$ be an invertible object of $Z(\cal F)$, where $\gamma_Y \co X \ot Y \to Y \ot X$ is a half-braiding. The forgetful functor $Z(\cal F) \to \cal F$ is a tensor functor \cite[Remark 7.13.4]{EGNO15}, so $X$ is invertible in $\cal F$, that is, $X \in \{1,\alpha,\alpha^2\}$. If $X = \alpha$, the isomorphism $\gamma_\rho$ gives $\alpha\rho \cong \rho\alpha \cong \alpha^2\rho$, which contradicts the distinctness of the six simple objects. The case $X = \alpha^2$ is excluded in the same way. If $X = 1$, then $\gamma$ is a tensor automorphism of the identity functor of $\cal F$, hence trivial, and $(X,\gamma)$ is the unit.
\end{proof}


%%%%%%%%%%%%%%%%%%%%%%%%%%%%%%%%%%%%%%%%%%%%%%%%%%%%%%%%%%%%%%%%%%%%%%%%%%%%%
\subsubsection{Dimensions and self-duality}

We use the labels $z \in \{0,b,a,c_0,c_1,c_2,d_1,\ldots,d_6\}$ of Theorem \ref{thm:zhu-inputs} for the simple $S$-modules $M_z$, with $M_0 = S$, and we write $h_z$ for their lowest conformal weights and $\dim M_z$ for their categorical dimensions, which are their Frobenius--Perron dimensions by (F4). Recall from (F3) that $V'|_S = S \oplus M_b$, that the two lattice extensions $H_\pm \cong V_{E_8}$ of $S$ restrict as
\begin{equation}
\label{eq:lagrangian-restriction}
H_\pm|_S \cong S \oplus M_b \oplus 2X_\pm,
\end{equation}
where the signs refer to the fixed lattice coordinates of (F3). The unordered pair of classes of $X_\pm$ is that of $M_a$ and $M_{c_0}$. By Theorem \ref{thm:HJ-main} and Corollary \ref{cor:S-consequences}(1), the intrinsic labels $a,c_0$ refer to the center in the convention of \cite{EvG11,ArhII}. The restrictions $L_\pm = H_\pm|_S$ are Lagrangian algebras of $\Rep(S)$ with $\cal F_\pm = \Rep(S)_{L_\pm}$.

\begin{lemma}
\label{lem:dimensions}
We have $\dim M_b = 3d+1$ and $\dim M_a = \dim M_{c_0} = \dim M_{c_1} = \dim M_{c_2} = 3d+2$. Every simple $S$-module is self-dual.
\end{lemma}

\begin{proof}
By Lemma \ref{zhu:lem:qdim-lattice}, the quantum dimensions of the restrictions to $S$ of the nine simple $V'$-modules are all equal to $1+\dim M_b$, so that $\dim M_{c_1} = \dim M_{c_2} = \dim M_a = \dim M_{c_0} = 1+\dim M_b$ by (F3). Since $L_+$ is a Lagrangian algebra, $(\operatorname{FPdim} L_+)^2 = \operatorname{FPdim}\Rep(S) = (\operatorname{FPdim}\cal F)^2$ \cite[Corollary 3.32 and Section 4]{DMNO13}, hence $1+\dim M_b+2(1+\dim M_b) = 9d+6$ and $\dim M_b = 3d+1$.

Let $z \mapsto z^*$ denote the duality on the simple $S$-modules. The dual of a simple object has the same twist and the same dimension, and by (F2) the twists $\theta_z = \e^{2\pi\i h_z}$ take the value $\e^{2\pi\i/3}$ only at $c_1$, the value $\e^{4\pi\i/3}$ only at $c_2$, and six distinct values $\e^{2\pi\i k/13}$, $k \in \{2,5,6,7,8,11\}$, at $d_1, \ldots, d_6$. Hence $c_1$, $c_2$ and the six $d_l$ are self-dual, and the duality permutes $\{b,a,c_0\}$, since these are the simple objects with twist $1$ other than the unit. A connected étale algebra in a fusion category is self-dual as an object: it is a Frobenius algebra by \cite[Section 2.2]{DMNO13} and \cite[Remark 3.4]{DMNO13}. Hence $L_\pm \cong L_\pm^*$ as objects. In each of $L_\pm$, the nonvacuum summand of multiplicity two is $X_\pm$, while the only nonvacuum summand of multiplicity one is $M_b$. Thus $X_\pm$ and $M_b$ are self-dual, and the assertion follows for both $M_a$ and $M_{c_0}$.
\end{proof}

\begin{remark}\normalfont
\label{rem:d-dimensions}
The dimensions of the six modules $M_{d_l}$ are not needed in this paper. By Theorem \ref{thm:HJ-main} they are equal to $3d$, see \cite[Section 2.2]{ArhII}, and the identity $1 + (3d+1)^2 + 4(3d+2)^2 + 6(3d)^2 = (9d+6)^2$ is consistent with Lemma \ref{lem:dimensions}.
\end{remark}


%%%%%%%%%%%%%%%%%%%%%%%%%%%%%%%%%%%%%%%%%%%%%%%%%%%%%%%%%%%%%%%%%%%%%%%%%%%%%
\subsubsection{Complex conjugation and the reverse category}
\label{subsec:conjugation}

The construction of the diagonal holomorphic extension of $S \ot S$ in \cite[Proposition 7.12]{ArhII} uses a braided equivalence $\Rep(S) \to \Rep(S)^{\mathrm{rev}}$, obtained there from the identification $\Rep(S) \cong Z(\cal H_3)$ and from the self-conjugacy of the even parts of the Haagerup subfactor \cite{JMS14}. Without this identification, the existence of such an equivalence is not automatic. The Drinfeld center of a fusion category need not be braided equivalent to its reverse: for the pointed category $\mathrm{Vec}_{\Z_3}^\omega$ with a suitable generator $[\omega]$ of $H^3(\Z_3,\mathbb T)$, the center $Z(\mathrm{Vec}_{\Z_3}^\omega)$ is a pointed modular category of rank nine whose twists are $1$ three times and $\e^{2\pi\i k/9}$ twice for each $k \in \{1,4,7\}$ \cite[Section 8.5]{EGNO15}. This multiset is not stable under complex conjugation, so that $Z(\mathrm{Vec}_{\Z_3}^\omega)$ and its reverse are not braided equivalent. The unitary structure of $S$ gives a weaker statement, which we record for later use.

For a $\Cbb$-linear category $\cal C$, we denote by $\ovl{\cal C}$ the complex conjugate category, with the same objects and morphisms and the conjugate scalar multiplication on the morphism spaces. If $\cal C$ is a ribbon fusion category, then so is $\ovl{\cal C}$, with the same braiding and twist morphisms, and the twist of a simple object of $\ovl{\cal C}$ is the complex conjugate of its twist in $\cal C$.

Recall that $S$ carries a structure of unitary vertex operator algebra \cite[Theorem 10.10]{Arh26}, with PCT operator $\theta \co S \to S$, an antilinear involution satisfying $\theta(a_{(n)}b) = \theta(a)_{(n)}\theta(b)$ and $\theta(\omega) = \omega$ \cite[Definition 2.1]{DoL14}. For an $S$-module $(M,Y_M)$, let $\overline M$ be the conjugate vector space, with its canonical antilinear map $m\mapsto\overline m$. Define $J(M)$ by
\[
a_{(n)}^{J(M)}\overline m = \overline{(\theta a)_{(n)}^M m}.
\]
Each displayed mode is complex linear in $\overline m$, and the assignment is complex linear in $a$. The original modes on $M$ are complex linear, and the conjugation maps and $\theta$ are antilinear. The module identities follow by conjugation, and the conformal weights are unchanged. On morphisms put $J(f)\overline m = \overline{f(m)}$. This is an antilinear functor to $\Rep(S)$, hence a linear functor to its conjugate category. We denote this latter functor by $\pi$ below.

\begin{prop}
\label{prop:conjugation}
The assignment $M \mapsto \pi(M)$ extends to an equivalence of ribbon categories $\pi \co \Rep(S) \to \ovl{\Rep(S)}^{\mathrm{rev}}$. In particular $\ovl{\Rep(S)}$ is braided equivalent to $\Rep(S)^{\mathrm{rev}}$. On the simple modules, $\pi(M_z) \cong M_z$ for every $z$.
\end{prop}

\begin{proof}
Let $M$, $N$, $P$ be $S$-modules and let $\cal Y$ be an intertwining operator of type $\binom{P}{M N}$, so that $\cal Y(m,x)n = \sum_{h \in \Cbb} m_{(h)}n x^{-h-1}$ with $m_{(h)}n \in P$. The conformal weights of the $S$-modules are rational by (F1) and (F2). The conjugated coefficients, defined by $\overline{\cal Y}(\overline m,x)\overline n = \sum_h \overline{m_{(h)}n}x^{-h-1}$, define an intertwining operator $\ovl{\cal Y}$ of type $\binom{\pi(P)}{\pi(M) \pi(N)}$, since the Jacobi identity and the $L(-1)$-derivative property have rational coefficients and since the modes of $S$ on $\pi(M)$ are those of $\theta(S) = S$. The map $\cal Y \mapsto \ovl{\cal Y}$ is a bijection between the spaces of intertwining operators, which is antilinear for the original scalar multiplications, hence linear onto the conjugate space. It is compatible with the P($z$)-tensor products of \cite{HL95,Hua05}, so that the vertex tensor category structure of $\Rep(S)$, which exists by (F1) and \cite[Theorem 4.6]{Hua08}, is transported through $\pi$ to a vertex tensor category structure on $\ovl{\Rep(S)}$, and the universal property of the P($z$)-tensor product identifies $\pi(M) \boxtimes \pi(N)$ with $\pi(M \boxtimes N)$. The associativity isomorphisms of \cite{Hua05} are characterized by the identity
\[
\cal Y_1(m,z_1)\cal Y_2(n,z_2)p = \cal Y_3(\cal Y_4(m,z_1-z_2)n,z_2)p
\]
for real numbers $z_1 > z_2 > z_1-z_2 > 0$, where the powers of $z_1$, $z_2$ and $z_1-z_2$ are the real ones. Both sides are transformed into the corresponding identity for the conjugate operators, so that $\pi$ is a tensor functor. The braiding of \cite{Hua08} is characterized by the skew-symmetry $\Omega(\cal Y)(n,z)m = \e^{zL(-1)}\cal Y(m,\e^{\pi\i}z)n$, in which $\e^{\pi\i}z$ means that each power $z^{h}$ is replaced by $\e^{\pi\i h}z^h$. In the conjugate scalar multiplication of $\pi(P)$, the multiplication by $\e^{\pi\i h}$ is the multiplication by $\e^{-\pi\i h}$ of $P$, since $h$ is real. Hence the braiding of $\Rep(S)$, viewed in $\ovl{\Rep(S)}$, is the inverse of the braiding transported through $\pi$, that is, $\pi$ is a braided tensor functor $\Rep(S) \to \ovl{\Rep(S)}^{\mathrm{rev}}$. Similarly the twist $\e^{2\pi\i L(0)}$ of $\pi(M)$, computed in the conjugate scalar multiplication, is the map $\e^{-2\pi\i L(0)}$ of $M$, which is the twist of $\ovl{\Rep(S)}^{\mathrm{rev}}$, so that $\pi$ is a ribbon functor. It is an equivalence, with quasi-inverse given by the same construction, since $\theta^2 = \mathrm{id}$.

For the action on the labels, $\pi(M)$ is a simple $S$-module with the same conformal weights as $M$, so that $\pi$ fixes the nine modules $M_0$, $M_{c_1}$, $M_{c_2}$, $M_{d_1}, \ldots, M_{d_6}$, whose lowest weights are pairwise distinct, and permutes $\{M_b, M_a, M_{c_0}\}$. By Lemma \ref{lem:dimensions}, $\pi$ fixes $M_b$. Let $\theta'$ be the PCT operator of $V'$ of \cite[Theorem 4.12]{DoL14}, for which $S$ is a unitary subalgebra with $\theta = \theta'|_S$ \cite[Theorem 10.10]{Arh26}. For a simple $V'$-module $W_g$, $g \in \Z_3 \times \Z_3$, the $V'$-module $\ovl{W_g}^{\theta'}$ is isomorphic to $W_{-g}$, since $\theta'$ maps $\e^\lambda$ to a multiple of $\e^{-\lambda}$, and $\pi(W_g|_S) = \ovl{W_g}^{\theta'}|_S = W_{-g}|_S$. By (F3), $W_{(1,1)}|_S \cong W_{(-1,-1)}|_S$ and $W_{(1,-1)}|_S \cong W_{(-1,1)}|_S$, so that $\pi$ fixes $M_a$ and $M_{c_0}$.
\end{proof}

\begin{remark}\normalfont
\label{rem:rev-not-automatic}
Proposition \ref{prop:conjugation} shows that $\Rep(S)^{\mathrm{rev}}$ is braided equivalent to $\ovl{\Rep(S)}$, not to $\Rep(S)$. A braided equivalence $\Rep(S) \cong \Rep(S)^{\mathrm{rev}}$ is equivalent to a braided equivalence $\Rep(S) \cong \ovl{\Rep(S)}$, that is, to the invariance of $\Rep(S)$ under complex conjugation of its structure constants, but the conjugation argument alone does not establish this extra equivalence. We retain the notation (R) for this property in Section \ref{subsec:hypothesis-R}. Theorem \ref{thm:net-diagonal} establishes it using the trace class theorem on the conformal-net side, where the identification of Theorem \ref{thm:net} with the unitary modular tensor category $Z(\cal H_3)$ is available, and Theorem \ref{thm:unitarity-equivalences} shows that it is equivalent to the Haagerup property of $S$.
\end{remark}


%%%%%%%%%%%%%%%%%%%%%%%%%%%%%%%%%%%%%%%%%%%%%%%%%%%%%%%%%%%%%%%%%%%%%%%%%%%%%
\subsection{The hypothesis \textup{(R)} and the action on the labels}
\label{subsec:hypothesis-R}

We consider the following hypothesis on the modular tensor category $\Rep(S)$ of \cite{Hua08}, which exists by (F1).
\[
\textup{(R)}\quad \text{There is a braided equivalence } \pi \co \Rep(S) \to \Rep(S)^{\mathrm{rev}}.
\]
Such a functor is also called a braid-reversed autoequivalence of $\Rep(S)$. Hypothesis (R) holds by Theorem \ref{thm:net-diagonal}(3), and also by Theorem \ref{thm:HJ-main} and Lemma \ref{lem:ZH3-self-reverse}, which recalls \cite[Proposition 7.12(1)]{ArhII}. We keep (R) as a named hypothesis in order to isolate the arguments which use it. The functor $\pi$ of (R) is not the functor $\pi$ of Proposition \ref{prop:conjugation}. By Proposition \ref{prop:conjugation}, (R) is equivalent to a braided equivalence between $\Rep(S)$ and its complex conjugate category. We write $\pi$ also for the induced permutation of the set $\{0,b,a,c_0,c_1,c_2,d_1,\ldots,d_6\}$ of labels of the simple $S$-modules, so that $\pi(M_z) \cong M_{\pi(z)}$.

\begin{lemma}
\label{lem:pi-labels}
Assume \textup{(I1)}, \textup{(F1)}--\textup{(F4)} and \textup{(R)}. Then $\pi$ is an equivalence of ribbon categories, where $\Rep(S)^{\mathrm{rev}}$ carries the twist $\theta^{-1}$ and the spherical structure of $\Rep(S)$. In particular $\dim M_{\pi(z)} = \dim M_z$ and $\theta_{\pi(z)} = \theta_z^{-1}$ for every label $z$. The permutation $\pi$ satisfies
\[
\begin{gathered}
\pi(0) = 0,\quad \pi(b) = b,\quad \pi(c_1) = c_2,\quad \pi(\{a,c_0\}) = \{a,c_0\},\\
\pi(d_1) = d_5,\quad \pi(d_2) = d_3,\quad \pi(d_4) = d_6,
\end{gathered}
\]
and it is an involution.
\end{lemma}

\begin{proof}
By (F4), $\Rep(S)$ is pseudo-unitary, and so is $\Rep(S)^{\mathrm{rev}}$, which has the same underlying fusion category. A pseudo-unitary fusion category has a unique spherical structure for which the categorical dimensions are the Frobenius--Perron dimensions \cite[Proposition 8.23 p.~34]{ENO05}. A tensor equivalence induces an isomorphism of Grothendieck rings mapping simple objects to simple objects, hence preserves Frobenius--Perron dimensions \cite[Section 3.3]{EGNO15}, and it transports spherical structures, so that $\pi$ transports the spherical structure of $\Rep(S)$ to the one of $\Rep(S)^{\mathrm{rev}}$. The twist of a braided category with a spherical structure is determined by the braiding and the spherical structure, and the twist of $\Rep(S)^{\mathrm{rev}}$ is $\theta^{-1}$. Hence $\pi$ is a ribbon equivalence, $\dim M_{\pi(z)} = \dim M_z$ and $\theta_{\pi(z)} = \theta_z^{-1}$.

By (F2), the twists $\theta_z = \e^{2\pi\i h_z}$ are $1$ for $z \in \{0,b,a,c_0\}$, $\e^{2\pi\i/3}$ for $c_1$, $\e^{4\pi\i/3}$ for $c_2$, and $\e^{2\pi\i k/13}$ with $k = 6, 11, 2, 5, 7, 8$ for $d_1, \ldots, d_6$, so that the twists of $c_1$, $c_2$, $d_1, \ldots, d_6$ are pairwise distinct and distinct from $1$. The relation $\theta_{\pi(z)} = \theta_z^{-1}$ gives $\pi(c_1) = c_2$, $\pi(c_2) = c_1$, and $\pi(d_l) = d_{l'}$ with $k' \equiv -k \bmod 13$, that is, $6 \leftrightarrow 7$, $11 \leftrightarrow 2$ and $5 \leftrightarrow 8$, which are the stated pairs. The unit is fixed. The simple objects with trivial twist are $0$, $b$, $a$, $c_0$, of dimensions $1$, $3d+1$, $3d+2$, $3d+2$ by Lemma \ref{lem:dimensions}, so that $\pi(b) = b$ and $\pi$ permutes $\{a,c_0\}$. Each orbit of $\pi$ on the labels has one or two elements, so the permutation is an involution.
\end{proof}

Lemma \ref{lem:pi-labels} allows either permutation on $\{a,c_0\}$, and the following construction of the diagonal extension applies in both cases. By Theorem \ref{thm:HJ-main}, Lemma \ref{lem:ZH3-self-reverse} and the triviality of the Brauer--Picard group of $\cal H_3$ \cite[Theorem 1.5 p.~5]{GrS12}, \cite[Theorem 1.1 p.~3]{ENO10}, a braided equivalence $\Rep(S) \to \Rep(S)^{\mathrm{rev}}$ is unique up to natural isomorphism and fixes $a$ and $c_0$. This is not used below.


%%%%%%%%%%%%%%%%%%%%%%%%%%%%%%%%%%%%%%%%%%%%%%%%%%%%%%%%%%%%%%%%%%%%%%%%%%%%%
\subsection{Diagonal holomorphic extensions of \texorpdfstring{$S \ot S$}{S x S}}
\label{subsec:diagonal-extension}

\begin{defi}
\label{def:diagonal-extension}
A \emph{diagonal holomorphic extension} of $S \ot S$ is a simple, strongly rational, holomorphic vertex operator algebra $W$ containing $S \ot S$ as a vertex subalgebra with the same conformal vector, such that
\[
\operatorname{Com}_W(S \ot \mathbf 1) = \mathbf 1 \ot S,
\quad
\operatorname{Com}_W(\mathbf 1 \ot S) = S \ot \mathbf 1.
\]
\end{defi}

Here holomorphic means that every $W$-module is a direct sum of copies of $W$. The central charge of such a $W$ is $16$. By (F1) and \cite[Theorem 4.5]{ABD04}, $S$ and $S \ot S$ are regular, and $\Rep(S \ot S)$ is braided equivalent to $\Rep(S) \boxtimes \Rep(S)$ \cite[Theorem 5.5 p.~48]{CKM22}, whose hypotheses, that both categories are locally finite and one of them is semisimple, hold since $\Rep(S)$ is a fusion category. We write $g'_1 = \mathfrak g' \ot \mathbf 1$, $g''_1 = \mathfrak g'' \ot \mathbf 1$, $g'_2 = \mathbf 1 \ot \mathfrak g'$, $g''_2 = \mathbf 1 \ot \mathfrak g''$ for the four simple ideals of $(S \ot S)_1 \cong \mathfrak{sl}_2^{\oplus 4}$, of levels $3$, $13$, $3$, $13$.

\begin{prop}
\label{prop:W-pi}
Assume \textup{(I1)}, \textup{(F1)}--\textup{(F4)} and \textup{(R)}, and put
\[
W_\pi \ov{\mathrm{def}}{=} \bigoplus_{z} M_z \ot M_{\pi(z)},
\]
the sum running over the twelve labels.
\begin{enumerate}
\item The $S \ot S$-module $W_\pi$ carries the structure of a simple $\Z$-graded vertex operator algebra of CFT type extending $S \ot S$, with the same conformal vector, and $W_\pi$ is a diagonal holomorphic extension of $S \ot S$.
\item $W_\pi$ is isomorphic to the lattice vertex operator algebra $V_{E_8 \oplus E_8}$.
\item The weight-one space is
\[
(W_\pi)_1 = (S \ot S)_1 \oplus \bigoplus_{z \in \{c_1,c_2,d_1,d_5,d_4,d_6\}} R_z \ot R_{\pi(z)},
\]
where $R_z$ is the lowest weight space of $M_z$. As a module over $(S \ot S)_1$, it does not depend on the choice of $\pi$ on $\{a,c_0\}$, it contains no irreducible $g'_1$-module nor $g'_2$-module of dimension four, no nonzero $(S \ot S)_1$-invariant vector, and the subspace of $g'_1 \oplus g'_2$-invariant vectors has dimension $78$.
\end{enumerate}
\end{prop}

\begin{proof}
1. We apply \cite[Theorem 5.10(1) p.~51]{CKM22}, which is \cite[Main Theorem 3(1) pp.~4--5]{CKM22}, with $U = V = S$ and $\cal U = \cal V = \Rep(S)$. Its hypotheses are that $S$ is simple and self-contragredient, which is (I1), that $\Rep(S)$ is closed under contragredients, semisimple, closed under submodules and quotients, and admits the vertex tensor category structure of Huang, Lepowsky and Zhang, which holds by (F1) and \cite{Hua08}, and that $\tau = \pi$ is a braid-reversed tensor equivalence whose twists satisfy $\theta_{\pi(M)} = \pm\pi(\theta_M^{-1})$, which holds with the sign $+$ by Lemma \ref{lem:pi-labels}. The conclusion is that $A = \bigoplus_z M_z' \ot \pi(M_z)$ carries the structure of a $\frac12\Z$-graded conformal vertex algebra extending $S \ot S$, which is simple since $\Rep(S)$ is rigid. By Lemma \ref{lem:dimensions}, every simple $S$-module is self-dual, so $M_z' \cong M_z$ and $A = W_\pi$ as an $S \ot S$-module. The conformal weights of $M_z \ot M_{\pi(z)}$ lie in $h_z+h_{\pi(z)}+\Z$, and $\e^{2\pi\i(h_z+h_{\pi(z)})} = \theta_z\theta_{\pi(z)} = 1$ by Lemma \ref{lem:pi-labels}, so $W_\pi$ is $\Z$-graded. Its weights are nonnegative, the weight-zero space is $S_0 \ot S_0 = \Cbb\mathbf 1$ since $h_z+h_{\pi(z)} > 0$ for $z \neq 0$ by (F1), and the graded pieces are finite dimensional, so $W_\pi$ is a vertex operator algebra of CFT type.

By \cite[Theorem 3.2 p.~10]{HKL15}, $W_\pi$ corresponds to a haploid commutative associative algebra $\cal A$ with trivial twist in $\Rep(S \ot S) \cong \Rep(S) \boxtimes \Rep(S)$, with underlying object $\bigoplus_z M_z \boxtimes M_{\pi(z)}$. The hypotheses of \cite[Theorem 5.12 p.~52]{CKM22} hold: $\Rep(S)$ is a braided fusion category, $W_\pi$ is a simple $\Z$-graded vertex operator algebra extension of $S \ot S$ in $\Rep(S) \boxtimes \Rep(S)$, the modules $M_z$ are distinct simple modules including $M_0 = S$, and $\dim\operatorname{Hom}_S(S,M_{\pi(z)}) = \delta_{z,0}$. Hence $\dim\cal A > 0$, and the category of $\cal A$-modules in $\Rep(S) \boxtimes \Rep(S)$ is semisimple \cite[Theorem 5.12 p.~52 and its proof]{CKM22}, so that $\cal A$ is separable \cite[Proposition 2.7 p.~5]{DMNO13} and is a connected étale algebra. By \cite[Corollary 5.13 p.~53]{CKM22}, whose hypotheses are that $S$ is strongly rational, which is (F1), and that $W_\pi$ is simple and of CFT type, the algebra $W_\pi$ is strongly rational. Its Frobenius--Perron dimension is
\[
\operatorname{FPdim}\cal A = \sum_z \dim M_z\dim M_{\pi(z)} = \sum_z (\dim M_z)^2 = \operatorname{FPdim}\Rep(S),
\]
by Lemma \ref{lem:pi-labels}, which is the square root of $\operatorname{FPdim}(\Rep(S) \boxtimes \Rep(S))$. Hence $\cal A$ is Lagrangian \cite[Definition 4.6 p.~17]{DMNO13}, and its category of local modules is trivial \cite[Remark 4.3 p.~16]{DMNO13}. By \cite[Theorem 3.4 p.~16]{HKL15}, the category of $W_\pi$-modules whose restriction to $S \ot S$ lies in $\Rep(S \ot S)$ is isomorphic to the category of local $\cal A$-modules. Since $S \ot S$ is regular, every weak $W_\pi$-module is a direct sum of simple ordinary $S \ot S$-modules \cite[Theorem 4.5]{ABD04}, and the submodule generated by one of them has finite length by the argument of \cite[proof of Theorem 12.3]{Arh26}, so every $W_\pi$-module is a direct sum of copies of $W_\pi$. Thus $W_\pi$ is holomorphic.

For the commutants, a vector of $W_\pi$ annihilated by the modes $a_{(n)} \ot \mathbf 1$, $a \in S$, $n \geq 0$, is annihilated by $L(0) \ot \mathbf 1$. On $M_z \ot M_{\pi(z)}$ the eigenvalues of $L(0) \ot \mathbf 1$ lie in $h_z+\N$, with $h_z > 0$ for $z \neq 0$ by (F1), so such a vector lies in $S_0 \ot S = \mathbf 1 \ot S$. Conversely $\mathbf 1 \ot S$ commutes with $S \ot \mathbf 1$. Hence $\operatorname{Com}_{W_\pi}(S \ot \mathbf 1) = \mathbf 1 \ot S$, and the same argument gives $\operatorname{Com}_{W_\pi}(\mathbf 1 \ot S) = S \ot \mathbf 1$.

3. The weight-one space of $M_z \ot M_{\pi(z)}$ is $R_z \ot R_{\pi(z)}$ if $h_z+h_{\pi(z)} = 1$, and it is zero otherwise, since $h_z+h_{\pi(z)} \in \Z_{\geq 0}$. By (F2) and Lemma \ref{lem:pi-labels}, $h_z+h_{\pi(z)} = 1$ exactly for the six labels $c_1$, $c_2$, $d_1$, $d_5$, $d_4$, $d_6$, since $\frac13+\frac23 = \frac{6}{13}+\frac{7}{13} = \frac{5}{13}+\frac{8}{13} = 1$, while $h_{d_2}+h_{d_3} = \frac{11}{13}+\frac{15}{13} = 2$, $h_b+h_b = 2$ and $h_a+h_{\pi(a)} = h_{c_0}+h_{\pi(c_0)} = 2$. This gives the decomposition, and shows that the choice of $\pi$ on $\{a,c_0\}$ does not affect $(W_\pi)_1$. By (F2), the lowest weight spaces, as $\mathfrak g' \oplus \mathfrak g''$-modules, are those of \cite[Table 2]{Arh26}, column (II):
\[
R_{c_1} = (0,2),\quad R_{c_2} = (0,0)+(0,4)+(1,1)+(1,3)+(2,2),\quad R_{d_1} = (0,2)+(1,1),
\]
\[
R_{d_4} = (0,0)+(1,1),\quad R_{d_5} = (0,4)+(1,3),\quad R_{d_6} = (0,2)+(1,3)+(2,0),
\]
where $(m,n)$ denotes $V(m) \ot V(n)$ and $V(m)$ is the irreducible $\mathfrak{sl}_2$-module of highest weight $m$. All the $\mathfrak g'$-weights occurring are at most $2$, so no $g'_k$-module of dimension four occurs in $(W_\pi)_1$, and $(S \ot S)_1$ is a sum of adjoint modules. A $(S \ot S)_1$-invariant vector of $R_z \ot R_{\pi(z)}$ would require $(0,0)$ to occur in both $R_z$ and $R_{\pi(z)}$, which fails for each of the three pairs, and $(S \ot S)_1$ is semisimple. The $g'_1 \oplus g'_2$-invariant vectors of $R_z \ot R_{\pi(z)}$ are the tensor products of the $\mathfrak g'$-invariant parts, of dimensions $3 \cdot 6$ for $(c_1,c_2)$, $3 \cdot 5$ for $(d_1,d_5)$ and $1 \cdot 3$ for $(d_4,d_6)$, and those of $(S \ot S)_1$ form $g''_1 \oplus g''_2$, of dimension $6$. The total is $6+2 \cdot 18+2 \cdot 15+2 \cdot 3 = 78$.

2. By part 1 and \cite[Theorem 2 p.~2]{DoM04b}, whose hypotheses are that $W_\pi$ is $C_2$-cofinite, holomorphic, of CFT type and of central charge $16$, $W_\pi$ is isomorphic to $V_{E_8 \oplus E_8}$ or to $V_{D_{16}^+}$. The exclusion of $V_{D_{16}^+}$ in \cite[proof of Proposition 7.12]{ArhII} uses only the three properties of $(W_\pi)_1$ established in part 3 and the levels $3$ of $g'_1$, $g'_2$. We recall it. If $W_\pi \cong V_{D_{16}^+}$, then $(W_\pi)_1 \cong \mathfrak{so}(E) \cong \Lambda^2E$ with $\dim E = 32$. The Dynkin index of $g'_k \subset \mathfrak{so}(E)$ is $3$, the vector representation has Dynkin index $2$, so $E|_{g'_k}$ has Dynkin index $6$ and is $6V(1) \oplus 20V(0)$ or $V(2) \oplus 2V(1) \oplus 25V(0)$. The second case gives $V(3) \subset V(2) \ot V(1) \subset \Lambda^2E$, excluded by part 3. Hence $E|_{g'_1 \oplus g'_2} = pV(1,1) \oplus (6-2p)(V(1,0) \oplus V(0,1)) \oplus (8+4p)V(0,0)$ for some $p \in \{0,1,2,3\}$, and the centralizer of $g'_1 \oplus g'_2$ in $\mathfrak{so}(E)$ is $\mathfrak{so}_p \oplus \mathfrak{so}_{8+4p} \oplus \mathfrak{sp}_{6-2p} \oplus \mathfrak{sp}_{6-2p}$, of dimension $70$, $86$, $127$ or $193$, which contradicts the value $78$ of part 3.
\end{proof}

The converse holds: a diagonal holomorphic extension forces (R).

\begin{lemma}
\label{lem:extension-to-R}
Assume \textup{(I1)} and \textup{(F1)}--\textup{(F4)}, and let $W$ be a diagonal holomorphic extension of $S \ot S$. Then there is a braided equivalence $\pi \co \Rep(S) \to \Rep(S)^{\mathrm{rev}}$ such that $W \cong \bigoplus_z M_z \ot M_{\pi(z)}$ as an $S \ot S$-module. In particular \textup{(R)} holds, $W_1$ has the properties of Proposition \ref{prop:W-pi}(3), and $W \cong V_{E_8 \oplus E_8}$.
\end{lemma}

\begin{proof}
Since $S \ot S$ is regular and $W$ has finite-dimensional weight spaces and nonnegative weights, $W$ is a finite direct sum of simple $S \ot S$-modules $M_z \ot M_{z'}$, because there are finitely many simple $S\ot S$-modules, their lowest weights form a finite set, and an infinite multiplicity would contradict finite-dimensionality of the corresponding weight space. Write $W = \bigoplus_{z \in I} M_z \ot V_z$, where $V_z \in \Rep(S)$ is the sum of the $M_{z'}$ paired with $M_z$ and $I$ is the set of labels with $V_z \neq 0$. A vector annihilated by the nonnegative modes of $S \ot \mathbf 1$ lies in $S_0 \ot V_0$, as in the proof of Proposition \ref{prop:W-pi}, so $\operatorname{Com}_W(S \ot \mathbf 1) = \mathbf 1 \ot V_0$ and the hypothesis gives $V_0 \cong S$. Likewise $\operatorname{Com}_W(\mathbf 1 \ot S) = \bigoplus_z M_z \ot \operatorname{Hom}_S(S,V_z)$, and the hypothesis gives $\dim\operatorname{Hom}_S(S,V_z) = \delta_{z,0}$. The hypotheses of \cite[Theorem 5.10(2) p.~51]{CKM22} are now satisfied for $U = V = S$: $\Rep(S)$ is a ribbon category, $W$ is a simple $\Z$-graded vertex operator algebra extension of $S \ot S$ of the form $\bigoplus_{z \in I} M_z \ot V_z$ with $M_0 = S$, and $\dim\operatorname{Hom}_S(S,V_z) = \delta_{z,0}$. Hence the full subcategories $\cal U_W$ and $\cal V_W$ of $\Rep(S)$ generated by the $M_z$, $z \in I$, and by the $V_z$ are ribbon subcategories, $\cal V_W$ is semisimple with the $V_z$ as distinct simple objects, and there is a braid-reversed equivalence $\tau \co \cal U_W \to \cal V_W$ with $\tau(M_z) \cong V_z'$.

By \cite[Theorem 3.2 p.~10]{HKL15}, $W$ corresponds to a connected commutative algebra $\cal A$ in $\Rep(S) \boxtimes \Rep(S)$, which is étale by \cite[Theorem 5.12 p.~52]{CKM22} and \cite[Proposition 2.7 p.~5]{DMNO13} as in the proof of Proposition \ref{prop:W-pi}. Since $W$ is holomorphic, its category of local modules is trivial by \cite[Theorem 3.4 p.~16]{HKL15}, so $\cal A$ is Lagrangian \cite[Remark 4.3 and Corollary 3.32]{DMNO13}, and $\operatorname{FPdim}\cal A = \operatorname{FPdim}\Rep(S) = \sum_z (\dim M_z)^2$, the sum over all twelve labels. On the other hand $\operatorname{FPdim}\cal A = \sum_{z \in I}\dim M_z\dim V_z = \sum_{z \in I}(\dim M_z)^2$, since the tensor equivalence $\tau$ preserves Frobenius--Perron dimensions and $\dim V_z' = \dim V_z$. Hence $I$ is the set of all labels, $\cal U_W = \Rep(S)$, and $\cal V_W$ is a fusion subcategory of $\Rep(S)$ of the same Frobenius--Perron dimension, hence $\cal V_W = \Rep(S)$. Thus $\pi \ov{\mathrm{def}}{=} \tau$ is a braided equivalence $\Rep(S) \to \Rep(S)^{\mathrm{rev}}$, and $V_z \cong \pi(M_z)' \cong M_{\pi(z)}$ by the self-duality of Lemma \ref{lem:dimensions}. The remaining assertions follow from parts 2 and 3 of Proposition \ref{prop:W-pi}, whose proofs only use the $S \ot S$-module structure of $W$ and the fact that $W$ is $C_2$-cofinite, holomorphic and of CFT type.
\end{proof}

\begin{remark}\normalfont
\label{rem:where-R-enters}
Under (I1) and (F1)--(F4), Proposition \ref{prop:W-pi} and Lemma \ref{lem:extension-to-R} show that (R) is equivalent to the existence of a diagonal holomorphic extension of $S \ot S$. The underlying $S\ot S$-module of $W_\pi$ depends only on the permutation of labels. Its multiplication also depends on the monoidal structure of $\pi$, and no uniqueness from the permutation alone is asserted. Two choices differing only on $\{a,c_0\}$ have the same weight-one $S_1\oplus S_1$-module. Theorem \ref{thm:net-diagonal} constructs a diagonal holomorphic extension of $S \ot S$ from the conformal net $\cal A_S$, using Proposition \ref{prop:sector-trace}, so that (R) is established under the inputs of that theorem. In the rest of this section, (R) and Theorem \ref{thm:net-diagonal} enter only through the existence of a diagonal holomorphic extension, which we take as the standing hypothesis.
\end{remark}


%%%%%%%%%%%%%%%%%%%%%%%%%%%%%%%%%%%%%%%%%%%%%%%%%%%%%%%%%%%%%%%%%%%%%%%%%%%%%
\subsection{The unitary structure}
\label{subsec:unitary-structure}

Let $W$ be a diagonal holomorphic extension of $S \ot S$. By Lemma \ref{lem:extension-to-R}, $W \cong V_\Gamma$ with $\Gamma = E_8 \oplus E_8$. We use the explicit model of $V_\Gamma$ of \cite[Section 7.5]{ArhII}, built with the bimultiplicative cocycle $\epsi(\alpha,\beta) = \epsi_8(\alpha^{(1)},\beta^{(1)})\epsi_8(\alpha^{(2)},\beta^{(2)})$, where $\epsi_8$ is the cocycle on $E_8$ specified there, and with the antilinear involution $\phi_1$ of \cite[Theorem 4.12]{DoL14}, which is the PCT operator of a unitary structure on $V_\Gamma$. For $\kappa(\lambda) = 2^{-\la\lambda,\varpi\ra}$, $\varpi = (\varpi_8,\varpi_8)$, $\varpi_8 = e_7+e_8$, let $\phi_\kappa = \phi_1 \circ \exp(h_\kappa(0))$, so that $\phi_\kappa(\e^\lambda) = \kappa(\lambda)\e^{-\lambda}$. Let $\mathfrak s_1, \ldots, \mathfrak s_4 \subset \mathfrak e_8 \oplus \mathfrak e_8 = (V_\Gamma)_1$ be the four commuting subalgebras isomorphic to $\mathfrak{sl}_2$ defined in \cite[Section 7.5]{ArhII}, spanned by the explicit $\mathfrak{sl}_2$-triples $(E_k,H_k,F_k)$, $F_k = -\phi_\kappa(E_k)$, with Dynkin indices $(2,1)$, $(6,7)$, $(1,2)$, $(7,6)$ in the two copies of $\mathfrak e_8$, so that they generate affine subalgebras of levels $3$, $13$, $3$, $13$. Let $A = L_{\mathfrak{sl}_2}(3,0) \ot L_{\mathfrak{sl}_2}(13,0) \subset S$ be the affine subalgebra generated by $S_1$, which is simple by Section \ref{sec:inputs}, let $\scr C = \operatorname{Com}_S(A)$, and let $\omega_{\scr C} = \omega-\omega_A$ be its conformal vector, of central charge $8-\frac95-\frac{13}5 = \frac{18}5$, where $\omega_A$ is the Sugawara vector of $A$. In $W$ we write $\omega_{\scr C \ot 1} = \omega_{\scr C} \ot \mathbf 1$, $\omega_{1 \ot \scr C} = \mathbf 1 \ot \omega_{\scr C}$, and similarly for $A$ and $S$.

\begin{prop}
\label{prop:unitary-structure}
Assume \textup{(I1)} and \textup{(F1)}--\textup{(F4)}, let $W$ be a diagonal holomorphic extension of $S \ot S$, and fix an isomorphism of vertex operator algebras $W \to V_\Gamma$.
\begin{enumerate}
\item After composing this isomorphism with an automorphism of $V_\Gamma$, the images of $g'_1$, $g''_1$, $g'_2$, $g''_2$ are $\mathfrak s_1$, $\mathfrak s_2$, $\mathfrak s_3$, $\mathfrak s_4$.
\item With this identification, the space $\cal G = \operatorname{Com}_W(A \ot A) \cap W_2$ has dimension $16$, and $\omega_{\scr C \ot 1}$ is one of four explicit vectors $u_0, \ldots, u_3 \in \cal G$, whose coordinates in a rational basis of $\cal G$ lie in $\Q(\sqrt{-6})$.
\item $\phi_\kappa$ is the PCT operator of a structure of unitary vertex operator algebra on $W$, and $S \ot \mathbf 1$, $\mathbf 1 \ot S$ and $S \ot S$ are unitary subalgebras for this structure, in the sense of \cite[Definition 5.22 p.~43]{CKLW18}.
\end{enumerate}
\end{prop}

\begin{proof}
Apply \cite[Lemmas 7.14--7.15 and Proposition 7.16]{ArhII}. The inputs to those proofs are the rank-$16$ lattice realization and the weight-one decomposition of Proposition \ref{prop:W-pi}, the levels $(3,13)$, the twelve top spaces and conformal weights in (F2), and positivity in (F1). The tensor-category identification $\mathrm{(H_J)}$ is not used after these data have been supplied.

The computation has three steps. The enumeration of embeddings of the four current algebras into $\mathfrak e_8\oplus\mathfrak e_8$ leaves the specified embedding. The common kernel of their relevant modes in weight two has dimension $16$, and the Virasoro and trace equations in that space have the four solutions $u_0,\ldots,u_3$. The explicit PCT operator $\phi_\kappa$ preserves the currents and these four vectors. Since it also fixes $\omega$ and the Sugawara vectors of the currents, it fixes $\omega_{1 \ot \scr C} = \omega-\omega_{A \ot 1}-\omega_{1 \ot A}-\omega_{\scr C \ot 1}$. It therefore preserves both copies of $S$: each copy is the commutant of the currents and coset conformal vector of the other copy, since positivity excludes a vector of zero opposite conformal weight in a nonvacuum summand. Stability under $L_1$ then gives unitary subalgebras by \cite[Proposition 5.23]{CKLW18}. The programs and their outputs form block C3 of Section \ref{sec:certificates}.
\end{proof}

\begin{remark}\normalfont
\label{rem:HT-structure}
The unitary structure of Proposition \ref{prop:unitary-structure}(3) is constructed from the lattice. When $W$ is the algebra of Theorem \ref{thm:net-diagonal}, it comes with the unitary structure of Henriques and Tener \cite{HeT25}, for which $S \ot S$ is a unitary subalgebra by construction, and Proposition \ref{prop:unitary-structure} is not needed for it. Two structures of unitary vertex operator algebra on a simple vertex operator algebra are unitarily isomorphic \cite[Proposition 5.19 pp.~40--41]{CKLW18}, so the conclusions of Theorem \ref{thm:complete-unitarity} do not depend on which structure is used.
\end{remark}


%%%%%%%%%%%%%%%%%%%%%%%%%%%%%%%%%%%%%%%%%%%%%%%%%%%%%%%%%%%%%%%%%%%%%%%%%%%%%
\subsection{The equivalent criteria}

\begin{thm}[Theorem C]
\label{thm:unitarity-equivalences}
Assume \textup{(I1)} and \textup{(F1)}--\textup{(F4)}, so that Theorem \ref{thm:net} is available. The following statements are equivalent. They all hold by Theorem \ref{thm:HJ-main}, and the implications \textup{(ii)}$\Rightarrow$\textup{(i)} and \textup{(iii)}$\Rightarrow$\textup{(iv)}$\Rightarrow$\textup{(v)}$\Rightarrow$\textup{(vi)}$\Rightarrow$\textup{(i)} are proved without Propositions \ref{prop:net-B} and \ref{prop:sector-trace}.
\begin{enumerate}
\item[\textup{(i)}] $\mathrm{(H_J)}$: $\Rep(S)$ is braided equivalent to $Z(\cal H_3)$.
\item[\textup{(ii)}] The fusion category $\cal F_x = \Rep(S)_{L_x}$ of \textup{(F4)} is unitarizable, for one, equivalently for both, $x \in \{+,-\}$.
\item[\textup{(iii)}] \textup{(R)}: there is a braided equivalence $\Rep(S) \to \Rep(S)^{\mathrm{rev}}$.
\item[\textup{(iv)}] $S \ot S$ admits a diagonal holomorphic extension in the sense of Definition \ref{def:diagonal-extension}.
\item[\textup{(v)}] $S$ is completely unitary, and the functor $\Pi \co \Rep^u(S) \to \Rep(\cal A_S)$ of Section \ref{subsec:complete-unitarity} is an equivalence of braided $C^*$-tensor categories.
\item[\textup{(vi)}] $\Rep(S)$ is braided equivalent to $\Rep(\cal A_S)$ or to $\Rep(\cal A_S)^{\mathrm{rev}}$.
\end{enumerate}
\end{thm}

\begin{proof}
(i) $\Rightarrow$ (iii). Composing a braided equivalence $\Rep(S) \to Z(\cal H_3)$, its reverse, and the equivalence $R$ of Lemma \ref{lem:ZH3-self-reverse} gives a braided equivalence $\Rep(S)^{\mathrm{rev}} \to \Rep(S)$, whose inverse is (iii).

(iii) $\Rightarrow$ (iv). This is Proposition \ref{prop:W-pi}: $W_\pi$ is a diagonal holomorphic extension of $S \ot S$.

(iv) $\Rightarrow$ (iii). This is Lemma \ref{lem:extension-to-R}.

(iv) $\Rightarrow$ (v). Let $W$ be a diagonal holomorphic extension. By Proposition \ref{prop:unitary-structure}, $W$ carries a unitary structure for which $S \ot \mathbf 1$ and $\mathbf 1 \ot S$ are unitary subalgebras. Since $W \cong V_\Gamma$ by Lemma \ref{lem:extension-to-R} and two unitary structures on $V_\Gamma$ are unitarily isomorphic \cite[Proposition 5.19 pp.~40--41]{CKLW18}, Theorem \ref{thm:complete-unitarity} gives (v).

(v) $\Rightarrow$ (vi). By (v), the forgetful functor $\Rep^u(S) \to \Rep(S)$ is a braided equivalence \cite[Remark 4.3 p.~18]{CGGH23}, and $\Pi$ is a braided equivalence onto $\Rep(\cal A_S)$ with the braiding of Connes fusion, which is the DHR braiding of Section \ref{sec:net} or its reverse by Remark \ref{rem:braiding-conventions}.

(vi) $\Rightarrow$ (i). By Theorem \ref{thm:net}(3), $\Rep(\cal A_S) \cong Z(\cal H_3)$ as braided categories, and $Z(\cal H_3)^{\mathrm{rev}} \cong Z(\cal H_3)$ by Lemma \ref{lem:ZH3-self-reverse}. Hence $\Rep(S) \cong Z(\cal H_3)$ in both cases of (vi).

(i) $\Rightarrow$ (ii). Under (i), $L_+$ and $L_-$ are Lagrangian algebras of $Z(\cal H_3)$. By \cite[Proposition 3.3]{ArhII}, the Lagrangian algebras of $Z(\cal H_3)$ are $L_1$, $L_2$, $L_3$, with categories of modules $\cal H_1$, $\cal H_2$, $\cal H_3$, which are unitary fusion categories. Hence $\cal F_\pm$ is equivalent to one of them.

(ii) $\Rightarrow$ (i). Apply Theorem \ref{zhu:thm:categorical}(4) for the given $x$.

(v) $\Rightarrow$ (ii). By (v) and \cite[Theorem 7.9 p.~53]{Gui19}, $\Rep^u(S)$ is a unitary modular tensor category, and the forgetful functor $\Rep^u(S) \to \Rep(S)$ is a braided equivalence. The algebra $L_x$ corresponds to a haploid commutative algebra in $\Rep^u(S)$, which is rigid since a connected étale algebra in a fusion category is rigid \cite[Remark 3.4 p.~7]{DMNO13}, hence equivalent to a $Q$-system $Q_x$ in the $C^*$-tensor category $\Rep^u(S)$ \cite[Theorem 3.2 p.~11]{CGGH23}. The category of $Q_x$-modules in $\Rep^u(S)$, with the tensor product over $Q_x$, is equivalent to $\cal F_x$ as a fusion category, since $Q_x \cong L_x$ as algebras. Every module over a $Q$-system is unitarizable and the unitary modules form a $C^*$-category \cite[Proposition 2.24 p.~26 and Theorem 2.28 p.~28]{Gui22}, and the unitary bimodules form a $C^*$-tensor category for the tensor product over the $Q$-system \cite[Theorem 3.13 p.~35]{Gui22}. For the commutative $Q$-system $Q_x$, a module is a bimodule through the braiding, and the modules form a full tensor subcategory of the bimodules \cite[Section 1]{KiO02}. Hence $\cal F_x$ is equivalent to a $C^*$-tensor category with simple unit and finitely many simple objects, that is, $\cal F_x$ is unitarizable.
\end{proof}

\begin{remark}\normalfont
\label{rem:status-equivalences}
Theorem \ref{thm:HJ-main} establishes these statements using the analytic diagonal extension. The implication (iv)$\Rightarrow$(v) applies to an arbitrary diagonal holomorphic extension and uses the computer-assisted Proposition \ref{prop:unitary-structure}. For the extension reconstructed from $\cal B$, the unitary structure and the unitary lattice realization are given by Theorem \ref{thm:net-diagonal}, and Proposition \ref{prop:unitary-structure} is not needed. An independent proof of (ii) or of (iii) would give $\mathrm{(H_J)}$ without Propositions \ref{prop:net-B} and \ref{prop:sector-trace}.
\end{remark}


%%%%%%%%%%%%%%%%%%%%%%%%%%%%%%%%%%%%%%%%%%%%%%%%%%%%%%%%%%%%%%%%%%%%%%%%%%%%%
\subsection{The lattice of the reconstructed extension}
\label{subsec:conditional}

\begin{cor}
\label{cor:HJ-conditional}
Assume \textup{(I1)} and \textup{(F1)}--\textup{(F4)}. The vertex operator algebra $W$ of Theorem \ref{thm:net-diagonal} is isomorphic to $V_{E_8 \oplus E_8}$, and the net $\cal B$ of Proposition \ref{prop:net-B} is isomorphic to $\cal A_{V_{E_8 \oplus E_8}}$. The categories $\cal F_+$ and $\cal F_-$ of \textup{(F4)} are equivalent to $\cal H_2$ and $\cal H_3$ in some order.
\end{cor}

\begin{proof}
By Theorem \ref{thm:net-diagonal}(1), $W$ is a diagonal holomorphic extension of $S \ot S$ in the sense of Definition \ref{def:diagonal-extension}. Hence $W \cong V_{E_8 \oplus E_8}$ by Lemma \ref{lem:extension-to-R}, and $\cal B \cong \cal A_{V_{E_8 \oplus E_8}}$ by Proposition \ref{prop:lattice-identification}. Let $\Phi_0 \co \Rep(S) \to Z(\cal H_3)$ be the braided equivalence of Section \ref{sec:center}. By (F4), $L_\pm$ is a Lagrangian algebra of $\Rep(S)$, so that $\Phi_0(L_\pm)$ is a Lagrangian algebra of $Z(\cal H_3)$. By (F3) and the choice of the labels in Section \ref{sec:center}, its underlying object is $0 \oplus b \oplus 2\Phi_0(X_\pm)$, where $\Phi_0(X_+) \not\cong \Phi_0(X_-)$ and $\{\Phi_0(X_+),\Phi_0(X_-)\} = \{a,c_0\}$. By \cite[Proposition 3.3]{ArhII}, $\Phi_0(L_+)$ and $\Phi_0(L_-)$ are $L_2$ and $L_3$ in some order, so that $\cal F_+$ and $\cal F_-$ are equivalent to $Z(\cal H_3)_{L_2} \cong \cal H_2$ and $Z(\cal H_3)_{L_3} \cong \cal H_3$ in some order.
\end{proof}


%%%%%%%%%%%%%%%%%%%%%%%%%%%%%%%%%%%%%%%%%%%%%%%%%%%%%%%%%%%%%%%%%%%%%%%%%%%%%
\section{Certificates and programs}
\label{sec:certificates}

The computations used in the proofs are listed below, with the statements they support. The certificate checks are written in Python~3 and use exact arithmetic over $\Q$ and over the fields specified in the programs. In particular, block A4 uses $\K = \Q(\sqrt{13},\sqrt{-3})$ and block C3 uses $\Q(\sqrt{-6})$ for the conformal-vector coordinates. The supplementary archives listed below record the scripts, outputs, checksums and replay instructions.
\begin{enumerate}
\item Block A1, \nolinkurl{a1_lattices.py} and \nolinkurl{a1b_extra.py}, outputs \nolinkurl{a1_lattices.out} and \nolinkurl{a1b_extra.out}: Lemma \ref{lem:common-sublattice} and the lattice facts used in Section \ref{subsec:involutions}. The first script works in the coordinates $(p,b)$ of $\mathfrak h$ and verifies the Gram matrices, the intersection $N = L_{\mathrm I} \cap L_{\mathrm{II}}$ on a box, the indices, the characters of $\sigma$ and $\tau$, the discriminant group $N'/N$ with its $2$-part and $39$-part, the isometry $\psi_N$ and the value $u = 22$, and the Weyl group actions $14$, $25$. The second script verifies the values of $(s,\cdot)$ and $(p,\cdot)$ on $L_{\mathrm{II}}'$ and $L_{\mathrm I}'$, the identities \eqref{eq:B-sets}, the isotropy statement of part 1 of the lemma, and, in the root coordinates of \cite[(4.1)]{DJX26}, the intersections \eqref{eq:intersections} on a box and the properties of $L_0$, $L_H$ and $E$ of part 5.
\item Block A4, the scripts \nolinkurl{a4_kfield.py}, \nolinkurl{a4_voa.py}, \nolinkurl{a4_data.py}, \nolinkurl{a4_common.py} and \nolinkurl{a4_smoke.py}, the six scripts \nolinkurl{a4_item1_lattice.py} to \nolinkurl{a4_item6_riesen.py}, the driver \nolinkurl{a4_run_all.sh} and the outputs: the independent verification of the finite identities of \cite{DJX26} described in Section \ref{subsec:DJX-verification}. The lattice vertex operator algebra $V_L$ is implemented from scratch in the basis and with the cocycle of \cite[Section 4]{DJX26}, with the mode formula of \cite[Appendix A]{DJX26} and without truncation in the conformal weight, independently of the implementation of \cite{Arh26}, which uses the Bourbaki basis of $E_6$ and a different cocycle. The field $\K$ is coded by hand in the basis $1, \sqrt{13}, \sqrt{-3}, \sqrt{-39}$.
\item Block C3, the directory \nolinkurl{run} with the programs \nolinkurl{R3} and \nolinkurl{R3-verif} of the archive \nolinkurl{Calculs-HU-v08} of \cite{ArhII}, three driver scripts, the scripts \nolinkurl{a3_c3_table_check.py} and \nolinkurl{a3_c3_compare.py}, and the directory \nolinkurl{out} with $69$ outputs and the timing file: the recomputation of \cite[Section 7.5]{ArhII} used in Section \ref{sec:algebraic-route}. The archived comparison with the original \nolinkurl{Calculs-HU-v08} covers $47$ regenerated data files, including the re-export of the portable certificate. The updated script \nolinkurl{a3_c3_compare.py} compares the shared files with the fixed public release of \cite{ArhII}, including its compressed secondary weight-two space; diagnostics specific to C3 are listed separately. The replay instructions distinguish these comparisons from the archived production runs. The script \nolinkurl{a3_c3_table_check.py} checks the input tables \nolinkurl{t12_table.json} and \nolinkurl{t12_W1.json} against the lowest weight spaces in Table \ref{zhu:table:blocks} and the exported Zhu data. These computations use the data of (F2).
\end{enumerate}
The archive \nolinkurl{Calculs-III-A1-A4-C3-v11.zip} contains the scripts and outputs of blocks A1, A4 and C3, with replay instructions and a SHA-256 manifest. The Zhu calculations and the two ideal-containment verifiers are in the separate supplementary archive \nolinkurl{Certificats-Zhu-III-v07.zip}, described in Section \ref{zhu:app:certificates}. Block C3 supports the alternative construction in Section \ref{sec:algebraic-route}. The main analytic proof uses the compatible unitary structure obtained by reconstruction. The weight-four certificates used for the $C_2$-cofiniteness result of \cite[Theorem 10.21]{Arh26} are distributed with the fixed supplementary release of \cite{Arh26}; their role is distinct from block C3.

The certificates are exact computations. They are not a formal verification of the implementations of the lattice vertex operators. The two implementations used in blocks A4 and C3 are independent of each other and of the one of \cite{Arh26}, and the archived comparisons report agreement on the quantities tested by both implementations.


\appendix
%%%%%%%%%%%%%%%%%%%%%%%%%%%%%%%%%%%%%%%%%%%%%%%%%%%%%%%%%%%%%%%%%%%%%%%%%%%%%
\section{Zhu algebra, simple modules and lattice branching}
\label{app:zhu}

We prove (F1)--(F3) of Theorem \ref{thm:zhu-inputs} and (F4) of Theorem \ref{thm:algebraic-F4} for the explicit $S$. The inputs from \cite{Arh26} are the lattice construction and the results (I1)--(I2) of Theorem \ref{thm:inputs-ArhI}, the generation of $A(S)$ by the current and weight-two classes in its Proposition~11.3, and its Proposition~11.4 on the affine simple-current extension $A^+$. For the exact relations we use the explicit $U_1$ and $U_2$ over $\K$ and $\omega=\omega_A+2y$ from that construction. The branching and block-survival arguments use the known lattice modules and their genus-one modular matrices. The categorical argument for (F4) uses the extension and recognition results of \cite{ArhII} stated for these inputs independently of $\mathrm{(H_J)}$.

%%%%%%%%%%%%%%%%%%%%%%%%%%%%%%%%%%%%%%%%%%%%%%%%%%%%%%%%%%%%%%%%%%%%%%%%%%%%%
\subsection{An integral descent lemma}
\label{zhu:sec:descent}

Throughout this section, $R$ is a discrete valuation ring with maximal ideal $\mathfrak m = \pi R$, residue field $k = R/\mathfrak m$ and field of fractions $K$. An $R$-algebra is finite if it is finitely generated as an $R$-module. A finite-dimensional algebra over a field $F$ is split semisimple if it is isomorphic to a finite product of full matrix algebras $\M_n(F)$. The following statement, and the use made of it below, follow the method of \cite[Lemma 9.6 p.~55]{Rie26b}. We give a complete proof.

\begin{lemma}
\label{zhu:lem:descent}
Let $\Lambda$ be a finite $R$-algebra such that $\Lambda/\mathfrak m\Lambda \cong \prod_{i \in I} \M_{n_i}(k)$ is split semisimple.
\begin{enumerate}
\item There exist a subset $I' \subset I$ and, for any algebraically closed field $\Omega$ containing $K$, a bijection between $I'$ and the set of simple factors of $\Lambda \ot_R \Omega$, such that the factor attached to $i$ is isomorphic to $\M_{n_i}(\Omega)$. In particular the $K$-algebra $\Lambda \ot_R K$ is semisimple.
\item Let $\epsi \in \Lambda$ be an idempotent and let $i \in I'$. The rank of $\epsi$ acting on the simple module of the factor of $\Lambda \ot_R \Omega$ attached to $i$ is equal to the rank of the image of $\epsi$ acting on the simple module of the factor $\M_{n_i}(k)$ of $\Lambda/\mathfrak m\Lambda$.
\item Let $z \in \Lambda$ and $i \in I'$, and assume that $z$ acts on the simple module of the factor of $\Lambda \ot_R \Omega$ attached to $i$ by a scalar $c \in R$. Then the image of $z$ in $\Lambda/\mathfrak m\Lambda$ acts on the simple module of the factor $\M_{n_i}(k)$ by the reduction of $c$.
\end{enumerate}
\end{lemma}

\begin{proof}
Let $T = \{x \in \Lambda : \pi^n x = 0 \text{ for some integer } n \geq 0\}$ be the torsion submodule. Since $R$ is central in $\Lambda$, the subspace $T$ is a two-sided ideal. The quotient $C \ov{\mathrm{def}}{=} \Lambda/T$ is a finitely generated torsion-free module over the principal ideal domain $R$, hence free of finite rank. Since $K$ is flat over $R$ and $T \ot_R K = 0$, the quotient map induces an isomorphism $\Lambda \ot_R K \to C \ot_R K$, and also $\Lambda \ot_R \Omega \cong C \ot_R \Omega$. The surjection $\Lambda \to C$ induces a surjective homomorphism of algebras $\Lambda/\mathfrak m\Lambda \to C/\mathfrak m C$. Its kernel is a two-sided ideal of $\prod_{i \in I} \M_{n_i}(k)$, hence of the form $\prod_{i \in I \setminus I'} \M_{n_i}(k)$ for some subset $I' \subset I$. Thus $C/\mathfrak m C \cong \prod_{i \in I'} \M_{n_i}(k)$, and the images $\epsi_i$, $i \in I'$, of the units of the factors are central orthogonal idempotents of $C/\mathfrak m C$ with sum $1$.

Let $\hat R$ be the $\mathfrak m$-adic completion of $R$. It is a complete discrete valuation ring with uniformizer $\pi$, residue field $k$, and its field of fractions $\hat K$ contains $K$. The $\hat R$-algebra $\hat C \ov{\mathrm{def}}{=} C \ot_R \hat R$ is a free $\hat R$-module of finite rank, it is complete for the $\pi$-adic topology, and $\hat C/\pi \hat C \cong C/\mathfrak m C$. We use the following lifting of idempotents. Let $Q$ be an $\hat R$-algebra which is free of finite rank as an $\hat R$-module, and let $\bar e$ be an idempotent of $Q/\pi Q$. Choose $x_0 \in Q$ lifting $\bar e$ and put $x_{n+1} = 3x_n^2 - 2x_n^3$. The identities $x_{n+1}^2 - x_{n+1} = (x_n^2 - x_n)^2(4x_n^2 - 4x_n - 3)$ and $x_{n+1} - x_n = -(x_n^2 - x_n)(2x_n - 1)$ show by induction that $x_n^2 - x_n \in \pi^{2^n} Q$ and that $(x_n)$ is a Cauchy sequence. Its limit $e$ is an idempotent of $Q$ lifting $\bar e$. Applying this to $Q = \hat C$ and $\bar e = \epsi_1$, then to the corner $(1-e_1)\hat C(1-e_1)$, which is a direct summand of $\hat C$ and reduces to $(1-\epsi_1)(C/\mathfrak m C)(1-\epsi_1)$, and so on, we obtain orthogonal idempotents $e_i \in \hat C$, $i \in I'$, with $\sum_i e_i = 1$, such that $e_i$ reduces to $\epsi_i$ (the last one is $1$ minus the sum of the others). For $i \neq j$, the $\hat R$-module $e_i \hat C e_j$ is a direct summand of $\hat C = \bigoplus_{r,s} e_r \hat C e_s$, hence free of finite rank, and its reduction modulo $\pi$ is $\epsi_i (C/\mathfrak m C) \epsi_j = 0$ since $\epsi_i$ is central. By Nakayama's lemma, $e_i \hat C e_j = 0$. Consequently each $e_i$ is central and $\hat C = \prod_{i \in I'} \hat C_i$ with $\hat C_i \ov{\mathrm{def}}{=} e_i \hat C e_i$ and $\hat C_i/\pi \hat C_i \cong \M_{n_i}(k)$.

Fix $i \in I'$ and put $n = n_i$. Lift a rank-one idempotent of $\M_n(k)$ to an idempotent $f \in \hat C_i$ as above. The left ideal $P \ov{\mathrm{def}}{=} \hat C_i f$ is a direct summand of $\hat C_i$, hence a free $\hat R$-module, and $P/\pi P \cong \M_n(k) E_{11} \cong k^n$. Thus $P$ is free of rank $n$, and $\hat C_i$ is free of rank $n^2$. Left multiplication gives a homomorphism of $\hat R$-algebras $\varphi \co \hat C_i \to \operatorname{End}_{\hat R}(P) \cong \M_n(\hat R)$ between free $\hat R$-modules of rank $n^2$, whose reduction modulo $\pi$ is the action of $\M_n(k)$ on $k^n$, an isomorphism. The determinant of $\varphi$ in bases is therefore a unit of the local ring $\hat R$, and $\varphi$ is an isomorphism. Hence $\hat C \cong \prod_{i \in I'} \M_{n_i}(\hat R)$ and $C \ot_R \hat K \cong \hat C \ot_{\hat R} \hat K \cong \prod_{i \in I'} \M_{n_i}(\hat K)$.

Put $B \ov{\mathrm{def}}{=} \Lambda \ot_R K \cong C \ot_R K$. The Jacobson radical $J(B)$ is a nilpotent ideal, so $J(B) \ot_K \hat K$ is a nilpotent ideal of the semisimple algebra $B \ot_K \hat K \cong C \ot_R \hat K$. Thus $J(B) = 0$ and $B$ is semisimple. Let $\Omega'$ be an algebraically closed field containing $\hat K$ and let $\ovl K$ be the algebraic closure of $K$ in $\Omega'$. Since $K$ has characteristic zero, $B \ot_K \ovl K$ is a product of matrix algebras over $\ovl K$, and its factors remain simple after extension to $\Omega'$. Since $B \ot_K \Omega' \cong \prod_{i \in I'} \M_{n_i}(\Omega')$, the algebra $B \ot_K \ovl K$ has factors of sizes $n_i$, $i \in I'$. Any algebraically closed field $\Omega$ containing $K$ contains an algebraic closure of $K$, and we fix a $K$-embedding $\iota \co \ovl K \to \Omega$. The simple factors of $B \ot_K \ovl K$ correspond bijectively, by extension of scalars, to those of $B \ot_K \Omega$ through $\iota$ and to those of $B \ot_K \Omega'$, and the latter are the factors $\hat C_i \ot_{\hat R} \Omega'$, $i \in I'$. This defines the bijection of the first part, which is proved.

For the second part, the image of $\epsi$ in $\hat C_i \cong \M_n(\hat R)$ is an idempotent matrix $E$. Its image $E \hat R^n$ is a direct summand of $\hat R^n$, hence free of some rank $r$. The rank of $E$ over $\hat K$, hence over $\Omega'$, is $r$, and the image of the reduction of $E$ modulo $\pi$ is $(E\hat R^n)/\pi(E \hat R^n) \cong k^r$. The rank of $\epsi$ on the simple module of a factor of $B \ot_K \ovl K$ does not change under the extensions of scalars to $\Omega$ and to $\Omega'$, which gives the claim. For the third part, $z$ acts by $c$ on the simple module of the corresponding factor of $B \ot_K \ovl K$, hence on that of $\hat C_i \ot_{\hat R} \Omega'$. Since $\hat R \to \Omega'$ is injective, the image of $z$ in $\hat C_i \cong \M_n(\hat R)$ is the scalar matrix $c$, whose reduction modulo $\pi$ is the scalar matrix $\bar c$.
\end{proof}

\begin{remark} \normalfont
\label{zhu:rem:descent-counterexamples}
The finiteness of $\Lambda$ cannot be removed, and the subset $I'$ can be proper. For the $R$-algebra $\Lambda = R[x]/(\pi x^2 - x)$, which is not finite over $R$, we have $\Lambda/\mathfrak m\Lambda \cong k$ while $\Lambda \ot_R K \cong K \times K$. For the finite $R$-algebra $\Lambda = R \times k$, we have $\Lambda/\mathfrak m\Lambda \cong k \times k$ while $\Lambda \ot_R K \cong K$. These are the two phenomena which a computation of a presentation modulo a prime cannot control by itself: the generic dimension may exceed the special one in the absence of a finiteness statement over $R$, and a special block may disappear generically. The first is excluded below by an explicit spanning set over $R$, and the second by the equality of the generic and special dimensions (Theorem \ref{zhu:thm:order}). The bijection of part 1 of Lemma \ref{zhu:lem:descent} depends on the choices of the embeddings of $\ovl K$, and only the sizes $n_i$, the ranks of the idempotents of $\Lambda$ and the central scalars in $R$ are transported by it. These are the only data used below. In particular the lemma alone does not distinguish two blocks with the same such data, as the blocks $a$ and $c_0$ of Section \ref{zhu:sec:modules}.
\end{remark}

\begin{remark} \normalfont
\label{zhu:rem:trace-form}
If the characteristic of $k$ does not divide any of the integers $n_i$, the semisimplicity of $\Lambda \ot_R K$ has a shorter proof. With the notation of the proof, the regular trace form $(x,y) \mapsto \Tr(L_{xy})$ of the free $R$-algebra $C$ has a Gram matrix with entries in $R$ whose reduction modulo $\mathfrak m$ is the Gram matrix of the regular trace form of $C/\mathfrak m C \cong \prod_{i \in I'} \M_{n_i}(k)$. The latter is nondegenerate, since on $\M_n(k)$ the regular trace is $n$ times the matrix trace. Hence the discriminant is a unit of $R$, the regular trace form of $C \ot_R K$ is nondegenerate, and $C \ot_R K$ is semisimple: if $x$ lies in its radical, then $xy$ is nilpotent for any $y$, so that $\Tr(L_{xy}) = 0$, and $x = 0$.
\end{remark}
%%%%%%%%%%%%%%%%%%%%%%%%%%%%%%%%%%%%%%%%%%%%%%%%%%%%%%%%%%%%%%%%
\subsection{The Morita corner of Zhu's algebra and its quiver presentation}
\label{zhu:sec:corner}

%%%%%%%%%%%%%%%%%%%%%%%%%%%%%%%%%%%%%%%%%%%%%%%%%%%%%%%%%%%%%%%%
\subsubsection{Matrix units and full idempotents}
\label{zhu:subsec:peirce}

The following elementary statement is the algebraic part of \cite[Lemma 8.2 and Proposition 8.3 pp.~34--35]{Rie26b}. We include its proof.

\begin{lemma}
\label{zhu:lem:peirce}
Let $\mathcal A$ be a unital algebra over a field $F$, let $(T_v)_{v \in \mathsf V}$ be a finite family of nonzero finite-dimensional $F$-vector spaces with fixed bases, let $\mathcal B = \prod_{v} \operatorname{End}_F(T_v)$ and let $\varphi \co \mathcal B \to \mathcal A$ be a unital homomorphism of algebras. Denote by $E^{(v)}_{ij}$ the matrix units of $\operatorname{End}_F(T_v)$ and put $e_v \ov{\mathrm{def}}{=} \varphi(E^{(v)}_{11})$, $e \ov{\mathrm{def}}{=} \sum_v e_v$ and $\mathcal D \ov{\mathrm{def}}{=} e \mathcal A e$.
\begin{enumerate}
\item The idempotent $e$ is full, that is $\mathcal A e \mathcal A = \mathcal A$.
\item The map
\[
\mathcal X \co \mathcal A \to \bigoplus_{v,w} \operatorname{Hom}_F(T_w, T_v) \ot_F e_v \mathcal D e_w,\quad b \mapsto \sum_{v,w,i,j} E^{v,w}_{ij} \ot \varphi(E^{(v)}_{1i})\, b\, \varphi(E^{(w)}_{j1}),
\]
where $E^{v,w}_{ij} \in \operatorname{Hom}_F(T_w,T_v)$ are the matrix units in the fixed bases, is an isomorphism of algebras, the right-hand side being equipped with the product $(F \ot d)(G \ot c) = FG \ot dc$ when the intermediate indices agree and $0$ otherwise.
\item The functor $M \mapsto eM$ is an equivalence between the categories of left $\mathcal A$-modules and of left $\mathcal D$-modules. In particular $\mathcal A$ is semisimple if and only if $\mathcal D$ is semisimple, and for a simple $\mathcal A$-module $M$ and any $v$, the dimension of $e_v M$ is the multiplicity of $T_v$ in $M$ considered as a $\mathcal B$-module through $\varphi$.
\end{enumerate}
\end{lemma}

\begin{proof}
1. We have $1 = \sum_{v,i} \varphi(E^{(v)}_{ii}) = \sum_{v,i} \varphi(E^{(v)}_{i1})\, e_v\, \varphi(E^{(v)}_{1i}) \in \mathcal A e \mathcal A$.

2. Since $e_v \varphi(E^{(v)}_{1i}) = \varphi(E^{(v)}_{1i})$ and $\varphi(E^{(w)}_{j1}) e_w = \varphi(E^{(w)}_{j1})$, the map $\mathcal X$ takes values in the indicated space. The linear map $\Psi(E^{v,w}_{ij} \ot d) \ov{\mathrm{def}}{=} \varphi(E^{(v)}_{i1})\, d\, \varphi(E^{(w)}_{1j})$ satisfies $\Psi(\mathcal X(b)) = \sum_{v,w,i,j} \varphi(E^{(v)}_{ii})\, b\, \varphi(E^{(w)}_{jj}) = b$ and $\mathcal X(\Psi(E^{v,w}_{ij} \ot d)) = E^{v,w}_{ij} \ot d$, since $\varphi(E^{(u)}_{1k}) \varphi(E^{(v)}_{i1}) = \delta_{uv}\delta_{ki} e_v$ and $d \in e_v \mathcal A e_w$. For the product, $\Psi(E^{v,w}_{ij} \ot d)\Psi(E^{w',u}_{kl} \ot c) = \delta_{ww'}\delta_{jk} \varphi(E^{(v)}_{i1})\, d c\, \varphi(E^{(u)}_{1l})$.

3. We show that $N \mapsto \mathcal A e \ot_{\mathcal D} N$ is a quasi-inverse. By part 1, write $1 = \sum_k x_k e y_k$ with $x_k, y_k \in \mathcal A$. For a left $\mathcal A$-module $M$, the map $\mu \co \mathcal A e \ot_{\mathcal D} eM \to M$, $a \ot m \mapsto am$, is surjective, since $m = \sum_k x_k e (y_k m)$. If $\sum_i a_i e \ot e m_i$ lies in its kernel, then
\[
\sum_i a_i e \ot e m_i = \sum_{k,i} x_k e\, (e y_k a_i e) \ot e m_i = \sum_k x_k e \ot e y_k \Big(\sum_i a_i e m_i\Big) = 0,
\]
because $e y_k a_i e \in \mathcal D$. Thus $\mu$ is an isomorphism, natural in $M$. For a left $\mathcal D$-module $N$, the decomposition $\mathcal A e = e\mathcal A e \oplus (1-e)\mathcal A e$ of right $\mathcal D$-modules gives $e(\mathcal A e \ot_{\mathcal D} N) = e \mathcal A e \ot_{\mathcal D} N = N$, naturally in $N$. Hence the two functors are inverse equivalences. An algebra is semisimple if and only if every short exact sequence of left modules splits, and an equivalence of abelian categories preserves split short exact sequences, so $\mathcal A$ is semisimple if and only if $\mathcal D$ is. The last assertion follows from $e_v M \cong \operatorname{Hom}_{\mathcal B}(T_v, M)$, since $e_v$ is the image of a rank-one idempotent of the factor $\operatorname{End}_F(T_v)$.
\end{proof}

%%%%%%%%%%%%%%%%%%%%%%%%%%%%%%%%%%%%%%%%%%%%%%%%%%%%%%%%%%%%%%%%
\subsubsection{The affine subalgebra}
\label{zhu:subsec:affine}

Recall that $A = L_{\mathfrak{sl}_2}(3,0) \ot L_{\mathfrak{sl}_2}(13,0)$ is the affine vertex subalgebra of $S$ generated by $S_1 = \mathfrak g = \mathfrak g' \oplus \mathfrak g''$ and that $A^+ = A \oplus J$ with $J = L(3,3) \ot L(13,13)$ is a vertex subalgebra of $S$ \cite[Proposition 11.4]{Arh26}. For $0 \leq r \leq 3$ and $0 \leq l \leq 13$, write $L(r,l) = L_{\mathfrak{sl}_2}(3,r) \ot L_{\mathfrak{sl}_2}(13,l)$ and $V(r,l)$ for the irreducible $\mathfrak g$-module with highest weight $(r,l)$, its lowest weight space. The lowest conformal weight of $L(r,l)$ is $h(r,l) = \frac{r(r+2)}{20} + \frac{l(l+2)}{60}$, and the fusion product of $J$ with $L(r,l)$ is $L(3-r,13-l)$. A direct computation gives
\begin{equation}
\label{zhu:eq:h-difference}
h(3-r,13-l) - h(r,l) = 4 - \frac{r+l}{2}.
\end{equation}
Since $A$ is rational and $J$ is a simple current of order two with integral lowest weight $h(3,13) = 4$, the simple $A^+$-modules are the modules $L(r,l) \oplus L(3-r,13-l)$ with $r+l$ even \cite[proof of Proposition 11.4]{Arh26}, which uses \cite[Theorem 3.4 p.~16]{HKL15}, and $A^+$ is rational \cite[p.~58]{Arh26}. By \eqref{zhu:eq:h-difference}, the lowest weight space of such a module is $V(r,l)$ if $r + l < 8$, and $V(r,l) \oplus V(3-r,13-l)$ if $r + l = 8$. We obtain fourteen simple $A^+$-modules, whose lowest weight spaces $T_0, \ldots, T_{13}$ are, as $\mathfrak g$-modules, $T_v = V(\lambda_v)$ for $0 \leq v \leq 11$ and
\[
T_{12} = V(1,7) \oplus V(2,6), \quad T_{13} = V(0,8) \oplus V(3,5),
\]
where $\lambda_0, \ldots, \lambda_{11}$ are the twelve pairs $(0,0)$, $(0,2)$, $(1,1)$, $(2,0)$, $(0,4)$, $(1,3)$, $(2,2)$, $(3,1)$, $(0,6)$, $(1,5)$, $(2,4)$, $(3,3)$. We call constituents the sixteen pairs $\lambda$ which occur in these lowest weight spaces. Since $A^+$ is rational, its Zhu algebra is semisimple and
\[
A(A^+) \cong \prod_{v=0}^{13} \operatorname{End}_\Cbb(T_v)
\]
\cite[Theorems 2.2.2 and 2.2.3 p.~261]{Zhu96}. By \cite[Proposition 11.4]{Arh26}, the inclusion $A^+ \subset S$ induces a unital homomorphism $\varphi \co A(A^+) \to A(S)$.

We apply Lemma \ref{zhu:lem:peirce} to $\mathcal A = A(S)$, $\mathcal B = A(A^+)$, with the bases of $V(r,l) = V(r) \ot V(l)$ given by the vectors $f_1^i f_2^j w$ of the highest weight vector $w$, the second factor varying fastest, and, for $v \in \{12, 13\}$, the first constituent listed first. We denote by $e_v \in A(S)$, $0 \leq v \leq 13$, the resulting idempotents and by
\[
\D \ov{\mathrm{def}}{=} e_{\mathsf V} A(S) e_{\mathsf V}, \qquad e_{\mathsf V} = \sum_{v=0}^{13} e_v,
\]
the corner. By Lemma \ref{zhu:lem:peirce}, $\D$ is Morita equivalent to $A(S)$.

%%%%%%%%%%%%%%%%%%%%%%%%%%%%%%%%%%%%%%%%%%%%%%%%%%%%%%%%%%%%%%%%
\subsubsection{Arrows and generation}
\label{zhu:subsec:arrows}

For $x \in S_1$ and $a \in S$, we have $[x] * [a] - [a] * [x] = [x_{(0)} a]$ in $A(S)$ \cite[Lemma 2.1.3 p.~248]{Zhu96}, \cite[(11.1)]{Arh26}. The class $[x]$ is the image under $\varphi$ of the class of $x$ in $A(A^+)$, which acts on $T_v$ by the zero mode, that is by the action $\rho_v(x)$ of $x \in \mathfrak g$. Hence $\mathcal X([x]) = \sum_v \rho_v(x) \ot e_v$, and $\mathcal X$ intertwines the derivation $\operatorname{ad}[x]$ of $A(S)$ with the natural action of $\mathfrak g$ on $\bigoplus_{v,w} \operatorname{Hom}(T_w,T_v)$, the factor $e_v\D e_w$ being acted upon trivially.

Let $z \in S_2$ be a highest weight vector of weight $\nu$ for $\mathfrak g$. Then $\mathcal X([z])$ is a highest weight vector of weight $\nu$. For constituents $\lambda \subset T_w$ and $\mu \subset T_v$, the space of highest weight vectors of weight $\nu$ in $\operatorname{Hom}(V(\lambda),V(\mu)) \cong V(\lambda)^* \ot V(\mu)$ has dimension at most one, by the Clebsch--Gordan rule for each factor of $\mathfrak g = \mathfrak{sl}_2 \oplus \mathfrak{sl}_2$. When it is nonzero, we say that $(\lambda, \mu)$ is admissible for $\nu$. A highest weight vector $F \neq 0$ of weight $\nu$ in $V(\lambda)^* \ot V(\mu)$ has a nonzero entry in row $1$ (the highest weight vector of $V(\mu)$) and column $(d_1,d_2)$, $d_k = (\nu_k - \mu_k + \lambda_k)/2$: for each factor $\mathfrak{sl}_2$, write a highest weight vector of the tensor product of two irreducible modules as $u = \sum_{i=0}^k c_i\, x_i \ot y_{k-i}$, where $(x_i)$ and $(y_j)$ are weight bases of the two factors, starting from their highest weight vectors, on which $e$ acts by $e x_i = \alpha_i x_{i-1}$ and $e y_j = \beta_j y_{j-1}$ with $\alpha_i, \beta_j \neq 0$ for $i, j \geq 1$ in the relevant range. The equation $eu = 0$ gives $c_{i+1}\alpha_{i+1} + c_i \beta_{k-i} = 0$ for $0 \leq i < k$, so that all the coefficients $c_i$ are nonzero if one of them is. We let $F^{\lambda,\mu}_\nu$ be the highest weight vector of weight $\nu$ normalized by the condition that this entry equals $1$. The component of $\mathcal X([z])$ in $\operatorname{Hom}(V(\lambda), V(\mu)) \ot e_v \D e_w$ is then $F^{\lambda,\mu}_\nu \ot x_{z,\lambda,\mu}$ for a unique element $x_{z,\lambda,\mu} \in e_v \D e_w$.

The space $S_2$ contains seventeen highest weight vectors for $\mathfrak g$ \cite[Proposition 11.3]{Arh26}. Seven of them lie in the affine subalgebra $A$, and their classes are polynomials in the classes of the currents. The ten others are the vectors $y$ and $x_C$ of type $(0,0)$ and the highest weight vectors of the primary types $(0,4)$, $(1,1)$, $(1,3)$, $(1,5)$, $(1,7)$, $(2,2)$, $(2,6)$, $(3,3)$. We call arrows the elements $x_{z,\lambda,\mu}$ attached to these ten vectors $z$ and to the pairs of constituents admissible for their weights. An arrow $x_{z,\lambda,\mu}$ goes from the vertex $w$ of $\lambda$ to the vertex $v$ of $\mu$. There are $515$ arrows. Let $\Gamma_Z$ be the quiver with vertices $0, \ldots, 13$ and these arrows, and $\Cbb\Gamma_Z$ its path algebra, with the convention that a path is written from left to right as a product in $\D$, so that a product $xy$ of arrows is nonzero in $\Cbb\Gamma_Z$ only if the source of $x$ is the target of $y$.

\begin{prop}
\label{zhu:prop:generation}
The homomorphism of algebras $\mathrm{ev} \co \Cbb\Gamma_Z \to \D$ which maps the vertex $v$ to $e_v$ and each arrow to the corresponding element of $\D$ is surjective.
\end{prop}

\begin{proof}
By \cite[Proposition 11.3]{Arh26}, $A(S)$ is generated by the classes of the six currents and of the seventeen highest weight vectors of $S_2$, and the classes of the seven affine highest weight vectors are polynomials in the classes of the currents. By Lemma \ref{zhu:lem:peirce}, $\mathcal X(A(S))$ is generated by the elements $\sum_v \rho_v(x) \ot e_v$, $x \in \mathfrak g$, and $\mathcal X([z]) = \sum F^{\lambda,\mu}_\nu \ot x_{z,\lambda,\mu}$ for the ten vectors $z$ above. A product of such generators is a sum of tensors $F \ot p$, where $F$ is a numerical matrix and $p$ a path in the arrows or a vertex idempotent. For $b \in e_v A(S) e_w$, the element $b = e_v b e_w$ is the coefficient of $E^{v,w}_{11}$ in $\mathcal X(b)$, which is therefore a linear combination of paths.
\end{proof}

%%%%%%%%%%%%%%%%%%%%%%%%%%%%%%%%%%%%%%%%%%%%%%%%%%%%%%%%%%%%%%%%
\subsubsection{Projected relations}
\label{zhu:subsec:projected}

Let $\{z_b\}_{0 \leq b \leq 119}$ be the basis of $S_2$ formed by the vectors $f_1(0)^i f_2(0)^j z$, where $z$ runs over the seventeen highest weight vectors \cite[Proposition 11.3]{Arh26}. A Zhu polynomial is a noncommutative polynomial $\Pi$ in symbols $Z_b$, $0 \leq b \leq 119$, and $J_x$, $x$ in a basis of $\mathfrak g$. Its value $\Pi([\cdot])$ in $A(S)$ is obtained by substituting $[z_b]$ and $[x]$. We define a formal matrix $\widehat{\mathcal X}(\Pi) \in \bigoplus_{v,w} \operatorname{Hom}(T_w,T_v) \ot \Cbb\Gamma_Z$ by substituting $\widehat{\mathcal X}(J_x) = \sum_v \rho_v(x) \ot e_v$, $\widehat{\mathcal X}(Z_b) = \sum F^{\lambda,\mu}_\nu \ot x_{z,\lambda,\mu}$ for the ten non-affine highest weight vectors $z = z_b$, the corresponding polynomials in the $\widehat{\mathcal X}(J_x)$ for the seven affine ones, and $\widehat{\mathcal X}(Z_{b'}) = \operatorname{ad}(\rho(f_1))^i \operatorname{ad}(\rho(f_2))^j \widehat{\mathcal X}(Z_b)$ when $z_{b'} = f_1(0)^i f_2(0)^j z_b$. By construction $(\mathrm{id} \ot \mathrm{ev})\widehat{\mathcal X}(\Pi) = \mathcal X(\Pi([\cdot]))$.

\begin{lemma}
\label{zhu:lem:projected}
Let $\Pi$ be a Zhu polynomial whose value in $A(S)$ is zero. Then every entry of $\widehat{\mathcal X}(\Pi)$ lies in the kernel of $\mathrm{ev}$. Moreover, for any arrow $x$ from a constituent $\lambda$ of $T_w$ to a constituent $\mu$ of $T_v$, the element $(h_\mu - h_\lambda)\, x + 2\, x_{y,\mu,\mu}\, x - 2\, x\, x_{y,\lambda,\lambda}$ lies in the kernel of $\mathrm{ev}$, where $h_\lambda = h(\lambda)$.
\end{lemma}

\begin{proof}
The first assertion follows from $(\mathrm{id} \ot \mathrm{ev})\widehat{\mathcal X}(\Pi) = \mathcal X(0) = 0$. For the second, write $\omega = \omega_A + \omega_C$ with $\omega_C = 2y$ \cite[proof of Proposition 11.4]{Arh26}. The class $[\omega]$ is central in $A(S)$ \cite[Theorem 2.1.1 p.~247]{Zhu96}. The class $[\omega_A]$ lies in the image of $A(A^+)$ and acts on $T_v$ by $L(0)$, that is by the scalar $h_\lambda$ on each constituent $\lambda$ of $T_v$ (both constituents of $T_{12}$, respectively $T_{13}$, have the same lowest weight). Since $y$ is a highest weight vector of weight $(0,0)$, $\mathcal X([y]) = \sum_\lambda \mathrm{id}_{V(\lambda)} \ot x_{y,\lambda,\lambda}$, the matrix $F^{\lambda,\lambda}_{(0,0)}$ being the identity. The component of $\mathcal X([\omega] * [z] - [z] * [\omega]) = 0$ on the line $F^{\lambda,\mu}_\nu$ gives the stated relation.
\end{proof}

For a Zhu polynomial $\Pi$ whose value is zero and which is a highest weight vector of weight $\eta$ for the adjoint action of $\mathfrak g$, we retain, for every pair $(\lambda,\mu)$ admissible for $\eta$, the entry of $\widehat{\mathcal X}(\Pi)$ in row $1$ and column $(d_1,d_2)$ of the block $\operatorname{Hom}(V(\lambda),V(\mu))$. These entries, together with the relations of the second part of Lemma \ref{zhu:lem:projected}, are the projected relations. Only the fact that they lie in the kernel of $\mathrm{ev}$ is used below. Completeness of these relations is not needed.
%%%%%%%%%%%%%%%%%%%%%%%%%%%%%%%%%%%%%%%%%%%%%%%%%%%%%%%%%%%%%%%%
\subsection{Exact relations and an integral upper presentation}
\label{zhu:sec:integral}

%%%%%%%%%%%%%%%%%%%%%%%%%%%%%%%%%%%%%%%%%%%%%%%%%%%%%%%%%%%%%%%%
\subsubsection{The field, the prime and the exact model}
\label{zhu:subsec:field}

Let $K = \mathbb{Q}(\sqrt{13}, \sqrt{-3})$, with $\mathbb{Q}$-basis $1, r, q, rq$, $r^2 = 13$, $q^2 = -3$. The algebra $S$ is defined over $K$ in the explicit model of $V' = V_{E_6 \oplus A_2}$ of \cite[Section 10.1]{Arh26}: the spaces $U_1$, $U_2$ are spanned by vectors with coordinates in $K$ in the monomial basis of $V'$, and all products below are computed exactly in this basis. The projected relations below are therefore elements of $K\Gamma_Z$, and their images under $\mathrm{ev} \co \Cbb\Gamma_Z \to \D$ vanish by Lemma \ref{zhu:lem:projected}.

Let $p = 8388571$. Then $13$ and $-3$ are squares modulo $p$, so that $p$ splits completely in $K$, and we fix the prime $\mathfrak p$ of the ring of integers $\mathcal O_K$ above the kernel of the homomorphism $\mathbb{Z}[r,q] \to \mathbb F_p$ sending $r$ and $q$ to the square roots $1649360$ and $2036966$ of $13$ and $-3$ modulo $p$ (the embedding used in \cite[Theorem 11.6]{Arh26}). This prime is unique since $p$ does not divide the index of $\mathbb{Z}[r,q]$ in $\mathcal O_K$, which divides a power of $2 \cdot 3 \cdot 13$. We denote by $R$ the localization of the ring of integers of $K$ at $\mathfrak p$. It is a discrete valuation ring with residue field $\mathbb F_p$, uniformizer $p$, and an element of $K$ whose four rational coordinates have denominators prime to $p$ lies in $R$. We write $x \mapsto \bar x$ for the reduction $R \to \mathbb F_p$.

%%%%%%%%%%%%%%%%%%%%%%%%%%%%%%%%%%%%%%%%%%%%%%%%%%%%%%%%%%%%%%%%
\subsubsection{The three families of relations over \texorpdfstring{$K$}{K}}
\label{zhu:subsec:families}

For $a, b \in S_2$ we put
\begin{align*}
B(a,b) &\ov{\mathrm{def}}{=} a_{(0)}b - \tfrac12 L(-1)(a_{(1)}b) + \tfrac14 L(-1)^2(a_{(2)}b) = \tfrac12\big(a_{(0)}b - b_{(0)}a\big),\\
N(a,b) &\ov{\mathrm{def}}{=} a_{(-1)}b - \tfrac12 L(-1)(a_{(0)}b) + \tfrac14 L(-1)^2(a_{(1)}b) - \tfrac1{12}L(-1)^3(a_{(2)}b) = \tfrac12\big(a_{(-1)}b + b_{(-1)}a\big),\\
c_\pm(a,b) &\ov{\mathrm{def}}{=} \tfrac12\big(a_{(1)}b \pm b_{(1)}a\big) \in S_2,
\end{align*}
the equalities following from the skew symmetry $b_{(n)}a = \sum_{j \geq 0} (-1)^{n+j+1} \frac{L(-1)^j}{j!} a_{(n+j)}b$ and from $L(-1)\mathbf 1 = 0$. By \cite[Lemma 2.1.1 p.~247 and (2.1.4) p.~248]{Zhu96}, in $A(S)$ we have, for $a, b \in S_2$, $v \in S$ homogeneous and $x \in S_1$,
\begin{equation}
\label{zhu:eq:zhu-formulas}
[L(-1)v] = -\operatorname{wt}(v)[v],\quad [x_{(-1)}v] = [v] * [x],\quad [B(a,b)] = [a]*[b] - [b]*[a] - [c_-(a,b)],
\end{equation}
\begin{equation}
\label{zhu:eq:zhu-formulas2}
[N(a,b)] = \tfrac12\big([a]*[b] + [b]*[a]\big) + [c_+(a,b)],
\end{equation}
see also \cite[(11.1)]{Arh26}. We consider, as in \cite[Theorem 11.6]{Arh26} and following \cite[Section 9]{Rie26b}, the $\mathfrak g$-equivariant linear maps
\begin{align*}
\Phi_3 &\co S_2 \oplus (\mathfrak g \ot S_2) \oplus \Lambda^2 S_2 \to V'_3,\quad (v, h \ot w, a \wedge b) \mapsto L(-1)v + h_{(-1)}w + B(a,b),\\
\Phi_{\mathrm{sym}} &\co \mathrm{Sym}^2 S_2 \to V'_4,\quad a \cdot b \mapsto N(a,b),\\
\Phi_{\mathrm{mix}} &\co S_3 \oplus (\mathfrak g \ot S_3) \oplus \mathrm{Sym}^2 S_2 \to V'_4,\quad (u, h \ot v, a \cdot b) \mapsto L(-1)u + h_{(-1)}v + N(a,b).
\end{align*}
A highest weight vector of weight $\eta$ of the domain of one of these maps which lies in its kernel gives, by \eqref{zhu:eq:zhu-formulas} and \eqref{zhu:eq:zhu-formulas2}, a Zhu polynomial which vanishes in $A(S)$ and which is a highest weight vector of weight $\eta$ for the adjoint action of $\mathfrak g$. For $\Phi_{\mathrm{mix}}$, the classes of the elements of $S_3$ are expressed through a $\mathfrak g$-equivariant section of $\Phi_3$.

\begin{prop}
\label{zhu:prop:exact-families}
The following statements hold, with all computations performed exactly over $K$.
\begin{enumerate}
\item The highest weight vectors of the domains of $\Phi_3$, $\Phi_{\mathrm{sym}}$ and $\Phi_{\mathrm{mix}}$ form spaces of dimensions $553$, $477$ and $1027$. The kernels of the three maps on these spaces have dimensions $458$, $70$ and $532$ over $K$, and their reductions modulo $\mathfrak p$ have the same dimensions. The image of $\Phi_3$ is $S_3$, of dimension $956$, with $95$ highest weight vectors.
\item The $1060$ resulting Zhu polynomials have all their coefficients in $R$.
\item The projected relations of Section \ref{zhu:subsec:projected}, computed from these polynomials and from the $515$ relations of the second part of Lemma \ref{zhu:lem:projected}, are $19981$ elements of $R\Gamma_Z$, each supported on a single pair of vertices.
\end{enumerate}
\end{prop}

\begin{proof}
This is the computation of items 3 and 4 of Section \ref{zhu:app:certificates}. The highest weight vectors of the formal domains are computed over $\mathbb Q$ as integer kernels of the action of $e_1$ and $e_2$, which has integer coefficients in the bases $f_1^i f_2^j$ of the irreducible $\mathfrak g$-modules. For each weight $\eta$, let $m_1, \ldots, m_h$ be the images of these highest weight vectors, which are exact vectors of $V'$ with coordinates in $K$. Their coordinates have denominators prime to $p$, which is checked when they are reduced modulo $\mathfrak p$. Choose, by Gaussian elimination modulo $\mathfrak p$, columns $C$ and rows $I$ such that the $r \times r$ minor $[m_i]_{i \in I, j \in C}$ is nonzero modulo $\mathfrak p$, where $r$ is the rank modulo $\mathfrak p$. This minor is nonzero over $K$, so that the rank over $K$ is at least $r$. For each $i \notin I$, the unique vector $c^{(i)}$ with $c^{(i)}_i = 1$, support in $I \cup \{i\}$ and $\sum_k c^{(i)}_k [m_k]_C = 0$ is computed exactly by restriction of scalars. The identity $\sum_k c^{(i)}_k m_k = 0$ is then verified exactly on all coordinates. This gives $h - r$ independent vectors of the kernel over $K$, so that the rank over $K$ is exactly $r$. The coefficients of $c^{(i)}$ lie in $R$ since the minor is a unit of $R$, and their reductions form a basis of the kernel modulo $\mathfrak p$. The coordinates of $c_\pm(a,b)$ in the basis $(z_b)$ of $S_2$ are found by exact solution of a $120 \times 120$ system and verified exactly. The section of $\Phi_3$ used for $\Phi_{\mathrm{mix}}$ is given by rational combinations of domain elements, chosen among the highest weight vectors whose images are independent modulo $\mathfrak p$, and their images span $S_3$ since their rank modulo $\mathfrak p$ is $956 = \dim S_3$ \cite[Theorem 10.21]{Arh26}. The projection uses the rational matrices $F^{\lambda,\mu}_\nu$ and $\rho_v(x)$. Finally, every denominator of every coefficient is checked to be prime to $p$.
\end{proof}

\begin{remark} \normalfont
\label{zhu:rem:no-rank-drop}
Part 1 answers the first obstruction recorded in \cite[Remark 11.7]{Arh26}: no rank drop occurs at $\mathfrak p$, and the reductions of the exact relations span the kernels of the reduced maps, which are the relations computed modulo $p$ in \cite[Theorem 11.6]{Arh26}. The polynomials themselves are not identical to those of \cite{Arh26}, since the bases of the kernels and the section of $\Phi_3$ are chosen differently, and this is the reason why the special fibre is recomputed in Section \ref{zhu:sec:special} from the reductions of the exact relations.
\end{remark}

%%%%%%%%%%%%%%%%%%%%%%%%%%%%%%%%%%%%%%%%%%%%%%%%%%%%%%%%%%%%%%%%
\subsubsection{The integral order}
\label{zhu:subsec:order}

Let $I_K \subset K\Gamma_Z$ be the two-sided ideal generated by the $19981$ projected relations of Proposition \ref{zhu:prop:exact-families}, and put
\[
I_R \ov{\mathrm{def}}{=} I_K \cap R\Gamma_Z, \qquad \Lambda \ov{\mathrm{def}}{=} R\Gamma_Z / I_R.
\]
Then $\Lambda$ is an $R$-algebra without $\mathfrak p$-torsion, $\Lambda \ot_R K = K\Gamma_Z/I_K$ since $I_K = K I_R$, and $\Lambda \ot_R \Cbb = \Cbb\Gamma_Z/\Cbb I_K$. By Lemma \ref{zhu:lem:projected} and Proposition \ref{zhu:prop:generation}, $\mathrm{ev}$ induces a surjective homomorphism of algebras $\Lambda \ot_R \Cbb \to \D$. The arrows are numbered by the generating vector $z$, in the order $y$, $x_C$, $(0,4)$, $(1,1)$, $(1,3)$, $(1,5)$, $(1,7)$, $(2,2)$, $(2,6)$, $(3,3)$ of their types, then by the source constituent and then by the target constituent, the constituents being ordered as in Section \ref{zhu:subsec:affine}. For reduction we compare paths first by length, with longer paths larger, and for equal length we take the reverse of the lexicographic order on the tuples of arrow indices. This order is compatible with left and right concatenation and has the descending chain condition. A rewriting rule is an element of $R\Gamma_Z$ of the form $w - \sum_u c_u u$, with $c_u \in R$, where every $u$ is shorter than $w$, or of the same length and lexicographically larger, and has the same endpoints as $w$. Its leading word is $w$.

\begin{lemma}
\label{zhu:lem:finiteness}
Let $\mathsf V' \subset \mathsf V$ be a set of vertices and $\mathcal E$ a set of arrows, let $\mathcal R$ be a finite set of rewriting rules, and assume:
\begin{enumerate}
\item for every vertex $v \notin \mathsf V'$, the idempotent $v$ lies in $I_R$, and for every arrow $x \in \mathcal E$, there is an element $x - \ell_x \in I_R$, where $\ell_x$ is an $R$-linear combination of vertices and of arrows with numerically larger indices than $x$,
\item every element of $\mathcal R$ lies in $I_R$,
\item every path of length $n$ in the arrows not in $\mathcal E$ between vertices of $\mathsf V'$ contains the leading word of a rule of $\mathcal R$.
\end{enumerate}
Then $\Lambda$ is spanned as an $R$-module by the images of the vertices of $\mathsf V'$ and of the paths of length at most $n - 1$ in the arrows not in $\mathcal E$ between vertices of $\mathsf V'$. In particular $\Lambda$ is a finite $R$-algebra.
\end{lemma}

\begin{proof}
By 1, every path through a vertex not in $\mathsf V'$ vanishes in $\Lambda$, and every arrow of $\mathcal E$ can be replaced by a combination of vertices and arrows with larger indices. Since the set of arrows is finite and totally ordered, finitely many replacements express every path as an $R$-combination of paths in the arrows not in $\mathcal E$, without increasing the length. If a path $w$ of length at least $n$ in these arrows contains the leading word $q$ of a rule $q - \sum c_u u$, write $w = w_1 q w_2$. Then $w \equiv \sum c_u w_1 u w_2$ modulo $I_R$, and each $w_1 u w_2$ is shorter than $w$, or of the same length and lexicographically larger. Since the set of paths of a fixed length is finite, a repetition of this step, together with the previous replacements, terminates and expresses $w$ modulo $I_R$ as an $R$-combination of shorter paths. Induction on the length concludes.
\end{proof}

\begin{thm}
\label{zhu:thm:order}
The following hold, with exact computations over $K$.
\begin{enumerate}
\item The hypotheses of Lemma \ref{zhu:lem:finiteness} hold with $n = 4$, $\mathsf V' = \{0, \ldots, 11\}$, a set $\mathcal E$ of $46$ arrows and a set $\mathcal R$ of $5659$ rules, $5221$ with leading words of length two and $438$ with leading words of length three. Consequently $\Lambda$ is a finite $R$-algebra.
\item Every leading coefficient used in the elimination is a unit of $R$, and every coefficient of the final eliminations and rules lies in $R$.
\item Let $J_K$ be the ideal of $K\Gamma_Z$ generated by the vertices $12$ and $13$, the elements $x - \ell_x$, $x \in \mathcal E$, and the rules of $\mathcal R$. Then $J_K = I_K$. All $131698$ overlap ambiguities and all $151$ inclusion ambiguities of $\mathcal R$ resolve over $K$, and the $762$ irreducible paths ($12$ vertices, $239$ arrows, $506$ paths of length two and $5$ of length three) form a basis of $\Lambda \ot_R K$.
\item $I_R$ is generated by the same elements as an ideal of $R\Gamma_Z$. The $R$-module $\Lambda$ is free with basis the $762$ irreducible paths, and $\Lambda/\mathfrak p\Lambda$ is the $\mathbb F_p$-algebra with this basis whose multiplication is given by the reductions of the rules.
\end{enumerate}
\end{thm}

\begin{proof}
This is the computation of item 5 of Section \ref{zhu:app:certificates}. It replays over $K$ the elimination of \cite[Theorem 11.6]{Arh26}: echelon forms on each pair of endpoints, with leading word the longest and then lexicographically smallest path, substitution of the linear eliminations and exclusion of the vertices $12$ and $13$ until stability, and completion by the overlaps of the rules. The completion has two stages: the $55444$ non-resolving overlaps of the $5105$ quadratic rules obtained after the eliminations produce, after echelon, the $5659$ rules of part 1, and all overlaps and inclusions of these rules resolve. Each division is by a leading coefficient, which is checked to be nonzero modulo $\mathfrak p$ with denominators prime to $p$, and the coefficients of the final rules and eliminations are checked to lie in $R$. The list of irreducible paths of each length is computed from the leading words. There are $239$, $506$, $5$ and $0$ irreducible paths of lengths $1$, $2$, $3$ and $4$, which gives hypothesis 3 of Lemma \ref{zhu:lem:finiteness} with $n = 4$. Hypotheses 1 and 2 of this lemma follow from the inclusion $J_K \subset I_K$ proved in the next paragraph and from part 2.

The inclusion $J_K \subset I_K$ holds by construction, since every element produced by the program of item 5 is an exact $K$-linear combination of products of paths with elements already known to lie in $I_K$. It is also certified independently by the programs of item 11 of Section \ref{zhu:app:certificates}, which express the vertices $12$ and $13$, the $46$ elements $x - \ell_x$ and the $5105$ quadratic rules as exact $K$-linear combinations of products of the generators with paths, and checks that each of the $554$ remaining rules reduces to zero modulo the span of the $5105$ certified rules and of their reduced overlaps. The inclusion $I_K \subset J_K$ is checked by the verifier of item 10 of Section \ref{zhu:app:certificates}: each of the $19981$ generators of $I_K$ reduces to zero modulo the rewriting data.

For part 3, substituting the eliminations and killing the vertices $12$ and $13$ gives an isomorphism $K\Gamma_Z/(e_{12}, e_{13}, x - \ell_x) \cong K\Gamma'$, where $\Gamma'$ is the quiver with the $12$ vertices of $\mathsf V'$ and the $239$ remaining arrows: the substitution is triangular for the order, no leading word contains an eliminated arrow, and the eliminations of the arrows at the vertices $12$ and $13$ reduce to zero (checked by the verifier of item 10 of Section \ref{zhu:app:certificates}). The path algebra $K\Gamma'$ is the tensor ring over $K^{\mathsf V'}$ of the bimodule spanned by the $239$ arrows, the order is a semigroup order with the descending chain condition, and every rule has its terms between the same endpoints as its leading word. By the diamond lemma \cite[Theorem 1.2 p.~181]{Ber78}, in the form of \cite[Theorem 6.1 p.~196]{Ber78} for tensor rings, the resolution of all ambiguities shows that the irreducible paths form a basis of $K\Gamma'/(\mathcal R) = K\Gamma_Z/J_K = \Lambda \ot_R K$.

For part 4, let $J_R$ be the ideal of $R\Gamma_Z$ generated by the same elements. The same reductions hold over $R$, since the rules are monic with coefficients in $R$, so that the diamond lemma over $R$ shows that $R\Gamma_Z/J_R$ is a free $R$-module with basis the irreducible paths. In particular it is torsion-free. Since $J_R \subset I_R$ and $(R\Gamma_Z/J_R) \ot_R K = K\Gamma_Z/J_K = \Lambda \ot_R K$, the kernel $I_R/J_R$ of the surjection $R\Gamma_Z/J_R \to \Lambda$ is a torsion submodule, hence zero. The last assertion follows by reduction modulo $\mathfrak p$.
\end{proof}

\begin{remark} \normalfont
\label{zhu:rem:order-free}
The unit pivots of Theorem \ref{zhu:thm:order}(2) allow reduction of the exact eliminations at $\mathfrak p$. The verifier of item 10 of Section \ref{zhu:app:certificates} compares the surviving vertices, active arrows and final rules with the modular computation on the reductions of the $19981$ relations. In particular, the modular computation of products of normal paths uses the multiplication of $\Lambda/\mathfrak p\Lambda$. The modular presentation from the reduced relations has dimension $762$, so its natural surjection onto $\Lambda/\mathfrak p\Lambda$ is an isomorphism.
\end{remark}
%%%%%%%%%%%%%%%%%%%%%%%%%%%%%%%%%%%%%%%%%%%%%%%%%%%%%%%%%%%%%%%%
\subsection{The special fibre}
\label{zhu:sec:special}

\begin{prop}
\label{zhu:prop:special}
Put $\D_p \ov{\mathrm{def}}{=} \Lambda/\mathfrak p\Lambda$.
\begin{enumerate}
\item The algebra $\D_p$ has the basis of the $762$ irreducible paths of Theorem \ref{zhu:thm:order}, and it is equal to the algebra presented by the reductions modulo $\mathfrak p$ of the $19981$ projected relations (Remark \ref{zhu:rem:order-free}).
\item The regular trace form of $\D_p$ is nondegenerate, its center has dimension $12$, a central element has twelve distinct eigenvalues in $\mathbb F_p$, and
\[
\D_p \cong \prod_{i=1}^{12} \M_{n_i}(\mathbb F_p), \qquad (n_i) = (1, 1, 2, 2, 2, 3, 5, 7, 11, 12, 12, 16).
\]
\item For each block, the ranks of the vertex idempotents on the simple module are those of Table \ref{zhu:table:blocks}.
\item Let $\lambda_v$ be the first constituent of $T_v$ and $h_v = h(\lambda_v)$, and put $Z \ov{\mathrm{def}}{=} \sum_{v=0}^{13} (h_v v + 2 x_{y,\lambda_v,\lambda_v}) \in R\Gamma_Z$. The image $\bar Z$ of $Z$ acts on the simple module of the block $i$ by the reduction of the rational number $h_i$ of Table \ref{zhu:table:blocks}.
\end{enumerate}
\end{prop}

\begin{proof}
Part 1 is part 4 of Theorem \ref{zhu:thm:order} and Remark \ref{zhu:rem:order-free}. Parts 2 and 3 are the computation of \cite[Theorem 11.6]{Arh26}, run on the reductions of the exact relations (item 6 of Section \ref{zhu:app:certificates}). By Remark \ref{zhu:rem:order-free} and the verifier of item 10 of Section \ref{zhu:app:certificates}, the surviving vertices, active arrows and final rules used by this program agree with the reductions of the exact presentation. The product on normal paths is therefore the product of $\D_p$. The regular trace form has rank $762$. If $x$ lies in the Jacobson radical, then $xy$ is nilpotent for every $y$, so that $\Tr(L_{xy}) = 0$, and nondegeneracy gives $x = 0$. Thus $\D_p$ is semisimple. A central element with twelve distinct eigenvalues in $\mathbb F_p$ gives twelve central idempotents by interpolation, so that the center is split. Each simple factor is a matrix algebra over a finite division ring, which is commutative \cite{Wed05}, hence equal to its center, which is $\mathbb F_p$ since the center of $\D_p$ is split. For part 4, the characteristic polynomial of left multiplication by $\bar Z$ on $\D_p$ factors as
\[
t\,(t - \tfrac13)(t - \tfrac{5}{13})^4(t - \tfrac{6}{13})^4(t-\tfrac{7}{13})^4(t - \tfrac{8}{13})^9(t - \tfrac23)^{25}(t - \tfrac{11}{13})^{49}(t - \tfrac{15}{13})^{256}(t-1)^{409}
\]
in $\mathbb F_p[t]$, where fractions denote their reductions modulo $p$ (item 7 of Section \ref{zhu:app:certificates}). Since $\bar Z$ is central, by the argument of the proof of Proposition \ref{zhu:prop:exact-h}, which uses only the relations of Lemma \ref{zhu:lem:projected}, and since $\D_p$ is semisimple, the kernel of $L_{\bar Z} - c$ is the sum of the blocks on which $\bar Z$ acts by $c$. For each root $c$, the program of item 7 of Section \ref{zhu:app:certificates} computes this kernel and the dimensions of its corners $e_v \ker(L_{\bar Z} - c) e_v$, which are the sums, over these blocks, of the squares of the ranks of $e_v$. For the nine roots $c \neq 1$, the kernel has dimension $n_i^2$ for a single block $i$, and the corner dimensions are the squares of the vertex ranks of the corresponding row of Table \ref{zhu:table:blocks}, which identifies the row. The three remaining blocks, of sizes $11$, $12$, $12$, carry the root $1$, of multiplicity $409 = 11^2 + 12^2 + 12^2$.
\end{proof}

\begin{table}[ht]
\centering
\footnotesize
\begin{tabular}{|c|c|p{8.2cm}|c|c|}
\hline
label & $n_i$ & lowest weight space as a $\mathfrak g$-module & $\dim$ & $h_i$\\
\hline
$0$ & $1$ & $(0,0)$ & $1$ & $0$\\
$c_1$ & $1$ & $(0,2)$ & $3$ & $1/3$\\
$c_2$ & $5$ & $(0,0)+(0,4)+(1,1)+(1,3)+(2,2)$ & $27$ & $2/3$\\
$b$ & $11$ & $2(0,2)+2(0,4)+(1,1)+2(1,3)+(1,5)+(2,2)+(2,4)+(3,1)$ & $80$ & $1$\\
$a$, $c_0$ & $12$ & $2(0,2)+(0,4)+(0,6)+2(1,1)+2(1,3)+(1,5)+(2,0)+(2,2)+(2,4)$ & $81$ & $1$\\
$d_1$ & $2$ & $(0,2)+(1,1)$ & $7$ & $6/13$\\
$d_4$ & $2$ & $(0,0)+(1,1)$ & $5$ & $5/13$\\
$d_5$ & $2$ & $(0,4)+(1,3)$ & $13$ & $7/13$\\
$d_6$ & $3$ & $(0,2)+(1,3)+(2,0)$ & $14$ & $8/13$\\
$d_2$ & $7$ & $(0,0)+(0,2)+(0,4)+(1,1)+(1,3)+(1,5)+(2,2)$ & $42$ & $11/13$\\
$d_3$ & $16$ & $2(0,0)+(0,2)+2(0,4)+(0,6)+2(1,1)+2(1,3)+2(1,5)+2(2,2)+(2,4)+(3,3)$ & $119$ & $15/13$\\
\hline
\end{tabular}
\caption{The twelve blocks of $\D_p$ and of $\Lambda \ot_R \ovl K$: matrix size $n_i$, lowest weight space $\bigoplus_v T_v^{m_{i,v}}$ as a $\mathfrak g$-module (where $(r,l)$ stands for $V(r,l)$), its dimension, and the eigenvalue $h_i$ of $Z$. The labels are those of \cite[Table 2]{Arh26}.}
\label{zhu:table:blocks}
\end{table}

%%%%%%%%%%%%%%%%%%%%%%%%%%%%%%%%%%%%%%%%%%%%%%%%%%%%%%%%%%%%%%%%
\subsection{Semisimplicity of Zhu's algebra and rationality}
\label{zhu:sec:rationality}

\begin{thm}
\label{zhu:thm:Z}
The algebra $\Lambda \ot_R K$ is semisimple and, for any algebraically closed field $\Omega$ containing $K$, $\Lambda \ot_R \Omega \cong \prod_{i=1}^{12} \M_{n_i}(\Omega)$ with the block sizes and vertex ranks of Proposition \ref{zhu:prop:special}. The corner $\D$ is semisimple and is a quotient of $\Lambda \ot_R \Cbb$. Zhu's algebra $A(S)$ is finite-dimensional and semisimple, that is, hypothesis $\mathrm{(Z)}$ of \cite[Theorem 11.1]{Arh26} holds.
\end{thm}

\begin{proof}
By Theorem \ref{zhu:thm:order}, $\Lambda$ is a finite $R$-algebra, free of rank $762$, and $\Lambda/\mathfrak p\Lambda = \D_p$ is split semisimple by Proposition \ref{zhu:prop:special}. Since $\Lambda$ has no torsion, the subset $I'$ of Lemma \ref{zhu:lem:descent} is the set of all twelve blocks, as $\dim_K \Lambda \ot_R K = 762 = \sum_i n_i^2$. Lemma \ref{zhu:lem:descent} gives the semisimplicity of $\Lambda \ot_R K$ and the decomposition of $\Lambda \ot_R \Omega$ into twelve blocks, with the same matrix sizes and vertex ranks as those of $\D_p$. By part 1 of Lemma \ref{zhu:lem:descent} with $\Omega = \Cbb$, $\Lambda \ot_R \Cbb \cong \prod_{i=1}^{12} \M_{n_i}(\Cbb)$, and the surjection $\Lambda \ot_R \Cbb \to \D$ of Section \ref{zhu:subsec:order} shows that $\D$ is a finite-dimensional semisimple algebra. By part 2 of Lemma \ref{zhu:lem:peirce}, $A(S)$ is finite-dimensional, and by part 3 it is semisimple.
\end{proof}

\begin{cor}
\label{zhu:cor:rational}
The vertex operator algebra $S$ is strongly rational. It has at most twelve simple modules up to isomorphism, and the lowest weight space of each simple $S$-module is, as a module over $A(A^+)$, isomorphic to $\bigoplus_v T_v^{m_{i,v}}$ for one of the rows $i$ of Table \ref{zhu:table:blocks}.
\end{cor}

\begin{proof}
The algebra $S$ is simple, of CFT type and self-contragredient \cite[Theorem 10.10]{Arh26}, and $C_2$-cofinite \cite[Theorem 10.21]{Arh26}. By Theorem \ref{zhu:thm:Z} and \cite[Corollary 5.19 p.~61]{McR21}, $S$ is rational, hence strongly rational since it has positive energy. By \cite[Theorem 2.2.2 p.~261]{Zhu96}, simple $S$-modules correspond to simple $A(S)$-modules through their lowest weight spaces, and by Lemma \ref{zhu:lem:peirce} and Theorem \ref{zhu:thm:Z} the simple $A(S)$-modules correspond to the blocks of $\D$, which form a subcollection of the twelve blocks of $\Lambda \ot_R \Cbb$. The vertex ranks give the multiplicities $m_{i,v}$ by part 3 of Lemma \ref{zhu:lem:peirce}. The class $[\omega]$ acts on the lowest weight space by $L(0)$. By the proof of Lemma \ref{zhu:lem:projected}, $\mathcal X([\omega]) = \sum_v h_v \mathrm{id}_{T_v} \ot e_v + 2\sum_\lambda \mathrm{id}_{V(\lambda)} \ot x_{y,\lambda,\lambda}$, so that $e_{\mathsf V}[\omega]e_{\mathsf V} = \sum_v e_v [\omega] e_v = \mathrm{ev}(Z)$, by the centrality of $[\omega]$. In particular $\mathrm{ev}(Z)$ is central in $\D$. The scalars by which $Z$ acts on the blocks of $\Lambda \ot_R \Cbb$, hence the lowest conformal weights, are determined in Proposition \ref{zhu:prop:exact-h} below.
\end{proof}
%%%%%%%%%%%%%%%%%%%%%%%%%%%%%%%%%%%%%%%%%%%%%%%%%%%%%%%%%%%%%%%%
\subsection{The twelve simple modules}
\label{zhu:sec:modules}

%%%%%%%%%%%%%%%%%%%%%%%%%%%%%%%%%%%%%%%%%%%%%%%%%%%%%%%%%%%%%%%%
\subsubsection{Exact conformal weights}
\label{zhu:subsec:exact-h}

\begin{prop}
\label{zhu:prop:exact-h}
In $\Lambda \ot_R K$ we have $\prod_{h \in H} (Z - h) = 0$, where
\[
H = \Big\{0, \frac13, \frac23, 1, \frac{5}{13}, \frac{6}{13}, \frac{7}{13}, \frac{8}{13}, \frac{11}{13}, \frac{15}{13}\Big\}.
\] Consequently, $Z$ acts on the simple module of the block $i$ of $\Lambda \ot_R \Omega$ by the rational number $h_i$ of Table \ref{zhu:table:blocks}. In particular, the lowest conformal weight of a simple $S$-module is one of the numbers $h_i$, and it is positive unless the module is $S$.
\end{prop}

\begin{proof}
The identity is checked exactly in the basis of irreducible paths of Theorem \ref{zhu:thm:order}, by computing $\prod_{h \in H}(Z-h)\cdot 1$ with $1 = \sum_{v=0}^{11} e_v$ (item 8 of Section \ref{zhu:app:certificates}). The element $Z$ is central in $\Lambda \ot_R K$: the vertices $12$ and $13$ vanish in $\Lambda$, and for an arrow $x$ between vertices $w, v \leq 11$, the relation of the second part of Lemma \ref{zhu:lem:projected} is $Zx - xZ$. Since $\Lambda \ot_R \Omega$ is semisimple, $Z$ acts on each block by a scalar, which is a root of $\prod_{h \in H}(t-h)$, hence an element of $H$. The elements of $H$ lie in $R$, since their denominators are prime to $p$. The image of $Z$ in $\Lambda/\mathfrak p\Lambda$ acts on the block $i$ by $\bar h_i$ (Proposition \ref{zhu:prop:special}), so that by part 3 of Lemma \ref{zhu:lem:descent} the scalar by which $Z$ acts on the block $i$ of $\Lambda \ot_R \Omega$ is an element of $H$ with reduction $\bar h_i$. The ten elements of $H$ have distinct reductions modulo $p$, which determines the scalar. The last assertion follows from Corollary \ref{zhu:cor:rational}.
\end{proof}

%%%%%%%%%%%%%%%%%%%%%%%%%%%%%%%%%%%%%%%%%%%%%%%%%%%%%%%%%%%%%%%%
\subsubsection{Restrictions of the lattice modules}
\label{zhu:subsec:restrictions}

The discriminant group $L^\circ/L$ of $L = E_6 \oplus A_2$ is $\Z_3 \times \Z_3$, and we denote by $W_{\epsi,\alpha}$, $\epsi,\alpha \in \Z_3$, the simple $V'$-modules as in \cite[Remark 11.2]{Arh26}, the first index referring to $E_6$ and the second to $A_2$. Their lowest conformal weights and lowest weight spaces are: $0$ and $\Cbb\mathbf 1$ for $W_{0,0} = V'$, $\frac13$ and a space of dimension $3$ for $W_{0,\pm1}$, $\frac23$ and a space of dimension $27$ for $W_{\pm1,0}$, $1$ and a space of dimension $81$ for $W_{\pm1,\pm1}$ and $W_{\pm1,\mp1}$. Moreover $\dim V'_1 = 86$ and $\dim S_1 = 6$.

\begin{prop}
\label{zhu:prop:restrictions}
The simple $S$-modules $M_0 = S$, $M_b$, $M_{c_1}$, $M_{c_2}$ with lowest weight spaces of the types $0$, $b$, $c_1$, $c_2$ of Table \ref{zhu:table:blocks} exist, and
\[
V'|_S \cong S \oplus M_b, \qquad W_{0,\pm1}|_S \cong M_{c_1}, \qquad W_{\pm1,0}|_S \cong M_{c_2}.
\]
Moreover each of the four modules $W_{\epsi,\alpha}$ with $\epsi\alpha \neq 0$ restricts to a simple $S$-module of conformal weight $1$ whose lowest weight space is of the type $a$ or $c_0$ of Table \ref{zhu:table:blocks}.
\end{prop}

\begin{proof}
By Corollary \ref{zhu:cor:rational}, every $V'$-module is a direct sum of simple $S$-modules, each of which has as lowest weight space one of the rows of Table \ref{zhu:table:blocks} and conformal weight $h_i$ by Proposition \ref{zhu:prop:exact-h}. A simple summand of $W_{\epsi,\alpha}$ has its lowest weight in the class of the conformal weights of $W_{\epsi,\alpha}$ modulo $\Z$, namely $\Z$, $\frac13 + \Z$ or $\frac23 + \Z$. The only rows with $h_i$ in these classes are the six rows $0$, $c_1$, $c_2$, $b$, $a$, $c_0$, with $h_i \in \{0, \frac13, \frac23, 1\}$. Hence the lowest weight space of every summand lies in degree at most one, and the multiplicities are determined by the dimensions of the spaces of $W_{\epsi,\alpha}$ in the degrees $0$, $\frac13$, $\frac23$ and $1$.

For $V'$, the space $V'_0 = \Cbb\mathbf 1$ gives one summand $S$, which is the vertex subalgebra $S$ itself. In degree one, $86 = 6 + 80\, m_b + 81\,(m_a + m_{c_0})$ forces $m_b = 1$ and $m_a = m_{c_0} = 0$. For $W_{0,\pm1}$, the only possible summands have conformal weight $\frac13$, and $3 = 3\,m_{c_1}$. For $W_{\pm1,0}$, the only possible summands have conformal weight $\frac23$, and $27 = 27\,m_{c_2}$. For $W_{\epsi,\alpha}$ with $\epsi\alpha \neq 0$, there is no vector of degree $0$, and $81 = 80\,m_b + 81\,(m_a + m_{c_0})$ forces a single summand of type $a$ or $c_0$. In each case the summand found is a simple $S$-module whose lowest weight space has the indicated type, which proves its existence. Finally $W_{-\epsi,-\alpha}$ is the contragredient module of $W_{\epsi,\alpha}$, and the restriction of a contragredient module is the contragredient of the restriction. The assertion $W_{\epsi,\alpha}|_S \cong W_{-\epsi,-\alpha}|_S$ is proved in Proposition \ref{zhu:prop:a-c0}.
\end{proof}

The last point requires an argument, since the rows $a$ and $c_0$ of Table \ref{zhu:table:blocks} cannot be distinguished by their conformal weights and lowest weight spaces. We use the genus-one modular matrices. For a strongly rational vertex operator algebra $U$ with simple modules $M_0 = U, \ldots, M_d$, let $s^U$ be the matrix of the transformation $\tau \mapsto -1/\tau$ of the trace functions \cite[Theorem 5.3.2 p.~299]{Zhu96}, which does not depend on the inserted vector. By \cite[Theorem 4.1 p.~6452]{DJX13}, which relies on the Verlinde formula of Huang \cite{Hua08b}, $s^U$ is symmetric, $(s^U)^2$ is the permutation matrix of contragredience and $s^U_{0,j} \neq 0$ for every $j$. If moreover the modules $M_i$, $i \neq 0$, have positive lowest weights, then $s^U_{i0}/s^U_{00} = \operatorname{qdim}_U M_i > 0$ \cite[Lemma 4.2 p.~6453 and its proof]{DJX13}. The lattice algebra $V'$ is strongly rational, its simple modules are the nine modules $W_g$, and they are simple currents of quantum dimension one \cite[Example 4.19 p.~6457]{DJX13}. Their contragredient modules are given by $W_g' \cong W_{-g}$, and their lowest conformal weights other than $0$ are positive.

\begin{lemma}
\label{zhu:lem:branching}
Let $U \subset V$ be strongly rational vertex operator algebras with the same conformal vector, let $(M_z)$ and $(W_g)$ be their simple modules, and let $\mathsf b_{gz}$ be the multiplicity of $M_z$ in $W_g|_U$, which is finite. Then $s^V \mathsf b = \mathsf b\, s^U$.
\end{lemma}

\begin{proof}
The multiplicity $\mathsf b_{gz}$ is at most the dimension of the homogeneous space of $W_g$ of weight $h_{M_z}$, which is finite. For $v \in U$ homogeneous for Zhu's grading $L[0]$, which is the same for $U$ and $V$, we have $T_{W_g}(v,\tau) = \sum_z \mathsf b_{gz} T_{M_z}(v,\tau)$ for the trace functions. Applying $\tau \mapsto -1/\tau$ \cite[Theorem 5.3.2 p.~299]{Zhu96} to both sides and comparing, we get $\sum_{z'} (s^V \mathsf b - \mathsf b s^U)_{g z'} T_{M_{z'}}(v,\tau) = 0$ for all such $v$. The trace functions $T_{M_{z'}}$ are linearly independent as functions of the insertion $v$ and of $\tau$ \cite[Theorem 5.3.1 p.~294 and its proof]{Zhu96}, which gives the claim.
\end{proof}

\begin{prop}
\label{zhu:prop:a-c0}
The four modules $W_{\epsi,\alpha}|_S$, $\epsi\alpha \neq 0$, are not all isomorphic. Consequently the six rows $0$, $b$, $c_1$, $c_2$, $a$, $c_0$ of Table \ref{zhu:table:blocks} are realized by simple $S$-modules. Moreover
\[
W_{1,1}|_S \cong W_{-1,-1}|_S \not\cong W_{1,-1}|_S \cong W_{-1,1}|_S,
\]
and these two simple modules are self-contragredient.
\end{prop}

\begin{proof}
Let $\mathsf b$ be the branching matrix of $S \subset V'$, with rows indexed by $g \in \Z_3 \times \Z_3$ and columns by the simple $S$-modules, and $\mathsf Q \ov{\mathrm{def}}{=} \mathsf b\, \mathsf b^T$, so that $\mathsf Q_{g,h}$ is the number of simple summands common to $W_g|_S$ and $W_h|_S$, counted with multiplicities. By Lemma \ref{zhu:lem:branching} and the symmetry of $s^S$ and $s^{V'}$, we have $\mathsf Q s^{V'} = \mathsf b\, s^S \mathsf b^T = s^{V'} \mathsf Q$. Since the modules $W_g$ have quantum dimension one, $s^{V'}_{g,0} = s^{V'}_{0,g} = s^{V'}_{00} \neq 0$ for every $g$. Suppose that the four modules $W_{\epsi,\alpha}|_S$ with $\epsi\alpha \neq 0$ are isomorphic. By Proposition \ref{zhu:prop:restrictions}, $\mathsf Q_{g,h} = 1$ for these four $g$ and $h$, $\mathsf Q_{(0,0),h} = 0$ for these $h$, and $\mathsf Q_{(0,0),(0,0)} = 2$. Then
\[
(\mathsf Q s^{V'})_{(0,0),(1,1)} = 2\, s^{V'}_{00}, \qquad (s^{V'}\mathsf Q)_{(0,0),(1,1)} = \sum_h s^{V'}_{(0,0),h}\mathsf Q_{h,(1,1)} = 4\, s^{V'}_{00},
\]
a contradiction. Hence both rows $a$ and $c_0$ are realized, and with Proposition \ref{zhu:prop:restrictions} the six rows are realized.

For the last assertion we use the full matrix $s^{V'}$. In the modular tensor category $\Rep(V')$ \cite{Hua08}, the simple objects $W_g$ are invertible, with twists $\vartheta_g = e^{2\pi i h_g}$, where $h_g$ is congruent to $\frac12(\lambda_g,\lambda_g)$ modulo $\Z$ for any $\lambda_g$ in the coset $g$. Their categorical dimensions are their quantum dimensions \cite[Theorem 3.10 p.~2142]{DLN15}, \cite[p.~5]{CKM22}, which are equal to one. For invertible simple objects of dimension one, the double braiding $c_{W_h,W_g}c_{W_g,W_h}$ is the scalar $\vartheta_{g+h}\vartheta_g^{-1}\vartheta_h^{-1} = e^{2\pi i(\lambda_g,\lambda_h)}$ by the balancing axiom, so that the unnormalized $S$-matrix of $\Rep(V')$ is $\tilde s_{g,h} = e^{2\pi i(\lambda_g,\lambda_h)}$. The normalized $S$-matrix of $\Rep(V')$ and the genus-one matrix $s^{V'}$ coincide up to a sign \cite[Theorem 3.10 p.~2142]{DLN15}. Allowing also the complex conjugate convention for the $S$-matrix of a modular category, we obtain $s^{V'}_{g,h} = \frac{\sigma}{3}\, \zeta^{3\eta(\lambda_g,\lambda_h)}$ for some signs $\sigma, \eta \in \{1,-1\}$, where $\zeta = e^{2\pi i/3}$ and $3(\lambda_g,\lambda_h) \equiv \epsi\epsi' + 2\alpha\alpha' \bmod 3$ for $g = (\epsi,\alpha)$, $h = (\epsi',\alpha')$.

Let $F$ be the set of the four elements $g = (\epsi,\alpha)$ with $\epsi\alpha \neq 0$, and for $g \in F$ let $[g] \subset F$ be the set of elements $h \in F$ with $W_h|_S \cong W_g|_S$. For $g, h \in F$, $\mathsf Q_{h,g}$ is $1$ if $h \in [g]$ and $0$ otherwise, and $\mathsf Q_{(\pm1,0),g} = \mathsf Q_{(0,\pm1),g} = 0$. For $k = (1,0)$, the row $k$ of $\mathsf Q$ is the indicator function of $\{(1,0),(-1,0)\}$ by Proposition \ref{zhu:prop:restrictions}. For $g = (\epsi,\alpha) \in F$, the entries $(k,g)$ of $\mathsf Q s^{V'}$ and $s^{V'} \mathsf Q$ are
\[
\frac{\sigma}{3}\big(\zeta^{\eta\epsi} + \zeta^{-\eta\epsi}\big) = -\frac{\sigma}{3}, \qquad \frac{\sigma}{3}\sum_{(\epsi',\alpha') \in [g]} \zeta^{\eta\epsi'}.
\]
A sum of $n_1$ numbers equal to $\zeta$ and $n_2$ numbers equal to $\zeta^{-1}$ is equal to $-1 = \zeta + \zeta^{-1}$ only if $n_1 = n_2 = 1$, since $\zeta$ and $\zeta^{-1}$ are linearly independent over $\mathbb{R}$. Hence $[g] = \{g, g^*\}$ with $g^* = (-\epsi,\alpha^*)$. The same computation with $k = (0,1)$, whose row in $\mathsf Q$ is the indicator function of $\{(0,1),(0,-1)\}$, gives $\alpha^* = -\alpha$, since $3(\lambda_{(0,1)},\lambda_{(\epsi',\alpha')}) \equiv 2\alpha' \bmod 3$. Thus $[g] = \{g,-g\}$ for every $g \in F$, which gives the stated isomorphisms, and self-contragredience follows since $W_{-g}$ is the contragredient of $W_g$ and $S$ and $V'$ have the same conformal vector. Item 9 of Section \ref{zhu:app:certificates} confirms the computation of the three partitions of $F$ into pairs.
\end{proof}

%%%%%%%%%%%%%%%%%%%%%%%%%%%%%%%%%%%%%%%%%%%%%%%%%%%%%%%%%%%%%%%%
\subsubsection{Survival of the six remaining blocks}
\label{zhu:subsec:survival}

The following argument is the one of \cite[Lemma 11.3 p.~61]{Rie26b}, adapted to $S$.

\begin{prop}
\label{zhu:prop:seven}
The algebra $S$ has a simple module which is not one of the six modules of Proposition \ref{zhu:prop:a-c0}.
\end{prop}

\begin{proof}
Assume the contrary. Then the branching matrix $\mathsf b$ of $S \subset V'$ has nine rows and six columns $0, b, c_1, c_2, a, c_0$, and by Propositions \ref{zhu:prop:restrictions} and \ref{zhu:prop:a-c0}, the row of $V'$ is $(1,1,0,0,0,0)$ and the other rows have a single entry $1$, in the columns $c_1$, $c_2$, $a$, $c_0$, each of these columns being reached. Hence $\ker \mathsf b = \Cbb (u_0 - u_b)$, where $(u_z)$ is the standard basis indexed by the six columns. By Lemma \ref{zhu:lem:branching}, $\mathsf b\, s^S (u_0 - u_b) = s^{V'} \mathsf b (u_0 - u_b) = 0$, so that $s^S (u_0 - u_b) = \kappa (u_0 - u_b)$ for some $\kappa \in \Cbb$. Since $S$ and $M_b$ are self-contragredient (the latter because $V'$ and $S$ are), $(s^S)^2 (u_0 - u_b) = u_0 - u_b$ and $\kappa = \pm 1$. Write $\delta_b = \operatorname{qdim}_S M_b > 0$. The coordinate $0$ of $s^S (u_0 - u_b) = \kappa (u_0 - u_b)$ gives $s^S_{00}(1 - \delta_b) = \kappa$. The entry $((0,0),0)$ of $s^{V'} \mathsf b = \mathsf b\, s^S$ gives $s^{V'}_{00} = s^S_{00}(1 + \delta_b)$, and $s^{V'}_{00} = \pm \frac13$ since $\sum_g s^{V'}_{0,g}s^{V'}_{g,0} = ((s^{V'})^2)_{00} = 1$ and $s^{V'}_{0,g} = s^{V'}_{g,0} = s^{V'}_{00}$ for the nine modules. Hence $\frac{1 - \delta_b}{1 + \delta_b} = \pm 3$, which forces $\delta_b \in \{-\frac12, -2\}$, a contradiction.
\end{proof}

\begin{thm}
\label{zhu:thm:twelve}
The algebra $S$ has exactly twelve simple modules up to isomorphism. Their lowest weight spaces and lowest conformal weights are given by the twelve rows of Table \ref{zhu:table:blocks}. All of them except $S$ have positive conformal weight. Moreover $\D \cong \Lambda \ot_R \Cbb$ and $\dim \D = 762$.
\end{thm}

\begin{proof}
By Proposition \ref{zhu:prop:seven}, some simple $S$-module $M$ has as lowest weight space one of the six rows $d_j$ of Table \ref{zhu:table:blocks}, with $h_M = \frac{k}{13}$ for some $k \in \{5,6,7,8,11,15\}$ by Proposition \ref{zhu:prop:exact-h}. The genus-one $T$-matrix of $S$ is the diagonal matrix with entries $e^{2\pi i(h - 1/3)}$, $h$ running over the conformal weights of the simple modules, since $c/24 = 1/3$. By Proposition \ref{zhu:prop:exact-h}, these entries are $39$-th roots of unity, and the entry of $M$ is $e^{2\pi i(k/13 - 1/3)} = \xi^{3k - 13}$ with $\xi = e^{2\pi i/39}$, a primitive $39$-th root of unity since $3k - 13$ is prime to $3$ and to $13$. Hence $T$ has order $39$. We follow the argument of \cite[Section 11.5 p.~62]{Rie26b}. By \cite[Theorem 3.10 p.~2142]{DLN15}, the representation of $\mathrm{SL}_2(\Z)$ defined by the trace functions of $S$ is a modular representation of the modular tensor category $\Rep(S)$, and by the Galois symmetry \cite[Theorem II p.~2124]{DLN15}, for every $\gamma \in \operatorname{Gal}(\mathbb{Q}(\xi)/\mathbb{Q})$, the matrix $\gamma^2(T)$ is conjugate to $T$ by a signed permutation matrix. In particular the diagonal of $\gamma^2(T)$ is a permutation of the diagonal of $T$. Let $\gamma(\xi) = \xi^m$ with $m$ prime to $39$. Then $m^2 \equiv 1 \bmod 3$, so that $13 m^2 \equiv 13 \bmod 39$ and $\gamma^2(\xi^{3k-13}) = \xi^{3k m^2 - 13} = e^{2\pi i(k'/13 - 1/3)}$ with $k' \equiv k m^2 \bmod 13$. The residues modulo $13$ of $5, 6, 7, 8, 11, 15$ are $5, 6, 7, 8, 11, 2$, the six non-squares modulo $13$, and $m^2$ runs over the six squares $1, 3, 4, 9, 10, 12$ modulo $13$ when $m$ runs over the units modulo $39$. Since the squares act transitively on the non-squares by multiplication, $k m^2 \bmod 13$ runs over all the non-squares. Hence the six numbers $e^{2\pi i(k'/13 - 1/3)}$, $k' \in \{5,6,7,8,11,15\}$, are diagonal entries of $T$. By Proposition \ref{zhu:prop:exact-h}, the rows $0$, $b$, $c_1$, $c_2$, $a$, $c_0$ have weights in $\frac13\Z$, so that their entries of $T$ are cube roots of unity, which are not primitive $39$-th roots of unity, and the six rows $d_j$ have the six distinct entries $e^{2\pi i(k'/13 - 1/3)}$. Hence each of these six values is the entry of a single row $d_j$. Hence the six rows $d_j$ are realized, and with Proposition \ref{zhu:prop:a-c0} all twelve blocks of $\Lambda \ot_R \Cbb$ survive in the quotient $\D$. Since this quotient is a surjection of semisimple algebras which is bijective on blocks, it is an isomorphism.
\end{proof}

\begin{cor}
\label{zhu:cor:dimension}
Zhu's algebra $A(S)$ is isomorphic to $\prod_{i=1}^{12} \operatorname{End}_\Cbb(R_i)$, where $R_i$ runs over the twelve lowest weight spaces of Table \ref{zhu:table:blocks}. In particular $\dim A(S) = 36625$.
\end{cor}

\begin{proof}
By Theorem \ref{zhu:thm:Z}, $A(S)$ is a finite-dimensional semisimple algebra over $\Cbb$, hence the product of the endomorphism algebras of its simple modules. By \cite[Theorem 2.2.2 p.~261]{Zhu96} and Theorem \ref{zhu:thm:twelve}, these are the twelve lowest weight spaces $R_i$, of dimensions $1$, $80$, $81$, $81$, $3$, $27$, $7$, $42$, $119$, $5$, $13$, $14$, and the sum of their squares is $36625$.
\end{proof}

\begin{remark} \normalfont
\label{zhu:rem:M-proved}
Theorem \ref{zhu:thm:twelve} and Propositions \ref{zhu:prop:restrictions} and \ref{zhu:prop:a-c0} prove hypothesis $\mathrm{(M)}$ of \cite[Remark 11.2]{Arh26}, with the restrictions $W_{0,0}|_S = S \oplus M_b$, $W_{0,\pm1}|_S \cong M_{c_1}$ (the module $P$ of loc. cit.), $W_{\pm1,0}|_S \cong M_{c_2}$ (the module $Q$), $W_{1,1}|_S \cong W_{-1,-1}|_S$ and $W_{1,-1}|_S \cong W_{-1,1}|_S$ (the modules $X_+$ and $X_-$), and the six modules $M_{d_l}$ with lowest weight spaces $R_{d_l}$ of \cite[Table 2]{Arh26} and lowest conformal weights $\frac{6}{13}$, $\frac{11}{13}$, $\frac{15}{13}$, $\frac{5}{13}$, $\frac{7}{13}$, $\frac{8}{13}$ for $l = 1, \ldots, 6$. These weights are congruent to $\frac{6l^2}{13}$ modulo $\Z$, as the twists of the objects $d_l$ of $Z(\cal H_3)$ in the labelling of \cite{EvG11}. The two blocks labelled $a$ and $c_0$ in Table \ref{zhu:table:blocks} have the same lowest weight and the same lowest weight space as an $S_1$-module. The computation identifies their corresponding $S$-modules with $X_+$ and $X_-$ as an unordered pair. The intrinsic labels in the center are distinguished by the Lagrangian algebras of \cite[Proposition 3.6 and Remark 3.7]{ArhII}.
\end{remark}
%%%%%%%%%%%%%%%%%%%%%%%%%%%%%%%%%%%%%%%%%%%%%%%%%%%%%%%%%%%%%%%%
\subsection{The representation category}
\label{zhu:sec:categorical}

%%%%%%%%%%%%%%%%%%%%%%%%%%%%%%%%%%%%%%%%%%%%%%%%%%%%%%%%%%%%%%%%
\subsubsection{Pseudo-unitarity}
\label{zhu:subsec:pseudo}

\begin{prop}
\label{zhu:prop:pseudo}
The category $\Rep(S)$ of $S$-modules is a pseudo-unitary modular tensor category. For every simple $S$-module $M$, the categorical dimension of $M$ and its Frobenius--Perron dimension are equal to the quantum dimension $\operatorname{qdim}_S M = s^S_{M,0}/s^S_{00} > 0$.
\end{prop}

\begin{proof}
By Corollary \ref{zhu:cor:rational} and \cite[Theorem 4.6]{Hua08}, $\Rep(S)$ is a modular tensor category. By Theorem \ref{zhu:thm:twelve}, every simple $S$-module other than $S$ has positive lowest conformal weight. Hence the quantum dimensions are positive \cite[Lemma 4.2 p.~6453]{DJX13}. The normalized $S$-matrix of $\Rep(S)$ and the genus-one matrix $s^S$ coincide up to a sign \cite[Theorem 3.10 p.~2142]{DLN15}, so that the categorical dimension $\tilde s_{i,0}/\tilde s_{0,0}$ of $M_i$ is $s^S_{i,0}/s^S_{00} = \operatorname{qdim}_S M_i$, see also \cite[p.~5]{CKM22}. Let $M_0 = S, \ldots, M_{11}$ be the simple $S$-modules and $N(i)$ the matrix of the fusion rules $N_{i,j}^k$. By \cite[Theorem 4.1 p.~6452]{DJX13}, the column $(s^S_{j,0})_j$ is an eigenvector of $N(i)$ for the eigenvalue $s^S_{i,0}/s^S_{0,0} = \operatorname{qdim}_S M_i$. Since this vector has positive entries and $N(i)$ has nonnegative entries, this eigenvalue is the Perron--Frobenius eigenvalue of $N(i)$, that is, the Frobenius--Perron dimension of $M_i$, as in \cite[Remarks 5.3 and 5.4 p.~6460]{DJX13}. Hence $\dim \Rep(S) = \sum_i (\operatorname{qdim}_S M_i)^2 = \operatorname{FPdim}\Rep(S)$.
\end{proof}

%%%%%%%%%%%%%%%%%%%%%%%%%%%%%%%%%%%%%%%%%%%%%%%%%%%%%%%%%%%%%%%%
\subsubsection{The Lagrangian algebras of the lattice extensions}
\label{zhu:subsec:lagrangian}

The arguments of this section follow the strategy outlined by Huang for the algebra of the other case \cite[Sections 5 and 6 pp.~7--9]{Hua26}, as carried out for Haagerup algebras in \cite[Sections 3 and 6]{ArhII}: Lagrangian algebras attached to the holomorphic lattice extensions, recognition of the fusion rules through the Drinfeld center, and the use of the second isotropic line for the associator of the pointed part. No statement of \cite{Hua26} is used.

The discriminant form of $E_6 \oplus A_2$ has two isotropic lines, generated by $(1,1)$ and $(1,-1)$. We put
\[
H_+ = V' \oplus W_{1,1} \oplus W_{-1,-1}, \qquad H_- = V' \oplus W_{1,-1} \oplus W_{-1,1}.
\]
Each of them is the lattice vertex operator algebra of an even unimodular overlattice of $E_6 \oplus A_2$, which is isometric to $E_8$, so that $H_\pm \cong V_{E_8}$ \cite[Section 2.4 and Lemma 4.1]{ArhII}. It is simple and self-contragredient, and $V'$ is strongly rational with the nine simple modules $W_g$ \cite{Don93}. By Propositions \ref{zhu:prop:restrictions} and \ref{zhu:prop:a-c0},
\[
H_x|_S \cong S \oplus M_b \oplus 2X_x, \qquad V'|_S \cong S \oplus M_b,
\]
where $X_+ = W_{1,1}|_S$ and $X_- = W_{1,-1}|_S$ are the two simple modules of types $a$ and $c_0$.

\begin{lemma}
\label{zhu:lem:qdim-lattice}
For every $g \in \Z_3 \times \Z_3$, we have $\operatorname{qdim}_S (W_g|_S) = \operatorname{qdim}_S(V'|_S) = 1 + \operatorname{qdim}_S M_b$.
\end{lemma}

\begin{proof}
By Proposition \ref{zhu:prop:restrictions}, $S$ occurs in $W_g|_S$ only for $g = 0$, with multiplicity one. The entry $(g,0)$ of the identity $s^{V'} b = b\, s^S$ of Lemma \ref{zhu:lem:branching} gives $s^{V'}_{g,0} = \sum_z b_{gz} s^S_{z,0} = s^S_{00}\operatorname{qdim}_S(W_g|_S)$, and $s^{V'}_{g,0} = s^{V'}_{00}$ for every $g$ (proof of Proposition \ref{zhu:prop:a-c0}).
\end{proof}

\begin{thm}
\label{zhu:thm:categorical}
Let $x \in \{+,-\}$.
\begin{enumerate}
\item The object $L_x = H_x|_S$ carries the structure of a Lagrangian algebra in $\Rep(S)$, and $\Rep(S)$ is braided equivalent to the Drinfeld center $Z(\cal F_x)$ of the fusion category $\cal F_x = \Rep(S)_{L_x}$ of $L_x$-modules.
\item The category $\cal F_x$ is pseudo-unitary and has exactly six simple objects up to isomorphism, three of which are invertible. There are an invertible object $\alpha$ of order three and a self-dual simple object $\rho$ such that the simple objects of $\cal F_x$ are $1$, $\alpha$, $\alpha^2$, $\rho$, $\alpha\rho$, $\alpha^2\rho$, and
\[
\rho\alpha \cong \alpha^2\rho, \qquad \rho^2 \cong 1 \oplus \rho \oplus \alpha\rho \oplus \alpha^2\rho.
\]
\item The fusion subcategory $\cal P_x$ generated by $\alpha$ is equivalent to $\mathrm{Vec}_{\Z_3}$ with trivial associator.
\item If $\cal F_x$ is equivalent to a unitary fusion category, then $\cal F_x$ is equivalent to $\cal H_2$ or $\cal H_3$, $\Rep(S)$ is braided equivalent to $Z(\cal H_3)$, $S$ is a Haagerup algebra, and $\mathrm{(H_J)}$ holds.
\end{enumerate}
\end{thm}

\begin{proof}
1. The algebra $H_x$ is an extension of $S$ with the same conformal vector and $\dim (H_x)_0 = 1$. Since $S$ is strongly rational and positive (Theorem \ref{zhu:thm:twelve}), $H_x$ corresponds to a haploid commutative associative algebra $L_x$ with trivial twist in $\Rep(S)$ \cite[Theorem 3.2 p.~10]{HKL15}, and the category of $H_x$-modules is isomorphic to the category $\Rep(S)^0_{L_x}$ of local $L_x$-modules \cite[Theorem 3.4 p.~16]{HKL15}. Since $H_x$ is simple and the categorical dimension of $L_x$ is positive by Proposition \ref{zhu:prop:pseudo}, the argument recalled in \cite[Section 2.4]{ArhII}, based on \cite[Lemma 1.20 p.~15 and Theorem 3.3 p.~17]{KiO02} and \cite[Proposition 2.7 p.~5]{DMNO13}, shows that $L_x$ is separable, hence a connected étale algebra. By Lemma \ref{zhu:lem:qdim-lattice} and Proposition \ref{zhu:prop:pseudo}, $\operatorname{FPdim}(L_x) = 3\,\operatorname{qdim}_S(V'|_S) = 3 s^{V'}_{00}/s^S_{00}$, with $|s^{V'}_{00}| = \frac13$ (proof of Proposition \ref{zhu:prop:seven}). Since $s^S$ is symmetric and $(s^S)^2$ is the permutation matrix of contragredience \cite[Theorem 4.1 p.~6452]{DJX13}, we have $\sum_i (s^S_{i,0})^2 = 1$ and $\operatorname{FPdim}\Rep(S) = \sum_i (s^S_{i,0}/s^S_{00})^2 = (s^S_{00})^{-2} = \operatorname{FPdim}(L_x)^2$. Hence $L_x$ is Lagrangian \cite[Definition 4.6 p.~17]{DMNO13}, and $\Rep(S)^0_{L_x} \cong \mathrm{Vec}$ \cite[Remark 4.3 p.~16]{DMNO13}, in accordance with the fact that $H_x$ is holomorphic. By \cite[Proposition 2.2]{ArhII}, which recalls \cite[Corollary 4.1 p.~16]{DMNO13}, there is a braided equivalence $\Phi \co \Rep(S) \to Z(\cal F_x)$, and $\Phi(L_x) \cong I(1)$ as algebras, where $I$ is the right adjoint of the forgetful functor $Z(\cal F_x) \to \cal F_x$ \cite[Lemma 3.1]{ArhII}.

2. By \cite[Lemma 3.1]{ArhII}, the number of simple objects of $\cal F_x$ is $\dim \operatorname{Hom}(L_x,L_x) = 1 + 1 + 2^2 = 6$. Since $L_x$ is Lagrangian, $\operatorname{FPdim}(\cal F_x) = \operatorname{FPdim}\Rep(S)/\operatorname{FPdim}(L_x) = \operatorname{FPdim}(L_x)$ \cite[Lemma 3.11 p.~10]{DMNO13}. The global dimension of $Z(\cal F_x)$ is $(\dim \cal F_x)^2$ \cite[Theorem 2.15]{ENO05}, and $\dim \cal F_x > 0$ \cite[Theorem 2.3]{ENO05}. Since a braided equivalence preserves global and Frobenius--Perron dimensions, Proposition \ref{zhu:prop:pseudo} gives $(\dim \cal F_x)^2 = \operatorname{FPdim}\Rep(S) = \operatorname{FPdim}(\cal F_x)^2$, so that $\cal F_x$ is pseudo-unitary. We equip $\cal F_x$ with its spherical structure for which the dimensions are the Frobenius--Perron dimensions \cite[Proposition 8.23 p.~34]{ENO05}. Then $K(\cal F_x) \ot \Cbb \cong \operatorname{End}(I(1)) \cong \operatorname{End}(L_x) \cong \Cbb \oplus \Cbb \oplus \M_2(\Cbb)$ \cite[Corollary 2.16 p.~7]{Ost15}, so that the Grothendieck ring is noncommutative.

By Lemma \ref{zhu:lem:qdim-lattice} and Proposition \ref{zhu:prop:pseudo}, $\operatorname{FPdim}(L_x) = 3(1 + \delta_b)$, where $\delta_b = \operatorname{qdim}_S M_b$. The subobject $A_{V'} = V'|_S$ of $L_x$ is a connected étale subalgebra \cite[Lemma 2.3]{ArhII}, applied to $S \subset V' \subset H_x$, since $V'$ is simple and $H_x$ is self-contragredient. By \cite[Proposition 2.2, part 3]{ArhII}, which recalls \cite[Theorem 4.10 p.~18]{DMNO13}, $\Phi(A_{V'})$ is the étale subalgebra of $I(1)$ attached to a fusion subcategory $\cal B_x$ of $\cal F_x$ with $\operatorname{FPdim}(A_{V'})\operatorname{FPdim}(\cal B_x) = \operatorname{FPdim}(\cal F_x)$, that is $\operatorname{FPdim}(\cal B_x) = 3$. A fusion category of prime Frobenius--Perron dimension is pointed \cite[Corollary 8.30 p.~37]{ENO05}. Hence the group $G$ of isomorphism classes of invertible objects of $\cal F_x$ contains a subgroup of order three, and $|G| \in \{3,6\}$ since $|G| \leq 6$. If $|G| = 6$, then $\cal F_x$ is pointed and $\operatorname{FPdim}(\cal F_x) = 6$, so that $\delta_b = 1$. Then $M_b$ is a simple current \cite[Proposition 4.17 p.~6456]{DJX13}, self-contragredient since $V'$ and $S$ are, so that $M_b \boxtimes M_b \cong S$ and the vertex operator of $V'$ maps $M_b \ot M_b$ into $S((z))$. The linear map $\sigma$ equal to $1$ on $S$ and $-1$ on $M_b$ is then an automorphism of $V'$ with fixed points $S$, and $S_1$ is the fixed-point subalgebra of an involution of $V'_1 \cong \mathfrak e_6 \oplus \mathfrak{sl}_3$. Such an involution preserves both ideals, and the fixed-point subalgebras of the involutions of $\mathfrak e_6$ and of $\mathfrak{sl}_3$ have dimension at least $36$ and $3$ \cite[Chapter 8]{Kac90}, \cite[Chapter X]{Hel78}, which contradicts $\dim S_1 = 6$. Hence $|G| = 3$, and $\cal B_x$ is the pointed part $\cal P_x$ of $\cal F_x$.

The proof of part 1 of \cite[Theorem 6.1]{ArhII}, consisting of its Lemmas 6.2--6.5 and Proposition 6.6, applies with the canonical spherical structure of a pseudo-unitary fusion category. Its inputs are the Frobenius--Perron character, the Ng--Schauenburg formula for indicators in a spherical category, the root-of-unity twists of the center, the dimension and divisibility results of \cite[Propositions 5.4 and 8.12]{ENO05}, the induction and formal-codegree results of \cite[Theorems 2.5 and 2.13]{Ost15}, and the arithmetic of roots of unity used in \cite[Section 6.3]{ArhII}. This part of recognition gives the fusion rules without invoking the classification of unitary fusion categories.

3. By \cite[Proposition 2.2, part 5]{ArhII}, the category of local modules $Z(\cal F_x)^0_{\Phi(A_{V'})}$ is braided equivalent to $Z(\cal P_x)$. On the other hand, $\Rep(S)^0_{A_{V'}}$ is isomorphic to $\Rep(V')$ \cite[Theorem 3.4 p.~16]{HKL15}, and this isomorphism is a braided equivalence for the braided tensor structure of $\Rep(V')$ given by \cite{Hua08} \cite[Theorem 3.65]{CKM24}, see also \cite[p.~5]{CKM22}. Hence $\Rep(V')$ is braided equivalent to $Z(\cal P_x)$, possibly with the reverse braiding, which has the same Lagrangian algebras. The two extensions $H_+$ and $H_-$ of $V'$ define, by the argument of part 1 applied to $V'$, two connected étale algebras in $\Rep(V')$ of Frobenius--Perron dimension $3$, which are Lagrangian since the nine simple $V'$-modules have quantum dimension one (proof of Proposition \ref{zhu:prop:a-c0}), so that $\operatorname{FPdim}\Rep(V') = 9$. These algebras are not isomorphic since their underlying objects $V' \oplus W_{1,1} \oplus W_{-1,-1}$ and $V' \oplus W_{1,-1} \oplus W_{-1,1}$ are not. By \cite[Proposition 2.2, part 2]{ArhII}, the category $\cal P_x$ has two inequivalent indecomposable module categories, and by \cite[Lemma 6.7, part 2]{ArhII} its associator is trivial.

4. Let $\cal F$ be a unitary fusion category equivalent to $\cal F_x$. It has six simple objects, three of which are invertible, a noncommutative Grothendieck ring, and by part 3 the fusion subcategory generated by its invertible objects is equivalent to $\mathrm{Vec}_{\Z_3}$ with trivial associator. By \cite[Theorem 6.1]{ArhII}, $\cal F$ is equivalent to $\cal H_2$ or $\cal H_3$. Since $Z(\cal H_2) \cong Z(\cal H_3)$ \cite[Section 2.1]{ArhII}, $\Rep(S) \cong Z(\cal F_x) \cong Z(\cal H_3)$. With Corollary \ref{zhu:cor:rational} and Theorem \ref{zhu:thm:twelve}, $S$ is a Haagerup algebra, and so is Jiao's algebra by \cite[Theorem 10.24]{Arh26}.
\end{proof}

\begin{remark}\normalfont
\label{zhu:rem:unitarity-gap}
Theorem \ref{zhu:thm:categorical} reduces identification of the center along this categorical route to unitarizability of $\cal F_x$. For the two categories attached to $S$, Theorem \ref{thm:HJ-main} supplies complete unitarity of $S$, and the implication \textup{(v)}$\Rightarrow$\textup{(ii)} of Theorem \ref{thm:unitarity-equivalences} gives the unitarizability of $\cal F_x$. The classification of \cite[Theorems 5.9 and 9.1]{Izu18} applies to $C^*$-fusion categories. The two nonunitary Galois conjugates in the endomorphism family of \cite[Theorem 3]{EvG17} have global dimension $3+3d^{-2}$, different from their Frobenius--Perron dimension $3+3d^2$, with $d=(3+\sqrt{13})/2$. Hence they are not pseudo-unitary and are not equivalent to $\cal F_x$.
\end{remark}

%%%%%%%%%%%%%%%%%%%%%%%%%%%%%%%%%%%%%%%%%%%%%%%%%%%%%%%%%%%%%%%%%%%%%%%%%%%%%
\subsection{Certificates and programs}
\label{zhu:app:certificates}

The archive \nolinkurl{Certificats-Zhu-III-v07.zip} contains the scripts, their inputs, provenance checksums and replay instructions. It separates computations that start from lattice vectors and their modes from checks that start with the projected rows and the exported rewriting data. Its README identifies the production inputs, the two ideal-containment checks, the full confluence check and the scope of each new replay. Historical production logs and new checks have separate directories. The computations are listed below with the statements they support. They are written in Python~3 with the libraries \texttt{python-flint} and \texttt{gmpy2}. Exact arithmetic in $K$ is done on quadruples of rationals in the basis $1, r, q, rq$ (module \texttt{kfast}), and exact linear algebra over $K$ by restriction of scalars to $\mathbb{Q}$ with pivots chosen modulo $\mathfrak p$ and an exact verification of every solution (module \texttt{kexact}). The explicit model of $V'$, the spaces $U_1$, $U_2$ and the exact lattice vertex operators are those of the archive of \cite{Arh26} (modules \texttt{lattice\_voa}, \texttt{exact\_engine}, \texttt{Kfield} and the data of \cite[Theorem 10.21]{Arh26}).
\begin{enumerate}
\item \texttt{z0\_highest.py} regenerates the lowering basis of the explicit $S_2$ and checks the highest vectors and the rank of the basis, using the exact lattice model of \cite{Arh26}.
\item \texttt{prod\_exact.py}, \texttt{bucket.py}: the exact products $a_{(n)}b$, $a, b \in S_2$, $0 \leq n \leq 2$, and the vectors $B(a,b)$ and $N(a,b)$ of Section \ref{zhu:subsec:families}, in the monomial basis of $V'$.
\item \texttt{z2x.py}, output \texttt{z2x.json}, \texttt{rel\_exact.pkl}: Proposition \ref{zhu:prop:exact-families}, parts 1 and 2 (the $1060$ Zhu polynomials, the kernel dimensions $458$, $70$, $532$ over $K$ and modulo $\mathfrak p$, the rank $956$ of $\Phi_3$, and the absence of $p$ in the denominators).
\item \texttt{z3x.py}, output \texttt{z3x.json}, \texttt{rows\_exact.pkl}: Proposition \ref{zhu:prop:exact-families}, part 3 (the $19981$ projected relations in $R\Gamma_Z$ and the $515$ arrows).
\item \texttt{z4x.py}, output \texttt{z4x.json}, \texttt{upper\_exact.pkl}: Theorem \ref{zhu:thm:order} (exact eliminations, rewriting rules, unit and integrality checks at $\mathfrak p$, irreducible paths, resolution of all ambiguities over $K$).
\item \texttt{z4\_upper\_algebra\_modp.py} of \cite{Arh26}, run on \texttt{rows\_exact\_reduced\_modp.pkl}, output \texttt{z4\_upper\_algebra\_modp.json}: Proposition \ref{zhu:prop:special}, parts 1--3.
\item The programs \texttt{specfib\_Z.py} and \texttt{specfib\_blocks.py}, with the outputs \texttt{specfib\_Z.log} and \texttt{specfib\_blocks.json}, prove part 4 of Proposition \ref{zhu:prop:special}: they compute the characteristic polynomial of $L_{\bar Z}$, the kernels of $L_{\bar Z} - c$ and the dimensions of their corners.
\item \texttt{exactZ.py}, output \texttt{exactZ.json}: Proposition \ref{zhu:prop:exact-h}.
\item \texttt{check\_patterns.py}, output \texttt{check\_patterns.json}: a check, with \texttt{sympy}, of the computation of the three partitions in the proof of Proposition \ref{zhu:prop:a-c0}, which is also done by hand there.
\item \texttt{verif\_exact.py}, output \texttt{verif\_exact.json}: an a posteriori verifier, written separately from \texttt{z4x.py}, with its own arithmetic routines in $K$ and its own reduction strategy (leftmost occurrence first, then shortest leading word). Starting from \texttt{rows\_exact.pkl} and from the output \texttt{upper\_exact.pkl} of \texttt{z4x.py}, it checks, in exact arithmetic, the compatibility of the eliminations and rules with the order and the endpoints, the reduction to zero of the $19981$ generators of $I_K$ (the inclusion $I_K \subset J_K$), the resolution of all overlap and inclusion ambiguities, the absence of eliminated arrows in the leading words, the counts of irreducible paths, the integrality at $\mathfrak p$, and the coincidence of the reductions modulo $\mathfrak p$ of the rules with the output of item 6. It does not check the inclusion $J_K \subset I_K$.
\item \texttt{membership\_0\_vertices\_E.py}, \texttt{membership\_a.py} and \texttt{membership\_b.py}, written separately from the production program to certify the reverse containment: a certificate of the inclusion $J_K \subset I_K$ which does not use the bookkeeping of \texttt{z4x.py}. With pivots chosen modulo $\mathfrak p$ and an exact verification of every identity, it shows that the vertices $12$ and $13$, the $46$ elements $x - \ell_x$ and the $5105$ quadratic rules lie in the $K$-span of the generators of $I_K$ multiplied by paths, after substitution of the elements already certified, and that each of the $554$ remaining rules reduces to zero modulo the span of the certified rules and of their reduced overlaps. These scripts use the rules of \texttt{z4x.py} as targets and share its arithmetic modules \texttt{kfast} and \texttt{Kfield}. The replay driver also contains a targeted check of the last step. It normalizes each target by the certified rules and accumulates exact overlap remainders in the corresponding endpoint corner until the targets reduce to zero.
\end{enumerate}
The products of item 2 and the data of the elimination are intermediate files which are recomputed by the programs. The certificates are exact computations. They are not a formal verification of the implementation of the lattice vertex operators, which is the one of \cite{Arh26}.

%%%%%%%%%%%%%%%%%%%%%%%%%%%%%%%%%%%%%%%%%%%%%%%%%%%%%%%%%%%%%%%%


\paragraph{AI statement.}
The author acknowledges the use of AI tools for language polishing, LaTeX editing, exploratory mathematical discussions and feedback, for drafting parts of the arguments, for reading and checking the statements of the cited literature, and for writing and running the computer programs of Section \ref{sec:certificates}. The author takes full responsibility for the content of the manuscript.

\paragraph{Competing interests} The author declares that he has no competing interests.

\paragraph{Data availability.} The ancillary files are \href{https://github.com/cedricarhancet/haagerup-c8/releases/download/article-III-v10-supplements-v01/Calculs-III-A1-A4-C3-v11.zip}{\nolinkurl{Calculs-III-A1-A4-C3-v11.zip}} and \href{https://github.com/cedricarhancet/haagerup-c8/releases/download/article-III-v10-supplements-v01/Certificats-Zhu-III-v07.zip}{\nolinkurl{Certificats-Zhu-III-v07.zip}}. They are distributed in the versioned release \href{https://github.com/cedricarhancet/haagerup-c8/releases/tag/article-III-v10-supplements-v01}{\texttt{article-III-v10-supplements-v01}}, with SHA-256 manifests and replay instructions. Sections \ref{sec:certificates} and \ref{zhu:app:certificates} describe their computational inputs and their scope. The dependencies are fixed to \href{https://github.com/cedricarhancet/haagerup-c8/releases/tag/article-I-v14-supplements-v01}{\texttt{article-I-v14-supplements-v01}} for \cite{Arh26} and \href{https://github.com/cedricarhancet/haagerup-c8/releases/tag/article-II-v14-supplements-v03}{\texttt{article-II-v14-supplements-v03}} for \cite{ArhII}; the version concordance identifies the manuscript sources and their hashes. The archived production outputs and the newly completed checks are identified separately.
\noindent\textit{Author's GitHub profile.} \url{https://github.com/cedricarhancet}


\begin{thebibliography}{AHKZ26}
\raggedright

\bibitem[ABD04]{ABD04}
T. Abe, G. Buhl and C. Dong.
\newblock Rationality, regularity, and $C_2$-cofiniteness.
\newblock Trans. Amer. Math. Soc. 356 (2004), no. 8, 3391--3402.

\bibitem[AHKZ26]{AHKZ26}
J. Albert, Y. Honda, J. Kaidi and Y. Zheng.
\newblock Haagerup symmetry in $(E_8)_1$?
\newblock Phys. Rev. Lett. 136 (2026), 091603. Page and equation references are to arXiv:2512.08225v2.

\bibitem[ArhI]{Arh26}
C. Arhancet.
\newblock Haagerup vertex operator algebras at central charge eight.
\newblock Preprint, 2026, \href{https://arxiv.org/abs/2611.01618}{arXiv:2611.01618}.

\bibitem[ArhII]{ArhII}
C. Arhancet.
\newblock Haagerup symmetries of the $E_8$ theory and conformal nets at central charge eight.
\newblock Preprint, 2026, \href{https://arxiv.org/abs/2611.03560}{arXiv:2611.03560}.

\bibitem[AsH99]{AsH99}
M. Asaeda and U. Haagerup.
\newblock Exotic subfactors of finite depth with Jones indices $(5+\sqrt{13})/2$ and $(5+\sqrt{17})/2$.
\newblock Comm. Math. Phys. 202 (1999), no. 1, 1--63.

\bibitem[BDH88]{BDH88}%
M. Baillet, Y. Denizeau and J.-F. Havet.
\newblock Indice d'une espérance conditionnelle.
\newblock Compos. Math. 66 (1988), no. 2, 199--236. \href{https://www.numdam.org/item/CM_1988__66_2_199_0/}{Numdam}.

\bibitem[Ber78]{Ber78}
G. M. Bergman.
\newblock The diamond lemma for ring theory.
\newblock Adv. in Math. 29 (1978), no. 2, 178--218.

\bibitem[Bis17a]{Bis17a}
M. Bischoff.
\newblock The relation between subfactors arising from conformal nets and the realization of quantum doubles.
\newblock In Proceedings of the 2014 Maui and 2015 Qinhuangdao conferences in honour of Vaughan F. R. Jones' 60th birthday, 15--24, Proc. Centre Math. Appl. Austral. Nat. Univ. 46, Austral. Nat. Univ., Canberra, 2017. Page references are to this published version of arXiv:1511.08931.

\bibitem[Bis17b]{Bis17b}
M. Bischoff.
\newblock Generalized orbifold construction for conformal nets.
\newblock Rev. Math. Phys. 29 (2017), no. 1, 1750002. Page references are to arXiv:1608.00253.

\bibitem[BKL15]{BKL15}
M. Bischoff, Y. Kawahigashi and R. Longo.
\newblock Characterization of 2D rational local conformal nets and its boundary conditions: the maximal case.
\newblock Doc. Math. 20 (2015), 1137--1184. Page references are to arXiv:1410.8848v2.

\bibitem[BKLR15]{BKLR15}
M. Bischoff, Y. Kawahigashi, R. Longo and K.-H. Rehren.
\newblock Tensor categories and endomorphisms of von Neumann algebras. With applications to quantum field theory.
\newblock SpringerBriefs in Mathematical Physics 3, Springer, Cham, 2015. Page references are to arXiv:1407.4793v3.

\bibitem[Buh02]{Buh02}
G. Buhl.
\newblock A spanning set for VOA modules.
\newblock J. Algebra 254 (2002), no. 1, 125--151.

\bibitem[CaM16]{CaM16}
S. Carnahan and M. Miyamoto.
\newblock Regularity of fixed-point vertex operator subalgebras.
\newblock Preprint, arXiv:1603.05645v4 (2018).

\bibitem[Car04]{Car04}
S. Carpi.
\newblock On the representation theory of Virasoro nets.
\newblock Comm. Math. Phys. 244 (2004), no. 2, 261--284. Numbering of arXiv:math/0306425.

\bibitem[CGGH23]{CGGH23}
S. Carpi, T. Gaudio, L. Giorgetti and R. Hillier.
\newblock Haploid algebras in $C^*$-tensor categories and the Schellekens list.
\newblock Comm. Math. Phys. 402 (2023), no. 1, 169--212. Page references are to arXiv:2211.12790v3.

\bibitem[CKLR19]{CKLR19}
T. Creutzig, S. Kanade, A. R. Linshaw and D. Ridout.
\newblock Schur--Weyl duality for Heisenberg cosets.
\newblock Transform. Groups 24 (2019), no. 2, 301--354.

\bibitem[CKLW18]{CKLW18}
S. Carpi, Y. Kawahigashi, R. Longo and M. Weiner.
\newblock From vertex operator algebras to conformal nets and back.
\newblock Mem. Amer. Math. Soc. 254 (2018), no. 1213. Page and equation references are to arXiv:1503.01260v4.

\bibitem[CKM22]{CKM22}
T. Creutzig, S. Kanade and R. McRae.
\newblock Gluing vertex algebras.
\newblock Adv. Math. 396 (2022), Paper No. 108174. Page references are to arXiv:1906.00119v2.

\bibitem[CKM24]{CKM24}
T. Creutzig, S. Kanade and R. McRae.
\newblock Tensor categories for vertex operator superalgebra extensions.
\newblock Mem. Amer. Math. Soc. 295 (2024), no. 1472. Statement numbers refer to arXiv:1705.05017v2.

\bibitem[CoM93]{CoM93}
D. H. Collingwood and W. M. McGovern.
\newblock Nilpotent orbits in semisimple Lie algebras.
\newblock Van Nostrand Reinhold, New York, 1993.

\bibitem[DJX13]{DJX13}
C. Dong, X. Jiao and F. Xu.
\newblock Quantum dimensions and quantum Galois theory.
\newblock Trans. Amer. Math. Soc. 365 (2013), no. 12, 6441--6469. Page references are to arXiv:1201.2738v1.

\bibitem[DJX26]{DJX26}
C. Dong, X. Jiao and F. Xu.
\newblock On the Haagerup vertex operator algebras.
\newblock Preprint, arXiv:2611.04778v1 (2026).

\bibitem[DLN15]{DLN15}
C. Dong, X. Lin and S.-H. Ng.
\newblock Congruence property in conformal field theory.
\newblock Algebra Number Theory 9 (2015), no. 9, 2121--2166.

\bibitem[DMNO13]{DMNO13}
A. Davydov, M. Müger, D. Nikshych and V. Ostrik.
\newblock The Witt group of non-degenerate braided fusion categories.
\newblock J. Reine Angew. Math. 677 (2013), 135--177. Page references are to arXiv:1009.2117v2.

\bibitem[DoL14]{DoL14}
C. Dong and X. Lin.
\newblock Unitary vertex operator algebras.
\newblock J. Algebra 397 (2014), 252--277. Page references are to arXiv:1308.2361v1.

\bibitem[DoL93]{DL93}
C. Dong and J. Lepowsky.
\newblock Generalized vertex algebras and relative vertex operators.
\newblock Progress in Mathematics 112, Birkh\"auser, Boston, 1993.

\bibitem[DoM04a]{DoM04}
C. Dong and G. Mason.
\newblock Rational vertex operator algebras and the effective central charge.
\newblock Int. Math. Res. Not. 2004, no. 56, 2989--3008. Page references are to arXiv:math/0201318.

\bibitem[DoM04b]{DoM04b}
C. Dong and G. Mason.
\newblock Holomorphic vertex operator algebras of small central charge.
\newblock Pacific J. Math. 213 (2004), no. 2, 253--266. Page references are to arXiv:math/0203005.

\bibitem[DoM97]{DoM97}
C. Dong and G. Mason.
\newblock On quantum Galois theory.
\newblock Duke Math. J. 86 (1997), no. 2, 305--321. Numbering of arXiv:hep-th/9412037.

\bibitem[Don93]{Don93}
C. Dong.
\newblock Vertex algebras associated with even lattices.
\newblock J. Algebra 161 (1993), no. 1, 245--265.

\bibitem[DoX06]{DoX06}
C. Dong and F. Xu.
\newblock Conformal nets associated with lattices and their orbifolds.
\newblock Adv. Math. 206 (2006), no. 1, 279--306. Page references are to arXiv:math/0411499v2.

\bibitem[EGNO15]{EGNO15}
P. Etingof, S. Gelaki, D. Nikshych and V. Ostrik.
\newblock Tensor categories.
\newblock Mathematical Surveys and Monographs 205, American Mathematical Society, Providence, RI, 2015.

\bibitem[ENO05]{ENO05}
P. Etingof, D. Nikshych and V. Ostrik.
\newblock On fusion categories.
\newblock Ann. of Math. (2) 162 (2005), no. 2, 581--642. Page references are to arXiv:math/0203060v11.

\bibitem[ENO10]{ENO10}
P. Etingof, D. Nikshych and V. Ostrik.
\newblock Fusion categories and homotopy theory.
\newblock Quantum Topol. 1 (2010), no. 3, 209--273. Page references are to arXiv:0909.3140v2.

\bibitem[ENO11]{ENO11}
P. Etingof, D. Nikshych and V. Ostrik.
\newblock Weakly group-theoretical and solvable fusion categories.
\newblock Adv. Math. 226 (2011), no. 1, 176--205. Page references are to arXiv:0809.3031v2.

\bibitem[EvG11]{EvG11}
D. E. Evans and T. Gannon.
\newblock The exoticness and realisability of twisted Haagerup--Izumi modular data.
\newblock Comm. Math. Phys. 307 (2011), no. 2, 463--512. Page and equation references are to arXiv:1006.1326v1.

\bibitem[EvG17]{EvG17}
D. E. Evans and T. Gannon.
\newblock Non-unitary fusion categories and their doubles via endomorphisms.
\newblock Adv. Math. 310 (2017), 1--43.

\bibitem[FHL93]{FHL93}
I. B. Frenkel, Y.-Z. Huang and J. Lepowsky.
\newblock On axiomatic approaches to vertex operator algebras and modules.
\newblock Mem. Amer. Math. Soc. 104 (1993), no. 494.

\bibitem[GrS12]{GrS12}
P. Grossman and N. Snyder.
\newblock Quantum subgroups of the Haagerup fusion categories.
\newblock Comm. Math. Phys. 311 (2012), no. 3, 617--643. Page references are to arXiv:1102.2631v2.

\bibitem[Gui19]{Gui19}
B. Gui.
\newblock Unitarity of the modular tensor categories associated to unitary vertex operator algebras, II.
\newblock Comm. Math. Phys. 372 (2019), no. 3, 893--950. Page references are to arXiv:1712.04931v6.

\bibitem[Gui21]{Gui21a}
B. Gui.
\newblock Categorical extensions of conformal nets.
\newblock Comm. Math. Phys. 383 (2021), no. 2, 763--839. Page references are to arXiv:1812.04470v8.

\bibitem[Gui22]{Gui22}
B. Gui.
\newblock Q-systems and extensions of completely unitary vertex operator algebras.
\newblock Int. Math. Res. Not. IMRN 2022, no. 10, 7550--7614. Page references are to arXiv:1908.03484v3.

\bibitem[Gui25]{Gui25}
B. Gui.
\newblock Comparison of extensions of unitary vertex operator algebras and conformal nets.
\newblock Preprint, arXiv:2505.03235v3.

\bibitem[Gui26]{Gui26}
B. Gui.
\newblock Unbounded field operators in categorical extensions of conformal nets.
\newblock Invent. Math., published online 2026. Page references are to arXiv:2001.03095v5.

\bibitem[GuL96]{GuL96}
D. Guido and R. Longo.
\newblock The conformal spin and statistics theorem.
\newblock Comm. Math. Phys. 181 (1996), no. 1, 11--35.

\bibitem[Haa94]{Haa94}
U. Haagerup.
\newblock Principal graphs of subfactors in the index range $4 < [M:N] < 3+\sqrt2$.
\newblock In: Subfactors (Kyuzeso, 1993), World Sci. Publ., River Edge, NJ, 1994, 1--38.

\bibitem[Hel78]{Hel78}
S. Helgason.
\newblock Differential geometry, Lie groups, and symmetric spaces.
\newblock Pure and Applied Mathematics 80. Academic Press, New York--London, 1978.

\bibitem[Hen19]{Hen19}
A. Henriques.
\newblock Loop groups and diffeomorphism groups of the circle as colimits.
\newblock Comm. Math. Phys. 366 (2019), no. 2, 537--565.

\bibitem[GuiF25]{GuiF25}
B. Gui.
\newblock On a Connes fusion approach to finite index extensions of conformal nets.
\newblock Preprint, arXiv:2112.15396v2, 2025.

\bibitem[HeT25]{HeT25}
A. G. Henriques and J. E. Tener.
\newblock Every conformal net has an associated unitary VOA.
\newblock Preprint, arXiv:2507.20735v3.

\bibitem[HeT26]{HeT26}
A. G. Henriques and J. E. Tener.
\newblock A genus zero functorial CFT associated to the vacuum sector of a conformal net.
\newblock Preprint, arXiv:2608.24437v1, 2026.

\bibitem[HKL15]{HKL15}
Y.-Z. Huang, A. Kirillov Jr. and J. Lepowsky.
\newblock Braided tensor categories and extensions of vertex operator algebras.
\newblock Comm. Math. Phys. 337 (2015), no. 3, 1143--1159. Page references are to arXiv:1406.3420v1.

\bibitem[Hua05]{Hua05}
Y.-Z. Huang.
\newblock Differential equations and intertwining operators.
\newblock Commun. Contemp. Math. 7 (2005), no. 3, 375--400.

\bibitem[Hua08]{Hua08}
Y.-Z. Huang.
\newblock Rigidity and modularity of vertex tensor categories.
\newblock Commun. Contemp. Math. 10 (2008), suppl. 1, 871--911.

\bibitem[Hua08a]{Hua08b}
Y.-Z. Huang.
\newblock Vertex operator algebras and the Verlinde conjecture.
\newblock Commun. Contemp. Math. 10 (2008), no. 1, 103--154.

\bibitem[Hua26b]{Hua26}
T.-C. Huang.
\newblock Haagerup conformal field theory at central charge eight.
\newblock Preprint, arXiv:2610.14441v2. The results are announced, with an outline of proof.

\bibitem[HuL95]{HL95}
Y.-Z. Huang and J. Lepowsky.
\newblock A theory of tensor products for module categories for a vertex operator algebra, I, II.
\newblock Selecta Math. (N.S.) 1 (1995), no. 4, 699--756, 757--786.

\bibitem[Izu01]{Izu01}
M. Izumi.
\newblock The structure of sectors associated with Longo--Rehren inclusions. II. Examples.
\newblock Rev. Math. Phys. 13 (2001), no. 5, 603--674.

\bibitem[Izu18]{Izu18}
M. Izumi.
\newblock The classification of $3^n$ subfactors and related fusion categories.
\newblock Quantum Topol. 9 (2018), no. 3, 473--562. Page references are to arXiv:1609.07604v2.

\bibitem[Jia26]{Jia26}
X. Jiao.
\newblock An explicit lattice construction for the Haagerup center.
\newblock Preprint, October 4, 2026, 54 pages, \url{https://www.researchgate.net/publication/415561299_An_explicit_lattice_construction_for_the_Haagerup_center}.

\bibitem[JMS14]{JMS14}
V. F. R. Jones, S. Morrison and N. Snyder.
\newblock The classification of subfactors of index at most 5.
\newblock Bull. Amer. Math. Soc. (N.S.) 51 (2014), no. 2, 277--327. Page references are to arXiv:1304.6141.

\bibitem[Kac90]{Kac90}
V. G. Kac.
\newblock Infinite-dimensional Lie algebras. Third edition.
\newblock Cambridge University Press, Cambridge, 1990.

\bibitem[KaL05]{KL05}
Y. Kawahigashi and R. Longo.
\newblock Noncommutative spectral invariants and black hole entropy.
\newblock Comm. Math. Phys. 257 (2005), no. 1, 193--225.

\bibitem[KiO02]{KiO02}
A. Kirillov Jr. and V. Ostrik.
\newblock On a $q$-analogue of the McKay correspondence and the ADE classification of $\widehat{\mathfrak{sl}}_2$ conformal field theories.
\newblock Adv. Math. 171 (2002), no. 2, 183--227. Page references are to arXiv:math/0101219v3.

\bibitem[KLM01]{KLM01}
Y. Kawahigashi, R. Longo and M. Müger.
\newblock Multi-interval subfactors and modularity of representations in conformal field theory.
\newblock Comm. Math. Phys. 219 (2001), no. 3, 631--669. Page references are to arXiv:math/9903104v2.

\bibitem[Li94]{Li94}
H. Li.
\newblock Symmetric invariant bilinear forms on vertex operator algebras.
\newblock J. Pure Appl. Algebra 96 (1994), no. 3, 279--297.

\bibitem[Lon03]{Lon03}
R. Longo.
\newblock Conformal subnets and intermediate subfactors.
\newblock Comm. Math. Phys. 237 (2003), no. 1-2, 7--30. Page references are to arXiv:math/0102196v2.

\bibitem[LoR95]{LoR95}
R. Longo and K.-H. Rehren.
\newblock Nets of subfactors.
\newblock Rev. Math. Phys. 7 (1995), no. 4, 567--597. Page references are to arXiv:hep-th/9411077v2.

\bibitem[Mas14]{Mas14}
G. Mason.
\newblock Lattice subalgebras of strongly regular vertex operator algebras.
\newblock In: Conformal field theory, automorphic forms and related topics, Contrib. Math. Comput. Sci. 8, Springer, Heidelberg, 2014, 31--53. Page references are to arXiv:1110.0544.

\bibitem[McR21]{McR21}
R. McRae.
\newblock On rationality for $C_2$-cofinite vertex operator algebras.
\newblock Preprint, arXiv:2108.01898v3.

\bibitem[McR22]{McR22}
R. McRae.
\newblock On semisimplicity of module categories for finite non-zero index vertex operator subalgebras.
\newblock Lett. Math. Phys. 112 (2022), Paper No. 25.

\bibitem[MoR24a]{MoR24a}
S. Möller and B. C. Rayhaun.
\newblock Equivalence relations on vertex operator algebras, I: Genus.
\newblock Preprint, arXiv:2408.07117v1 (2024).

\bibitem[Ost15]{Ost15}
V. Ostrik.
\newblock Pivotal fusion categories of rank 3.
\newblock Mosc. Math. J. 15 (2015), no. 2, 373--396. Page references are to arXiv:1309.4822v1.

\bibitem[Reu23]{Reu23}
D. Reutter.
\newblock Uniqueness of unitary structure for unitarizable fusion categories.
\newblock Comm. Math. Phys. 397 (2023), no. 1, 37--52. Page references are to arXiv:1906.09710v2.

\bibitem[Rie26a]{Rie26}
A. Riesen.
\newblock The Haagerup VOA at $c = 8$.
\newblock Research announcement, arXiv:2610.13933v1 (October 9, 2026).

\bibitem[Rie26b]{Rie26b}
A. Riesen.
\newblock The Haagerup VOA at central charge eight.
\newblock Working manuscript \texttt{haagerup\_readable\_proof.pdf}, release \texttt{snapshot-2026-10-09}, \url{https://github.com/Doctorfm98/Haagerup_VOA/releases/tag/snapshot-2026-10-09}. Section, page and statement numbers refer to this release.

\bibitem[RTT22]{RTT22}
C. Raymond, Y. Tanimoto and J. E. Tener.
\newblock Unitary vertex algebras and Wightman conformal field theories.
\newblock Comm. Math. Phys. 395 (2022), no. 1, 299--330. Page references are to arXiv:2203.10795v1.

\bibitem[Ten19]{Ten19}
J. E. Tener.
\newblock Representation theory in chiral conformal field theory: from fields to observables.
\newblock Selecta Math. (N.S.) 25 (2019), no. 5, Paper No. 76. Page references are to arXiv:1810.08168v2.

\bibitem[Ten24]{Ten24}
J. E. Tener.
\newblock Fusion and positivity in chiral conformal field theory.
\newblock Geom. Funct. Anal. 34 (2024), no. 4, 1226--1296. Page references are to arXiv:1910.08257v3.

\bibitem[Wed05]{Wed05}
J. H. M. Wedderburn.
\newblock A theorem on finite algebras.
\newblock Trans. Amer. Math. Soc. 6 (1905), no. 3, 349--352.

\bibitem[Wei06]{Wei06}
M. Weiner.
\newblock Conformal covariance and positivity of energy in charged sectors.
\newblock Comm. Math. Phys. 265 (2006), no. 2, 493--506.

\bibitem[Zhu96]{Zhu96}
Y. Zhu.
\newblock Modular invariance of characters of vertex operator algebras.
\newblock J. Amer. Math. Soc. 9 (1996), no. 1, 237--302.

\end{thebibliography}

\end{document}
